% The real-rooting time of heat-evolved random polynomials.
% Assembled on 21 October 2026 from the companion drafts in ../paper/ (see ANNALS-PLAN.md there):
%   Part I   = paper-I.tex   lines 140-546  (Sections 2-5 and the proof of the main theorem)
%   Part II  = paper-II.tex  lines 183-670  (Sections 2.1-2.5) and 834-1091 (Section 3)
%   Part III = paper-III.tex lines 155-476  (Section 2) and 864-1598 (Section 4)
% The introduction, the questions section and the bibliography are new or merged. Cross-paper citations
% were converted to internal references by script; the sentences that referred to material now in the
% companion papers [A], [B], [C] were rewritten by hand and are marked "% MERGE:" in the source.
\documentclass[11pt]{amsart}
\usepackage[margin=1.1in]{geometry}
\usepackage{amsmath,amssymb,amsthm}
\usepackage[colorlinks=true,linkcolor=blue,citecolor=blue,urlcolor=blue]{hyperref}

\newtheorem{theorem}{Theorem}[section]
\newtheorem{corollary}[theorem]{Corollary}
\newtheorem{proposition}[theorem]{Proposition}
\newtheorem{lemma}[theorem]{Lemma}
\theoremstyle{definition}
\newtheorem{definition}[theorem]{Definition}
\newtheorem{hypothesis}[theorem]{Hypothesis}
\newtheorem{question}[theorem]{Question}
\theoremstyle{remark}
\newtheorem{remark}[theorem]{Remark}
\numberwithin{equation}{section}

\newcommand{\C}{\mathbb{C}}
\newcommand{\R}{\mathbb{R}}
\newcommand{\N}{\mathbb{N}}
\newcommand{\Q}{\mathbb{Q}}
\newcommand{\Z}{\mathbb{Z}}
\newcommand{\D}{\mathbb{D}}
\newcommand{\E}{\mathbb{E}}
\newcommand{\PP}{\mathbb{P}}
\newcommand{\eps}{\varepsilon}
\newcommand{\He}{\mathrm{He}}
\DeclareMathOperator{\re}{Re}
\DeclareMathOperator{\im}{Im}
\DeclareMathOperator{\dist}{dist}
\DeclareMathOperator{\sech}{sech}
\DeclareMathOperator{\arcosh}{arcosh}

\title[The real-rooting time of heat-evolved random polynomials]{The real-rooting time of heat-evolved random polynomials}
\author{B. Baggins}
\date{Draft, October 2026}
\subjclass[2020]{Primary 26C10, 60G55; Secondary 30B20, 30C15, 33C45, 35K05, 47B10, 60G15}
\keywords{Heat flow, real zeros, Weyl polynomials, random power series, de Bruijn--Newman constant, Hermite polynomials, linear combinations of orthogonal polynomials, collision of zeros, Kepler's equation, Pfaffian point processes, Fredholm determinants}

\begin{document}

\begin{abstract}
Let $G_{n,t}=\exp(-\frac{t}{2n}\partial_z^2)G_n$ be a random Weyl polynomial of degree $n$ under the backward heat flow, with independent coefficients of a law that is not a point mass and has finite variance. Hall, Ho, Jalowy and Kabluchko showed that the empirical zero distribution collapses onto the real line at $t=1$. We show that the first time at which all zeros of $G_{n,t}$ are real, a finite-degree analogue of the de Bruijn--Newman constant, converges in distribution to $T_*=R_*^{-2}$, where $R_*$ is the smallest modulus of a nonreal zero of the random power series $\sum_j\eta_jw^j$ with coefficients of the same law; the convergence is almost sure under a coupling of the degrees that shares the top coefficients, whereas under the standard coupling the set of limit points of the real-rooting time is almost surely $[1,\infty]$. The proof rests on an exact count of the nonreal zeros of analytic combinations $\sum_{j\le n}a_jh_{n-j}$ of orthonormal Hermite polynomials: if the real power series $a$ has radius of convergence greater than $1$, no zeros on the unit circle and simple zeros in the unit disk, then for all large $n$ the combination has exactly as many nonreal zeros as $a$ has in the disk. Its method expresses the polynomial and its derivative in the bulk through a $2\times2$ matrix applied to $(h_n,h_{n-1})$, built from only $O(\log n)$ steps of the backward three-term recurrence, with the remaining terms controlled by a Christoffel--Darboux bound. For Gaussian coefficients the law of $T_*$ is a Fredholm determinant, $\PP(T_*>t)\sim\frac4{3\pi}t^{-3/2}$ as $t\to\infty$, and $(t-1)\log\PP(T_*\le t)\to-\pi^2/12$ as $t\downarrow1$, half of the constant for complex coefficients. At the scale $1/n$ the last pair of nonreal zeros collides inside the real bulk, at relative position $\cos\theta_*$, where $R_*e^{\pm i\theta_*}$ are the nonreal zeros of minimal modulus, and the real-rooting time is within $O(1/n)$ of $T_*$. In a window of size $1/n$ the polynomial is, up to a positive factor and $O(1/n)$, the model $\zeta\cos(\zeta+c)+\lambda\sin(\zeta+c)$, whose threshold for real-rootedness is given by Kepler's equation; the phase $c$ rotates with $n$, so $n(\Lambda_n-T_*)$ oscillates over an explicit interval and has an explicit limit law jointly with $T_*$.
\end{abstract}

\maketitle

\section{Introduction}\label{sec:intro}

\subsection{Heat-evolved Weyl polynomials and their real-rooting time}
Let $\xi_0,\xi_1,\dots$ be independent standard real Gaussian random variables and let
\begin{equation}\label{eq:weyl}
G_n(z)=\sum_{k=0}^{n}\xi_k\frac{(\sqrt n\,z)^k}{\sqrt{k!}}
\end{equation}
be the real Weyl polynomial of degree $n$, whose zeros are asymptotically uniformly distributed in the unit disk. Let it evolve under the backward heat flow,
\begin{equation}\label{eq:heat}
G_{n,t}=\exp\Bigl(-\frac{t}{2n}\,\partial_z^2\Bigr)G_n=\sum_{i\ge0}\frac1{i!}\Bigl(-\frac{t}{2n}\Bigr)^{i}\partial_z^{2i}G_n,\qquad t\ge0,
\end{equation}
a finite sum. Hall, Ho, Jalowy and Kabluchko \cite[Theorem~2.1]{HHJK} proved (for much more general coefficients) that the empirical distribution of the zeros of $G_{n,t}$ converges in probability to the uniform distribution on the ellipse $\{x^2/(1+t)^2+y^2/(1-t)^2\le1\}$ for $0\le t<1$, and to the semicircle law on $[-2\sqrt t,2\sqrt t]$ for $t\ge1$. Macroscopically the zeros thus become real at time $1$.

The operators $\exp(-\frac{s}{2n}\partial_z^2)$, $s\ge0$, preserve real-rootedness (Lemma~\ref{lem:heat}), so the set of times at which $G_{n,t}$ has only real zeros is an interval $[\Lambda_n^{\rm std},\infty)$ (we put $\Lambda_n^{\rm std}=0$ if $G_n\equiv0$). The \emph{real-rooting time} $\Lambda_n^{\rm std}$ is a finite-degree analogue of the de Bruijn--Newman constant \cite{dB,Ne,RT}; the superscript refers to the \emph{standard coupling} of the degrees, in which one sequence $(\xi_k)$ serves for every $n$. Weak convergence of empirical measures cannot see a bounded number of nonreal zeros, so the macroscopic picture says nothing about $\Lambda_n^{\rm std}$ beyond $\PP(\Lambda_n^{\rm std}>t)\to1$ for each $t<1$, which Proposition~\ref{prop:early} sharpens to an exponential bound. This paper identifies the limit law of $\Lambda_n^{\rm std}$ (Theorem~\ref{cor:weyl}), and then describes, at the scale $1/n$, how the last pair of nonreal zeros becomes real (Theorems~\ref{thm:collision-intro}--\ref{thm:jointlaw-intro}).

The results on the limit law hold for general coefficients. From now on let $\mu$ be a Borel probability measure on $\R$ that is not a point mass and has $\int x^2\,d\mu(x)<\infty$, and let the $\xi_k$ in \eqref{eq:weyl} be independent with law $\mu$. The standard Gaussian law remains the main example, and Theorems~\ref{thm:collision-intro}--\ref{thm:jointlaw-intro}, \ref{thm:sharp-intro} and~\ref{thm:std} and Proposition~\ref{prop:early} concern only this case. Let $\eta_0,\eta_1,\dots$ be independent random variables with law $\mu$ and
\begin{equation}\label{eq:kac}
q(w)=\sum_{j\ge0}\eta_jw^j ,
\end{equation}
a random power series that almost surely has radius of convergence $1$; we write $\D=\{|w|<1\}$. For standard Gaussian coefficients we call $q$ the \emph{Gaussian power series}. Almost surely $q$ has nonreal zeros in $\D$ (Proposition~\ref{prop:nonreal}), and we let
\begin{equation}\label{eq:Rstar}
R_*=\min\{|w|:\ w\in\D\setminus\R,\ q(w)=0\},\qquad T_*=R_*^{-2}\in(1,\infty).
\end{equation}
Let $\mathcal E$ be the set of $t>1$ for which $q$ has a zero on the circle $\{|w|=t^{-1/2}\}$ with positive probability. It is countable, and it is empty if $\mu$ has no atoms (Lemma~\ref{lem:circle}).

\begin{theorem}[Real-rooting time of heat-evolved Weyl polynomials]\label{cor:weyl}
Let $\Lambda_n^{\rm std}$ be the real-rooting time of $G_{n,t}$. Then $\Lambda_n^{\rm std}\to T_*$ in distribution as $n\to\infty$. Moreover, for every $t>1$ with $t\notin\mathcal E$, in particular for every $t>1$ if $\mu$ has no atoms, the number of nonreal zeros of $G_{n,t}$ converges in distribution to the number of nonreal zeros of $q$ in the disk $\{|w|<t^{-1/2}\}$, and
\[
\lim_{n\to\infty}\PP\bigl(G_{n,t}\text{ has only real zeros}\bigr)=\PP(T_*\le t)=\PP\bigl(q\text{ has no zero in }\{|w|<t^{-1/2},\ \im w>0\}\bigr).
\]
\end{theorem}

\begin{proposition}[The limit law]\label{prop:law}
Almost surely $1<T_*<\infty$. The atoms of the law of $T_*$ lie in $\mathcal E$; in particular, if $\mu$ has no atoms, the distribution function of $T_*$ is continuous. For standard Gaussian coefficients,
\[
\PP(T_*>t)=\frac4{3\pi}\,t^{-3/2}\bigl(1+o(1)\bigr)\quad(t\to\infty),\qquad\PP(T_*>t)\le\varkappa(t^{-1/2})\quad(t>1),
\]
where $\varkappa(r)=\frac1\pi\bigl[\frac{1+r^2}{1-r^2}\arctan r-\operatorname{artanh}r\bigr]=\frac4{3\pi}r^3\bigl(1+r^2+O(r^4)\bigr)$ is the expected number of zeros of $q$ in $\{|w|<r,\ \im w>0\}$. In particular $\PP(T_*>t)>0$ for every $t$.
\end{proposition}

So the last nonreal pair of zeros survives macroscopic collapse by a random amount of time, governed by the nonreal zero of $q$ closest to the origin. The limit law depends on $\mu$: if $\mu$ has a bounded density $f$ that is continuous at $0$, then $t^{3/2}\PP(T_*>t)\to\frac83f(0)^2\int x^2\,d\mu$ (Theorem~\ref{thm:tail}), and for Rademacher coefficients even $T_*\le4$ (Remark~\ref{rem:nonuniversal}). Theorem~\ref{cor:weyl} is proved in Section~\ref{sec:heat}, as a consequence of an almost sure statement under another coupling of the degrees (Theorem~\ref{thm:heat}), which in turn rests on a deterministic theorem about linear combinations of Hermite polynomials.

\subsection{Exact counts of nonreal zeros}\label{sec:intro-count}
Throughout, $\He_k$ are the probabilists' Hermite polynomials and
\[
h_k=\frac{\He_k}{\sqrt{k!}}
\]
are the orthonormal polynomials for the standard Gaussian measure $(2\pi)^{-1/2}e^{-x^2/2}\,dx$.

\begin{definition}\label{def:admissible}
A power series $a(w)=\sum_{j\ge0}a_jw^j$ is \emph{admissible} if its coefficients are real, its radius of convergence is greater than $1$, $a(0)\neq0$, $a$ has no zeros on the unit circle, and all zeros of $a$ in $\D$ are simple.
\end{definition}

An admissible $a$ has finitely many zeros $w_1,\dots,w_k$ in $\D$; they are nonzero, and the set is closed under complex conjugation. Let $\zeta_i=w_i+w_i^{-1}$. The Joukowski map $w\mapsto w+w^{-1}$ maps $\D\setminus\{0\}$ conformally onto $\C\setminus[-2,2]$, and maps $(-1,1)\setminus\{0\}$ onto $\R\setminus[-2,2]$, so the $\zeta_i$ are distinct points off $[-2,2]$ and $\zeta_i$ is real exactly when $w_i$ is.

\begin{theorem}[Exact zero count]\label{thm:main}
Let $a$ be admissible, with zeros $w_1,\dots,w_k$ in $\D$, and let
\[
Q_n(x)=\sum_{j=0}^{n}a_j\,h_{n-j}(x).
\]
For every sufficiently small $\eps>0$ there is $n_0$ such that for every $n\ge n_0$ all zeros of the polynomial $z\mapsto Q_n(\sqrt n\,z)$ are simple, and they are:
\begin{enumerate}
\item[(a)] for each $i$, exactly one zero in the disk $|z-\zeta_i|<\eps$, which is real if $w_i$ is real and nonreal if $w_i$ is nonreal;
\item[(b)] $n-k$ further zeros, all real and lying in $(-2-\eps,2+\eps)$.
\end{enumerate}
In particular, for all large $n$ the number of nonreal zeros of $Q_n$ equals the number of nonreal zeros of $a$ in $\D$.
\end{theorem}

When $a$ is a polynomial the theorem concerns fixed finite combinations of consecutive orthonormal Hermite polynomials. The case of interest here is when $a$ is a genuine power series, possibly with radius of convergence arbitrarily close to $1$; then $Q_n$ involves every $h_k$, $0\le k\le n$, and the coefficients $a_j$ may decay at an exponential rate arbitrarily close to $1$. Exterior asymptotics place one zero near each $\zeta_i$; the difficulty is to exclude pairs of nonreal zeros approaching the real bulk, which they cannot detect. Theorem~\ref{thm:main} is proved in Part~I (Sections~\ref{sec:prelim}--\ref{sec:proof}).

The heat flow enters through the following identity. With $q$ as in \eqref{eq:kac}, define the \emph{reversed} Weyl polynomials and their heat evolutions
\begin{equation}\label{eq:reversed}
P_n(z)=\sum_{k=0}^{n}\eta_{n-k}\frac{(\sqrt n\,z)^k}{\sqrt{k!}},\qquad P_{n,t}=\exp\Bigl(-\frac{t}{2n}\partial_z^2\Bigr)P_n .
\end{equation}
For each $n$ the vector $(\eta_n,\dots,\eta_0)$ has the law of $(\xi_0,\dots,\xi_n)$, so $(P_{n,t})_{t\ge0}$ has the law of $(G_{n,t})_{t\ge0}$; only the coupling between different degrees differs. In \eqref{eq:reversed} the top coefficients $\eta_0,\eta_1,\dots$ are shared by all degrees. Since $\exp(-\frac12\partial_x^2)x^k=\He_k(x)$, one finds (Lemma~\ref{lem:heat}) that for $t>0$
\begin{equation}\label{eq:heat-hermite}
P_{n,t}(z)=t^{n/2}\sum_{j=0}^{n}\eta_j\,t^{-j/2}\,h_{n-j}\Bigl(\sqrt{\tfrac nt}\,z\Bigr).
\end{equation}
Up to the factor $t^{n/2}$ and the dilation $z\mapsto z/\sqrt t$, this is the polynomial $Q_n(\sqrt n\,\cdot)$ of Theorem~\ref{thm:main} for $a(w)=q(w/\sqrt t)$, which has radius of convergence $\sqrt t>1$ when $t>1$. The weights $1/\sqrt{k!}$ of the Weyl polynomials enter here: under the heat flow every real polynomial becomes a combination of the $h_{n-j}$, but for other coefficient profiles the coefficients of this combination depend on $n$; such profiles, and a Laguerre analogue of the heat flow, are treated in the companion paper \cite{CompC}.

\begin{theorem}[Almost sure real-rooting time]\label{thm:heat}
Let $\Lambda_n=\inf\{t\ge0:\ P_{n,t}\text{ has only real zeros}\}$, with $\Lambda_n=0$ if $P_n\equiv0$. Almost surely all zeros of $q$ in $\D\setminus\{0\}$ are simple and $q$ has infinitely many nonreal zeros in $\D$, so that $R_*$ and $T_*$ in \eqref{eq:Rstar} are defined, and
\[
\lim_{n\to\infty}\Lambda_n=T_* .
\]
Moreover, fix $t>1$ with $t\notin\mathcal E$. Almost surely, for every sufficiently small $\eps>0$ (depending on $q$) and all sufficiently large $n$: all zeros of $P_{n,t}$ are simple; for each nonreal zero $u$ of $q$ with $|u|<t^{-1/2}$ exactly one zero of $P_{n,t}$ lies within distance $\eps$ of $tu+u^{-1}$; and all other zeros of $P_{n,t}$ are real.
\end{theorem}

\subsection{The final collision}\label{sec:intro-collision}
In Sections~\ref{sec:intro-collision}--\ref{sec:intro-standard} the coefficients are standard Gaussian. Let $u_*=R_*e^{i\theta_*}$, $0<\theta_*<\pi$, be a nonreal zero of $q$ of modulus $R_*$. For fixed $t\in(1,T_*)$ it produces the nonreal zero of $P_{n,t}$ near $tu_*+u_*^{-1}$ of Theorem~\ref{thm:heat}, and as $t\uparrow T_*$ this point tends to $2R_*^{-1}\cos\theta_*$, an interior point of the limiting bulk $[-2\sqrt{T_*},2\sqrt{T_*}]=[-2R_*^{-1},2R_*^{-1}]$. So heuristically the last pair of zeros becomes real inside the bulk, at relative position $\cos\theta_*$, and not at its edge. The next theorem shows that this is so, to within $O(1/n)$ in time and in place.

\begin{theorem}[The final collision]\label{thm:collision-intro}
Let $q$ be the Gaussian power series and $\Lambda_n$ the real-rooting time of $P_{n,t}$. Almost surely $q$ has, up to conjugation, exactly one nonreal zero $u_*=R_*e^{i\theta_*}$, $0<\theta_*<\pi$, of modulus $R_*$, and there are $K<\infty$ and $n_0$, depending on $q$, such that for every $n\ge n_0$:
\begin{enumerate}
\item[(a)] $|\Lambda_n-T_*|\le K/n$;
\item[(b)] all zeros of $P_{n,\Lambda_n}$ are real; exactly one of them, $z_n$, is multiple, it is a double zero, and $|z_n-2R_*^{-1}\cos\theta_*|\le K/n$;
\item[(c)] for $T_*-K^{-1}\le t<\Lambda_n$ the polynomial $P_{n,t}$ has exactly two nonreal zeros, and for every $t\ge T_*-K^{-1}$ with $t\ne\Lambda_n$ all zeros of $P_{n,t}$ are simple.
\end{enumerate}
\end{theorem}

Theorem~\ref{thm:collision-intro} is Corollary~\ref{cor:collision-gauss}, with the $\delta$ of Theorem~\ref{thm:collision} written as $K^{-1}$ after enlarging $K$. It is the Gaussian case of a deterministic result, Theorem~\ref{thm:collision}, which holds for every real power series $q$ with $q(0)\ne0$, one simple conjugate pair of nonreal zeros of minimal modulus, no real zero of that modulus and simple real zeros inside it. For a fixed $q$ the collision is a real zero of multiplicity $2$ or $3$; for random coefficients a triple zero is excluded almost surely, because $P_{n,t}=\partial_zP_{n,t}=\partial_z^2P_{n,t}=0$ imposes three conditions on the three bottom coefficients of $P_n$, while only the two parameters $t$ and $z$ are free. For the Weyl polynomials the theorem gives tightness of $n(\Lambda_n^{\rm std}-\tilde T_{*,n})$, where $\tilde T_{*,n}$ has the law of $T_*$ and is defined under an explicit coupling (Corollary~\ref{cor:collision-weyl}). Theorem~\ref{thm:collision-intro} and Corollary~\ref{cor:collision-weyl} hold also for independent identically distributed coefficients whose law has a density and finite variance (Remark~\ref{rem:collision-laws}). The constants $K$ and $n_0$ depend on $q$; in the companion paper \cite{CompA} they are bounded by polynomials in condition parameters of $q$, which gives a power tail for $\sup_nn|\Lambda_n-T_*|$ on the event $\{T_*\ge1+\eps\}$. Theorem~\ref{thm:collision-intro} does not say whether $n(\Lambda_n-T_*)$ converges; the next theorem shows that in general it does not.

\subsection{The first-order law}\label{sec:intro-firstorder}
The next order is not a constant: $n(\Lambda_n-T_*)$ oscillates with $n$. In a window of size $1/n$ in time and in space around the collision, $P_{n,t}$ is, after a positive factor and up to $O(1/n)$, the one-phase model $\zeta\cos(\zeta+c)+\lambda\sin(\zeta+c)$, where $\lambda$ is the rescaled time and $c=c_n$ is, up to a constant, the phase of the pair $(h_n,h_{n-1})$ at $2\sqrt n\cos\theta_*$. This model has a real multiple zero only at $\lambda=-\cos^2s$, where $s-\sin s\cos s=c$, a form of Kepler's equation; and the phase advances by $-(\theta_*-\sin\theta_*\cos\theta_*)$ per degree, given by the same Kepler function.

Let $u_*=R_*e^{i\theta_*}$ ($0<\theta_*<\pi$) be the nonreal zero of $q$ of least modulus, $\gamma_*=u_*q'(u_*)$ and $\rho_*=u_*q''(u_*)/q'(u_*)$, and put
\begin{gather*}
W=\frac1{2(\im u_*)^2},\qquad\tau_c=\frac{\re(u_*)\im(\rho_*)}{2|u_*|^2\im u_*},\\
c_n^\natural=\arg\gamma_*-\frac\pi4-\Bigl(n+\frac12\Bigr)\theta_*+\frac n2\sin2\theta_*-\frac14\sin2\theta_*\re\rho_* .
\end{gather*}
For $c\in\R$ let $\Psi(c)\in[0,\pi]$ be the root of Kepler's equation
\begin{equation}\label{eq:kepler-intro}
\Psi-\sin\Psi=2\dist(c,\pi\Z),
\end{equation}
let $\tilde c\in[-\frac\pi2,\frac\pi2)$ be congruent to $c$ modulo $\pi$, and let $\nu=(\theta_*-\sin\theta_*\cos\theta_*)/\pi$.

\begin{theorem}[The first-order law]\label{thm:firstorder-intro}
Let $q$ be the Gaussian power series, $\Lambda_n$ the real-rooting time of $P_{n,t}$, and $z_n$ the multiple zero of $P_{n,\Lambda_n}$ of Theorem~\ref{thm:collision-intro}. Almost surely
\[
n(\Lambda_n-T_*)=\tau_c-W\cos\Psi(c_n^\natural)+O(n^{-2/3}),
\]
with error $O(1/n)$ along every sequence of degrees with $\dist(c_n^\natural,\pi\Z)$ bounded away from $0$, and
\[
n\Bigl(z_n-\frac{2\cos\theta_*}{R_*}\Bigr)=-\frac{\re(u_*)\re(\rho_*)}{2|u_*|^2}+\operatorname{sgn}(\tilde c_n^\natural)
\frac{\sin\Psi(c_n^\natural)}{2\im u_*}+\re(u_*)\,n(\Lambda_n-T_*)+O(n^{-1/3}).
\]
Moreover, almost surely $\nu$ is irrational and the set of limit points of $n(\Lambda_n-T_*)$ is the interval $[\tau_c-W,\tau_c+W]$, of length $1/(\im u_*)^2$.
\end{theorem}

\begin{theorem}[The joint law]\label{thm:jointlaw-intro}
Let $q$ be the Gaussian power series and $\Lambda_n$ the real-rooting time of $P_{n,t}$. Then
\[
\bigl(T_*,\,n(\Lambda_n-T_*)\bigr)\longrightarrow\bigl(T_*,\,\tau_c-W\cos\Psi(U)\bigr)\qquad\text{in distribution},
\]
where $U$ is uniform on $[0,\frac\pi2]$ and independent of $q$; the same holds with $T_*$ replaced by any random element of a Polish space that is a measurable function of $q$. Given $q$, the limit has the density $(\pi W)^{-1}\bigl((W+s)/(W-s)\bigr)^{1/2}$ in $s=\tau-\tau_c\in(-W,W)$. For the real-rooting time $\Lambda_n^{\rm std}$ of the heat-evolved Weyl polynomials $G_{n,t}$ and the threshold $\tilde T_{*,n}$ of the coupling of Corollary~\ref{cor:collision-weyl}, $\bigl(\tilde T_{*,n},n(\Lambda_n^{\rm std}-\tilde T_{*,n})\bigr)\to\bigl(T_*,\tau_c-W\cos\Psi(U)\bigr)$ in distribution.
\end{theorem}

Theorem~\ref{thm:firstorder-intro} follows from Theorem~\ref{thm:firstorder} and Proposition~\ref{prop:phase}, applied on the almost sure event of Lemma~\ref{lem:moduli}, together with Theorem~\ref{thm:collision-intro} and Theorem~\ref{thm:jointlaw}; Theorem~\ref{thm:jointlaw-intro} is Theorem~\ref{thm:jointlaw}(b) and Corollary~\ref{cor:firstorder-weyl}. Theorem~\ref{thm:firstorder} holds for every deterministic real power series under the conditions (G1)--(G3) of Section~\ref{sec:collision}, with an exact phase $c_n$ in place of $c_n^\natural$. The exponents $\frac23$ and $\frac13$ come from the cusps of the threshold curve at $c_n\in\pi\Z$, where the model has a triple zero, and they cannot be improved from the window expansion alone (Remark~\ref{rem:sharp}). For $q=1+bw+aw^2$ with $b^2<4a$ the upper end $\tau_c+W$ equals $a/2$, which shows that the elementary bound $\Lambda_n\le a\sqrt{n/(n-1)}$ is sharp to first order along a subsequence when $\nu$ is irrational (Remark~\ref{rem:benchmark}).

\subsection{The law of \texorpdfstring{$T_*$}{T*} near \texorpdfstring{$t=1$}{t=1}}\label{sec:intro-quant}
For standard Gaussian coefficients the law of $T_*$ has an exact form, and $-\log\PP(T_*\le t)$ is determined to first order as $t\downarrow1$.

\begin{theorem}[The law of $T_*$ near $t=1$]\label{thm:sharp-intro}
Let the coefficients be standard Gaussian. For $0<r<1$ and $m,n\ge0$ let $\mathcal G_{mn}(r)=\frac2\pi\int\im(z^m\bar z^n)\,|1-z^2|^{-1}dA(z)$, the integral over $\{|z|<r,\ \im z>0\}$. Then the matrix $X(r)=\bigl(m\mathcal G_{m',m-1}(r)-(m+1)\mathcal G_{m',m+1}(r)\bigr)_{m,m'\in\{0,2,4,\dots\}}$, in which the first term is $0$ for $m=0$, is a trace class operator on $\ell^2(\{0,2,4,\dots\})$, and $\PP(T_*\le t)=\det\bigl(I-X(t^{-1/2})\bigr)$ for every $t>1$. Moreover
\[
\PP(T_*\le t)\ge\prod_{k\ge0}\bigl(1-t^{-2k-1/2}\bigr)\quad(t>1),\qquad
\PP(T_*\le t)\le\exp\Bigl(-\frac{\pi^2/12-o(1)}{t-1}\Bigr)\quad(t\downarrow1),
\]
so that $(t-1)\log\PP(T_*\le t)\to-\pi^2/12$ as $t\downarrow1$.
\end{theorem}

The first statement is Proposition~\ref{prop:fredholm}, which rests on the Pfaffian structure found by Matsumoto and Shirai \cite{MS} and an identity of Rains \cite{Ra}; the upper bound is Theorem~\ref{thm:sharp}, and the lower bound, valid for every $t>1$, is Theorem~\ref{thm:lower}. The constant $\pi^2/12$ is half of the constant $\pi^2/6$ for the Gaussian power series with complex coefficients \cite{PV}. The lower bound comes from a comparison, for $r=t^{-1/2}$, of the symmetric part of $X(r)$ with the diagonal operator $R=\operatorname{diag}(r^{2m+1})_{m\in\{0,2,4,\dots\}}$; here $\det(I-R)=\prod_{k\ge0}(1-r^{4k+1})$, which has half the exponent of the product $\prod_{k\ge1}(1-r^{2k})$ of the complex case.

Two quantitative questions are left to companion papers. A rate of convergence in Theorem~\ref{cor:weyl}, of order $n^{-\alpha}$ uniformly on $[1+\eps,\infty)$ for a small explicit $\alpha$ that does not depend on $\eps$, is obtained in \cite{CompA} from a version of Theorem~\ref{thm:main} with an explicit onset. Times $t=t_n$ that tend to $1$ are the subject of \cite{CompB}, which shows that if $t_n\le2$ and $(t_n-1)n^{1/3}/(\log n)^{5/4}\to\infty$, then with probability tending to $1$ the polynomial $P_{n,t_n}$ has at least as many nonreal zeros as $q$ has in $\{|w|<t_n^{-1/2}\}$, with any excess confined to two small disks about the edges $\pm2\sqrt{t_n}$.

\subsection{The standard coupling}\label{sec:intro-standard}
In the standard coupling of Theorem~\ref{cor:weyl} the top coefficients $\xi_n,\xi_{n-1},\dots$, which play the role of $\eta_0,\eta_1,\dots$ in \eqref{eq:heat-hermite} and so make up the symbol, are renewed as $n$ grows, whereas in the reversed coupling of Theorem~\ref{thm:heat} they are shared by all degrees. For each $n$, $\Lambda_n^{\rm std}$ has the law of $\Lambda_n$, so $\Lambda_n^{\rm std}\to T_*$ in distribution; but it has no almost sure limit.

\begin{theorem}[No almost sure limit under the standard coupling]\label{thm:std}
Let $\xi_0,\xi_1,\dots$ be independent standard Gaussian variables.
\begin{enumerate}
\item[(a)] Almost surely the set of limit points of $(\Lambda_n^{\rm std})_{n\ge1}$ in $[0,\infty]$ is $[1,\infty]$. In particular, almost surely $\liminf_n\Lambda_n^{\rm std}=1$ and $\limsup_n\Lambda_n^{\rm std}=\infty$.
\item[(b)] Let $\mathcal N\subset\N$ be an infinite set that does not depend on the $\xi_k$. Almost surely, for all $1<t_1<t_2<\infty$ we have $t_1<\Lambda_n^{\rm std}\le t_2$ for infinitely many $n\in\mathcal N$, and the set of limit points of $(\Lambda_n^{\rm std})_{n\in\mathcal N}$ in $[0,\infty]$ is $[1,\infty]$.
\item[(c)] $\Lambda_n^{\rm std}$ does not converge in probability to any random variable with values in $[0,\infty]$.
\end{enumerate}
\end{theorem}

\begin{proposition}[Real-rootedness before time $1$ is exponentially rare]\label{prop:early}
Let $\xi_0,\xi_1,\dots$ be independent standard Gaussian variables, let $0<s<1$ and put
\[
\beta_s^2=\tfrac14\sqrt s\,(1-\sqrt s)\log(1/s),\qquad\gamma_s=\frac{\beta_s^2}{128(1-s^2)}.
\]
There is an explicit $n_s$, defined in \eqref{eq:std-onset} (for instance $n_s=25\,051$ for $s=\frac12$), such that for all $n\ge n_s$
\begin{equation}\label{eq:std-early}
\PP\bigl(G_{n,s}\text{ has only real zeros}\bigr)\le e^{-0.8(n+1)}+13\,e^{-\gamma_sn/4}.
\end{equation}
Consequently $\PP(G_{n,s}\text{ has only real zeros})\le C_se^{-c_sn}$ for all $n\ge1$, with $c_s=\gamma_s/4$ and $C_s=14\,e^{\gamma_sn_s/4}$. The event depends only on $(\xi_0,\dots,\xi_n)$, so these bounds hold in every coupling of the degrees; in particular $\PP(\Lambda_n^{\rm std}\le s)=\PP(\Lambda_n\le s)\le C_se^{-c_sn}$.
\end{proposition}

In (b) the set $\mathcal N$ must be fixed in advance (Remark~\ref{rem:std-quantifier}). Theorem~\ref{thm:std} shows that the convergence in distribution in Theorem~\ref{cor:weyl} cannot be improved to convergence in probability. Its proof plants a prescribed top block of coefficients at a sparse sequence of degrees and tests real-rootedness at the zeros of $h_{n-1}$ (Lemma~\ref{lem:interlace}); Proposition~\ref{prop:early}, which gives $\liminf_n\Lambda_n^{\rm std}\ge1$, combines Mehler's formula with the hole-probability argument of Sodin and Tsirelson \cite{ST}. Both results are proved in Section~\ref{sec:standard}.

\subsection{Ideas of the proofs}\label{sec:idea}
\emph{The exact count (Part~I).} Put $x=\sqrt n\,z$. Away from $[-2,2]$, classical ratio asymptotics give $Q_n(\sqrt n z)/h_n(\sqrt n z)\to a(r(z))$, where $r$ is the inverse of the Joukowski map; Rouch\'e's theorem then yields one zero near each $\zeta_i$ (Section~\ref{sec:exterior}). This exterior information is far from an exact count. A pair of nonreal zeros at distance $o(1)$ from $[-2,2]$ is invisible to it, and so is any statement about zero distributions. To count exactly we deform $a$ continuously, keeping its zeros in $\D$ fixed, until only a polynomial $p_{\mathrm{in}}$ of degree $k$ with these zeros remains. For the endpoint, Shohat's theorem on quasi-orthogonal polynomials (Lemma~\ref{lem:quasi}) gives $n-k$ real zeros inside $(-2,2)$, and the exterior argument gives the remaining $k$. Along the deformation the number of nonreal zeros can change only if two real zeros collide. Everything therefore rests on excluding \emph{real double zeros in the bulk}, uniformly along the deformation (Section~\ref{sec:bulk}).

For this we write, for $|z|\le L$ with some $L>2$,
\[
\begin{pmatrix}Q_n(x)\\ Q_n'(x)/\sqrt n\end{pmatrix}=H_{n,M}(z)\begin{pmatrix}h_n(x)\\ h_{n-1}(x)\end{pmatrix}+\begin{pmatrix}E(x)\\ E'(x)/\sqrt n\end{pmatrix},
\]
where $H_{n,M}$ is obtained by expressing the first $M+1$ terms $h_n,\dots,h_{n-M}$ through $h_n$ and $h_{n-1}$ by the backward recurrence, and $E$ is the sum of the remaining terms. The key points are these:
\begin{itemize}
\item Only $M=M_n\asymp\log n$ backward steps are used. Over this range the backward recurrence is a small perturbation of the free recurrence $f_j=zf_{j-1}-f_{j-2}$, with growth at most $e^{\kappa j}$ on $[-2\cosh\kappa,2\cosh\kappa]$ (Lemma~\ref{lem:transfer}). Hence $H_{n,M_n}$ converges to a matrix $H_a(z)$ whose determinant is $a(w)a(w^{-1})$ for $z=w+w^{-1}$, which equals $|a(e^{i\theta})|^2>0$ on $[-2,2]$.
\item The tail $E$ is not expanded at all. By the Christoffel--Darboux formula, every $h_k(x)$ with $k<n$ is at most $O(\sqrt n)$ times $\|(h_n(x),h_{n-1}(x))\|$ in the bulk. If $|a_j|\le CR_0^{-j}$ and $M_n=\lceil\ell\log n\rceil$, then $\sum_{j>M_n}|a_j|=O(n^{-\ell\log R_0})$, so the tail is negligible as soon as $\ell\log R_0>\frac12$.
\end{itemize}
Consequently $(Q_n,Q_n'/\sqrt n)$ is bounded below by a constant times $\|(h_n,h_{n-1})\|$, which never vanishes. The truncation is essential. Expressing all $n$ terms through $(h_n,h_{n-1})$ produces coefficients of size at least $e^{n/2}n^{-1/4}$ at the edge of the bulk, and for radius of convergence below $\sqrt e$ the resulting matrix does not even stay bounded (Remark~\ref{rem:full}).

\emph{The limit law (Part~II).} Theorem~\ref{thm:main} at a fixed time $t$, applied to the symbol $q(\cdot/\sqrt t)$, gives the zeros of $P_{n,t}$ for a single power series $q$ (Theorem~\ref{thm:fixedq}); since the heat flow does not increase the number of nonreal zeros, this already pins down $\Lambda_n$ to within any neighbourhood of $T_*$. Probability enters only through three facts about the random series: its zeros in $\D\setminus\{0\}$ are simple, it has no zero on a fixed circle, and it has nonreal zeros (Section~\ref{sec:heat}). The tail of $T_*$ is the probability of a nonreal zero in a small disk about the origin, a first-moment computation. For Gaussian coefficients the gap probability $\PP(T_*\le t)$ is a Fredholm determinant, by the Pfaffian structure of the zeros; the hole exponent $\pi^2/12$ is obtained from above by a comparison with the complex case and from below by an operator inequality for the symmetric part of the kernel. For the standard coupling, a prescribed block of top coefficients is planted at a sparse sequence of degrees, and real-rootedness is tested by signs at the zeros of $h_{n-1}$ (Section~\ref{sec:standard}).

\emph{The final collision (Part~III).} The backward heat flow never increases the number of nonreal zeros (Laguerre, P\'olya), and a real zero of multiplicity $m\ge2$ is preceded, just before, by $2\lfloor m/2\rfloor$ nonreal zeros nearby, as for the heat polynomials. Theorem~\ref{thm:main} at the two fixed times $T_*\pm\delta$ gives the counts $2$ and $0$, so there is exactly one collision in between, and no other multiple zero at any time $t\ge T_*-\delta$. To locate it we rerun the bulk estimate of Section~\ref{sec:bulk} at rate $1/n$, uniformly in $t$ near $T_*$. The limiting head matrix now degenerates: $\det H_{a^{(t)}}(y)=|q(t^{-1/2}e^{i\theta})|^2$ for $y=2\cos\theta$ vanishes only at $(t,y)=(T_*,2\cos\theta_*)$, and its smallest singular value grows linearly away from that point. Since the head converges at rate $1/n$ (Lemma~\ref{lem:head-rate}) and the tail is $O(1/n)$, every real multiple zero lies within $O(1/n)$ of the collision point.

\emph{The first-order law.} At the collision point the whole limiting head matrix vanishes, so the window of size $1/n$ sees the $O(1/n)$ corrections. A second-order expansion of the backward recurrence gives an envelope with an explicit first-order term, and the pair $(h_n,h_{n-1})$ near $2\sqrt n\cos\theta_*$ is written through a complex carrier controlled by the Hermite differential equation. Together they show that $Q^{(t)}_n$ is, up to a positive factor, the model $f_{\lambda,c}(\zeta)=\zeta\cos(\zeta+c)+\lambda\sin(\zeta+c)$ plus $O(1/n)$ in $C^2$. An exact identity, $|(f,\partial_\zeta f)|^2=(J(s)-c)^2+(\lambda+\cos^2s)^2$ with $s=\zeta+c$ and $J(s)=s-\sin s\cos s$, shows that the model has exactly one real multiple zero. A degree argument on a fixed box, with a homotopy in $c$ that also covers the cusps, produces a multiple zero of the perturbed model, and the uniqueness in Theorem~\ref{thm:collision} identifies it with the final collision. The Plancherel--Rotach asymptotics convert the phase to closed form, and equidistribution of $n\nu$ modulo $1$ gives the limit points and the joint law.

\subsection{Related work}
Zeros of linear combinations of orthogonal polynomials have a long history. Shohat \cite{Sh} proved that a combination of $K+1$ consecutive orthogonal polynomials of degree $n$ has at least $n-K$ real zeros in the convex hull of the support of the measure (Lemma~\ref{lem:quasi}); Beardon and Driver \cite{BD} located these zeros between the zeros of the $n$th polynomial. Tur\'an \cite{Tu} showed that all zeros of $\sum_{i\le n}b_iH_i$, in the physicists' normalisation, lie in the strip $|\im z|\le\frac12\bigl(1+\max_{i<n}|b_i|/|b_n|\bigr)$, so that the Hermite coefficients control the distance of the zeros from the real axis, a point of view he introduced for entire functions whose zeros lie in a horizontal strip. Iserles and Saff \cite{IS} proved that $\sum_kq_kH_k$ is real-rooted whenever $\sum_kq_kx^k$ is, together with results of the same type for other classical families; this is the real-rooted case of the heat-flow fact in Lemma~\ref{lem:heat}(b). Closest to Theorem~\ref{thm:main} is the work of Peherstorfer \cite{Pe}. For monic orthogonal polynomials $p_k$ of a fixed measure of Szeg\H{o} class on $[-1,1]$, he considered combinations $\sum_{j=0}^{m(n)}\mu_{j,n}p_{n-j}$ whose length $m(n)$ may grow with $n-m(n)\to\infty$. Under a separation condition on the zeros of an associated polynomial and a summability condition on the coefficients, he proved that these combinations eventually have exactly $n-l(n)$ simple zeros in $(-1,1)$, where $l(n)$ is the number of zeros of the associated polynomial $\sum_j2^j\mu_{j,n}z^{m(n)-j}$ outside the unit circle. Theorem~\ref{thm:main} is a result of the same type, but it does not follow from that theorem as stated. The scaled polynomials $z\mapsto h_k(\sqrt n z)$, $0\le k\le n$, are orthogonal with respect to the $n$-dependent weight $e^{-nz^2/2}$ on the whole real line, not with respect to a fixed measure on a compact interval. Moreover, the combinations $Q_n$ have full length $n$, so after any truncation a tail remains that must be controlled separately. We do not know whether the method of \cite{Pe} extends to this setting.
% TODO(author): the full text of [Pe] could not be consulted while this draft was written; compare its proof
% with Sections 3--5 before submission.
Dur\'an \cite{Du1} studies fixed finite combinations $\sum_{j=0}^{K}\gamma_j\widetilde H_{n-j}$ of Hermite polynomials through the polynomial $P(x)=\sum_j\gamma_jx^{K-j}$: in the standard normalisation the combination is real-rooted for all large $n$ even when $P$ has nonreal zeros, whereas in the Appell normalisation $H_n/(2^nn!)$ it has, for all large $n$, exactly as many nonreal zeros as $P$ \cite[Theorems~1.2 and~1.3]{Du1}. For Laguerre polynomials he proves exact counts of the same kind in three of four normalisations, in two of them under the condition that $P$ has no nonreal zeros in the closed unit disk, while in the monic normalisation the combination is eventually real-rooted \cite[Theorems~1.3--1.6]{Du3}; for polynomial symbols these are the Laguerre counterparts of Theorem~\ref{thm:main}, and the Laguerre case of our method in \cite{CompC} is to be compared with them. In \cite{Du2} he shows that every positive measure has a normalisation of its orthogonal polynomials for which all zeros of fixed finite combinations are eventually real and simple. The normalisation matters: relative to monic Hermite polynomials, a fixed combination $\sum_ja_jh_{n-j}$ has the $n$-dependent coefficients $a_j\sqrt{n!/(n-j)!}$. The exterior part of Theorem~\ref{thm:main} is a consequence of ratio asymptotics, which are available also for varying recurrence coefficients; see \cite{KVA} for the corresponding zero distributions. From a spectral point of view, $Q_n(\sqrt n\,\cdot)$ is the characteristic polynomial of a rank-one, non-symmetric perturbation of the scaled Jacobi matrix (Lemma~\ref{lem:colleague}). Results on spectral clustering for such perturbations \cite{GSC} allow eigenvalues near the interval and do not address their reality. The parametrisation $z=w+w^{-1}$ and the role of an analytic function of $w$ resemble the Jost function theory of Jacobi matrices \cite{DS}, where the perturbations are self-adjoint. The exact count extends to the orthonormal polynomials of varying weights $e^{-nV}$ with $V$ a one-cut regular polynomial, by the Riemann--Hilbert asymptotics of \cite{DKMVZ}, and to Laguerre polynomials; this is carried out in \cite{CompC}.

The intensity of the complex zeros of the Gaussian power series was computed by Shepp and Vanderbei \cite{SV}; the zeros form a Pfaffian point process, found by Forrester \cite{Fo} via truncated orthogonal matrices and by Matsumoto and Shirai \cite{MS}; the real zeros go back to Kac \cite{Ka}. Theorem~\ref{cor:weyl} expresses $\PP(T_*\le t)$ as a gap probability of this process, and Theorem~\ref{thm:sharp-intro} is the real-coefficient counterpart of the hole probability of Peres and Vir\'ag \cite{PV}. The heat-flow model and its normalisation are those of \cite{HHJK}; for the heat flow applied to Gaussian analytic functions themselves see \cite{HHJK2}. That the backward heat flow does not increase the number of nonreal zeros of a real polynomial goes back to Laguerre and P\'olya \cite{Ob}; for real-rooted polynomials it is also in \cite{IS}.

The real-rooting time is a finite-degree analogue of the de Bruijn--Newman constant $\Lambda$ \cite{dB,Ne}, the real-rooting time of the backward heat flow started from a Fourier transform attached to the Riemann zeta function, for which Rodgers and Tao proved $\Lambda\ge0$ \cite{RT} and Polymath~15 gave an upper bound \cite{PM15}. Collisions of zeros under this flow are central there. Csordas, Smith and Varga \cite{CSV} bound $\Lambda$ from below through close pairs of zeros (Lehmer pairs), and the local form of a collision at a multiple zero (our Lemma~\ref{lem:localform}) is \cite[Proposition~3.1(ii)]{PM15} for the de Bruijn--Newman family; the polynomials that describe it are the heat polynomials of Rosenbloom and Widder \cite{RW}. The carrier $\cos(\zeta+c)$ of the window model of Section~\ref{sec:firstorder} is the arithmetic-progression equilibrium state used by Rodgers and Tao, and the model $f_{\lambda,c}$ adds one zero to it; it is an exact trajectory of the backward heat flow, and its threshold $\lambda_c$ is the real-rooting time of that flow. For $c\in\pi\Z$ and $c\in\frac\pi2+\pi\Z$ the threshold is classical: up to a factor, $f_{\lambda,0}$ and $f_{\lambda,\pi/2}$ are Dini functions of orders $\pm\frac12$, whose zeros are all real exactly when $\lambda\ge-1$, respectively $\lambda\ge0$ \cite[\S15.25, p.~482]{Wat}. We have not found the general threshold curve, a Kepler equation, in the literature.

\subsection{Organisation}
Part~I proves Theorem~\ref{thm:main}: Section~\ref{sec:prelim} collects facts about Hermite and Chebyshev polynomials, Section~\ref{sec:transfer} bounds the backward recurrence over logarithmically many steps, also at rate $1/n$, Section~\ref{sec:bulk} excludes real double zeros in the bulk, Section~\ref{sec:exterior} treats the exterior, and Section~\ref{sec:proof} completes the proof. Part~II treats the heat flow: Section~\ref{sec:heat} proves Theorems~\ref{thm:heat} and~\ref{cor:weyl}, Proposition~\ref{prop:law}, the tail of $T_*$ and, for Gaussian coefficients, Theorem~\ref{thm:sharp-intro}; Section~\ref{sec:standard} proves Theorem~\ref{thm:std} and Proposition~\ref{prop:early}. Part~III treats the final collision: Section~\ref{sec:collision} proves the deterministic collision theorem, Theorem~\ref{thm:collision}, and its Gaussian form, Corollary~\ref{cor:collision-gauss}, which gives Theorem~\ref{thm:collision-intro}, and ends with the Weyl polynomials (Corollary~\ref{cor:collision-weyl}); Section~\ref{sec:firstorder} proves the first-order law (Theorem~\ref{thm:firstorder}) and the joint law (Theorem~\ref{thm:jointlaw}), and hence Theorems~\ref{thm:firstorder-intro} and~\ref{thm:jointlaw-intro}. Section~\ref{sec:remarks} contains remarks and open questions.

The companion papers extend this work in three directions that are not needed here: explicit onsets for Theorem~\ref{thm:main}, a rate of convergence in Theorem~\ref{cor:weyl} and explicit constants in Theorem~\ref{thm:collision-intro} \cite{CompA}; the nonreal zeros at times $t_n\downarrow1$ \cite{CompB}; and exact counts for varying weights and Laguerre polynomials, with the real-rooting times of other coefficient profiles and of a Laguerre flow \cite{CompC}.

\clearpage
\part{Exact counts of nonreal zeros}

\section{Preliminaries}\label{sec:prelim}

\subsection{Hermite polynomials}
The polynomials $h_k$ satisfy, with $h_{-1}=0$,
\begin{equation}\label{eq:recurrence}
x\,h_k(x)=\sqrt{k+1}\,h_{k+1}(x)+\sqrt k\,h_{k-1}(x),\qquad h_k'(x)=\sqrt k\,h_{k-1}(x)\qquad(k\ge0).
\end{equation}
For $m\ge1$ let $J_m$ be the $m\times m$ symmetric tridiagonal matrix with zero diagonal and entries $(J_m)_{k-1,k}=(J_m)_{k,k-1}=\sqrt k$, $1\le k\le m-1$ (rows and columns indexed by $0,\dots,m-1$). Let $e_0,\dots,e_{m-1}$ be the standard basis and $p_m(x)=(h_0(x),\dots,h_{m-1}(x))^T$. By \eqref{eq:recurrence},
\begin{equation}\label{eq:jacobi}
(xI-J_m)\,p_m(x)=\sqrt m\,h_m(x)\,e_{m-1}.
\end{equation}

\begin{lemma}\label{lem:zeros}
For $m\ge1$ the zeros of $h_m$ are the eigenvalues of $J_m$. They are real and simple, and for $m\ge2$ they lie in $(-2\sqrt{m-1},2\sqrt{m-1})$. The polynomials $h_m$ and $h_{m-1}$ have no common zero.
\end{lemma}

\begin{proof}
By \eqref{eq:jacobi} and $p_m(x)\neq0$, every zero of $h_m$ is an eigenvalue of $J_m$; as $h_m$ has $m$ zeros counted with multiplicity and $J_m$ is symmetric with nonzero off-diagonal entries, hence has $m$ simple real eigenvalues, the zeros of $h_m$ are exactly these eigenvalues. Each row of $J_m$ has absolute row sum at most $\sqrt{m-2}+\sqrt{m-1}<2\sqrt{m-1}$, which bounds $\|J_m\|$. Finally, if $h_m(x)=h_{m-1}(x)=0$, then \eqref{eq:recurrence} gives $h_{m-2}(x)=\dots=h_0(x)=0$, contradicting $h_0=1$.
\end{proof}

\begin{lemma}[Christoffel--Darboux]\label{lem:cd}
For $n\ge1$ and $x=\sqrt n\,z$,
\[
\sum_{k=0}^{n-1}h_k(x)^2=n\bigl(h_n(x)^2-z\,h_n(x)h_{n-1}(x)+h_{n-1}(x)^2\bigr).
\]
Consequently, if $L>0$ and $z$ is real with $|z|\le L$, then $\sum_{k=0}^{n-1}h_k(x)^2\le n(1+\tfrac L2)\bigl(h_n(x)^2+h_{n-1}(x)^2\bigr)$.
\end{lemma}

\begin{proof}
This is the confluent Christoffel--Darboux formula $\sum_{k<n}h_k^2=\sqrt n(h_n'h_{n-1}-h_nh_{n-1}')$ \cite[Section~3.2, (3.2.4)]{Sz}, since $h_n'=\sqrt n\,h_{n-1}$ and $h_{n-1}'=\sqrt{n-1}\,h_{n-2}=xh_{n-1}-\sqrt n\,h_n$ by \eqref{eq:recurrence}, and $\sqrt n\,x=nz$.
% TODO(author): the pinpoint [Sz, (3.2.4)] (confluent Christoffel--Darboux formula) was not checked against the book;
% in the fold of 21 October 2026 it replaced a direct three-line induction (F_{n+1}=F_n+h_n^2 for
% F_n=n(h_n^2+h_{n-1}^2)-sqrt(n) x h_n h_{n-1}), which can be restored if the pinpoint is in doubt.
The inequality follows from $|zh_nh_{n-1}|\le\frac{|z|}2(h_n^2+h_{n-1}^2)$.
\end{proof}

\begin{lemma}[Gaussian quadrature {\cite[Theorems~3.4.1--3.4.2]{Sz}}]\label{lem:quadrature}
Let $x_1<\dots<x_{n+1}$ be the zeros of $h_{n+1}$. There are weights $\lambda_i>0$ such that $\int R\,d\gamma=\sum_i\lambda_iR(x_i)$ for every polynomial $R$ of degree at most $2n+1$, where $\gamma$ is the standard Gaussian measure.
\end{lemma}

\begin{lemma}\label{lem:resolvent}
Let $n\ge1$ and $z\in\C\setminus[-2,2]$. Then $h_n(\sqrt n z)\neq0$, and the vector $\rho_n(z)=\bigl(h_k(\sqrt nz)/h_n(\sqrt nz)\bigr)_{k=0}^{n-1}$ equals $(zI-n^{-1/2}J_n)^{-1}e_{n-1}$. In particular
\[
\sum_{k=0}^{n-1}\Bigl|\frac{h_k(\sqrt nz)}{h_n(\sqrt nz)}\Bigr|^2\le\dist(z,[-2,2])^{-2}.
\]
\end{lemma}

\begin{proof}
By Lemma~\ref{lem:zeros} the zeros of $h_n(\sqrt n\,\cdot)$ lie in $(-2,2)$. Dividing \eqref{eq:jacobi} by $\sqrt n\,h_n(\sqrt nz)$ gives $(zI-n^{-1/2}J_n)\rho_n(z)=e_{n-1}$. Since $n^{-1/2}J_n$ is symmetric with spectrum in $[-2,2]$, the resolvent has norm at most $\dist(z,[-2,2])^{-1}$.
\end{proof}

\subsection{Chebyshev polynomials and the Joukowski map}
Let $U_l$ be the Chebyshev polynomials of the second kind, $U_l(\cos\theta)=\sin((l+1)\theta)/\sin\theta$, extended by the recurrence $U_{l+1}(y)=2yU_l(y)-U_{l-1}(y)$ to $U_{-1}=0$ and $U_{-2}=-1$. For $w\ne0,\pm1$ and $y=\frac12(w+w^{-1})$,
\begin{equation}\label{eq:chebyshev-joukowski}
U_l(y)=\frac{w^{l+1}-w^{-l-1}}{w-w^{-1}}\qquad(l\ge-2),
\end{equation}
since both sides satisfy the same recurrence and agree for $l=-1,0$. The polynomial $z\mapsto U_l(z/2)$ is the characteristic polynomial of the $l\times l$ path matrix with zero diagonal and unit off-diagonal entries.

\begin{lemma}\label{lem:chebyshev}
Let $\kappa>0$, $l\ge0$ and $y\in\R$ with $|y|\le\cosh\kappa$. Then
\[
|U_l(y)|\le U_l(\cosh\kappa)=\frac{\sinh((l+1)\kappa)}{\sinh\kappa}\le(l+1)e^{l\kappa}.
\]
\end{lemma}

\begin{proof}
On $[-1,1]$, $|U_l|\le l+1=U_l(1)$. All zeros of $U_l$ lie in $(-1,1)$ and its leading coefficient is positive, so $U_l$ is increasing on $[1,\infty)$; by parity, $|U_l(y)|\le U_l(|y|)\le U_l(\cosh\kappa)$ for $1\le|y|\le\cosh\kappa$. Finally $\sinh((l+1)\kappa)/\sinh\kappa=\sum_{i=0}^{l}e^{(l-2i)\kappa}\le(l+1)e^{l\kappa}$.
\end{proof}

Let $r(z)=\frac12\bigl(z-\sqrt{z^2-4}\bigr)$, with the branch for which $r(z)\sim1/z$ at infinity. It is the inverse of the Joukowski map $w\mapsto w+w^{-1}$ from $\D\setminus\{0\}$ onto $\C\setminus[-2,2]$, and $|r(z)|<1$ there.

\subsection{Two facts about real polynomials}

\begin{lemma}\label{lem:constancy}
Let $(P_s)_{s\in[0,1]}$ be real polynomials of degree $n$ whose coefficients depend continuously on $s$. If every real zero of every $P_s$ is simple, then the number of nonreal zeros of $P_s$, counted with multiplicity, does not depend on $s$.
\end{lemma}

\begin{proof}
It suffices to show that the number $N(s)$ of real zeros is locally constant. Fix $s_0$, and let $x_1,\dots,x_r$ be the real zeros of $P_{s_0}$, which are simple. The zeros of all $P_s$ with $s$ near $s_0$ lie in a fixed disk $\{|z|\le\varrho\}$, because the leading coefficients stay away from $0$. Choose disjoint closed disks $B_i$ centred at $x_i$ that contain no other zero of $P_{s_0}$. By Rouch\'e's theorem, for $s$ close to $s_0$ each $B_i$ contains exactly one zero of $P_s$, which is real because $B_i$ is symmetric about $\R$ and the zeros of $P_s$ are symmetric. Also $P_{s_0}$ has no zero on the compact set $[-\varrho,\varrho]\setminus\bigcup_i B_i$, hence neither has $P_s$ for $s$ close to $s_0$. Thus $N(s)=r$ near $s_0$.
\end{proof}

\begin{lemma}[Shohat's theorem]\label{lem:quasi}
Let $0\le k\le n$ and let $P=\sum_{j=0}^{k}c_jh_{n-j}$ with real $c_j$ and $c_0\ne0$. Then $P$ has at least $n-k$ distinct real zeros in the open interval between the smallest and the largest zero of $h_{n+1}$, and in particular in $(-2\sqrt n,2\sqrt n)$.
\end{lemma}

That a combination of $k+1$ consecutive orthogonal polynomials of degree $n$ has at least $n-k$ real zeros in the convex hull of the support of the measure is a classical theorem of Shohat \cite{Sh}; see also \cite{BD} and \cite[Lemma~3.1]{Du2}. We include the short proof by Gaussian quadrature, which also places the zeros between the extreme zeros of $h_{n+1}$.

\begin{proof}
$P$ has degree $n$ and is orthogonal in $L^2(\gamma)$ to all polynomials of degree less than $n-k$. With the notation of Lemma~\ref{lem:quadrature}, for every such polynomial $R$ the product $PR$ has degree at most $2n-k-1\le2n+1$, so
\begin{equation}\label{eq:discrete-orth}
\sum_{i=1}^{n+1}\lambda_iP(x_i)R(x_i)=0 .
\end{equation}
Not all $P(x_i)$ vanish, since $P\ne0$ has degree $n$. Consider the values $P(x_i)\ne0$ in increasing order of $i$, and suppose they have $\sigma<n-k$ sign changes. For each pair of consecutive nonzero values of opposite sign, at nodes $x_a<x_b$, choose a point $\tau\in(x_a,x_b)$ that is not a node, and let $R$ be $\pm$ the product of the factors $x-\tau$ over all these pairs. Then $R$ has degree $\sigma<n-k$, and the sign can be chosen so that $P(x_i)R(x_i)>0$ whenever $P(x_i)\ne0$, contradicting \eqref{eq:discrete-orth}. Hence $\sigma\ge n-k$. Each sign change gives a zero of $P$ in the corresponding open interval $(x_a,x_b)$, and these intervals are disjoint. By Lemma~\ref{lem:zeros}, $|x_i|<2\sqrt n$.
\end{proof}

\section{The backward recurrence over logarithmically many steps}\label{sec:transfer}

Fix $n\ge2$ and put $x=\sqrt n\,z$ and $f_j(z)=h_{n-j}(\sqrt nz)$ for $0\le j\le n$. By \eqref{eq:recurrence},
\begin{equation}\label{eq:backward}
f_j=\frac{z\,f_{j-1}-\sqrt{1-\frac{j-2}{n}}\,f_{j-2}}{\sqrt{1-\frac{j-1}n}}\qquad(2\le j\le n).
\end{equation}
Define polynomials $A_{n,j},B_{n,j}$ in $z$, for $0\le j\le n$, by $(A_{n,0},B_{n,0})=(1,0)$, $(A_{n,1},B_{n,1})=(0,1)$ and the recurrence \eqref{eq:backward}. Then
\begin{equation}\label{eq:AB}
h_{n-j}(\sqrt nz)=A_{n,j}(z)\,h_n(\sqrt nz)+B_{n,j}(z)\,h_{n-1}(\sqrt nz)\qquad(0\le j\le n).
\end{equation}
As $n\to\infty$ with $j$ fixed, the coefficients of \eqref{eq:backward} tend to those of the free recurrence $f_j=zf_{j-1}-f_{j-2}$, whose solutions with the same initial values are $-U_{j-2}(z/2)$ and $U_{j-1}(z/2)$. By induction on $j$, the coefficients of the polynomials $A_{n,j},B_{n,j}$ converge, so
\begin{equation}\label{eq:fixed-j}
A_{n,j}(z)\to A_j(z):=-U_{j-2}(z/2),\qquad B_{n,j}(z)\to B_j(z):=U_{j-1}(z/2)
\end{equation}
uniformly on compact subsets of $\C$, for every fixed $j\ge0$.

For $0\le r<n$ let $\beta_r=1-r/n$, and for $j\ge1$ let $d_{n,j}=\prod_{r=1}^{j-1}\sqrt{\beta_r}$. For $m\ge1$ and positive numbers $b_1,\dots,b_{m-1}$ let $\Pi_m(z;b_1,\dots,b_{m-1})$ be the characteristic polynomial of the $m\times m$ symmetric tridiagonal matrix with zero diagonal whose off-diagonal entries have squares $b_1,\dots,b_{m-1}$; put $\Pi_0=1$.

\begin{lemma}[Continuant formulas]\label{lem:continuant}
For $2\le j\le n$,
\[
B_{n,j}=\frac{\Pi_{j-1}(z;\beta_1,\dots,\beta_{j-2})}{d_{n,j}},\qquad A_{n,j}=-\frac{\Pi_{j-2}(z;\beta_2,\dots,\beta_{j-2})}{d_{n,j}}.
\]
\end{lemma}

\begin{proof}
Put $d_{n,0}=d_{n,1}=1$. If $X_j$ satisfies \eqref{eq:backward}, then $N_j=d_{n,j}X_j$ satisfies $N_j=zN_{j-1}-\beta_{j-2}N_{j-2}$ for $j\ge2$ (with $\beta_0=1$). Expanding a tridiagonal determinant along its last row gives $\Pi_m=z\Pi_{m-1}-b_{m-1}\Pi_{m-2}$. For $B$ the initial values are $N_0=0$, $N_1=1=\Pi_0$, $N_2=z=\Pi_1$; for $A$ they are $N_0=1$, $N_1=0$, $N_2=-1=-\Pi_0$, $N_3=-z=-\Pi_1(z)$. Both formulas follow by induction.
\end{proof}

\begin{lemma}[Defect expansion]\label{lem:defect}
Write $b_e=1-\delta_e$ for the edges $e=1,\dots,m-1$ of the path on $m$ vertices. Then
\[
\Pi_m(z;1-\delta_1,\dots,1-\delta_{m-1})=\sum_{S}\ \prod_{e\in S}\delta_e\prod_{C}U_{|C|}(z/2),
\]
where $S$ runs over the matchings of the path (sets of pairwise disjoint edges, including $S=\emptyset$), and $C$ runs over the connected components of the path with the endpoints of the edges of $S$ removed, $|C|$ denoting the number of vertices of $C$.
\end{lemma}

\begin{proof}
Expanding the determinant over permutations, only products of disjoint transpositions along edges contribute, so $\Pi_m=\sum_{S'}(-1)^{|S'|}\prod_{e\in S'}b_e\,z^{m-2|S'|}$, the sum over matchings $S'$. Expand $\prod_{e\in S'}(1-\delta_e)=\sum_{S\subseteq S'}(-1)^{|S|}\prod_{e\in S}\delta_e$ and exchange the sums. For a fixed matching $S$ the matchings $S'\supseteq S$ are $S\cup S''$ with $S''$ a matching of the path with the vertices of $S$ removed. The two signs $(-1)^{|S|}$ cancel, and $\sum_{S''}(-1)^{|S''|}z^{m-2|S|-2|S''|}$ is the product of the unit-weight path polynomials $U_{|C|}(z/2)$ of the components.
\end{proof}

\begin{lemma}[Transfer bound]\label{lem:transfer}
Let $\kappa>0$ and $c_\kappa=\coth\kappa+\frac12$. For $n\ge2$, $0\le j\le n/2$ and real $z$ with $|z|\le2\cosh\kappa$,
\[
|A_{n,j}(z)|+|B_{n,j}(z)|\le2(j+1)\,e^{j\kappa}\,e^{c_\kappa j^2/n}.
\]
The limits in \eqref{eq:fixed-j} satisfy $|A_j(z)|+|B_j(z)|\le2(j+1)e^{j\kappa}$ for the same $z$.
\end{lemma}

\begin{proof}
For $j\le1$ the left side is $1$. Let $2\le j\le n/2$. Each polynomial $\Pi_m$ in Lemma~\ref{lem:continuant} has $m\le j-1$ and edge squares $1-\delta_e$ with $0\le\delta_e\le j/n$. In Lemma~\ref{lem:defect} every $\delta_e\ge0$, and by Lemma~\ref{lem:chebyshev}, $|U_l(z/2)|\le U_l(\cosh\kappa)$ with $U_l(\cosh\kappa)>0$. Hence $|\Pi_m(z;1-\delta)|$ is at most the same expansion with each $U_l(z/2)$ replaced by $U_l(\cosh\kappa)$. That expansion is nondecreasing in every $\delta_e$, so it is at most its value with every $\delta_e$ replaced by $j/n$. With $\beta=1-j/n\in[\frac12,1)$, the constant-edge path is $\sqrt\beta$ times the unit path, so
\[
|\Pi_m(z;1-\delta)|\le\Pi_m(2\cosh\kappa;\beta,\dots,\beta)=\beta^{m/2}U_m\Bigl(\frac{\cosh\kappa}{\sqrt\beta}\Bigr)\le(m+1)e^{m\kappa'},
\]
where $\cosh\kappa'=\cosh\kappa/\sqrt\beta$. Since $\beta^{-1/2}-1\le1-\beta$ for $\beta\in[\frac12,1]$ and $\arcosh'\le1/\sinh\kappa$ on $[\cosh\kappa,\infty)$, we get $\kappa'\le\kappa+\coth\kappa\cdot j/n$. For the denominators, $\log(1-u)\ge-2u$ for $u\in[0,\frac12]$ gives $d_{n,j}\ge\exp(-\sum_{r<j}r/n)\ge e^{-j^2/(2n)}$. Using $m=j-1$ for $B_{n,j}$ and $m=j-2$ for $A_{n,j}$,
\[
|A_{n,j}|+|B_{n,j}|\le(2j-1)\,e^{(j-1)\kappa}\exp\Bigl(\coth\kappa\cdot\frac{j^2}{n}+\frac{j^2}{2n}\Bigr),
\]
which implies the claim. The bound for the limits follows from Lemma~\ref{lem:chebyshev}.
\end{proof}

\begin{lemma}[The backward recurrence at rate $1/n$]\label{lem:head-rate}
Let $\kappa>0$, put $\kappa_1=\min(\kappa,1)$ and $c_1=15/\kappa_1$, and let $j\ge0$, $n\ge1$ and $z\in\R$ with $|z|\le2\cosh\kappa$.
\begin{enumerate}
\item[(a)] If $(j+1)^2\le\kappa_1n$, then
\[
|A_{n,j}(z)-A_j(z)|+|B_{n,j}(z)-B_j(z)|\le c_1\,\frac{j^3e^{(j-1)\kappa}}{n}.
\]
\item[(b)] If $(j+2)^2\le\kappa_1n$, then, with $\omega_{n,j}=\sqrt{1-j/n}$,
\[
|\omega_{n,j}A_{n,j+1}(z)-A_{j+1}(z)|+|\omega_{n,j}B_{n,j+1}(z)-B_{j+1}(z)|\le(c_1+3)\,\frac{(j+1)^3e^{j\kappa}}{n}.
\]
\end{enumerate}
\end{lemma}

\begin{proof}
\emph{A first-order defect bound.} Let $m\ge2$ and $0\le\delta_e\le\bar\delta$ for $e=1,\dots,m-1$, where $m\bar\delta\le\kappa_1$. We claim that
\begin{equation}\label{eq:defect-rate}
\bigl|\Pi_m(z;1-\delta_1,\dots,1-\delta_{m-1})-U_m(z/2)\bigr|\le\frac{e^{4/3}}{\kappa_1}\,m(m+1)\,\bar\delta\,e^{m\kappa}.
\end{equation}
By Lemma~\ref{lem:defect} the difference is the sum over the nonempty matchings $S$ in that lemma, and by Lemma~\ref{lem:chebyshev}, $|U_l(z/2)|\le U_l(\cosh\kappa)$. So the difference is at most $\Phi(\delta)-\Phi(0)$, where
\[
\Phi(\delta)=\sum_S\prod_{e\in S}\delta_e\prod_CU_{|C|}(\cosh\kappa)=\Pi_m(2\cosh\kappa;1-\delta_1,\dots,1-\delta_{m-1})
\]
is a polynomial in $(\delta_e)$ with nonnegative coefficients. Hence $\Phi(\delta)-\Phi(0)\le g(\bar\delta)-g(0)$, where $g(s)=\Phi(s,\dots,s)$ is a polynomial in $s$ with nonnegative coefficients; in particular $(g(s)-g(0))/s$ is nondecreasing on $(0,\infty)$. With $s_1=\kappa_1/m\in[\bar\delta,\frac12]$ this gives $g(\bar\delta)-g(0)\le(\bar\delta/s_1)\,g(s_1)=(m\bar\delta/\kappa_1)\,g(s_1)$. As in the proof of Lemma~\ref{lem:transfer}, $g(s)=(1-s)^{m/2}U_m(\cosh\kappa/\sqrt{1-s})$, and $\cosh\kappa'=\cosh\kappa/\sqrt{1-s_1}$ gives $\kappa'\le\kappa+s_1\coth\kappa$. Moreover $\kappa_1\coth\kappa\le\frac43$: for $\kappa\le1$ because $\kappa\coth\kappa\le1+\kappa^2/3$ (compare the Taylor series of $\kappa\cosh\kappa$ and $(1+\kappa^2/3)\sinh\kappa$ termwise), and for $\kappa>1$ because $\coth\kappa<\coth1<\frac43$. So $m\kappa'\le m\kappa+\frac43$, and $g(s_1)\le(m+1)e^{m\kappa'}\le e^{4/3}(m+1)e^{m\kappa}$ by Lemma~\ref{lem:chebyshev}, which proves \eqref{eq:defect-rate}.

(a) For $j\le1$ the left side is $0$. Let $j\ge2$. Then $j^2\le n$, $n\ge9$ and $j\le\sqrt n\le n/2$, so Lemma~\ref{lem:continuant} applies.

\emph{Denominators.} As in the proof of Lemma~\ref{lem:transfer}, $1\ge d_{n,j}\ge e^{-j^2/(2n)}\ge e^{-1/2}$, so $0\le d_{n,j}^{-1}-1\le e^{j^2/(2n)}-1\le\frac{j^2}{2n}e^{1/2}\le\frac{j^2}n$.

\emph{The entry $B$.} By Lemma~\ref{lem:continuant}, $B_{n,j}=\Pi_{j-1}(z;\beta_1,\dots,\beta_{j-2})/d_{n,j}$, with edge defects $\delta_e=e/n\le\bar\delta:=(j-2)/n$. For $j=2$, $\Pi_1(z)=z=U_1(z/2)$. For $j\ge3$ we have $(j-1)\bar\delta\le j^2/n\le\kappa_1$, and \eqref{eq:defect-rate} with $m=j-1$ gives $|\Pi_{j-1}-U_{j-1}(z/2)|\le(e^{4/3}/\kappa_1)j^3e^{(j-1)\kappa}/n$. Since $|U_{j-1}(z/2)|\le je^{(j-1)\kappa}$ (Lemma~\ref{lem:chebyshev}),
\[
|B_{n,j}-B_j|\le\frac{|\Pi_{j-1}-U_{j-1}(z/2)|}{d_{n,j}}+|U_{j-1}(z/2)|\,\bigl(d_{n,j}^{-1}-1\bigr)\le\Bigl(\frac{e^{11/6}}{\kappa_1}+1\Bigr)\frac{j^3e^{(j-1)\kappa}}n .
\]

\emph{The entry $A$.} Here $A_{n,j}=-\Pi_{j-2}(z;\beta_2,\dots,\beta_{j-2})/d_{n,j}$, with $m=j-2$ and edge defects $(e+1)/n\le(j-2)/n$, $1\le e\le j-3$. For $j\le3$, $\Pi_{j-2}=U_{j-2}(z/2)$; for $j\ge4$, $(j-2)\bar\delta=(j-2)^2/n\le\kappa_1$, so \eqref{eq:defect-rate} applies. With $|U_{j-2}(z/2)|\le(j-1)e^{(j-2)\kappa}$ the same computation gives $|A_{n,j}-A_j|\le(e^{11/6}/\kappa_1+1)j^3e^{(j-2)\kappa}/n$. Adding the two bounds and using $2(e^{11/6}+\kappa_1)\le15$ gives (a).

(b) Since $0\le1-\omega_{n,j}\le j/n$, for $X\in\{A,B\}$ we have $|\omega_{n,j}X_{n,j+1}-X_{j+1}|\le|X_{n,j+1}-X_{j+1}|+\frac jn|X_{j+1}|$, and $|A_{j+1}|+|B_{j+1}|=|U_{j-1}(z/2)|+|U_j(z/2)|\le(2j+1)e^{j\kappa}$ by Lemma~\ref{lem:chebyshev}. Part (a) at $j+1$ and $j(2j+1)\le3(j+1)^3$ give (b).
\end{proof}

The rate $1/n$ cannot be improved: by \eqref{eq:backward}, $A_{n,2}-A_2=1-(1-1/n)^{-1/2}\sim-1/(2n)$.

\section{The bulk: no real double zeros}\label{sec:bulk}

We work with families of power series, as needed for the deformation in Section~\ref{sec:proof}.

\begin{hypothesis}\label{hyp:H}
$(a_s)_{s\in[0,1]}$ is a family of real power series $a_s(w)=\sum_{j\ge0}a_j(s)w^j$ such that:
\begin{enumerate}
\item[(H1)] there are $R_0>1$ and $C_0<\infty$ with $|a_j(s)|\le C_0R_0^{-j}$ for all $j\ge0$ and $s\in[0,1]$;
\item[(H2)] each coefficient $a_j(s)$ depends continuously on $s$;
\item[(H3)] $a_0(s)$ does not depend on $s$, and $a_0(s)\neq0$;
\item[(H4)] $a_s(w)\neq0$ for all $s\in[0,1]$ and $|w|=1$.
\end{enumerate}
\end{hypothesis}

By (H1) and (H2) the map $(s,w)\mapsto a_s(w)$ is continuous on $[0,1]\times\{|w|<R_0\}$. For such a family put
\[
Q_n^{s}(x)=\sum_{j=0}^{n}a_j(s)\,h_{n-j}(x),\qquad v_n(x)=\begin{pmatrix}h_n(x)\\h_{n-1}(x)\end{pmatrix}.
\]
By Lemma~\ref{lem:zeros}, $v_n(x)\neq0$ for every real $x$.

\subsection{The head matrix and its limit}
For $0\le M\le n-1$, $\omega_{n,j}=\sqrt{1-j/n}$ and a real sequence $a=(a_j)$ define
\[
G_{n,j}(z)=\begin{pmatrix}A_{n,j}(z)&B_{n,j}(z)\\ \omega_{n,j}A_{n,j+1}(z)&\omega_{n,j}B_{n,j+1}(z)\end{pmatrix},\qquad H_{n,M}[a](z)=\sum_{j=0}^{M}a_j\,G_{n,j}(z).
\]
Let $Q_{n,M}=\sum_{j=0}^{M}a_jh_{n-j}$. By \eqref{eq:recurrence}, $Q_{n,M}'(x)/\sqrt n=\sum_{j=0}^{M}a_j\omega_{n,j}h_{n-j-1}(x)$, so \eqref{eq:AB} gives, for $x=\sqrt nz$,
\begin{equation}\label{eq:head}
\begin{pmatrix}Q_{n,M}(x)\\Q_{n,M}'(x)/\sqrt n\end{pmatrix}=H_{n,M}[a](z)\,v_n(x).
\end{equation}
Note that $M+1\le n$, so no index below $0$ occurs. Correspondingly let $G_j(z)=\begin{pmatrix}A_j&B_j\\A_{j+1}&B_{j+1}\end{pmatrix}$ with $A_j,B_j$ from \eqref{eq:fixed-j}.

\begin{lemma}[The limiting matrix]\label{lem:limit}
Let $a=(a_j)$ be real with $|a_j|\le C_0R_0^{-j}$, $R_0>1$, and let $\kappa_0>0$ with $e^{\kappa_0}<R_0$. For real $z$ with $|z|\le2\cosh\kappa_0$ the series
\[
\mathcal A(z)=\sum_{j\ge0}a_jA_j(z)=a_0-\sum_{j\ge2}a_jU_{j-2}(z/2),\qquad\mathcal B(z)=\sum_{j\ge0}a_jB_j(z)=\sum_{j\ge1}a_jU_{j-1}(z/2)
\]
converge absolutely and uniformly, and
\[
H_a(z):=\sum_{j\ge0}a_jG_j(z)=\begin{pmatrix}\mathcal A(z)&\mathcal B(z)\\-\mathcal B(z)&\mathcal A(z)+z\mathcal B(z)\end{pmatrix},\qquad\det H_a(z)=\mathcal A^2+z\mathcal A\mathcal B+\mathcal B^2 .
\]
If $w$ lies on the unit circle, or is real with $e^{-\kappa_0}\le|w|\le e^{\kappa_0}$, then $z=w+w^{-1}$ is real with $|z|\le2\cosh\kappa_0$, and $\det H_a(z)=a(w)\,a(w^{-1})$. In particular $\det H_a(2\cos\theta)=|a(e^{i\theta})|^2$.
\end{lemma}

\begin{proof}
Absolute and uniform convergence follows from Lemma~\ref{lem:transfer} and $e^{\kappa_0}<R_0$. From $A_{j+1}=-U_{j-1}(z/2)=-B_j$ and $B_{j+1}=U_j(z/2)=zU_{j-1}(z/2)-U_{j-2}(z/2)=zB_j+A_j$ (valid for all $j\ge0$ with our conventions) we get the second row. Fix $w$ as stated. The sequence $A_j+wB_j$ satisfies the free recurrence with initial values $1,w$, so $A_j+wB_j=w^j$. Summing against $a_j$ (both series converge absolutely), $\mathcal A+w\mathcal B=a(w)$, and likewise $\mathcal A+w^{-1}\mathcal B=a(w^{-1})$. Multiplying, $(\mathcal A+w\mathcal B)(\mathcal A+w^{-1}\mathcal B)=\mathcal A^2+z\mathcal A\mathcal B+\mathcal B^2$. For $w=e^{i\theta}$, $a(w^{-1})=\overline{a(w)}$ because the coefficients are real.
\end{proof}

\subsection{The bulk estimate}

\begin{proposition}[No real double zeros in the bulk]\label{prop:bulk}
Let $(a_s)$ satisfy (H1), (H2) and (H4). There are $L>2$, $c>0$ and $n_1$ such that for all $n\ge n_1$, all $s\in[0,1]$ and all real $x$ with $|x|\le L\sqrt n$,
\[
Q_n^s(x)^2+\frac{Q_n^s{}'(x)^2}{n}\ \ge\ c^2\bigl(h_n(x)^2+h_{n-1}(x)^2\bigr)>0 .
\]
In particular, $Q_n^s$ has no multiple real zero in $[-L\sqrt n,L\sqrt n]$.
\end{proposition}

\begin{proof}
\emph{Choice of constants.} Let $R_0,C_0$ be as in (H1), and fix $\kappa_0>0$ with $e^{\kappa_0}<R_0$. By Lemma~\ref{lem:limit} and the Weierstrass test, the map $(s,z)\mapsto H_{a_s}(z)$ is continuous on the set $[0,1]\times[-2\cosh\kappa_0,\,2\cosh\kappa_0]$. By (H4) and compactness, $\mu=\min\{|a_s(w)|:s\in[0,1],\,|w|=1\}>0$, and by Lemma~\ref{lem:limit}, $\det H_{a_s}(z)\ge\mu^2$ for $z\in[-2,2]$. By uniform continuity there is $\kappa\in(0,\kappa_0]$ such that
\begin{equation}\label{eq:det-lower}
\det H_{a_s}(z)\ge\mu^2/2\qquad\text{for all }s\in[0,1]\text{ and real }|z|\le L:=2\cosh\kappa .
\end{equation}
Let $\Gamma=\sup\{\|H_{a_s}(z)\|:s\in[0,1],\,|z|\le L\}<\infty$. Fix $\ell>1/(2\log R_0)$ and put $M_n=\lceil\ell\log n\rceil$. Below, $n$ is so large that $M_n+2\le n/2$ and $(M_n+1)^2\le n$.

\emph{Convergence of the head.} We claim that
\begin{equation}\label{eq:head-conv}
\delta_n=\sup\bigl\{\|H_{n,M_n}[a_s](z)-H_{a_s}(z)\|:\ s\in[0,1],\ z\in[-L,L]\bigr\}\to0 .
\end{equation}
By Lemma~\ref{lem:transfer}, for $j\le M_n$ and $|z|\le L$ both $\|G_{n,j}(z)\|$ and $\|G_j(z)\|$ are at most $C_1(j+2)e^{(j+1)\kappa}$, where $C_1=4e^{c_\kappa}$; here we bound the operator norm by the sum of the absolute values of the entries and use $\omega_{n,j}\le1$. For any fixed $J$ and all $n$ with $M_n>J$,
\[
\|H_{n,M_n}[a_s]-H_{a_s}\|\le\sum_{j\le J}C_0\|G_{n,j}-G_j\|+2C_0C_1\sum_{j>J}(j+2)e^{(j+1)\kappa}R_0^{-j}.
\]
The second sum is small for large $J$, since $e^{\kappa}<R_0$. For fixed $j$, $G_{n,j}\to G_j$ uniformly on $[-L,L]$ by \eqref{eq:fixed-j} and $\omega_{n,j}\to1$. This proves \eqref{eq:head-conv}.

\emph{The head is uniformly invertible.} For a real $2\times2$ matrix $X$ the smallest singular value is $|\det X|/\|X\|$, and $|\det X-\det Y|\le\|X-Y\|(\|X\|+\|Y\|)$ (write the determinant as a bilinear form in the columns). By \eqref{eq:det-lower} and \eqref{eq:head-conv}, for large $n$ we have $\delta_n\le1$ and $\det H_{n,M_n}[a_s](z)\ge\mu^2/2-\delta_n(2\Gamma+1)\ge\mu^2/4$, hence
\begin{equation}\label{eq:head-lower}
\|H_{n,M_n}[a_s](z)\,u\|\ge2c\,\|u\|\qquad(u\in\R^2),\qquad 2c:=\frac{\mu^2}{4(\Gamma+1)},
\end{equation}
for all $s\in[0,1]$ and $|z|\le L$.

\emph{The tail.} Let $E=Q_n^s-Q_{n,M_n}$ with $Q_{n,M_n}$ built from $a_s$. Then $E(x)=\sum_{j>M_n}a_j(s)h_{n-j}(x)$ and $E'(x)/\sqrt n=\sum_{j>M_n}a_j(s)\omega_{n,j}h_{n-j-1}(x)$, where only indices below $n$ occur, $0\le\omega_{n,j}\le1$ and $h_{-1}=0$. By Cauchy--Schwarz and Lemma~\ref{lem:cd}, for $x=\sqrt nz$ with $|z|\le L$,
\[
\Bigl\|\Bigl(E(x),\frac{E'(x)}{\sqrt n}\Bigr)\Bigr\|\le\sqrt2\Bigl(\sum_{j>M_n}a_j(s)^2\Bigr)^{1/2}\Bigl(\sum_{k<n}h_k(x)^2\Bigr)^{1/2}\le\eps_n\|v_n(x)\|,
\]
with $\eps_n=\sqrt2\,C_0(1-R_0^{-2})^{-1/2}R_0^{-M_n}\sqrt{n(1+L/2)}$. Since $R_0^{-M_n}\le n^{-\ell\log R_0}$ and $\ell\log R_0>\frac12$, we have $\eps_n\to0$.

\emph{Conclusion.} By \eqref{eq:head}, \eqref{eq:head-lower} and the tail bound, $\|(Q_n^s(x),Q_n^s{}'(x)/\sqrt n)\|\ge(2c-\eps_n)\|v_n(x)\|\ge c\|v_n(x)\|$ once $\eps_n\le c$, and $v_n(x)\neq0$ by Lemma~\ref{lem:zeros}.
\end{proof}

\begin{remark}[Why the head must be short]\label{rem:full}
Expressing the whole sum $Q_n$ through $(h_n,h_{n-1})$, that is, using all of \eqref{eq:AB} for $0\le j\le n$, is badly conditioned. At $z=2$ write $b_j=B_{n,j}(2)$, so that $\omega_{n,j-1}b_j=2b_{j-1}-\omega_{n,j-2}b_{j-2}$ for $2\le j\le n$ by \eqref{eq:backward}. Induction gives $b_j\ge b_{j-1}/\omega_{n,j-1}>0$ for $1\le j\le n$: if $b_{j-2}\le\omega_{n,j-2}b_{j-1}$, then $\omega_{n,j-2}b_{j-2}\le b_{j-1}$, so $\omega_{n,j-1}b_j\ge b_{j-1}$. Hence, by the bound $n!\le e\sqrt n\,(n/e)^n$,
\[
B_{n,n}(2)\ge\prod_{i=1}^{n-1}\frac1{\omega_{n,i}}=\Bigl(\frac{n^{n-1}}{(n-1)!}\Bigr)^{1/2}=\Bigl(\frac{n^{n}}{n!}\Bigr)^{1/2}\ge\frac{e^{n/2}}{(e^2n)^{1/4}} .
\]
For $a(w)=\sum_j t^{-j/2}w^j$ with $1<t<e$, which is admissible with radius $\sqrt t$, every term of $\sum_{j\le n}a_jB_{n,j}(2)$ is nonnegative, and the sum is at least $t^{-n/2}B_{n,n}(2)\to\infty$. So the full backward coefficients are unbounded at the edge of the bulk, although the physical values $Q_n(2\sqrt n)$ and $h_n(2\sqrt n)$ are not affected. The proof above avoids the problem by expanding only $O(\log n)$ terms and estimating the rest directly through Lemma~\ref{lem:cd}.
\end{remark}

\section{The exterior}\label{sec:exterior}

\begin{proposition}[Exterior asymptotics]\label{prop:exterior}
Let $(a_s)$ satisfy (H1). For every compact set $K\subset\C\setminus[-2,2]$,
\[
\sup_{s\in[0,1],\,z\in K}\Bigl|\frac{Q_n^s(\sqrt nz)}{h_n(\sqrt nz)}-a_s(r(z))\Bigr|\to0\qquad(n\to\infty).
\]
\end{proposition}

\begin{proof}
Let $\ell^2=\ell^2(\{0,1,2,\dots\})$ with basis $(e_i)$. Let $T_\infty$ be the free Jacobi operator, $T_\infty e_i=e_{i-1}+e_{i+1}$ (with $e_{-1}=0$), and let $T_n e_i=\tau_{n,i-1}e_{i-1}+\tau_{n,i}e_{i+1}$, where $\tau_{n,i}=\sqrt{1-(i+1)/n}$ for $0\le i\le n-2$ and $\tau_{n,i}=0$ otherwise. Thus $T_n$ is the direct sum of $n^{-1/2}J_n$, with the order of the basis reversed, and the zero operator. Both operators are self-adjoint with norm at most $2$.

Fix $z\in K$ and let $d_K=\dist(K,[-2,2])$ and $\varrho=\max_K|r|<1$. By Lemma~\ref{lem:resolvent}, $u=(z-T_n)^{-1}e_0$ has coordinates $u_i=h_{n-1-i}(\sqrt nz)/h_n(\sqrt nz)$ for $i\le n-1$ and $u_i=0$ for $i\ge n$. A direct check shows that $y=(z-T_\infty)^{-1}e_0$ has coordinates $y_i=r(z)^{i+1}$: indeed $zy_0-y_1=1$ and $zy_i-y_{i-1}-y_{i+1}=0$ because $r+r^{-1}=z$. By the second resolvent identity,
\[
\|u-y\|=\|(z-T_n)^{-1}(T_n-T_\infty)y\|\le d_K^{-1}\|(T_n-T_\infty)y\|.
\]
The operator $T_n-T_\infty$ is tridiagonal with zero diagonal and off-diagonal entries $-\epsilon_{n,i}$, where $0\le\epsilon_{n,i}=1-\tau_{n,i}\le\min(1,(i+1)/n)$. Since $|y_i|\le\varrho^{i+1}$,
\[
\|(T_n-T_\infty)y\|\le2\Bigl(\sum_{i\ge0}\epsilon_{n,i}^2\varrho^{2(i+1)}\Bigr)^{1/2}\le\frac2n\Bigl(\sum_{i\ge0}(i+1)^2\varrho^{2(i+1)}\Bigr)^{1/2}.
\]
Now $Q_n^s(\sqrt nz)/h_n(\sqrt nz)=a_0(s)+\sum_{j=1}^{n}a_j(s)u_{j-1}$ and $a_s(r(z))=a_0(s)+\sum_{j\ge1}a_j(s)y_{j-1}$. Since $u_i=0$ for $i\ge n$, their difference is $\sum_{j\ge1}a_j(s)(u_{j-1}-y_{j-1})$, which is bounded by
\[
\Bigl(\sum_{j\ge1}a_j(s)^2\Bigr)^{1/2}\|u-y\|\le\frac{2C_0}{(1-R_0^{-2})^{1/2}\,d_K\,n}\Bigl(\sum_{i\ge0}(i+1)^2\varrho^{2(i+1)}\Bigr)^{1/2}
\]
by (H1). The right side does not depend on $s\in[0,1]$ or $z\in K$ and tends to $0$.
\end{proof}

\begin{lemma}[Colleague matrix]\label{lem:colleague}
Let $a_0,\dots,a_n$ be real with $a_0\neq0$, let $Q(x)=\sum_{j=0}^{n}a_jh_{n-j}(x)$, and put
\[
b=(a_n,a_{n-1},\dots,a_1)^T\in\R^n,\qquad C=n^{-1/2}J_n-a_0^{-1}e_{n-1}b^T .
\]
Then
\[
\det(zI-C)=\frac{\sqrt{n!}}{a_0\,n^{n/2}}\,Q(\sqrt nz).
\]
Consequently every zero $z$ of $Q(\sqrt n\,\cdot)$ satisfies $|z|\le2+|a_0|^{-1}\bigl(\sum_{j=1}^{n}a_j^2\bigr)^{1/2}$.
\end{lemma}

\begin{proof}
Since $\det(xI-J_n)=\He_n(x)=\sqrt{n!}\,h_n(x)$ and, by \eqref{eq:jacobi}, $\operatorname{adj}(xI-J_n)e_{n-1}=\He_n(x)(xI-J_n)^{-1}e_{n-1}=\sqrt{(n-1)!}\,p_n(x)$ (first where $h_n(x)\ne0$, then everywhere as a polynomial identity), the matrix determinant lemma gives
\[
\det\bigl(xI-J_n+\tfrac{\sqrt n}{a_0}e_{n-1}b^T\bigr)=\sqrt{n!}\,h_n(x)+\frac{\sqrt n}{a_0}\sqrt{(n-1)!}\;b^Tp_n(x)=\frac{\sqrt{n!}}{a_0}Q(x).
\]
Put $x=\sqrt nz$ and divide by $n^{n/2}$. The bound follows from $\|C\|\le\|n^{-1/2}J_n\|+\|b\|/|a_0|$ and Lemma~\ref{lem:zeros}.
\end{proof}

\begin{lemma}[Quantitative exterior lemma]\label{lem:ext-quant}
Let $a(w)=\sum_{j\ge0}a_jw^j$ be a real power series with $|a_j|\le C_0R_0^{-j}$ for all $j\ge0$, where $C_0<\infty$ and $R_0>1$, and let $Q_n=\sum_{j=0}^{n}a_jh_{n-j}$. Let $0<d<1$ and $\delta>0$, and let $B$ be a closed disk with $B\cap\R=\emptyset$ and $B\subset\{|w|\le1-d\}$. Suppose that $|a|\ge\delta$ on $\partial B$ and that $a$ has $m\ge1$ zeros in the interior of $B$, counted with multiplicity. If
\begin{equation}\label{eq:n-ext}
n>n_{\rm ext}:=\frac{2\sqrt2\,C_0}{(1-R_0^{-2})^{1/2}\,\delta\,d^{7/2}},
\end{equation}
then the polynomial $z\mapsto Q_n(\sqrt n\,z)$ has exactly $m$ zeros, counted with multiplicity, in the domain $\mathcal J(\operatorname{int}B)$, where $\mathcal J(w)=w+w^{-1}$, and all of them are nonreal. In particular $Q_n$ has a nonreal zero.
\end{lemma}

Unlike Proposition~\ref{prop:exterior}, the lemma has an explicit onset, which depends on $a$ only through $C_0$, $R_0$, $\delta$ and $d$. This is what is needed when the symbol changes with $n$, as in the standard coupling of Section~\ref{sec:standard}.

\begin{proof}
\emph{A distance bound.} For $0<|w|<1$,
\begin{equation}\label{eq:joukowski-dist}
\dist\bigl(\mathcal J(w),[-2,2]\bigr)\ge\frac{(1-|w|)^2}{|w|}.
\end{equation}
Indeed $|\mathcal J(w)\mp2|=|w\mp1|^2/|w|$, so $|\mathcal J(w)-2|+|\mathcal J(w)+2|=2(|w|+|w|^{-1})$. For $y\in[-2,2]$ we have $|y-2|+|y+2|=4$, and the triangle inequality gives $2|\mathcal J(w)-y|\ge2(|w|+|w|^{-1})-4=2(1-|w|)^2/|w|$.

\emph{Proof of the lemma.} Since $B\cap\R=\emptyset$, the disk $B$ lies in $\D\setminus\{0\}$, where $\mathcal J$ is conformal. Hence $\Omega=\mathcal J(\operatorname{int}B)$ is a Jordan domain bounded by $\mathcal J(\partial B)$, and $K=\mathcal J(B)=\overline\Omega$ is a compact subset of $\C\setminus[-2,2]$. Moreover $K\cap\R=\emptyset$, because $\im \mathcal J(w)=\im w\,(1-|w|^{-2})\ne0$ for $w\in\D\setminus\R$. Let $F(z)=a(r(z))$ for $z\in\C\setminus[-2,2]$. Since $r(\mathcal J(w))=w$ for $w\in\D\setminus\{0\}$ and $\mathcal J$ is conformal near $B$, $F$ has exactly $m$ zeros in $\Omega$, counted with multiplicity, and $|F|\ge\delta$ on $\partial\Omega=\mathcal J(\partial B)$.

The function $(1-x)^2/x$ is decreasing on $(0,1)$, so by \eqref{eq:joukowski-dist}, $\dist(K,[-2,2])\ge d^2/(1-d)\ge d^2$; and $\varrho:=\max_K|r|=\max_B|w|\le1-d$. With $x=\varrho^2$ we have $1-x\ge1-(1-d)^2\ge d$, so
\[
\sum_{i\ge0}(i+1)^2\varrho^{2(i+1)}=\frac{x(1+x)}{(1-x)^3}\le\frac{2}{d^{3}}.
\]
The constant family $a_s=a$ satisfies (H1). So the last display of the proof of Proposition~\ref{prop:exterior}, with these bounds for $\dist(K,[-2,2])$ and $\varrho$, gives for $n>n_{\rm ext}$
\[
\sup_{z\in K}\Bigl|\frac{Q_n(\sqrt nz)}{h_n(\sqrt nz)}-F(z)\Bigr|\le\frac{2C_0}{(1-R_0^{-2})^{1/2}\,d^2\,n}\cdot\frac{\sqrt2}{d^{3/2}}=\frac{n_{\rm ext}\,\delta}{n}<\delta\le|F|\qquad\text{on }\partial\Omega .
\]
Both functions are analytic near $K$, since the zeros of $h_n(\sqrt n\,\cdot)$ lie in $(-2,2)$ (Lemma~\ref{lem:resolvent}). By Rouch\'e's theorem $Q_n(\sqrt n\,\cdot)/h_n(\sqrt n\,\cdot)$ has exactly $m$ zeros in $\Omega$, and since $h_n(\sqrt n\,\cdot)$ has no zeros in $K$, so has $Q_n(\sqrt n\,\cdot)$. They are nonreal because $\Omega\cap\R=\emptyset$.
\end{proof}

\section{Proof of the exact count}\label{sec:proof}

Let $a$ be admissible with radius of convergence $R>1$ and zeros $w_1,\dots,w_k$ in $\D$. These are finitely many, since $a$ is analytic and not identically zero near the closed disk and has no zeros on its boundary. They are nonzero because $a(0)=a_0\ne0$, and they are closed under conjugation.

\emph{The deformation.} Put $p_{\mathrm{in}}(w)=a_0\prod_{i=1}^{k}(1-w/w_i)$, a real polynomial of degree $k$, and $p_{\mathrm{out}}=a/p_{\mathrm{in}}$. Then $p_{\mathrm{out}}$ is analytic in $\{|w|<R\}$ (its singularities at the $w_i$ are removable, the zeros being simple), has real Taylor coefficients, satisfies $p_{\mathrm{out}}(0)=1$, and has no zeros in $\{|w|\le1\}$. Let
\[
a_s(w)=p_{\mathrm{in}}(w)\,p_{\mathrm{out}}(sw)=\sum_{j\ge0}a_j(s)\,w^j,\qquad s\in[0,1].
\]
For $s=1$ this is $a$, and for $s=0$ it is the polynomial $p_{\mathrm{in}}$. The family satisfies Hypothesis~\ref{hyp:H}. For (H1), fix $R_0\in(1,R)$; then $|a_s(w)|\le\max_{|w|\le R_0}|p_{\mathrm{in}}|\cdot\max_{|u|\le R_0}|p_{\mathrm{out}}(u)|=:C_0$ for $|w|=R_0$, and Cauchy's estimates apply. The coefficients $a_j(s)$ are polynomials in $s$, and $a_0(s)=a_0$. For $|w|\le1$ we have $|sw|\le1$, so $p_{\mathrm{out}}(sw)\ne0$. Hence, for every $s$, \emph{the zeros of $a_s$ in $\{|w|\le1\}$ are exactly $w_1,\dots,w_k$, all simple.} Let $Q_n^s=\sum_{j=0}^{n}a_j(s)h_{n-j}$; it is a real polynomial of degree $n$ with leading coefficient $a_0/\sqrt{n!}$, and its coefficients are continuous in $s$. Put
\[
\widehat Q_{n,s}(z)=Q_n^s(\sqrt n\,z).
\]

\emph{Geometry.} Let $L>2$ and $n_1$ be given by Proposition~\ref{prop:bulk}, and let $\varrho_0=2+|a_0|^{-1}C_0(1-R_0^{-2})^{-1/2}$. By (H1) and Lemma~\ref{lem:colleague}, all zeros of every $\widehat Q_{n,s}$ lie in $\{|z|\le\varrho_0\}$. Let $\delta_0>0$ be the minimum of the distances $\dist(\zeta_i,[-2,2])$, $|\zeta_i-\zeta_{i'}|$ ($i\ne i'$) and $|\im\zeta_i|$ (for nonreal $\zeta_i$); if $k=0$ put $\delta_0=1$. Let
\[
0<\eps<\min\bigl(1,\ (L-2)/2,\ \delta_0/3\bigr).
\]
Let $D_i$ be the closed disk of radius $\eps$ about $\zeta_i$, and $U=\{z:\dist(z,[-2,2])<\eps\}$. The disks $D_i$ are pairwise disjoint and disjoint from $\overline U$; those with nonreal centre do not meet $\R$, and those with real centre are symmetric about $\R$. Moreover $U\cap\R=(-2-\eps,2+\eps)\subset(-L,L)$. Let
\[
K=\{|z|\le\varrho_0+|\zeta_1|+\dots+|\zeta_k|+1\}\setminus\Bigl(U\cup\bigcup_i\operatorname{int}D_i\Bigr),
\]
a compact subset of $\C\setminus[-2,2]$ that contains every circle $\partial D_i$.

\emph{Uniform exterior control.} Let $F_s(z)=a_s(r(z))$. Since $r$ is a bijection of $\C\setminus[-2,2]$ onto $\D\setminus\{0\}$ with $r(\zeta_i)=w_i$ and $r'\neq0$, the zeros of $F_s$ are exactly $\zeta_1,\dots,\zeta_k$, all simple, for every $s$. The function $(s,z)\mapsto F_s(z)$ is continuous on $[0,1]\times K$ and has no zeros there, so $m=\min_{[0,1]\times K}|F_s(z)|>0$. By Proposition~\ref{prop:exterior} there is $n_2$ such that $|\widehat Q_{n,s}(z)/h_n(\sqrt nz)-F_s(z)|<m$ on $K$ for all $n\ge n_2$ and $s\in[0,1]$. Since $h_n(\sqrt n\,\cdot)$ has no zeros off $[-2,2]$ (Lemma~\ref{lem:resolvent}), we conclude for $n\ge n_2$ and every $s$:
\begin{enumerate}
\item[(i)] $\widehat Q_{n,s}$ has no zeros in $K$, and none in $\{|z|>\varrho_0\}$; so all its zeros lie in $U\cup\bigcup_i\operatorname{int}D_i$;
\item[(ii)] by Rouch\'e's theorem on $\partial D_i$, $\widehat Q_{n,s}$ has exactly one zero in $\operatorname{int}D_i$, counted with multiplicity. It is real when $\zeta_i$ is real, because $D_i$ is symmetric about $\R$ and $\widehat Q_{n,s}$ has real coefficients. It is nonreal when $\zeta_i$ is nonreal.
\end{enumerate}
Let $n_0\ge\max(n_1,n_2,k)$. For $n\ge n_0$ and every $s$, the real zeros of $\widehat Q_{n,s}$ in $U$ lie in $(-L,L)$ and are simple by Proposition~\ref{prop:bulk}, and the real zeros in the disks $D_i$ are simple by (ii). So all real zeros of every $\widehat Q_{n,s}$ are simple.

\emph{Counting.} Fix $n\ge n_0$. By Lemma~\ref{lem:constancy}, the number $N(s)$ of nonreal zeros of $\widehat Q_{n,s}$ does not depend on $s$. At $s=0$ we have $Q_n^0=\sum_{j=0}^{k}a_j(0)h_{n-j}$, because $a_s$ is the polynomial $p_{\mathrm{in}}$ of degree $k$ for $s=0$, and $a_0(0)=a_0\ne0$. By Lemma~\ref{lem:quasi}, $\widehat Q_{n,0}$ has at least $n-k$ distinct real zeros in $(-2,2)\subset U$. Together with the $k$ zeros in the disks $D_i$, which are disjoint from $U$, these account for all $n$ zeros. So the zeros of $\widehat Q_{n,0}$ in $U$ are $n-k$ simple real zeros, and $N(0)$ is the number of nonreal $w_i$. Hence so is $N(1)$.

At $s=1$, by (ii) the disks $D_i$ with nonreal centre already contain $N(1)$ nonreal zeros. So $\widehat Q_{n,1}(z)=Q_n(\sqrt nz)$ has no nonreal zeros in $U$. Its $n-k$ zeros in $U$ are therefore real, lie in $(-2-\eps,2+\eps)$, and are simple by Proposition~\ref{prop:bulk}, and its zeros in the disks $D_i$ are simple by (ii). This proves Theorem~\ref{thm:main} for every $\eps$ as above. \qed

\begin{remark}\label{rem:hypotheses}
The conditions of Definition~\ref{def:admissible} on the zeros of $a$ are used as follows. If $a$ has a zero $e^{i\theta}$ on the unit circle, then $\det H_a(2\cos\theta)=0$ and Proposition~\ref{prop:bulk} fails near $2\cos\theta$. If $a$ has a multiple real zero in $\D$, then Rouch\'e's theorem only gives a cluster of zeros near the corresponding $\zeta_i$, whose reality is not decided by the leading-order asymptotics; multiple nonreal zeros are allowed in the version of Theorem~\ref{thm:main} with an explicit onset in \cite{CompA}. % MERGE: was Theorem thm:quant of Paper I Neither situation is covered by Theorem~\ref{thm:main}; Section~\ref{sec:collision} treats the first one quantitatively for the family $a=q(\cdot/\sqrt t)$ of \eqref{eq:heat-hermite}, in which a conjugate pair of zeros crosses the unit circle.
\end{remark}


\clearpage
\part{The limit law}

\section{Heat flow and random power series}\label{sec:heat}

Theorem~\ref{thm:main} is deterministic, and so is its consequence for a fixed power series (Theorem~\ref{thm:fixedq}). Probability enters only through three facts about the random series $q$: its zeros in $\D\setminus\{0\}$ are simple (Lemma~\ref{lem:resample}), it has no zero on a fixed circle (Lemma~\ref{lem:circle}), and it has nonreal zeros (Proposition~\ref{prop:nonreal}).

\subsection{The heat flow}\label{subsec:heatflow}

\begin{lemma}\label{lem:heat}
Let $P(z)=\sum_{k=0}^{n}c_k(\sqrt nz)^k/\sqrt{k!}$ with real coefficients $c_k$, not all zero, let $d=\max\{k:c_k\ne0\}$, and for $t\ge0$ let $P_t=\exp(-\frac t{2n}\partial_z^2)P$.
\begin{enumerate}
\item[(a)] For $t>0$, $P_t(z)=\sum_{k=0}^{n}c_k\,t^{k/2}h_k\bigl(\sqrt{n/t}\,z\bigr)$.
\item[(b)] For a real polynomial $p$ let $Z(p)$ be the number of nonreal zeros of $p$, counted with multiplicity. If $p\not\equiv0$ is a real polynomial and $\alpha\ge0$, then $\exp(-\alpha\partial_z^2)p$ has the degree and the leading coefficient of $p$, and $Z(\exp(-\alpha\partial_z^2)p)\le Z(p)$. In particular, if $P_t$ has only real zeros, then so has $P_{t+s}$ for every $s\ge0$.
\item[(c)] The set $\{t\ge0:P_t\text{ has only real zeros}\}$ is a closed interval $[\Lambda,\infty)$ with $\Lambda<\infty$.
\end{enumerate}
\end{lemma}

\begin{proof}
(a) Comparing the generating function $e^{xs-s^2/2}=\sum_k\He_k(x)s^k/k!$ with
\[
\exp\bigl(-\tfrac12\partial_x^2\bigr)e^{xs}=e^{-s^2/2}e^{xs}
\]
gives $\He_k(x)=\exp(-\frac12\partial_x^2)x^k$. Substituting $x=\sqrt{n/t}\,z$, so that $\frac12\partial_x^2=\frac t{2n}\partial_z^2$, gives
\[
\exp\Bigl(-\frac t{2n}\partial_z^2\Bigr)(\sqrt nz)^k=t^{k/2}\He_k\bigl(\sqrt{n/t}\,z\bigr).
\]

(b) The degree and the leading coefficient are preserved, since $\partial^2$ lowers the degree. The inequality $Z(\exp(-\alpha\partial^2)p)\le Z(p)$ is a classical theorem of Laguerre and P\'olya: if $\varphi$ is an entire function of the Laguerre--P\'olya class, such as $\varphi(x)=e^{-\alpha x^2}=\lim_m(1-\frac\alpha mx^2)^m$, then $\varphi(\partial)$ does not increase the number of nonreal zeros of a real polynomial \cite{Ob}. For real-rooted $p$ the real-rootedness of $\exp(-\alpha\partial^2)p$, equivalently of $\sum_kq_k\He_k$ when $\sum_kq_kx^k$ is real-rooted, is also proved by Iserles and Saff \cite{IS}.
% TODO(author): pinpoint in [Ob] (Laguerre's theorem for phi(D), phi in the Laguerre--Polya class) not checked; no
% referee of idea 106 has seen the book. Give chapter and Satz, or cite a checked source for the Hermite--Poulain--
% Laguerre theorem, then delete this comment. Nothing depends on it: the sketch that follows is a complete outline.
For a real-rooted polynomial $\varphi$ this reduces, factor by factor, to $Z(p-cp')\le Z(p)$ for real $c\ne0$, which follows from Rolle's theorem applied to $e^{-x/c}p(x)$; the case $\varphi(x)=e^{-\alpha x^2}$ follows by the lower semicontinuity below, since the polynomials $(1-\frac\alpha m\partial^2)^mp$, all of the degree of $p$, converge coefficientwise to $\exp(-\alpha\partial^2)p$. The last statement of (b) is the case $p=P_t$, $\alpha=\frac s{2n}$, since $P_{t+s}=\exp(-\frac s{2n}\partial_z^2)P_t$.

\emph{Lower semicontinuity.} If real polynomials $p_m$ of degree $d$ converge coefficientwise to a polynomial $p$ of degree $d$, then $Z(p)\le\liminf_mZ(p_m)$. Indeed, around each nonreal zero of $p$, of multiplicity $\mu$, choose a closed disk that does not meet $\R$ or the other zeros of $p$; by Hurwitz's theorem it contains exactly $\mu$ zeros of $p_m$ for large $m$, all nonreal.

(c) By (b) the set is upward closed. It is closed by the lower semicontinuity above, because the coefficients of $P_t$ are polynomials in $t$ and $P_t$ has the degree and the leading coefficient of $P$. It is nonempty: by (a), $t^{-d/2}P_t(\sqrt{t/n}\,x)=\sum_{k\le d}c_kt^{(k-d)/2}h_k(x)\to c_dh_d(x)$ coefficientwise as $t\to\infty$. Since $h_d$ has $d$ simple real zeros, the argument of Lemma~\ref{lem:constancy} shows that $P_t$ is real-rooted for large $t$.
\end{proof}

Let $(\eta_j)_{j\ge0}$ be any real sequence and define $P_n$, $P_{n,t}$ by \eqref{eq:reversed}. Applied to $P_n$ (with $c_k=\eta_{n-k}$) and substituting $k=n-j$, part (a) gives \eqref{eq:heat-hermite}:
\begin{equation}\label{eq:PQ}
P_{n,t}(z)=t^{n/2}\,Q_n^{(t)}\bigl(\sqrt n\,z/\sqrt t\bigr),\qquad Q_n^{(t)}=\sum_{j=0}^{n}\eta_jt^{-j/2}h_{n-j},
\end{equation}
which is the polynomial of Theorem~\ref{thm:main} for $a=a^{(t)}$, where $a^{(t)}(w)=q(w/\sqrt t)$. If $P_n\not\equiv0$, the real-rooting time $\Lambda_n=\inf\{t\ge0:P_{n,t}\text{ has only real zeros}\}$ is finite, and by Lemma~\ref{lem:heat}(c) the polynomial $P_{n,t}$ has only real zeros exactly when $t\ge\Lambda_n$. If $P_n\equiv0$ we put $\Lambda_n=0$.

\subsection{A fixed power series}

Theorem~\ref{thm:heat} is the almost sure form of a deterministic statement about one power series. Let $(\eta_j)_{j\ge0}$ be real, not all zero, and let $q(w)=\sum_j\eta_jw^j$ have radius of convergence $\rho_q\in(0,\infty]$. Let $Z_q$ be the set of zeros of $q$ in $\{0<|w|<\rho_q\}$ and put
\begin{equation}\label{eq:Rstar-det}
R_*=\inf\{|u|:\ u\in Z_q\setminus\R\}\in(0,\rho_q],\qquad T_*=R_*^{-2}\in[0,\infty),
\end{equation}
with $R_*=\rho_q$ if $Z_q\subset\R$, and $\infty^{-2}=0$. The infimum is positive because the zeros of $q\not\equiv0$ do not accumulate at $0$. When $\rho_q=1$ and $q$ has a nonreal zero in $\D$, this agrees with \eqref{eq:Rstar}.

\begin{theorem}[A fixed power series]\label{thm:fixedq}
Let $q$, $\rho_q$, $Z_q$, $R_*$, $T_*$ be as above, and let $\Lambda_n$ be the real-rooting time of $P_{n,t}$.
\begin{enumerate}
\item[(a)] Let $t>0$ with $\sqrt t\,\rho_q>1$ be such that $q$ has no zero on $\{|w|=t^{-1/2}\}$ and all zeros of $q$ in $\{0<|w|<t^{-1/2}\}$ are simple. For every sufficiently small $\eps>0$ there is $n_0$ such that for $n\ge n_0$: all zeros of $P_{n,t}$ are simple; for each $u\in Z_q$ with $|u|<t^{-1/2}$, exactly one zero of $P_{n,t}$ lies within distance $\eps$ of $tu+u^{-1}$, and it is real if and only if $u$ is real; all other zeros of $P_{n,t}$ are real and lie in $(-2\sqrt t-\eps,2\sqrt t+\eps)$. In particular the number of nonreal zeros of $P_{n,t}$ equals the number of nonreal zeros of $q$ in $\{|w|<t^{-1/2}\}$.
\item[(b)] If every real $u\in Z_q$ with $|u|<R_*$ is a simple zero of $q$, then $\limsup_{n}\Lambda_n\le T_*$.
\item[(c)] If $R_*<\rho_q$, then $\liminf_n\Lambda_n\ge T_*$.
\end{enumerate}
In particular, if $q$ has a nonreal zero in its disk of convergence and its nonzero real zeros of modulus less than $R_*$ are simple, then $\Lambda_n\to T_*$.
\end{theorem}

\begin{proof}
\emph{Reduction to $\eta_0\ne0$.} Let $m_0=\min\{j:\eta_j\ne0\}$ and $\tilde q(w)=q(w)/w^{m_0}$, and for $n>m_0$ let $n'=n-m_0$ and let $\widetilde P_{n'}$, $\widetilde P_{n',t}$ be the polynomials \eqref{eq:reversed} built from the coefficients $\tilde\eta_j=\eta_{m_0+j}$ of $\tilde q$. Since $\eta_{n-k}=0$ for $k>n'$, we have $P_n(z)=\sum_{k\le n'}\eta_{n-k}(\sqrt nz)^k/\sqrt{k!}=\widetilde P_{n'}(\lambda z)$ with $\lambda=\sqrt{n/n'}$. For every polynomial $p$, $\exp(-s\partial_z^2)[p(\lambda\,\cdot)](z)=[\exp(-s\lambda^2\partial^2)p](\lambda z)$; with $s=t/(2n)$ this gives
\begin{equation}\label{eq:reduction}
P_{n,t}(z)=\widetilde P_{n',t}\bigl(\sqrt{n/n'}\,z\bigr)\qquad(n>m_0,\ t\ge0).
\end{equation}
So $\Lambda_n(q)=\Lambda_{n-m_0}(\tilde q)$, and the zeros of $P_{n,t}$ are those of $\widetilde P_{n',t}$ multiplied by $\sqrt{n'/n}=1+O(1/n)$. Since $Z_{\tilde q}=Z_q$ and $\rho_{\tilde q}=\rho_q$, the numbers $R_*,T_*$ and all hypotheses are the same for $q$ and $\tilde q$. Hence (b) and (c) for $q$ follow from (b) and (c) for $\tilde q$. For (a), the points $tu+u^{-1}$ ($u\in Z_q$, $|u|<t^{-1/2}$) are finitely many, as $t^{-1/2}<\rho_q$; they are distinct and off $[-2\sqrt t,2\sqrt t]$, because $u\mapsto tu+u^{-1}=\sqrt t(w+w^{-1})$, $w=\sqrt tu$, is injective on $\{0<|u|<t^{-1/2}\}$ with values off that interval. Take $\eps$ smaller than a third of their mutual distances and of their distances to $[-2\sqrt t,2\sqrt t]$, and apply (a) for $\tilde q$ with $\eps/2$. The zeros of $\widetilde P_{n',t}$ then lie in a bounded set, so the dilation by $\sqrt{n'/n}$ moves each of them by less than $\eps/2$ for large $n$; it preserves reality and simplicity, and it maps the $\eps/2$-neighbourhoods into the $\eps$-neighbourhoods, which are disjoint. So (a) holds for $q$. From now on $\eta_0\ne0$.

(a) The series $a^{(t)}(w)=q(w/\sqrt t)$ has real coefficients, radius of convergence $\sqrt t\rho_q>1$, $a^{(t)}(0)=\eta_0\ne0$ and no zero on the unit circle. Its zeros in $\D$ are the points $\sqrt t\,u$ with $u\in Z_q$, $|u|<t^{-1/2}$, and these are simple. So $a^{(t)}$ is admissible. By \eqref{eq:PQ}, $z$ is a zero of $P_{n,t}$ if and only if $z/\sqrt t$ is a zero of $Q_n^{(t)}(\sqrt n\,\cdot)$, with the same multiplicity, and $\sqrt t\bigl(\sqrt tu+(\sqrt tu)^{-1}\bigr)=tu+u^{-1}$. Theorem~\ref{thm:main}, applied to $a^{(t)}$ with $\eps/\sqrt t$ in place of $\eps$, gives (a) after the dilation $z\mapsto\sqrt t\,z$.

(b) The zeros of $q$ in $\{|w|<\rho_q\}$ are isolated, so the set $M=\{|u|:u\in Z_q\}$ is countable. Let $t>T_*$, $t>0$, with $t^{-1/2}\notin M$. Then $t^{-1/2}<R_*\le\rho_q$, so $\sqrt t\rho_q>1$; $q$ has no zero on $\{|w|=t^{-1/2}\}$; and every $u\in Z_q$ with $|u|<t^{-1/2}$ is real, since $|u|<R_*$, hence simple. By (a), $P_{n,t}$ has only real zeros for all large $n$, that is, $\Lambda_n\le t$. Such $t$ are dense in $(T_*,\infty)$.

(c) Since $R_*<\rho_q$, the closed disk $\{|w|\le\min(R_*+1,(R_*+\rho_q)/2)\}$ lies in the disk of convergence and contains finitely many zeros of $q$, so the infimum $R_*$ is attained at a nonreal zero $u_*$. Let $\rho_q^{-2}<t<T_*$. Then $a^{(t)}$ has radius of convergence $\sqrt t\rho_q>1$, so by Cauchy's estimates the constant family $a_s=a^{(t)}$ satisfies (H1) of Hypothesis~\ref{hyp:H} for some $R_0>1$. The point $w_*=\sqrt tu_*$ is a nonreal zero of $a^{(t)}$ in $\D$, so $\zeta_*=w_*+w_*^{-1}$ is nonreal: $\im\zeta_*=\im w_*\,(1-|w_*|^{-2})\ne0$. The function $a^{(t)}\circ r$ is analytic on $\C\setminus[-2,2]$, vanishes at $\zeta_*$, and is not identically zero, since it tends to $\eta_0$ at infinity. Choose a closed disk $B$ centred at $\zeta_*$, disjoint from $\R$, on whose boundary $a^{(t)}\circ r$ has no zero. By Proposition~\ref{prop:exterior}, $Q_n^{(t)}(\sqrt nz)/h_n(\sqrt nz)\to a^{(t)}(r(z))$ uniformly on $B$, so by Rouch\'e's theorem this quotient has a zero in $B$ for all large $n$. As $h_n(\sqrt n\,\cdot)$ has no zeros off $[-2,2]$ (Lemma~\ref{lem:resolvent}), it is a zero of $Q^{(t)}_n(\sqrt n\,\cdot)$, and it is nonreal. By \eqref{eq:PQ}, $P_{n,t}$ has a nonreal zero, so $\Lambda_n>t$ by Lemma~\ref{lem:heat}(c). Hence $\liminf\Lambda_n\ge t$ for every $t<T_*$ close to $T_*$.
\end{proof}

\begin{remark}\label{rem:fixedq}
(i) No condition on circles is needed for the limit: (b) uses only times outside the countable set $\{t:t^{-1/2}\in M\}$, and (c) uses none. (ii) Simplicity is needed only for the real zeros of modulus less than $R_*$; for independent coefficients with a common law as in Lemma~\ref{lem:resample} it holds almost surely by that lemma. (iii) $T_*\le1$ is possible when $\rho_q>1$. For $q(w)=1+w^2/4$ we have $R_*=2$, so $\Lambda_n\to\frac14$; indeed $Q^{(t)}_n=h_n+(4t)^{-1}h_{n-2}$, and the zeros $\pm2i\sqrt t$ of $a^{(t)}$ lie in $\D$ exactly when $t<\frac14$. So $T_*>1$ in Proposition~\ref{prop:law} is a property of series of radius $1$. (iv) If $q$ has no nonreal zero in its disk of convergence and $\rho_q<\infty$, only (b) applies, and gives $\limsup\Lambda_n\le\rho_q^{-2}$. For $t<\rho_q^{-2}$ the series $a^{(t)}$ has radius less than $1$, and Proposition~\ref{prop:exterior} does not apply.
\end{remark}

\subsection{Random coefficients}

Throughout this subsection $\eta_0,\eta_1,\dots$ are independent random variables with a common law $\mu$ that is not a point mass and satisfies $\int\log^+|x|\,d\mu(x)<\infty$, and $q(w)=\sum_j\eta_jw^j$. Put $\gamma_\mu=\sum_{x\in\R}\mu(\{x\})^2$; then $\gamma_\mu\le\max_x\mu(\{x\})<1$. Since $\sum_{k\ge1}\PP(X>\delta k)\le\delta^{-1}\E X$ for $X=\log^+|\eta_0|$, we have $\sum_k\PP(|\eta_k|>e^{\delta k})<\infty$ for every $\delta>0$, and the Borel--Cantelli lemma gives, almost surely, $|\eta_k|\le e^{\delta k}$ for all large $k$, for every $\delta>0$ (it suffices to take $\delta=1/m$, $m\in\N$), so that $\rho_q\ge1$. Since $\mu\ne\delta_0$, there is $b>0$ with $\PP(|\eta_0|\ge b)>0$, and then $|\eta_k|\ge b$ for infinitely many $k$ by the second Borel--Cantelli lemma; hence $q\not\equiv0$ and $\rho_q=1$ almost surely. The events used below are measurable on $\{\rho_q\ge1\}$. For instance, $q$ has a multiple zero in a compact set $K\subset\D$ exactly when $\min_K(|q|+|q'|)=0$, and this minimum may be taken over a countable dense subset of $K$.

\begin{lemma}[Simple zeros]\label{lem:resample}
Let $\eta_0,\eta_1,\dots$ be independent with a common law $\mu$ that is not a point mass and has $\int\log^+|x|\,d\mu(x)<\infty$. Almost surely every zero of $q$ in $\D\setminus\{0\}$ is simple.
\end{lemma}

\begin{proof}
\emph{A deterministic claim.} Let $g\not\equiv0$ be analytic in $\D$, let $K\subset\D\setminus\{0\}$ be compact with $r_K=\min_K|u|$ and $\rho_K=\max_K|u|$, and let $\delta>0$ with $e^\delta\rho_K<1$. Then there is $k_0$ such that for every $k\ge k_0$ and every $c\in\C$ with $0<|c|\le e^{\delta k}$, the function $g+cw^k$ has no multiple zero in $K$.

To see this, let $K'=\{u:\dist(u,K)\le\tau\}$ with $\tau>0$ so small that $K'\subset\D\setminus\{0\}$. In $K'$ the function $g$ has finitely many zeros $z_1,\dots,z_p$, of multiplicities $m_1,\dots,m_p$, and $g=h\prod_i(w-z_i)^{m_i}$ with $h$ analytic and zero-free on $K'$. Let $m=\sum_im_i$ and $M=\max_im_i$, fix $s>0$ with $2s<\tau$ and $2s<|z_i-z_{i'}|$ for $i\ne i'$, and put $C_1=\max_K|h'/h|+m/s$ and $c_1=\min_K|h|\cdot\min(s,1)^m$. Let $u\in K$. If $|u-z_i|\ge s$ for all $i$, then $|g'(u)/g(u)|\le C_1$. If $|u-z_i|<s$ for some $i$, then $|u-z_{i'}|>s$ for $i'\ne i$, so $|g'(u)/g(u)-m_i/(u-z_i)|\le C_1$ and $|g(u)|\ge c_1|u-z_i|^{m_i}$.

Now let $u_1\in K$ be a multiple zero of $g+cw^k$ with $c\ne0$. Then $g(u_1)=-cu_1^k\ne0$ and $g'(u_1)=-kcu_1^{k-1}$, so $g'(u_1)/g(u_1)=k/u_1$, whose modulus lies in $[k/\rho_K,k/r_K]$ and so exceeds $k$. Let $k>C_1$. Then the first case is excluded, so $|u_1-z_i|<s$ for some $i$, and $m_i/|u_1-z_i|\le k/r_K+C_1\le2k/r_K$. Hence $|u_1-z_i|\ge r_K/(2k)$ and $|g(u_1)|\ge c_1(r_K/2k)^M$. But $|g(u_1)|=|c||u_1|^k\le(e^\delta\rho_K)^k$, which is smaller than $c_1(r_K/2k)^M$ for all large $k$. This proves the claim.

\emph{Resampling.} Fix a compact $K\subset\D\setminus\{0\}$ and $\delta>0$ with $e^\delta\max_K|u|<1$, and let $A$ be the event that $\rho_q\ge1$ and $q$ has a multiple zero in $K$. Let $\eta'_0,\eta'_1,\dots$ be independent with law $\mu$, and independent of $\eta=(\eta_j)_{j\ge0}$. Let $\eta^{(k)}$ be $\eta$ with its $k$-th coordinate replaced by $\eta'_k$. Since the coordinates are independent, $\eta^{(k)}$ has the law of $\eta$, and the series built from $\eta^{(k)}$ is $q+(\eta'_k-\eta_k)w^k$. Let $\varphi(x)=\mu(\{x\})$. Then $\PP(\eta'_k=\eta_k\mid\eta)=\varphi(\eta_k)$, and $\eta^{(k)}=\eta$ on $\{\eta'_k=\eta_k\}$, so
\begin{equation}\label{eq:resample}
\PP(A)=\PP(\eta^{(k)}\in A)=\E\bigl[\mathbf 1_A(\eta)\,\varphi(\eta_k)\bigr]+\PP\bigl(\eta^{(k)}\in A,\ \eta'_k\ne\eta_k\bigr).
\end{equation}
\emph{First term.} Let $\eps>0$. There are $J$ and an event $A_J$ determined by $\eta_0,\dots,\eta_{J-1}$ with $\PP(A\triangle A_J)<\eps$. For $k\ge J$ the variable $\eta_k$ is independent of $A_J$, so $\E[\mathbf 1_{A_J}\varphi(\eta_k)]=\PP(A_J)\sum_x\mu(\{x\})^2=\gamma_\mu\PP(A_J)$. Since $0\le\varphi\le1$, the first term is at most $\gamma_\mu\PP(A)+2\eps$ for $k\ge J$.
\emph{Second term.} By the Borel--Cantelli arguments above, applied with $\delta/2$ to $\eta$ and to $\eta'$, almost surely $\rho_q\ge1$, $q\not\equiv0$, and $|\eta_k|+|\eta'_k|\le e^{\delta k}$ for all large $k$. On this event the deterministic claim, applied to $g=q$ and $c=\eta'_k-\eta_k$, shows that the indicator of $\{\eta^{(k)}\in A,\ \eta'_k\ne\eta_k\}$ vanishes for all large $k$. By dominated convergence the second term tends to $0$ as $k\to\infty$.

Hence $\PP(A)\le\gamma_\mu\PP(A)+2\eps$ for every $\eps>0$, and $\PP(A)=0$ because $\gamma_\mu<1$. Exhausting $\D\setminus\{0\}$ by the compact sets $\{1/m\le|w|\le1-1/m\}$ proves the lemma.
\end{proof}

No density, continuity or symmetry of the laws is used: the lemma covers Rademacher coefficients and all other nondegenerate atomic laws with a logarithmic moment. The limit $k\to\infty$ is essential: for random polynomials of fixed degree with an atomic coefficient law, a double zero in $\D\setminus\{0\}$ can have positive probability, for instance $1-4w+4w^2=(1-2w)^2$ when the law charges $1$, $-4$ and $4$. When $\mu(\{0\})>0$ the zero of $q$ at $w=0$ can be multiple; it is removed by \eqref{eq:reduction}.

\begin{lemma}[Zeros on circles]\label{lem:circle}
Let $\eta_0,\eta_1,\dots$ be as in Lemma~\ref{lem:resample}.
\begin{enumerate}
\item[(a)] If $\mu$ has no atoms, then for every $r\in(0,1)$, almost surely $q$ has no zero on the circle $\{|w|=r\}$.
\item[(b)] The set $\mathcal E$ of $t>1$ for which $\PP\bigl(q\text{ has a zero on }\{|w|=t^{-1/2}\}\bigr)>0$ is countable. By (a), it is empty if $\mu$ has no atoms.
\end{enumerate}
\end{lemma}

\begin{proof}
(a) Let $g(w)=\sum_{j\ge1}\eta_jw^j$, which is independent of $\eta_0$. Almost surely $\rho_q\ge1$ and not all $\eta_j$, $j\ge1$, vanish. Then $q$ has a zero on $\{|w|=r\}$ if and only if $-\eta_0\in g(\{|w|=r\})\cap\R$. The function $\theta\mapsto\im g(re^{i\theta})=\sum_{j\ge1}\eta_jr^j\sin(j\theta)$ is real-analytic, $2\pi$-periodic and not identically zero, so it has finitely many zeros in $[0,2\pi)$, and $g(\{|w|=r\})\cap\R$ is finite. Since $\eta_0$ is independent of $g$ and has no atoms, Fubini's theorem shows that $-\eta_0$ lies in this finite set with probability $0$.

(b) Fix $b\in(0,1)$. Almost surely $q\not\equiv0$ and $\rho_q\ge1$, so the number $N_b$ of zeros of $q$ in $\{|w|\le b\}$ is finite. Let $C_r$ be the event that $q$ has a zero on $\{|w|=r\}$. For distinct $r_1,\dots,r_m\in(0,b]$ each zero lies on at most one of the circles $\{|w|=r_i\}$, so
\[
\sum_{i=1}^m\PP(C_{r_i})\le\E\min(m,N_b)=o(m)\qquad(m\to\infty),
\]
by dominated convergence. So for each $\eps>0$ only finitely many $r\in(0,b]$ have $\PP(C_r)\ge\eps$. Letting $\eps\downarrow0$ and $b\uparrow1$ shows that only countably many $r\in(0,1)$ have positive probability.
\end{proof}

\begin{proposition}[Nonreal zeros]\label{prop:nonreal}
Let $\eta_0,\eta_1,\dots$ be independent with a common law $\mu$ that is not a point mass and has $\int x^2\,d\mu(x)<\infty$. Almost surely $q$ has infinitely many nonreal zeros in $\D$.
\end{proposition}

The proof looks at $q$ near the boundary point $i$, where a central limit theorem applies. We use the following standard facts. Let $\Omega\subset\C$ be a domain and $H(\Omega)$ the space of analytic functions on $\Omega$ with the topology of locally uniform convergence, a Polish space. By Montel's theorem, sets of functions that are uniformly bounded on $\Omega$ are relatively compact in $H(\Omega)$. The Borel $\sigma$-field of $H(\Omega)$ is generated by the evaluations at a countable dense subset of $\Omega$, so the law of a random element of $H(\Omega)$ is determined by its finite-dimensional distributions. We also use Prokhorov's theorem, the portmanteau theorem for closed sets, Slutsky's lemma ($X_N\Rightarrow X$ and $d(X_N,Y_N)\to0$ in probability imply $Y_N\Rightarrow X$), and the fact that $X_N^{(l)}\Rightarrow X^{(l)}$ for $l=1,\dots,L$, with $X_N^{(1)},\dots,X_N^{(L)}$ independent for each $N$, implies joint convergence to independent copies; see \cite[Chapter~1]{Bi2}. The multivariate Lindeberg--Feller theorem follows from the one-dimensional one by the Cram\'er--Wold device \cite[Theorems~27.2 and~29.4]{Bi}. Finally, a Gaussian analytic function $F$ on $\Omega$ with covariance kernel $K(s,s')=\E F(s)\overline{F(s')}$, $K(s,s)>0$, has a first intensity of zeros equal to $\frac1{4\pi}\Delta\log K(s,s)$ with respect to Lebesgue measure (the Edelman--Kostlan formula \cite[Section~2.4]{HKPV}).

\begin{proof}[Proof of Proposition~\ref{prop:nonreal}]
Let $m=\E\eta_0$ and $\sigma^2=\operatorname{Var}\eta_0>0$; dividing $q$ by $\sigma$ we may assume $\sigma=1$. Almost surely $\rho_q\ge1$, as shown at the beginning of this subsection, since $\log^+|x|\le x^2$. Let $\Omega_1=\{|s-1|<\frac12\}$ and $\Omega_2=\{|s-1|<\frac34\}$. For $s\in\Omega_2$ we have $\re s>\frac14$ and $|s|<\frac74$, so $|1-\eps s|^2\le1-\eps/4$ for $0<\eps\le\eps_0=4/49$. For such $\eps$ put
\[
w_\eps(s)=i(1-\eps s)=\eps\im s+i(1-\eps\re s),\qquad f_\eps(s)=\sqrt\eps\,q\bigl(w_\eps(s)\bigr)\qquad(s\in\Omega_2).
\]
The points $w_\eps(s)$, $s\in\Omega_2$, lie in $\D$ and are nonreal.

\emph{Step 1: the mean.} Write $q(w)=m/(1-w)+q^0(w)$, where $q^0$ has the centred coefficients $\eta_j-m$, and $f_\eps=d_\eps+f^0_\eps$ accordingly. Since $|1-w_\eps(s)|\ge\sqrt2-\frac74\eps\ge1$, the deterministic function $d_\eps(s)=\sqrt\eps\,m/(1-w_\eps(s))$ tends to $0$ uniformly on $\Omega_2$ as $\eps\to0$.

\emph{Step 2: finite-dimensional distributions.} For $s,s'\in\Omega_2$,
\[
\E f^0_\eps(s)\overline{f^0_\eps(s')}=\frac{\eps}{1-(1-\eps s)(1-\eps\bar s')}\to\frac1{s+\bar s'},\qquad\E f^0_\eps(s)f^0_\eps(s')=\frac{\eps}{1+(1-\eps s)(1-\eps s')}\to0 .
\]
Fix $s_1,\dots,s_p\in\Omega_2$. The vector $(f^0_\eps(s_l))_{l\le p}$ is the sum over $k$ of the independent centred vectors $X_k=\sqrt\eps(\eta_k-m)(w_\eps(s_l)^k)_{l\le p}$. The sum over $k>K_\eps=\lceil\eps^{-2}\rceil$ has coordinates of variance at most $\eps\sum_{k>K_\eps}(1-\eps/4)^k\le4(1-\eps/4)^{K_\eps}\to0$. For the sum over $k\le K_\eps$ we have $|X_k|\le\sqrt{p\eps}\,|\eta_k-m|$, so for $\tau>0$, with $\ell_\eps=\E\bigl[(\eta_0-m)^2\mathbf 1_{\{|\eta_0-m|>\tau/\sqrt{p\eps}\}}\bigr]$, the Lindeberg sum is
\[
\sum_{k}\E\bigl[|X_k|^2\mathbf 1_{\{|X_k|>\tau\}}\bigr]\le\ell_\eps\sum_{l\le p}\eps\sum_k|w_\eps(s_l)|^{2k}\le4p\,\ell_\eps\to0 .
\] By the Lindeberg--Feller theorem, the finite-dimensional distributions of $f^0_\eps$ converge to those of a centred complex Gaussian process $F$ on $\Omega_2$ with $\E F(s)\overline{F(s')}=K(s,s'):=1/(s+\bar s')$ and $\E F(s)F(s')=0$.

\emph{Step 3: tightness.} For analytic $f$ on $\Omega_2$ and $s\in\Omega_1$ the disk $B(s,\frac14)$ lies in $\Omega_2$, and $|f(s)|^2\le\frac{16}\pi\int_{B(s,1/4)}|f|^2\,dA$. Since $\E|f^0_\eps(s)|^2=\eps/(1-|1-\eps s|^2)\le4$ on $\Omega_2$ and $\Omega_2$ has area $\frac{9\pi}{16}$,
\[
\E\sup_{\Omega_1}|f^0_\eps|^2\le\frac{16}\pi\int_{\Omega_2}\E|f^0_\eps|^2\,dA\le36 .
\]
By Markov's inequality, Montel's theorem and Prokhorov's theorem, the laws of $f^0_\eps$ on $H(\Omega_1)$ are tight. Every subsequential limit has the finite-dimensional distributions of $F$, so $F$ has a version in $H(\Omega_1)$, a Gaussian analytic function with kernel $K$, and $f^0_\eps\Rightarrow F$ in $H(\Omega_1)$ as $\eps\to0$. By Step 1, also $f_\eps\Rightarrow F$.

\emph{Step 4: the limit has zeros.} Here $K(s,s)=1/(2\re s)>0$ and $\frac1{4\pi}\Delta\log K(s,s)=\frac1{4\pi(\re s)^2}>0$. By the Edelman--Kostlan formula, $F$ has a zero in $B=\{|s-1|<\frac14\}$ with positive probability $p_B$. Also $F\not\equiv0$ almost surely, since $F(1)$ is a nondegenerate Gaussian variable.

\emph{Step 5: independent scales.} Let $N\ge4$ be an integer, $L\ge1$, and $\eps_l=N^{-2l}$ for $l=1,\dots,L$. Split $f^0_{\eps_l}=f^{\rm b}_l+f^{\rm r}_l$, where $f^{\rm b}_l$ keeps the indices $N^{2l-1}<k\le N^{2l+1}$. These blocks of indices are disjoint, so $f^{\rm b}_1,\dots,f^{\rm b}_L$ are independent. As in Step 3,
\[
\E\sup_{\Omega_1}|f^{\rm r}_l|^2\le9\sup_{\Omega_2}\eps_l\Bigl(\sum_{k\le N^{2l-1}}1+\sum_{k>N^{2l+1}}(1-\eps_l/4)^k\Bigr)\le9\Bigl(\frac2N+4e^{-N/4}\Bigr).
\]
Let $N\to\infty$ with $L$ fixed. Then $f^0_{\eps_l}\Rightarrow F$ and $\sup_{\Omega_1}|f^{\rm r}_l|\to0$ in probability, so $f^{\rm b}_l\Rightarrow F$ by Slutsky's lemma, for each $l$. By independence, $(f^{\rm b}_1,\dots,f^{\rm b}_L)\Rightarrow(F_1,\dots,F_L)$ in $H(\Omega_1)^L$, where $F_1,\dots,F_L$ are independent copies of $F$. By Slutsky's lemma and Step 1, also $(f_{\eps_1},\dots,f_{\eps_L})\Rightarrow(F_1,\dots,F_L)$.

\emph{Step 6: counting.} The set $\mathcal C=\{g\in H(\Omega_1):\ g\text{ has no zero in }B,\text{ or }g|_B\equiv0\}$ is closed, by Hurwitz's theorem. The regions $w_{\eps_l}(B)$ lie in the annuli $\{\frac34\eps_l<|w-i|<\frac54\eps_l\}$, which are disjoint for $N\ge2$, and they consist of nonreal points of $\D$. If $q\not\equiv0$ and $q$ has fewer than $j$ nonreal zeros in $\D$, then $f_{\eps_l}\notin\mathcal C$ for at most $j-1$ values of $l$. The set $\mathcal C_{L,j}$ of $(g_1,\dots,g_L)$ with $g_l\in\mathcal C$ for at least $L-j+1$ values of $l$ is a finite union of finite intersections of closed sets, hence closed. Let $V_j$ be the event that $q$ has fewer than $j$ nonreal zeros in $\D$. By the portmanteau theorem, and since $\PP(F\in\mathcal C)=1-p_B$,
\[
\PP(V_j)\le\limsup_{N\to\infty}\PP\bigl((f_{\eps_1},\dots,f_{\eps_L})\in\mathcal C_{L,j}\bigr)\le\PP\bigl((F_1,\dots,F_L)\in\mathcal C_{L,j}\bigr)=\PP\bigl(\mathrm{Bin}(L,p_B)<j\bigr).
\]
The left side depends neither on $N$ nor on $L$, and the right side tends to $0$ as $L\to\infty$. So $q$ has at least $j$ nonreal zeros in $\D$ almost surely, for every $j$.
\end{proof}

No zero--one law is used: the probability of the fixed event is bounded by a binomial tail for every $L$. For standard Gaussian coefficients the central limit step is not needed, since $f_\eps$ is exactly Gaussian. In that case Feldheim's law of large numbers for zeros of stationary Gaussian analytic functions in a strip \cite[Theorem~1]{Fe}, applied to $q(\tanh u)/\cosh u$ (whose covariance kernel is $\sech(u-\bar v)$), gives another proof.

The convergence of rescaled random power series with independent coefficients to a Gaussian analytic function near the boundary of the disk is the universality principle of Ledoan, Merkli and Starr \cite{LMS}, extended by Shirai \cite{Shi}; Steps~2 and~3 are an instance of it, at a single boundary point and with independent scales. For symmetric coefficient laws, Nazarov, Nishry and Sodin \cite[Theorem~1.3]{NNS} prove a stronger statement under weaker assumptions: if $\limsup_j|\eta_j|^{1/j}=1$ and $\sum_j\eta_j^2=\infty$ almost surely, then almost surely $q$ takes every complex value $b$ infinitely often in $\D$, with $\sum(1-|w|)=\infty$ over the solutions of $q(w)=b$. Whether this yields nonreal zeros without a second moment is discussed in Remark~\ref{rem:logmoment}.

We now prove Theorems~\ref{thm:heat} and~\ref{cor:weyl}. Let $\eta_0,\eta_1,\dots$ be independent with a common law $\mu$ that is not a point mass and has $\int x^2\,d\mu<\infty$. Then $\int\log^+|x|\,d\mu<\infty$, since $\log^+|x|\le x^2$, so the results above apply.

\begin{proof}[Proof of Theorem~\ref{thm:heat}]
Work on the almost sure event on which $\rho_q=1$, every zero of $q$ in $\D\setminus\{0\}$ is simple (Lemma~\ref{lem:resample}), and $q$ has a nonreal zero in $\D$ (Proposition~\ref{prop:nonreal}). The zeros of $q$ in each disk $\{|w|\le b\}$, $b<1$, are finitely many, so the infimum \eqref{eq:Rstar-det} is attained, it agrees with the minimum \eqref{eq:Rstar}, and $0<R_*<1=\rho_q$. Hence $1<T_*<\infty$, and Theorem~\ref{thm:fixedq} gives $\Lambda_n\to T_*$.

Now fix $t>1$ with $t\notin\mathcal E$, where $\mathcal E$ is the countable set of Lemma~\ref{lem:circle}(b). Then almost surely $q$ also has no zero on $\{|w|=t^{-1/2}\}$. On this event the hypotheses of Theorem~\ref{thm:fixedq}(a) hold, since $\sqrt t\rho_q=\sqrt t>1$ and all zeros of $q$ in $\{0<|w|<t^{-1/2}\}$ are simple, and Theorem~\ref{thm:fixedq}(a) gives the description of the zeros of $P_{n,t}$.
\end{proof}

\begin{proof}[Proof of Theorem~\ref{cor:weyl}]
For each $n$ the real-rooting time of a polynomial is a measurable function of its coefficients, namely the infimum of the rational $t\ge0$ at which the polynomial is real-rooted (Lemma~\ref{lem:heat}(c)). Since $(\eta_n,\dots,\eta_0)$ and $(\xi_0,\dots,\xi_n)$ have the same law, $\Lambda_n^{\rm std}$ has the law of $\Lambda_n$, which converges almost surely, hence in distribution, to $T_*$ by Theorem~\ref{thm:heat}. Likewise, for $t>1$ with $t\notin\mathcal E$, the number of nonreal zeros of $G_{n,t}$ has the law of the number $N_t^{(n)}$ of nonreal zeros of $P_{n,t}$, which by Theorem~\ref{thm:heat} is almost surely eventually equal to
\[
N_t=\#\{u\in\D\setminus\R:\ q(u)=0,\ |u|<t^{-1/2}\}.
\]
In particular $\PP(G_{n,t}\text{ real-rooted})=\PP(N_t^{(n)}=0)\to\PP(N_t=0)$. Finally $\{N_t=0\}=\{R_*\ge t^{-1/2}\}=\{T_*\le t\}$, and by the symmetry of the zeros under conjugation this is the event that $q$ has no zero in $\{|w|<t^{-1/2},\ \im w>0\}$.
\end{proof}

\begin{remark}[Without a second moment]\label{rem:logmoment}
Let the $\eta_j$ be independent with a common law $\mu$ that is not a point mass, and assume only $\int\log^+|x|\,d\mu<\infty$. Lemmas~\ref{lem:resample} and~\ref{lem:circle} and Theorem~\ref{thm:fixedq} still give the following. Almost surely, on the event that $q$ has a nonreal zero in $\D$, $\Lambda_n\to T_*$; off this event $R_*=\rho_q=1$, and $\limsup\Lambda_n\le1$ by Theorem~\ref{thm:fixedq}(b). For every $t>1$ with $t\notin\mathcal E$, almost surely the number of nonreal zeros of $P_{n,t}$ is eventually the number of nonreal zeros of $q$ in $\{|w|<t^{-1/2}\}$. So the statements of Theorem~\ref{cor:weyl} for fixed $t$ hold, with $\PP(T_*\le t)$ read as the probability that $q$ has no nonreal zero in $\{|w|<t^{-1/2}\}$. We do not know whether $q$ has nonreal zeros almost surely when only $\int\log^+|x|\,d\mu<\infty$ is assumed. For symmetric laws, \cite[Theorem~1.3]{NNS} gives $\sum_{q(w)=0}(1-|w|)=\infty$ almost surely under this assumption alone, since then $\rho_q=1$ and $\sum_j\eta_j^2=\infty$ almost surely; so it would suffice to show that the real zeros satisfy $\sum(1-|x|)<\infty$.
\end{remark}

\subsection{The limit law}

By Theorem~\ref{cor:weyl}, $\PP(T_*>t)$ is the limit of $\PP(G_{n,t}\text{ is not real-rooted})$ for $t\notin\mathcal E$. Its decay as $t\to\infty$ is governed by the probability that $q$ has a nonreal zero in a small disk about the origin, which requires the first two coefficients to be small.

\begin{theorem}[The tail of $T_*$]\label{thm:tail}
Let $\eta_0,\eta_1,\dots$ be independent. Suppose that $\eta_0$ and $\eta_1$ have densities $f_0,f_1$ that are bounded on $\R$ and continuous at $0$, that $\E\eta_2^2<\infty$, and that $\sup_{j\ge3}\E\eta_j^2<\infty$. Then $q$ is almost surely analytic in $\D$, and
\[
\lim_{r\downarrow0}\ r^{-3}\,\PP\bigl(q\text{ has a nonreal zero in }\{|w|<r\}\bigr)=\frac83\,f_0(0)f_1(0)\,\E[\eta_2^2].
\]
In particular, if the $\eta_j$ are identically distributed with a bounded density $f$ that is continuous at $0$, and $\E\eta_0^2<\infty$, then
\[
\PP(T_*>t)=\frac83\,f(0)^2\,\E[\eta_0^2]\;t^{-3/2}\bigl(1+o(1)\bigr)\qquad(t\to\infty),
\]
read as $o(t^{-3/2})$ if $f(0)=0$. For standard Gaussian coefficients the constant is $\frac4{3\pi}$, and for the uniform law on a symmetric interval it is $\frac29$.
\end{theorem}

\begin{proof}
Let $A_r$ be the event that $q$ has a nonreal zero in $\{|w|<r\}$. It is the union, over the closed disks $B$ with rational centres and radii contained in $\{|w|<r\}\setminus\R$, of the events $\{\min_B|q|=0\}$, hence measurable. Since $\sup_{j\ge2}\E|\eta_j|<\infty$, the Borel--Cantelli lemma gives $\limsup_j|\eta_j|^{1/j}\le1$ almost surely, so $q$ is analytic in $\D$. In the identically distributed case $\{T_*>t\}=A_{t^{-1/2}}$, the series $q$ having nonreal zeros by Theorem~\ref{thm:heat}; and the standard Gaussian density has $f(0)^2=1/(2\pi)$, while the uniform law on $[-b,b]$ has $f(0)^2\E\eta_0^2=(2b)^{-2}\cdot b^2/3=\frac1{12}$. So it suffices to prove the first statement.

Fix $0<r\le\frac14$ and put
\[
F_r(z)=r^{-2}q(rz)=X+Yz+\eta_2z^2+h_r(z),\qquad X=\frac{\eta_0}{r^2},\quad Y=\frac{\eta_1}{r},\quad h_r(z)=\sum_{j\ge3}\eta_jr^{j-2}z^j .
\]
Then $A_r$ is the event that $F_r$ has a nonreal zero in $\D$. Let $\omega'=(\eta_2,\eta_3,\dots)$ and $\Gamma=\frac{27}8\sum_{j\ge3}|\eta_j|(\frac38)^{j-3}$; by Minkowski's inequality $\|\Gamma\|_2\le\frac{27}5\sup_{j\ge3}\|\eta_j\|_2<\infty$. For $|z|\le\frac32$ we have $r^{j-2}|z|^j\le r\,4^{-(j-3)}\cdot\frac{27}8(\frac32)^{j-3}$, so
\begin{equation}\label{eq:hr}
\sup_{|z|\le3/2}|h_r(z)|\le r\Gamma .
\end{equation}
For fixed $\omega'$ let $D_r(\omega')$ be the set of $(x,y)\in\R^2$ such that $x+yz+\eta_2z^2+h_r(z)$ has a nonreal zero in $\D$. Since $(\eta_0,\eta_1)$ is independent of $\omega'$ and has the density $f_0(u)f_1(v)$, the substitution $(\eta_0,\eta_1)=(r^2x,ry)$ gives
\begin{equation}\label{eq:tail-cond}
\PP(A_r\mid\omega')=r^3\int_{D_r(\omega')}f_0(r^2x)f_1(ry)\,dx\,dy .
\end{equation}

\emph{A box.} Let $x,y,c\in\R$, and let $h$ be analytic near $\{|z|\le\frac32\}$, with real Taylor coefficients and $|h|\le\eps$ there. Suppose $F(z)=x+yz+cz^2+h(z)$ has a nonreal zero in $\D$. Then $F$ has at least two zeros in $\D$, a conjugate pair. If $|x|>|y|+|c|+\eps$, then $|F-x|<|x|$ on $|z|=1$, and Rouch\'e's theorem gives no zero in $\D$; so $|x|\le|y|+|c|+\eps$. On $|z|=\frac32$ we then have $|x+yz|\ge\frac32|y|-|x|\ge\frac12|y|-|c|-\eps$ and $|cz^2+h(z)|\le\frac94|c|+\eps$. If $|y|>\frac{13}2|c|+4\eps$, then $|x+yz|>|cz^2+h|$ on $|z|=\frac32$, so $F$ has at most one zero in $\{|z|<\frac32\}$, which is impossible. Hence $|y|\le\frac{13}2|c|+4\eps$ and $|x|\le\frac{15}2|c|+5\eps$. By \eqref{eq:hr}, with $\eps=r\Gamma\le\Gamma$,
\begin{equation}\label{eq:box}
D_r(\omega')\subset\mathcal B(\omega'):=\bigl[-\tfrac{15}2|\eta_2|-5\Gamma,\ \tfrac{15}2|\eta_2|+5\Gamma\bigr]\times\bigl[-\tfrac{13}2|\eta_2|-4\Gamma,\ \tfrac{13}2|\eta_2|+4\Gamma\bigr],
\end{equation}
a rectangle of area at most $240(|\eta_2|+\Gamma)^2$.

\emph{The limit region.} For $c\in\R$ let $D_0(c)$ be the set of $(x,y)$ such that $F_0(z)=x+yz+cz^2$ has a nonreal zero in $\D$. For $c\ne0$ the zeros of $F_0$ are nonreal if and only if $y^2<4cx$, and then they have modulus $\sqrt{x/c}$; so $D_0(c)=\{(x,y):0<x/c<1,\ y^2<4cx\}$, and for $c>0$ its area is $\int_0^c4\sqrt{cx}\,dx=\frac83c^2$, as also for $c<0$ by the symmetry $(x,y,c)\mapsto(-x,-y,-c)$. Also $D_0(0)=\emptyset$.

\emph{Pointwise convergence.} Fix $\omega'$ with $\Gamma<\infty$ and put $c=\eta_2$. We claim that $\mathbf 1_{D_r(\omega')}(x,y)\to\mathbf 1_{D_0(c)}(x,y)$ as $r\downarrow0$ for almost every $(x,y)\in\R^2$. Fix $(x,y)$; by \eqref{eq:hr}, $F_r\to F_0$ uniformly on $\{|z|\le\frac32\}$, where now $F_r(z)=x+yz+cz^2+h_r(z)$. Let first $c\ne0$, and exclude the null set $\{y^2=4cx\}\cup\{x=c\}$.
(i) If $(x,y)\in D_0(c)$, then $F_0$ has a simple nonreal zero $z_0\in\D$. Take a closed disk $B\ni z_0$ in $\D\setminus\R$ on whose boundary $F_0$ does not vanish and which contains no other zero of $F_0$. By Rouch\'e's theorem $F_r$ has a zero in $B$ for small $r$, so $(x,y)\in D_r(\omega')$.
(ii) If $y^2>4cx$, then $F_0$ has two simple real zeros $u_1\ne u_2$. For each $u_i$ with $|u_i|<\frac54$ take a closed disk $B_i$ centred at $u_i$, of radius less than $\frac14$, such that these disks are disjoint and contain no other zero of $F_0$, and let $K=\overline\D\setminus\bigcup_i\operatorname{int}B_i$. Then $F_0$ does not vanish on $K\cup\bigcup_i\partial B_i$, so once $\sup_{|z|\le3/2}|F_r-F_0|$ is less than the minimum of $|F_0|$ there, $F_r$ has no zero in $K$ and exactly one zero in each $B_i$, which is real because $F_r$ is real on $\R$ and $B_i$ is symmetric about $\R$. So all zeros of $F_r$ in $\overline\D$ are real, and $(x,y)\notin D_r(\omega')$.
(iii) If $y^2<4cx$ and $x/c>1$, then $F_0$ has no zero in $\overline\D$, hence neither has $F_r$ for small $r$.
These cases exhaust the complement of the null set. If $c=0$, exclude the null set $\{|x|=|y|\}$. Then $|F_0|\ge||x|-|y||>0$ on $|z|=1$, so for small $r$ the function $F_r$ has as many zeros in $\D$ as $F_0$, at most one, and a single zero of a function that is real on $\R$ is real. So $(x,y)\notin D_r(\omega')$ for small $r$, and $D_0(0)=\emptyset$.

\emph{Conclusion.} For fixed $\omega'$ with $\Gamma<\infty$, the integrand $\mathbf 1_{D_r(\omega')}(x,y)f_0(r^2x)f_1(ry)$ in \eqref{eq:tail-cond} is dominated by $\|f_0\|_\infty\|f_1\|_\infty\mathbf 1_{\mathcal B(\omega')}$, by \eqref{eq:box}, and converges almost everywhere to $f_0(0)f_1(0)\mathbf 1_{D_0(\eta_2)}$, by the continuity of $f_0,f_1$ at $0$. By dominated convergence, $r^{-3}\PP(A_r\mid\omega')\to\frac83f_0(0)f_1(0)\eta_2^2$. Moreover $r^{-3}\PP(A_r\mid\omega')\le240\|f_0\|_\infty\|f_1\|_\infty(|\eta_2|+\Gamma)^2$, which is integrable. Dominated convergence in $\omega'$ gives $r^{-3}\PP(A_r)\to\frac83f_0(0)f_1(0)\E[\eta_2^2]$.
\end{proof}

The exponent $\frac32$ has a transparent source: $(\eta_0,\eta_1)$ must lie in a region of size $r^2\times r$, and $\eta_2$ is unconstrained. Only the densities of $\eta_0,\eta_1$ at $0$ and the second moment of $\eta_2$ enter the constant.

For Gaussian coefficients there is also an exact bound, valid for all $t>1$, through the expected number of nonreal zeros.

\begin{lemma}[Expected number of pairs]\label{lem:pairs}
Let $\eta_0,\eta_1,\dots$ be independent standard Gaussian variables, and let $k_r$ be the number of zeros of $q$ in $\{|w|<r,\ \im w>0\}$. For $0\le r<1$,
\[
\E k_r=\varkappa(r):=\frac1\pi\Bigl[\frac{1+r^2}{1-r^2}\arctan r-\operatorname{artanh}r\Bigr]=\frac4{3\pi}\,r^3\bigl(1+r^2+O(r^4)\bigr)\qquad(r\to0).
\]
Consequently $\PP(T_*>t)\le\varkappa(t^{-1/2})$ for every $t>1$.
\end{lemma}

\begin{proof}
The zeros of $q$ in the upper half-disk have the first intensity
\[
\rho_c(z)=\frac{|z-\bar z|}{\pi\,|1-z^2|\,(1-|z|^2)^2}
\]
with respect to Lebesgue measure: this is the limiting intensity of the complex zeros of real Gaussian Kac polynomials computed by Shepp and Vanderbei \cite[Theorem~1.2]{SV}, and the case $n=1$ of the Pfaffian formula of Matsumoto and Shirai \cite[Theorem~2.7]{MS}. In polar coordinates $z=\varrho e^{i\theta}$ we have $|1-\varrho^2e^{2i\theta}|^2=(1+\varrho^2)^2-4\varrho^2\cos^2\theta$, and with $u=\cos\theta$,
\[
\int_0^{\pi}\frac{\sin\theta\,d\theta}{|1-\varrho^2e^{2i\theta}|}=2\int_0^1\frac{du}{\sqrt{(1+\varrho^2)^2-4\varrho^2u^2}}=\frac1\varrho\arcsin\frac{2\varrho}{1+\varrho^2}=\frac2\varrho\arctan\varrho\qquad(0<\varrho<1).
\]
Hence $\E k_r=\frac4\pi\int_0^r\varrho\arctan\varrho\,(1-\varrho^2)^{-2}\,d\varrho$. Integrating by parts with $\frac{d}{d\varrho}\frac1{2(1-\varrho^2)}=\frac{\varrho}{(1-\varrho^2)^2}$, and using $\frac1{(1-\varrho^2)(1+\varrho^2)}=\frac12\bigl(\frac1{1-\varrho^2}+\frac1{1+\varrho^2}\bigr)$,
\[
\E k_r=\frac4\pi\Bigl[\frac{\arctan r}{2(1-r^2)}-\frac{\operatorname{artanh}r+\arctan r}4\Bigr]=\varkappa(r).
\]
Multiplying the Taylor series gives $\pi\varkappa(r)=\frac43(r^3+r^5)+\frac{152}{105}(r^7+r^9)+O(r^{11})$. Finally, with $r=t^{-1/2}$, $\{T_*>t\}=\{R_*<r\}=\{k_r\ge1\}$, and Markov's inequality gives $\PP(T_*>t)\le\E k_r$.
\end{proof}

By Theorem~\ref{thm:tail}, the bound of Lemma~\ref{lem:pairs} is asymptotically sharp as $t\to\infty$. The intensity $\rho_c$ can also be obtained from the Rice formula \cite[Theorem~6.2]{AW}, applied to the Gaussian field $(\re q,\im q)$ on $\D\setminus\R$.

\begin{proof}[Proof of Proposition~\ref{prop:law}]
That $1<T_*<\infty$ almost surely was shown in the proof of Theorem~\ref{thm:heat}. For $t>1$, $\{T_*=t\}$ is contained in the event that $q$ has a zero on $\{|w|=t^{-1/2}\}$, which is null for $t\notin\mathcal E$. Together with $\PP(T_*\le1)=0$ this shows that the atoms of the law of $T_*$ lie in $\mathcal E$, which is empty if $\mu$ has no atoms (Lemma~\ref{lem:circle}). For standard Gaussian coefficients the asymptotics are Theorem~\ref{thm:tail}, and the bound is Lemma~\ref{lem:pairs}. The asymptotics give $\PP(T_*>t)>0$ for large $t$, hence for every $t$, as $\PP(T_*>t)$ is nonincreasing in $t$.
\end{proof}

\begin{remark}[The law of $T_*$ is not universal]\label{rem:nonuniversal}
The convergence $\Lambda_n\to T_*$ holds for every law as in Theorem~\ref{thm:heat}, but the law of $T_*$ depends on $\mu$, also after normalisation: $T_*$ does not change when $q$ is multiplied by a nonzero constant. If the support of $\mu$ lies in $\{x:b_1\le|x|\le b_2\}$ with $0<b_1\le b_2$, then $|q(w)|\ge b_1-b_2|w|/(1-|w|)>0$ for $|w|<b_1/(b_1+b_2)$, so $R_*\ge b_1/(b_1+b_2)$ and $T_*\le(1+b_2/b_1)^2$ almost surely. For Rademacher coefficients, $\pm1$ with probability $\frac12$ each, this gives $T_*\le4$, whereas for Gaussian coefficients $\PP(T_*>t)>0$ for every $t$. Among laws with a bounded density that is continuous at $0$, the tail constant $\frac83f(0)^2\E[\eta_0^2]$ of Theorem~\ref{thm:tail} also depends on the law: it is $\frac4{3\pi}$ for centred Gaussian coefficients and $\frac29$ for coefficients uniform on a symmetric interval.
\end{remark}

\subsection{The law of \texorpdfstring{$T_*$}{T*} near \texorpdfstring{$t=1$}{t=1}}\label{subsec:fredholm}

Throughout this subsection the $\eta_j$ are independent standard Gaussian variables. For $0<r<1$ let $\D_r=\{|w|<r\}$ and $\D_r^+=\{w\in\D_r:\im w>0\}$, and let $k_r$ be the number of zeros of $q$ in $\D_r^+$, as in Lemma~\ref{lem:pairs}; so $\PP(T_*\le t)=\PP(k_{t^{-1/2}}=0)$ for $t>1$. By \cite[Theorem~2.7]{MS} the zeros of $q$ in $\D^+=\{w\in\D:\im w>0\}$ form a Pfaffian point process: their correlation functions with respect to Lebesgue measure are
\begin{equation}\label{eq:ms-kernel}
\rho_n(z_1,\dots,z_n)=\frac1{(\pi i)^n}\prod_{j=1}^n\frac1{|1-z_j^2|}\,\operatorname{Pf}\bigl(K^c(z_i,z_j)\bigr)_{i,j=1}^n,\qquad
K^c(z,w)=\begin{pmatrix}F(z,w)&F(z,\bar w)\\F(\bar z,w)&F(\bar z,\bar w)\end{pmatrix},
\end{equation}
where $F(x,y)=(x-y)/(1-xy)^2$, $(K^c(z_i,z_j))_{i,j}$ is the antisymmetric $2n\times2n$ matrix with the $2\times2$ blocks $K^c(z_i,z_j)$, and $\operatorname{Pf}B=b_{12}$ for $n=1$. For $n=1$ this is the intensity $\rho_c$ of Lemma~\ref{lem:pairs}. For a trace class operator $T$ we write $\operatorname{tr}T$ and $\det(I+T)$ for its trace and its Fredholm determinant \cite[Section~14.5]{vN}, $\|T\|$ for its norm and $\|T\|_p$ for its Schatten $p$-norm; $2\N_0=\{0,2,4,\dots\}$.

\begin{proposition}[The law of $T_*$ as a Fredholm determinant]\label{prop:fredholm}
For $0<r<1$ and $m,n\ge0$ let $\mathcal G_{mn}(r)=\frac2\pi\int_{\D_r^+}\im(z^m\bar z^n)\,|1-z^2|^{-1}dA(z)$, and for $m,m'\in2\N_0$ let $X_{m,m'}(r)=m\,\mathcal G_{m',m-1}(r)-(m+1)\,\mathcal G_{m',m+1}(r)$, where the first term is $0$ for $m=0$. Then $X(r)$ is a trace class operator on $\ell^2(2\N_0)$, $\operatorname{tr}X(r)=\varkappa(r)$, and $\E\,s^{k_r}=\det\bigl(I-(1-s)X(r)\bigr)$ for every $s\in\C$. In particular $\PP(T_*\le t)=\det\bigl(I-X(t^{-1/2})\bigr)$ for every $t>1$.
\end{proposition}

\begin{proof}
Fix $r$ and put $\lambda=1-s$.

\emph{Step 1: a Fredholm Pfaffian.} Let $d(z)$ be a square root of $(\pi i|1-z^2|)^{-1}$ and $\mathbb K(z,w)=d(z)d(w)K^c(z,w)$. Since $\operatorname{Pf}(DBD^T)=\det D\operatorname{Pf}B$ with $D=\operatorname{diag}(d(z_1),d(z_1),\dots,d(z_n),d(z_n))$, \eqref{eq:ms-kernel} reads $\rho_n=\operatorname{Pf}(\mathbb K(z_i,z_j))$. The kernel is bounded on $\{|z|\le r\}$: there $|1-z^2|\ge1-r^2$ and $|1-zw|\ge1-r^2$, so the entries of $\mathbb K$ are at most $C_r=2r/(\pi(1-r^2)^3)$. By Hadamard's inequality and $\operatorname{Pf}(B)^2=\det B$, an antisymmetric $2n\times2n$ matrix with entries bounded by $C$ has $|\operatorname{Pf}B|\le(2n)^{n/2}C^n$. Hence $\E\binom{k_r}n=\frac1{n!}\int_{(\D_r^+)^n}\rho_n\le(\pi r^2C_r)^n(2n)^{n/2}/n!$, so $\E(1+\delta)^{k_r}=\sum_n\delta^n\E\binom{k_r}n<\infty$ for every $\delta>0$, and by dominated convergence
\begin{equation}\label{eq:fpf}
\E\,s^{k_r}=\sum_{n\ge0}\frac{(-\lambda)^n}{n!}\int_{(\D_r^+)^n}\operatorname{Pf}\bigl(\mathbb K(z_i,z_j)\bigr)_{i,j\le n}\,dA(z_1)\cdots dA(z_n)\qquad(s\in\C),
\end{equation}
the Fredholm Pfaffian $\operatorname{Pf}(J-\lambda\mathbb K)$ of \cite[(8.3)]{Ra}. In particular $\E s^{k_r}$ is entire in $s$.

\emph{Step 2: truncation and the square.} Expanding, $F(x,y)=\sum_{m,n\ge0}A_{mn}x^my^n$ with $A_{m,m-1}=m=-A_{m-1,m}$ for $m\ge1$ and $A_{mn}=0$ if $|m-n|\ne1$. So $K^c(z,w)=\Phi(z)A\Phi(w)^T$, where $\Phi(z)$ is the $2\times\infty$ matrix with rows $(z^m)_{m\ge0}$ and $(\bar z^m)_{m\ge0}$. For odd $M$ let $A^{(M)}=(A_{mn})_{m,n\le M}$, let $\Phi_M$ be the corresponding $2\times(M+1)$ matrix, $K^{(M)}(z,w)=d(z)d(w)\Phi_M(z)A^{(M)}\Phi_M(w)^T$, and let $f_M(s)$ be the series \eqref{eq:fpf} with $K^{(M)}$ in place of $\mathbb K$. Since $\sum_{m,n}|A_{mn}|r^{m+n}=2r/(1-r^2)^2$, the entries of $K^{(M)}$ are bounded by $C_r$ and $K^{(M)}\to\mathbb K$ uniformly on $(\D_r^+)^2$; so $f_M(s)\to\E s^{k_r}$ as $M\to\infty$, by dominated convergence.

The kernel $K^{(M)}$ is bounded and antisymmetric, $K^{(M)}(w,z)^T=-K^{(M)}(z,w)$. Let $J=\bigl(\begin{smallmatrix}0&1\\-1&0\end{smallmatrix}\bigr)$. The identity $\operatorname{Pf}(J+K)^2=\det(I+J^{-1}K)$ of \cite[Lemma~8.1, (8.6)]{Ra}, with $K=-\lambda K^{(M)}$ and $J^{-1}=-J$, gives
\begin{equation}\label{eq:rains}
f_M(s)^2=\sum_{n\ge0}\frac{\lambda^n}{n!}\int_{(\D_r^+)^n}\sum_{a\in\{1,2\}^n}\det\Bigl(\bigl(JK^{(M)}(z_i,z_j)\bigr)_{a_ia_j}\Bigr)_{i,j\le n}\,dA(z_1)\cdots dA(z_n),
\end{equation}
the Fredholm determinant of the kernel $\lambda JK^{(M)}$ on $\D_r^+\times\{1,2\}$. Rains states (8.6) without analytic hypotheses. For a bounded antisymmetric $2\times2$ kernel $K$ on a space of finite measure both sides are absolutely convergent series, by Hadamard's inequality, and (8.6) follows by integrating over the $n$-th power of the space the pointwise identity
\[
\sum_{S\sqcup S'=\{1,\dots,n\}}\operatorname{Pf}K(S)\operatorname{Pf}K(S')=\sum_{a\in\{1,2\}^n}\det\Bigl(\bigl(J^{-1}K(y_i,y_j)\bigr)_{a_ia_j}\Bigr)_{i,j\le n},
\]
where $K(S)=(K(y_i,y_j))_{i,j\in S}$. This identity is the coefficient of $w_1\cdots w_n$ in $\operatorname{Pf}(J_n+N)^2=\det(J_n+N)=\det(I+J_n^{-1}N)$, where $J_n=\operatorname{diag}(J,\dots,J)$ and $N=(\sqrt{w_iw_j}\,K(y_i,y_j))_{i,j\le n}$, whose Pfaffian is $\sum_{S}\prod_{i\in S}w_i\operatorname{Pf}K(S)$. Now $JK^{(M)}(z,w)=u(z)v(w)$ with $u(z)=d(z)J\Phi_M(z)$ and $v(w)=d(w)A^{(M)}\Phi_M(w)^T$, and $d(w)^2(\Phi_M(w)^TJ\Phi_M(w))_{mn}=(w^m\bar w^n-\bar w^mw^n)/(\pi i|1-w^2|)$. By the Cauchy--Binet formula and Andr\'eief's identity, \eqref{eq:rains} equals $\det(I+\lambda\int_{\D_r^+}v\,u\,dA)=\det(I+\lambda A^{(M)}\mathcal G^{(M)})$, where $\mathcal G^{(M)}=(\mathcal G_{mn}(r))_{m,n\le M}$.

\emph{Step 3: parity.} The map $z\mapsto-\bar z$ preserves $\D_r^+$, $dA$ and $|1-z^2|$, and $\im((-\bar z)^m(-z)^n)=-(-1)^{m+n}\im(z^m\bar z^n)$; so $\mathcal G_{mn}=0$ when $m+n$ is even, and so is $A_{mn}$. Order the indices $0,\dots,M$ as (even, odd), two groups of $(M+1)/2$. Then $A^{(M)}=\bigl(\begin{smallmatrix}0&A_{eo}\\-A_{eo}^T&0\end{smallmatrix}\bigr)$ and $\mathcal G^{(M)}=\bigl(\begin{smallmatrix}0&\mathcal G_{eo}\\-\mathcal G_{eo}^T&0\end{smallmatrix}\bigr)$, so $A^{(M)}\mathcal G^{(M)}=-\operatorname{diag}(A_{eo}\mathcal G_{eo}^T,A_{eo}^T\mathcal G_{eo})$ and, by $\det(I-BC)=\det(I-CB)$ and transposition, $f_M(s)^2=\det(I-\lambda X^{(M)})^2$ with $X^{(M)}=A_{eo}\mathcal G_{eo}^T$. Here $X^{(M)}_{m,m'}=\sum_nA_{mn}\mathcal G_{m'n}=m\mathcal G_{m',m-1}-(m+1)\mathcal G_{m',m+1}$ for even $m,m'\le M-1$, since both neighbours $m\pm1$ of $m$ are at most $M$: $X^{(M)}$ is the compression of $X(r)$ to the even indices at most $M-1$. Since $|\mathcal G_{mn}|\le r^{m+n+2}/(1-r^2)$, we have $|X_{m,m'}|\le(2m+1)r^{m+m'+1}/(1-r^2)$ and $\sum_{m,m'}|X_{m,m'}|<\infty$. A matrix unit has trace norm $1$, so $X(r)$ is trace class, $X^{(M)}\to X(r)$ in trace norm, and $\det(I-\lambda X^{(M)})\to\det(I-\lambda X(r))$ \cite[Lemmas~14.38 and~14.39]{vN}. Hence $f^2=g^2$ for the entire functions $f(s)=\E s^{k_r}$ and $g(s)=\det(I-(1-s)X(r))$. Since $f(1)=g(1)=1$, the factor $f+g$ of $(f-g)(f+g)\equiv0$ is not identically zero, so $f=g$: the sign is fixed by the value $1$ at $s=1$. The coefficients of $s-1$ give $\operatorname{tr}X(r)=\E k_r=\varkappa(r)$, and $s=0$ gives $\PP(k_r=0)=\det(I-X(r))$.
\end{proof}

\begin{remark}[Numerical evaluation]\label{rem:fredholm-num}
In polar coordinates $\mathcal G_{mn}(r)=\frac2\pi\int_0^r\varrho^{m+n+1}S_{m-n}(\varrho)\,d\varrho$ with $S_k(\varrho)=\int_0^\pi\sin(k\theta)\,|1-\varrho^2e^{2i\theta}|^{-1}d\theta$. Here $S_k=0$ for even $k$, $S_{-k}=-S_k$, $S_1(\varrho)=2\varrho^{-1}\arctan\varrho$ (proof of Lemma~\ref{lem:pairs}), and for odd $k\ge1$, with $x=\varrho^2$,
\[
(k+1)x\,S_{k+2}-k(1+x^2)\,S_k+(k-1)x\,S_{k-2}=-2(1-x),
\]
as one finds by integrating $D\,g'+2x\sin(2\theta)\,g=0$, where $D=|1-xe^{2i\theta}|^2$ and $g=D^{-1/2}$, against $\cos k\theta$ over $[0,\pi]$. Solved as a boundary value problem (forward recursion is unstable), with Gauss--Legendre quadrature in $\varrho$ and truncation at $m\le14/(1-r)$, this evaluates $-\log\PP(T_*\le t)$ to about $10^{-7}$ (absolute; tested by a $25\%$ change of the truncation) for $t-1\ge0.0026$. For example, numerically $\PP(T_*\le100/81)=0.10868062$ ($r=0.9$) and $-\log\PP(T_*\le10000/9801)=36.180955$ ($r=0.99$). The law of $k_r$ given by $\det(I-(1-s)X(r))$ agrees with Monte Carlo simulations of $q$ for $0.5\le r\le0.9$.
\end{remark}

\begin{theorem}[The hole probability from above]\label{thm:sharp}
For standard Gaussian coefficients,
\[
\limsup_{t\downarrow1}\,(t-1)\log\PP(T_*\le t)\le-\frac{\pi^2}{12}.
\]
That is, for every $c<\pi^2/12$ there is $t_0>1$ such that $\PP(T_*\le t)\le e^{-c/(t-1)}$ for $1<t\le t_0$.
\end{theorem}

The constant is half of the constant $\pi^2/6$ for the Gaussian power series with complex coefficients \cite{PV}, and it is attained, by Theorem~\ref{thm:lower} below. The factor $\frac12$ arises in the upper bound because the kernel of the determinant is, up to an error of Hilbert--Schmidt norm $O(\log\frac1{t-1})$, the reproducing kernel of the even functions in the Bergman space of $\D$ (Lemma~\ref{lem:even-bergman}), and in the lower bound because the symmetric part of $X(r)$ is at most $\operatorname{diag}(r^{2m+1})_{m\in2\N_0}$ (Lemma~\ref{lem:qform} and Step~1 of the proof of Theorem~\ref{thm:lower}).

\emph{The operators.} Fix $r\in(0,1)$ and put $\varsigma=\log t=-2\log r$ and $\ell_r=\log\frac1{1-r}$. For $m\in2\N_0$ let $\psi_m(z)=mz^{m-1}-(m+1)z^{m+1}$, and define $U:\ell^2(2\N_0)\to L^2(\D_r)$ and $V:L^2(\D_r)\to\ell^2(2\N_0)$, with Lebesgue measure on $\D_r$, by
\[
(Uc)(z)=\sum_{m\in2\N_0}c_mz^m,\qquad(Vf)_m=\int_{\D_r}v_m(z)f(z)\,dA(z),\qquad v_m(z)=\frac{\operatorname{sgn}(\im z)\,\overline{\psi_m(z)}}{\pi i\,|1-z^2|}.
\]
As $|1-z^2|\ge1-r^2$ on $\D_r$ and $\sum_m(2m+1)^2r^{2m}<\infty$, both are Hilbert--Schmidt. Splitting $\D_r$ into $\D_r^+$ and its conjugate gives $\mathcal G_{mn}(r)=\frac1{\pi i}\int_{\D_r}\operatorname{sgn}(\im z)\,z^m\bar z^n|1-z^2|^{-1}dA(z)$, hence $X(r)=VU$. The operator $\mathcal L=UV$ on $L^2(\D_r)$ has the kernel $\sum_mz^mv_m(w)$; summing with $\sum_{m\in2\N_0}mx^m=2x^2(1-x^2)^{-2}$ and $\sum_{m\in2\N_0}(m+1)x^m=(1+x^2)(1-x^2)^{-2}$, $x=z\bar w$,
\begin{equation}\label{eq:Lkernel}
\mathcal L(z,w)=\frac{\operatorname{sgn}(\im w)}{2\pi i\,|1-w^2|}\Bigl[\frac{z-\bar w}{(1-z\bar w)^2}-\frac{z+\bar w}{(1+z\bar w)^2}\Bigr]\qquad(z,w\in\D_r).
\end{equation}
$X(r)$ and $\mathcal L$ are trace class, as products of two Hilbert--Schmidt operators. For $\mu\ne0$ and $k\ge1$, $U$ maps $\ker(VU-\mu)^k$ into $\ker(UV-\mu)^k$, as $(UV-\mu)^kU=U(VU-\mu)^k$, and injectively, as $Uc=0$ and $(VU-\mu)^kc=0$ give $(-\mu)^kc=0$; likewise $V$ maps back injectively. So $X(r)$ and $\mathcal L$ have the same nonzero eigenvalues $\lambda_k$, repeated according to algebraic multiplicity, and $\sum_k|\lambda_k|<\infty$ \cite[Proposition~14.21]{vN}. We use three classical facts.

(i) $\PP(T_*\le t)=\det(I-X(r))=\prod_k(1-\lambda_k)$, by Proposition~\ref{prop:fredholm} and \cite[Theorem~14.44]{vN}.

(ii) $\operatorname{tr}\mathcal L^j=\sum_k\lambda_k^j$ for $j\ge1$, and $\operatorname{tr}\mathcal L=\operatorname{tr}X(r)=\varkappa(r)$. For $j=1$ this is Lidskii's theorem \cite[Theorem~14.34]{vN}, applied to $\mathcal L$ and to $X(r)$. For $j\ge2$, from $I-\nu^j\mathcal L^j=\prod_{\omega^j=1}(I-\omega\nu\mathcal L)$, the multiplicativity of $\det$ \cite[Lemma~14.40]{vN} and \cite[Theorem~14.44]{vN}, $\det(I-\nu^j\mathcal L^j)=\prod_k(1-\nu^j\lambda_k^j)$ for $\nu\in\C$; compare the coefficients of $\nu^j$.

(iii) For $p\ge2$,
\begin{equation}\label{eq:weyl-ineq}
\textstyle\sum_k|\lambda_k|^p\le\|\mathcal L\|_p^p .
\end{equation}
This is Weyl's inequality between the eigenvalues and the singular values \cite{We}; for compact operators see \cite[Chapter~1]{Si}.
% TODO(author): pinpoint in [Si] for Weyl's inequality sum |lambda_k|^p <= ||A||_p^p, A compact, not checked; no
% referee of idea 106 has seen the book. It is believed to be Theorem 1.15 of the 2005 edition, in Chapter 1
% ("Preliminaries"). Check it and cite the theorem number, then delete this comment.

The comparison operator $\mathcal K$ is the integral operator on $L^2(\D_r)$ with the kernel
\[
K_{\rm ev}(z,w)=\frac1{2\pi}\Bigl[\frac1{(1-z\bar w)^2}+\frac1{(1+z\bar w)^2}\Bigr]=\frac1\pi\sum_{m\in2\N_0}(m+1)z^m\bar w^m .
\]
The $z^m$ are orthogonal in $L^2(\D_r)$, with $\|z^m\|^2=\pi r^{2m+2}/(m+1)$, so $\mathcal Kz^m=r^{2m+2}z^m$ for even $m$, and $\mathcal K$ vanishes on the orthogonal complement of these functions. So $\mathcal K$ is self-adjoint, $\|\mathcal K\|\le1$, its nonzero eigenvalues are $e^{-(2k+1)\varsigma}$, $k\ge0$, and
\begin{equation}\label{eq:traces}
\operatorname{tr}\mathcal K^j=\|\mathcal K\|_j^j=\textstyle\sum_{k\ge0}e^{-(2k+1)j\varsigma}=1/(2\sinh(j\varsigma))\qquad(j\ge1).
\end{equation}

\begin{lemma}[The even Bergman kernel]\label{lem:even-bergman}
$\|\mathcal L-\mathcal K\|_2\le1+\ell_r$.
\end{lemma}

\begin{proof}
Let $\mathcal R=\mathcal L-\mathcal K$, with kernel $R=\mathcal L-K_{\rm ev}$. As $\mathcal L(\bar z,\bar w)=\overline{\mathcal L(z,w)}$ and $K_{\rm ev}(\bar z,\bar w)=\overline{K_{\rm ev}(z,w)}$, $\|\mathcal R\|_2^2=2\int_{\D_r^+}\int_{\D_r}|R(z,w)|^2\,dA(z)\,dA(w)$. Let $w\in\D_r^+$ and $\tau=2\im w/|1-w^2|$. From $-i(z-\bar w)=-i(z-w)+2\im w$ and $i(z+\bar w)=i(z+w)+2\im w$,
\[
R(z,w)=\frac1{2\pi}\Bigl[\Bigl(\frac{-i(z-w)}{|1-w^2|}-(1-\tau)\Bigr)\frac1{(1-z\bar w)^2}+\Bigl(\frac{i(z+w)}{|1-w^2|}-(1-\tau)\Bigr)\frac1{(1+z\bar w)^2}\Bigr].
\]
Here $0\le\tau\le1$, as $|1-w^2|^2=(1-|w|^2)^2+4(\im w)^2$, and $|z\mp w|\le|1\mp z\bar w|$, as $|1-z\bar w|^2-|z-w|^2=(1-|z|^2)(1-|w|^2)$. Expanding in powers of $z\bar w$ and using the orthogonality of the $z^k$, $\int_{\D_r}|1\mp z\bar w|^{-2}dA(z)=\sum_{k\ge0}\pi r^{2k+2}|w|^{2k}/(k+1)\le\pi r\sum_{k\ge0}r^{k+1}/(k+1)\le\pi\ell_r$ and $\int_{\D_r}|1\mp z\bar w|^{-4}dA(z)\le\sum_{k\ge0}\pi(k+1)|w|^{2k}=\pi(1-|w|^2)^{-2}$. With $(a_1+\dots+a_4)^2\le4(a_1^2+\dots+a_4^2)$,
\[
\int_{\D_r}|R(z,w)|^2\,dA(z)\le\frac2\pi\Bigl[\frac{\ell_r}{|1-w^2|^2}+\frac{(1-\tau)^2}{(1-|w|^2)^2}\Bigr].
\]
To integrate over $w$, substitute $w=\tanh(\zeta/2)$, $\zeta=u+iv$, which maps the strip $\{0<v<\frac\pi2\}$ conformally onto $\D^+$. Then $w=(\sinh u+i\sin v)/(\cosh u+\cos v)$, $|dw/d\zeta|=(\cosh u+\cos v)^{-1}$, $|1-w^2|=2(\cosh u+\cos v)^{-1}$, $1-|w|^2=2\cos v\,(\cosh u+\cos v)^{-1}$ and $\tau=\sin v$. Also $|w|^2=(\cosh u-\cos v)/(\cosh u+\cos v)\ge\tanh^2(u/2)$, so $w\in\D_r^+$ forces $|u|<2\operatorname{artanh}r$. Hence $\int_{\D_r^+}|1-w^2|^{-2}dA(w)\le\int_{|u|<2\operatorname{artanh}r}\int_0^{\pi/2}\frac14\,dv\,du=\frac\pi2\operatorname{artanh}r$ and
\[
\int_{\D_r^+}\frac{(1-\tau)^2\,dA(w)}{(1-|w|^2)^2}\le\int_{|u|<2\operatorname{artanh}r}\int_0^{\pi/2}\frac{(1-\sin v)^2}{4\cos^2v}\,dv\,du=\Bigl(2-\frac\pi2\Bigr)\operatorname{artanh}r .
\]
So $\|\mathcal R\|_2^2\le2\operatorname{artanh}r\,(\ell_r+\frac{4-\pi}\pi)\le(\log2+\ell_r)(\ell_r+\frac{4-\pi}\pi)\le(1+\ell_r)^2$.
\end{proof}

\begin{lemma}\label{lem:log-ineq}
Let $p\ge2$ be an integer and $C_p=\max(2e,1+H_{p-1})$, where $H_{p-1}=\sum_{j<p}\frac1j$. For every $\lambda\in\C$, $\log|1-\lambda|\le-\re\sum_{j=1}^{p-1}\lambda^j/j+C_p|\lambda|^p$.
\end{lemma}

\begin{proof}
Let $h(\lambda)$ be the left side minus the sum. If $|\lambda|\ge1$, then $h\le\log(1+|\lambda|)+\sum_{j<p}|\lambda|^j/j\le|\lambda|+H_{p-1}|\lambda|^{p-1}\le(1+H_{p-1})|\lambda|^p$. If $|\lambda|<1$ and $|1-\lambda|<1/(ep)$, then $h<-1-\log p+H_{p-1}\le0$, as $H_{p-1}\le1+\log p$. If $|\lambda|<1$ and $|1-\lambda|\ge1/(ep)$, then $h=-\re\sum_{j\ge p}\lambda^j/j$. With $\sigma_j=\sum_{i=p}^j\lambda^i$ we have $|\sigma_j|\le2|\lambda|^p/|1-\lambda|$, and summation by parts gives $|\sum_{j\ge p}\lambda^j/j|=|\sum_{j\ge p}\sigma_j(\frac1j-\frac1{j+1})|\le2|\lambda|^p/(p|1-\lambda|)\le2e|\lambda|^p$.
\end{proof}

\begin{proof}[Proof of Theorem~\ref{thm:sharp}]
Let $t>1$ and $r=t^{-1/2}$, so $r^2=e^{-\varsigma}$. By Lemma~\ref{lem:even-bergman} and \eqref{eq:traces}, $\|\mathcal R\|_2\le1+\ell_r$ and $\|\mathcal K\|_2^2=1/(2\sinh2\varsigma)\le1/(4\varsigma)$; for $0<\varsigma\le1$, $1-r\ge\varsigma/4$, so $\ell_r\le\log(4/\varsigma)$.

\emph{Traces.} For $j=1$, $\operatorname{tr}\mathcal L=\varkappa(r)$ by (ii), and the closed form of $\varkappa$ in Lemma~\ref{lem:pairs} gives $\varkappa(r)=\frac1{2(1-r^2)}+O(\ell_r)$, so $\varsigma\operatorname{tr}\mathcal L\to\frac12$ as $\varsigma\to0$. For $j\ge2$, expand $\operatorname{tr}(\mathcal K+\mathcal R)^j$ into $2^j$ words. By cyclicity a word with one factor $\mathcal R$ has trace $\operatorname{tr}(\mathcal K^{j-1}\mathcal R)$, of modulus at most $\|\mathcal K\|_2\|\mathcal R\|_2$ as $\|\mathcal K\|\le1$; a word with at least two factors $\mathcal R$ has trace norm at most $\|\mathcal R\|_2^2\max(1,\|\mathcal R\|)^{j-2}\le(1+\|\mathcal R\|_2)^j$. With \eqref{eq:traces}, $|\operatorname{tr}\mathcal L^j-1/(2\sinh(j\varsigma))|\le j(1+\ell_r)/(2\sqrt\varsigma)+(4+2\ell_r)^j$. So $\varsigma\operatorname{tr}\mathcal L^j\to\frac1{2j}$ as $\varsigma\to0$, for every $j\ge1$.

\emph{Schatten norms.} Let $p\ge2$. By \eqref{eq:weyl-ineq}, the triangle inequality for $\|\cdot\|_p$, $\|\mathcal R\|_p\le\|\mathcal R\|_2$ and \eqref{eq:traces},
\[
\varsigma\sum_k|\lambda_k|^p\le\varsigma\bigl(\|\mathcal K\|_p+\|\mathcal R\|_2\bigr)^p=\frac{\varsigma}{2\sinh(p\varsigma)}\Bigl(1+\|\mathcal R\|_2\bigl(2\sinh(p\varsigma)\bigr)^{1/p}\Bigr)^p\to\frac1{2p}\qquad(\varsigma\to0).
\]

\emph{Conclusion.} By (i), Lemma~\ref{lem:log-ineq} and (ii), for every integer $p\ge2$,
\[
-\log\PP(T_*\le t)=-\sum_k\log|1-\lambda_k|\ge\sum_{j=1}^{p-1}\frac{\re\operatorname{tr}\mathcal L^j}j-C_p\sum_k|\lambda_k|^p .
\]
Multiplying by $\varsigma$ and letting $t\downarrow1$, $\liminf_{t\downarrow1}\varsigma\,(-\log\PP(T_*\le t))\ge\sum_{j=1}^{p-1}\frac1{2j^2}-\frac{C_p}{2p}$. As $C_p=O(\log p)$, letting $p\to\infty$ gives $\frac12\zeta(2)=\frac{\pi^2}{12}$ on the right. Finally $0<\varsigma=\log t\le t-1$ and $-\log\PP(T_*\le t)\ge0$.
\end{proof}

\begin{remark}[An explicit bound]\label{rem:explicit}
For $p=2$, $\log|1-\lambda|=\frac12\log(1-2\re\lambda+|\lambda|^2)\le-\re\lambda+\frac12|\lambda|^2$, and the proof gives, for every $t>1$ and $r=t^{-1/2}$, $\PP(T_*\le t)\le\exp(-\varkappa(r)+\frac12\|\mathcal L\|_2^2)$. Here $\|\mathcal L\|_2^2=\operatorname{tr}(U^*U\,VV^*)=\sum_{m\in2\N_0}\|z^m\|^2\|v_m\|^2$ is a series of one-dimensional integrals. Numerically, $(t-1)(\varkappa(r)-\frac12\|\mathcal L\|_2^2)$ increases with $r$ on a grid of step $0.0025$ in $r\in[0.8,0.9995]$, from $0.2796$ at $r=0.8$; so, numerically, $\PP(T_*\le t)\le e^{-0.279/(t-1)}$ for $1.001\le t\le1.5625$. In the other direction, Theorem~\ref{thm:lower} below gives $\PP(T_*\le t)\ge\prod_{k\ge0}(1-t^{-2k-1/2})$ for every $t>1$.
\end{remark}

\emph{The lower bound.} The lower bound that matches Theorem~\ref{thm:sharp} comes from a comparison of the symmetric part of $X(r)$ with a diagonal operator. It rests on an identity for the quadratic form of $X(r)$, which follows from Green's formula.

\begin{lemma}[The quadratic form of $X(r)$]\label{lem:qform}
Let $0<r<1$ and $\sigma_r(\theta)=(1+r^2)\sin\theta\,|1-r^2e^{2i\theta}|^{-1}$ for $0\le\theta\le\pi$. Then $0\le\sigma_r\le1$, and for every real $c\in\ell^2(2\N_0)$, with $f(z)=\sum_{m\in2\N_0}c_mz^m$,
\begin{equation}\label{eq:qform}
\sum_{m,m'\in2\N_0}c_m\,X_{m,m'}(r)\,c_{m'}=\frac r\pi\int_0^\pi|f(re^{i\theta})|^2\sigma_r(\theta)\,d\theta-\frac1\pi\int_{-r}^rf(x)^2\,dx .
\end{equation}
\end{lemma}

\begin{proof}
Both sides of \eqref{eq:qform} are continuous in $c\in\ell^2(2\N_0)$: the left side because $X(r)$ is bounded, the right side because $f$ depends linearly on $c$ and $\sup_{|z|\le r}|f(z)|\le\sum_m|c_m|r^m\le\|c\|(1-r^2)^{-1/2}$, by the Cauchy--Schwarz inequality. So it suffices to take $c$ finitely supported; then $f$ is an even polynomial with real coefficients.

Put $\psi=\sum_mc_m\psi_m$. As $\psi_m$ has real coefficients, $m\,\im(z^{m'}\bar z^{m-1})-(m+1)\im(z^{m'}\bar z^{m+1})=\im(z^{m'}\overline{\psi_m(z)})$, so the definitions of $X_{m,m'}(r)$ and $\mathcal G_{mn}(r)$ give
\[
\sum_{m,m'}c_m\,X_{m,m'}(r)\,c_{m'}=\frac2\pi\int_{\D_r^+}\frac{\im\bigl(f(z)\overline{\psi(z)}\bigr)}{|1-z^2|}\,dA(z).
\]
Let $s(z)=\sqrt{1-z^2}$ be the principal branch. It is holomorphic and has no zero on $\D$, since $\re(1-z^2)=1-x^2+y^2>0$ there, and $s>0$ on $(-1,1)$. For holomorphic $g$ we have $(1-z^2)g'-zg=s\,(sg)'$, as $ss'=-z$; with $g=z^m$ the left side is $\psi_m$, so $\psi=s\,u'$ with $u=sf$. Put $v=f/s$. As $|1-z^2|=s\bar s$ and $v$ is holomorphic,
\[
\frac{f\,\overline\psi}{|1-z^2|}=\frac fs\,\overline{u'}=v\,\overline{u'}=\partial_{\bar z}(v\bar u).
\]
The function $v\bar u$ is smooth on $\D$, which contains the closure of $\D_r^+$. Green's formula $\int_\Omega\partial_{\bar z}F\,dA=\frac1{2i}\oint_{\partial\Omega}F\,dz$, with $\partial\Omega$ oriented counterclockwise, therefore applies to $F=v\bar u$ and $\Omega=\D_r^+$ (the corners at $\pm r$ are harmless). The boundary consists of the segment from $-r$ to $r$ and the arc $z=re^{i\theta}$, $0\le\theta\le\pi$, so the left side of \eqref{eq:qform} equals
\[
\frac2\pi\im\Bigl(\frac1{2i}\int_{-r}^r(v\bar u)(x)\,dx+\frac1{2i}\int_0^\pi(v\bar u)(re^{i\theta})\,ire^{i\theta}\,d\theta\Bigr).
\]
On $[-r,r]$, $s>0$ and $f$ is real, so $v\bar u=f^2$, and the first term contributes $-\frac1\pi\int_{-r}^rf(x)^2\,dx$. On the arc, $v\bar u=|f|^2\,\bar s/s$, so the second term contributes $\frac r\pi\int_0^\pi|f(re^{i\theta})|^2\,\im\bigl(e^{i\theta}\bar s/s\bigr)\,d\theta$. Finally, at $z=re^{i\theta}$ we have $\bar s^2=1-r^2e^{-2i\theta}$ and $|s|^2=|1-r^2e^{2i\theta}|$, so
\[
e^{i\theta}\,\frac{\bar s}{s}=\frac{e^{i\theta}\,\bar s^2}{|s|^2}=\frac{e^{i\theta}-r^2e^{-i\theta}}{|1-r^2e^{2i\theta}|}=\frac{(1-r^2)\cos\theta+i\,(1+r^2)\sin\theta}{|1-r^2e^{2i\theta}|}.
\]
This number has modulus $1$, and its imaginary part is $\sigma_r(\theta)$, which proves \eqref{eq:qform}; since $\sin\theta\ge0$ for $0\le\theta\le\pi$, it also gives $0\le\sigma_r\le1$.
\end{proof}

\begin{theorem}[The hole probability from below]\label{thm:lower}
For standard Gaussian coefficients and every $t>1$,
\[
\PP(T_*\le t)\ \ge\ \prod_{k\ge0}\bigl(1-t^{-2k-1/2}\bigr)\ \ge\ \bigl(1-t^{-1/2}\bigr)\exp\Bigl(-\frac{\pi^2}{12\log t}\Bigr).
\]
Consequently $\lim_{t\downarrow1}(t-1)\log\PP(T_*\le t)=-\pi^2/12$.
\end{theorem}

\begin{proof}
Let $r=t^{-1/2}$, so that $r^{4k+1}=t^{-2k-1/2}=e^{-(2k+1/2)\varsigma}$ with $\varsigma=\log t$. For $n\ge1$ let $X_n$ be the compression of $X(r)$ to the indices $m\in\{0,2,\dots,2n-2\}$, a real $n\times n$ matrix, and put
\[
S_n=\tfrac12\bigl(X_n+X_n^T\bigr),\qquad B_n=\tfrac12\bigl(X_n^T-X_n\bigr),\qquad R_n=\operatorname{diag}\bigl(r^{2m+1}\bigr)_{m\in\{0,2,\dots,2n-2\}},
\]
so that the entries of $R_n$ are $r^{4k+1}$, $0\le k<n$.
By Step~3 of the proof of Proposition~\ref{prop:fredholm} ($X_n$ is the matrix $X^{(M)}$ of that step with $M=2n-1$), $X_n\to X(r)$ in trace norm and $\det(I-X_n)\to\det(I-X(r))=\PP(T_*\le t)$ as $n\to\infty$.

\emph{Step 1: $S_n\le R_n$.} Let $c\in\R^n$, indexed by $m\in\{0,2,\dots,2n-2\}$, and $f(z)=\sum_mc_mz^m$. For real $c$ we have $c^TS_nc=c^TX_nc$, which is the left side of \eqref{eq:qform} for $c$ extended by zeros. In \eqref{eq:qform} we have $\sigma_r\le1$ (Lemma~\ref{lem:qform}), the last term is not positive, and $\frac1\pi\int_0^\pi e^{i(m-m')\theta}\,d\theta=\delta_{mm'}$ for even $m-m'$; hence
\[
c^TS_nc\le\frac r\pi\int_0^\pi|f(re^{i\theta})|^2\,d\theta=\sum_mr^{2m+1}c_m^2=c^TR_nc .
\]
As $S_n$ and $R_n$ are real symmetric, $S_n\le R_n$. In particular $H_n:=I-S_n\ge I-R_n\ge(1-r)I$ is positive definite.

\emph{Step 2: $\det(I-X_n)\ge\det H_n$.} This is the inequality of Ostrowski and Taussky \cite{OT}, $|\det M|\ge\det\frac12(M+M^*)$ when $\frac12(M+M^*)$ is positive definite, in the case we need; we include the short proof. Since $X_n=S_n-B_n$, we have $I-X_n=H_n+B_n=H_n^{1/2}(I+C_n)H_n^{1/2}$ with $C_n=H_n^{-1/2}B_nH_n^{-1/2}$, so $\det(I-X_n)=\det H_n\cdot\det(I+C_n)$. The matrix $C_n$ is real and antisymmetric, so $iC_n$ is Hermitian and the eigenvalues of $C_n$ are purely imaginary. The characteristic polynomial of $C_n$ has real coefficients, so its nonzero eigenvalues, repeated according to multiplicity, come in pairs $\pm i\kappa_j$ with $\kappa_j>0$. Hence $\det(I+C_n)=\prod_j(1+i\kappa_j)(1-i\kappa_j)=\prod_j(1+\kappa_j^2)\ge1$.

\emph{Step 3: $\det H_n\ge\det(I-R_n)$.} Put $P=I-R_n$. Then $0<P\le H_n$, so $P^{-1/2}H_nP^{-1/2}\ge I$ has all its eigenvalues at least $1$, and $\det H_n=\det P\cdot\det\bigl(P^{-1/2}H_nP^{-1/2}\bigr)\ge\det P$.

\emph{Step 4: the bounds.} By Steps~2 and~3, $\det(I-X_n)\ge\det(I-R_n)=\prod_{k<n}(1-r^{4k+1})\ge\prod_{k\ge0}(1-r^{4k+1})$, since every factor lies in $(0,1)$. Letting $n\to\infty$ gives the first inequality. For the second, let $\phi(u)=-\log(1-e^{-u})$ for $u>0$. Then $\phi$ is positive and decreasing, and $\int_0^\infty\phi(u)\,du=\sum_{j\ge1}\int_0^\infty\frac{e^{-ju}}j\,du=\sum_{j\ge1}\frac1{j^2}=\frac{\pi^2}6$. As $\phi$ is decreasing, $\phi\bigl((2k+\frac12)\varsigma\bigr)\le\frac1{2\varsigma}\int_{(2k-3/2)\varsigma}^{(2k+1/2)\varsigma}\phi(u)\,du$ for $k\ge1$, and these intervals are disjoint subsets of $(0,\infty)$. Hence
\[
-\log\prod_{k\ge1}\bigl(1-t^{-2k-1/2}\bigr)=\sum_{k\ge1}\phi\bigl((2k+\tfrac12)\varsigma\bigr)\le\frac1{2\varsigma}\int_0^\infty\phi(u)\,du=\frac{\pi^2}{12\log t},
\]
and the factor with $k=0$ is $1-t^{-1/2}$.

\emph{The limit.} By the bounds, $(t-1)\log\PP(T_*\le t)\ge(t-1)\log(1-t^{-1/2})-\frac{\pi^2}{12}\cdot\frac{t-1}{\log t}$. As $t\downarrow1$, $(t-1)/\log t\to1$, and $(t-1)\log(1-t^{-1/2})\to0$ because $1-t^{-1/2}\sim\frac12(t-1)$. So $\liminf_{t\downarrow1}(t-1)\log\PP(T_*\le t)\ge-\pi^2/12$, and Theorem~\ref{thm:sharp} gives the reverse inequality for the $\limsup$.
\end{proof}

\begin{remark}\label{rem:lower}
(a) For the comparison operator $R=\operatorname{diag}(r^{2m+1})_{m\in2\N_0}$, $\det(I-R)=\prod_{k\ge0}(1-r^{4k+1})=\exp\bigl(-(\pi^2/12+o(1))/(t-1)\bigr)$ as $t\downarrow1$, half the exponent of the product $\prod_{k\ge1}(1-r^{2k})$ that gives the hole probability for the Gaussian power series with complex coefficients \cite{PV}. A halving of the same kind occurs for the gap probabilities of non-Hermitian random matrices: at maximal non-Hermiticity, the probability that a disk about the origin contains no eigenvalue of the quaternion ($\beta=4$) Ginibre ensemble is obtained from the product for the complex ($\beta=2$) ensemble by skipping every second factor \cite{APS}.

(b) The proof gives more. The compressions $S_n$ of $S=\frac12(X(r)+X(r)^T)$ converge to $S$ in trace norm, as $\sum_{m,m'}|X_{m,m'}(r)|<\infty$ (proof of Proposition~\ref{prop:fredholm}), so Steps~1--2 give $\PP(T_*\le t)\ge\det(I-S)$. This bound is much closer: numerically, at $r=0.99$, $\log\PP(T_*\le t)=-36.181$ and $\log\det(I-S)=-36.269$, against $\sum_{k\ge0}\log(1-r^{4k+1})=-42.090$. The bound of Theorem~\ref{thm:lower} is sharp in the exponent but not for moderate $t$: at $t=100/81$ ($r=0.9$) it gives $0.0112$, against $\PP(T_*\le t)=0.1087$ (Remark~\ref{rem:fredholm-num}).

(c) For $0\le s\le1$ the same argument applies to $(1-s)X(r)$: Step~1 gives $(1-s)S_n\le(1-s)R_n$, and $I-(1-s)R_n$ is positive definite. With Proposition~\ref{prop:fredholm} this gives $\E s^{k_r}=\det\bigl(I-(1-s)X(r)\bigr)\ge\prod_{k\ge0}\bigl(1-(1-s)r^{4k+1}\bigr)$. So the generating function of $k_r$ dominates, on $[0,1]$, that of a sum of independent Bernoulli variables with parameters $r^{4k+1}$, $k\ge0$; whether $k_r$ itself is such a sum is Question~\ref{q:bernoulli}.
\end{remark}



\section{The standard coupling}\label{sec:standard}

Throughout this section $\xi_0,\xi_1,\dots$ are independent standard real Gaussian variables, $G_n$ and $G_{n,t}$ are given by \eqref{eq:weyl} and \eqref{eq:heat}, and
\[
\Lambda_n^{\rm std}=\inf\{t\ge0:\ G_{n,t}\text{ has only real zeros}\}.
\]
In this coupling all degrees share the bottom coefficients $\xi_0,\xi_1,\dots$, while in the reversed coupling \eqref{eq:reversed} they share the top coefficients $\eta_0,\eta_1,\dots$; $\Lambda_n$ is the real-rooting time of $P_{n,t}$. As in the proof of Theorem~\ref{cor:weyl}, $\Lambda_n^{\rm std}$ and $\Lambda_n$ have the same law for each $n$, so $\Lambda_n^{\rm std}\to T_*$ in distribution, and only the joint law of different degrees differs. Theorem~\ref{thm:heat}, an almost sure statement, is specific to the reversed coupling: under the standard coupling there is no almost sure limit, and not even a limit in probability (Theorem~\ref{thm:std}).


\begin{remark}\label{rem:std-quantifier}
In Theorem~\ref{thm:std}(b) the order of the quantifiers matters: the set $\mathcal N$ is fixed before the sample point is drawn. For every sample point in the event of Theorem~\ref{thm:std}(a) the set $\{n:\Lambda_n^{\rm std}\in(2,3)\}$ is infinite, and along it all limit points lie in $[2,3]$.
\end{remark}

\subsection{The standard coupling in Hermite coordinates}
For real numbers $a_0,\dots,a_n$, or a real polynomial or power series $a(w)=\sum_ja_jw^j$, write
\[
Q_n[a]=\sum_{j=0}^{n}a_j\,h_{n-j}.
\]
Let $b_n(u)=\sum_{j=0}^{n}\xi_{n-j}u^j$ and $a^{(n,t)}(w)=b_n(w/\sqrt t)$, so that $a^{(n,t)}_j=\xi_{n-j}t^{-j/2}$. Lemma~\ref{lem:heat}(a) with $c_k=\xi_k$ and $k=n-j$ gives, for $t>0$,
\begin{equation}\label{eq:std-hermite}
G_{n,t}(z)=t^{n/2}\,Q_n\bigl[a^{(n,t)}\bigr]\Bigl(\sqrt{\tfrac nt}\,z\Bigr).
\end{equation}
This is \eqref{eq:PQ} with the top block $\xi_n,\xi_{n-1},\dots$ of the single sequence in place of $\eta_0,\eta_1,\dots$, so the symbol $a^{(n,t)}$ is renewed as $n$ grows. In the Hermite variable $y=z/\sqrt t$ the zeros of $G_{n,t}$ are those of $y\mapsto Q_n[a^{(n,t)}](\sqrt n\,y)$; in particular $G_{n,t}$ has only real zeros if and only if $Q_n[a^{(n,t)}]$ has. Almost surely $\xi_n\ne0$ for every $n$, and on this event Lemma~\ref{lem:heat}(c) gives, for every $t\ge0$,
\begin{equation}\label{eq:std-interval}
\{\Lambda_n^{\rm std}\le t\}=\{G_{n,t}\text{ has only real zeros}\},\qquad\{\Lambda_n^{\rm std}>t\}=\{G_{n,t}\text{ has a nonreal zero}\}.
\end{equation}
Only these one-sided descriptions are used below; we never approximate $\Lambda_n^{\rm std}$ by a function of the top block.

\subsection{An interlacing criterion}

\begin{lemma}[Interlacing criterion]\label{lem:interlace}
Let $n\ge2$ and $Q=\sum_{j=0}^{n}a_jh_{n-j}$ with real $a_j$ and $a_0\ne0$. Let $x_1<\dots<x_{n-1}$ be the zeros of $h_{n-1}$, and put $\alpha_j(x)=h_{n-j}(x)/h_n(x)$.
\begin{enumerate}
\item[(a)] $h_n(x_i)\ne0$, $\alpha_1(x_i)=0$, $\alpha_2(x_i)=-\sqrt{n/(n-1)}$, and $|\alpha_j(x_i)|\le\sqrt n$ for $1\le j\le n$.
\item[(b)] For $\kappa>0$ and $0\le j\le n/2$, $|\alpha_j(x_i)|\le2(j+1)e^{j\kappa+c_\kappa j^2/n}$, where $c_\kappa=\coth\kappa+\frac12$.
\item[(c)] If $\bigl|\sum_{j=2}^{n}a_j\alpha_j(x_i)\bigr|<|a_0|$ for every $i$, then $Q$ has $n$ simple real zeros, one in each of the intervals $(-\infty,x_1),(x_1,x_2),\dots,(x_{n-1},\infty)$. This holds in particular if, for some $\kappa>0$,
\begin{equation}\label{eq:interlace-crit}
|a_2|\sqrt{\frac n{n-1}}+\sum_{j=3}^{n}|a_j|\,\beta_{n,j}<|a_0|,\qquad
\beta_{n,j}=\begin{cases}\min\bigl\{\sqrt n,\ 2(j+1)e^{j\kappa+c_\kappa j^2/n}\bigr\},& j\le n/2,\\[2pt] \sqrt n,& j>n/2.\end{cases}
\end{equation}
\end{enumerate}
\end{lemma}

\begin{proof}
(a) By Lemma~\ref{lem:cd} at $x=x_i$, where $h_{n-1}$ vanishes, $\sum_{k<n}h_k(x_i)^2=n\,h_n(x_i)^2$. Since $h_0=1$, this gives $h_n(x_i)\ne0$ and $|h_k(x_i)|\le\sqrt n\,|h_n(x_i)|$ for $k<n$, that is, $|\alpha_j(x_i)|\le\sqrt n$ for $1\le j\le n$. Clearly $\alpha_1(x_i)=0$, and the recurrence \eqref{eq:recurrence} with $k=n-1$, which at $x_i$ reads $0=\sqrt n\,h_n(x_i)+\sqrt{n-1}\,h_{n-2}(x_i)$, gives $\alpha_2(x_i)=-\sqrt{n/(n-1)}$.

(b) By \eqref{eq:AB} with $z=x_i/\sqrt n$ and $h_{n-1}(x_i)=0$, we have $\alpha_j(x_i)=A_{n,j}(x_i/\sqrt n)$. By Lemma~\ref{lem:zeros}, $|x_i|<2\sqrt n$, so Lemma~\ref{lem:transfer} applies with $|z|<2\le2\cosh\kappa$.

(c) By \eqref{eq:recurrence} and (a), $h_{n-1}'(x_i)=\sqrt{n-1}\,h_{n-2}(x_i)=-\sqrt n\,h_n(x_i)$. The zeros of $h_{n-1}$ are simple and its leading coefficient is positive, so $\operatorname{sign}h_{n-1}'(x_i)=(-1)^{n-1-i}$, and hence $\operatorname{sign}h_n(x_i)=(-1)^{n-i}$. Since $\alpha_1(x_i)=0$, the hypothesis says that $Q(x_i)/h_n(x_i)=a_0+\sum_{j\ge2}a_j\alpha_j(x_i)$ has the sign of $a_0$, so $\operatorname{sign}Q(x_i)=(-1)^{n-i}\operatorname{sign}a_0$. The leading coefficient of $Q$ is $a_0/\sqrt{n!}$, so $Q(x)$ has the sign of $a_0$ for large positive $x$ and the sign of $(-1)^na_0$ for large negative $x$. Hence $Q$ changes sign on each of the $n$ intervals listed. As $Q$ has degree $n$, this accounts for all its zeros, and they are simple. Under \eqref{eq:interlace-crit} the hypothesis holds by (a) and (b).
\end{proof}

\begin{remark}\label{rem:interlace}
Sign tests at the zeros of $h_{n-1}$ are the classical route to the interlacing of the zeros of consecutive orthogonal polynomials (cf.\ \cite[Section~3.3, Theorem~3.3.2]{Sz}), and Lemma~\ref{lem:interlace}(c) is a test of this kind. What makes it effective for the random symbols below is the bound $|\alpha_j|\le\sqrt n$ of part (a), which is uniform in $j$ and comes from the Christoffel--Darboux formula. In the application $a_j=\xi_{n-j}t^{-j/2}$, where $t>1$ may be close to $1$. A geometric bound $|\alpha_j(x_i)|\le C^j$ alone would bring in the term $|\xi_0|(C/\sqrt t)^n$, which is unbounded unless $t>C^2$; with (a) and (b) the test works for every $t>1$ (Step~4 of the proof of Lemma~\ref{lem:std-planted}).
\end{remark}

\subsection{Every point of \texorpdfstring{$[1,\infty]$}{[1,infinity]} is a limit point}

\begin{lemma}\label{lem:std-planted}
Let $\mathcal N\subset\N$ be an infinite set that does not depend on the $\xi_k$, and let $1<t_1<t_2<\infty$. Almost surely $t_1<\Lambda_n^{\rm std}\le t_2$ for infinitely many $n\in\mathcal N$.
\end{lemma}

\begin{proof}
All constants below depend only on $t_1$ and $t_2$.

\emph{Step 1: parameters.} Put
\[
\lambda=t_1^{1/4},\qquad\rho=t_1^{-1/2},\qquad\theta_1=\lambda\rho=t_1^{-1/4},\qquad\sigma=(t_1t_2)^{-1/4},\qquad p(u)=\sigma^2+u^2 .
\]
Then $\lambda>1$, $\theta_1<1$ and $t_2^{-1/2}<\sigma<\rho<1$, and $p$ has the simple zeros $\pm i\sigma$. Let
\[
2\delta_2=\sigma^2-t_2^{-1}>0,\qquad\eps_0=\tfrac12\min\{\sigma,\ \rho-\sigma\},\qquad2\delta_1=\eps_0\sigma,\qquad d=1-\sqrt{t_1}\,(\sigma+\eps_0).
\]
Since $\sigma+\eps_0\le(\sigma+\rho)/2$, we have $\frac12(1-\sqrt{t_1}\,\sigma)\le d<1$, so $0<d<1$. Let $M_n=\lceil3\log n/\log t_1\rceil$, so that $\lambda^{M_n}=t_1^{M_n/4}\ge n^{3/4}$. Since $\lambda/\sqrt{t_2}<t_1^{1/4}t_1^{-1/2}=t_1^{-1/4}<1$, we can choose $\kappa>0$ with $\theta_2:=\lambda e^{\kappa}/\sqrt{t_2}<1$; let $c_\kappa=\coth\kappa+\frac12$ and $S=\sum_{j\ge3}(j+1)\theta_2^j$. Finally choose $\eps\in(0,1]$ with
\[
\eps\Bigl(3+\frac{\theta_1^3}{1-\theta_1}\Bigr)\le\frac{\delta_1}2,\qquad\eps\,(3+3S)\le\delta_2 .
\]

\emph{Step 2: the events.} For $n$ with $M_n+2\le n/2$ let $E_n$ be the event that
\[
|\xi_n-\sigma^2|<\eps,\qquad|\xi_{n-1}|<\eps,\qquad|\xi_{n-2}-1|<\eps,\qquad|\xi_{n-j}|\le\eps\lambda^j\quad(3\le j\le M_n),
\]
which depends only on the top block $(\xi_{n-j})_{0\le j\le M_n}$, and let $F_n=\{\max_{k\le n}|\xi_k|\le\lambda^{M_n}\}$.

(i) By independence,
\[
\PP(E_n)\ge\pi_0:=\PP(|\xi_0-\sigma^2|<\eps)\,\PP(|\xi_0|<\eps)\,\PP(|\xi_0-1|<\eps)\prod_{j\ge3}\PP(|\xi_0|\le\eps\lambda^j)>0,
\]
because every factor is positive and $\sum_{j\ge3}\PP(|\xi_0|>\eps\lambda^j)\le\sum_{j\ge3}2e^{-\eps^2\lambda^{2j}/2}<\infty$, as $\lambda>1$.

(ii) Since $\PP(|\xi_k|>2\sqrt{\log k})\le2k^{-2}$ for $k\ge2$, the first Borel--Cantelli lemma gives an almost surely finite $C$ with $\max_{k\le n}|\xi_k|\le C+2\sqrt{\log n}$ for all $n\ge2$. As $\lambda^{M_n}\ge n^{3/4}$, almost surely $F_n$ holds for all large $n$.

(iii) Since $n-M_n\to\infty$, we can choose $n_1<n_2<\cdots$ in $\mathcal N$ such that $M_n+2\le n/2$ for all $n\ge n_1$ and $n_{m+1}-M_{n_{m+1}}>n_m$ for all $m$. The event $E_{n_m}$ depends only on the $\xi_k$ with $n_m-M_{n_m}\le k\le n_m$. These blocks of indices are disjoint, so the events $E_{n_m}$ are independent, and $\sum_m\PP(E_{n_m})=\infty$ by (i). By the second Borel--Cantelli lemma, almost surely $E_{n_m}$ occurs for infinitely many $m$.

\emph{Step 3: the coefficients on $E_n\cap F_n$.} Let $n\ge n_1$, and assume that $E_n\cap F_n$ occurs.

(a) $|\xi_{n-j}|\le2\lambda^j$ for $0\le j\le n$. For $j\le2$ this follows from $\sigma^2+\eps\le2$, $\eps\le2\lambda$ and $1+\eps\le2$; for $3\le j\le M_n$ from $\eps\le1$; and for $j>M_n$ from $F_n$, which gives $|\xi_{n-j}|\le\lambda^{M_n}<\lambda^j$. Hence $|a^{(n,t_1)}_j|\le2(\lambda/\sqrt{t_1})^j=2\,t_1^{-j/4}$ for all $j$.

(b) For $|u|\le\rho$, using $\lambda\rho=\theta_1$ and $\rho<1$,
\[
|b_n(u)-p(u)|\le|\xi_n-\sigma^2|+|\xi_{n-1}|+|\xi_{n-2}-1|+\eps\sum_{j=3}^{M_n}\theta_1^j+\sum_{j>M_n}\theta_1^j\le\eps\Bigl(3+\frac{\theta_1^3}{1-\theta_1}\Bigr)+\frac{\theta_1^{M_n+1}}{1-\theta_1}.
\]
Let $N_1\ge n_1$ be such that $\theta_1^{M_n+1}/(1-\theta_1)\le\delta_1/2$ for $n\ge N_1$. Then $|b_n-p|\le\delta_1$ on $\{|u|\le\rho\}$ for $n\ge N_1$.

\emph{Step 4: $\Lambda_n^{\rm std}\le t_2$.} Let $N_2\ge n_1$ be such that for all $n\ge N_2$
\[
e^{c_\kappa M_n^2/n}\le\tfrac32\qquad\text{and}\qquad\Bigl(\sqrt{\tfrac n{n-1}}-1\Bigr)+\frac{n^{-1/4}}{1-\lambda/\sqrt{t_2}}<\delta_2 .
\]
Let $n\ge N_2$ and $a=a^{(n,t_2)}$, so that $a_j=\xi_{n-j}t_2^{-j/2}$. On $E_n\cap F_n$ we have $|a_0|>\sigma^2-\eps$ and $|a_2|<(1+\eps)t_2^{-1}$. For $3\le j\le M_n$, since $M_n\le n/2$,
\[
|a_j|\,\beta_{n,j}\le\eps\lambda^jt_2^{-j/2}\cdot2(j+1)e^{j\kappa}\cdot\tfrac32=3\eps(j+1)\theta_2^j ,
\]
and for $M_n<j\le n$, by $F_n$, $|a_j|\le\lambda^{M_n}t_2^{-j/2}\le(\lambda/\sqrt{t_2})^j$ and $\beta_{n,j}\le\sqrt n$. Hence the left side of \eqref{eq:interlace-crit} is at most
\[
(1+\eps)\,t_2^{-1}\sqrt{\tfrac n{n-1}}+3\eps S+\tau_n,\qquad\tau_n=\frac{\sqrt n\,(\lambda/\sqrt{t_2})^{M_n+1}}{1-\lambda/\sqrt{t_2}}\le\frac{n^{-1/4}}{1-\lambda/\sqrt{t_2}},
\]
where we used $\lambda/\sqrt{t_2}<t_1^{-1/4}$ and $t_1^{M_n/4}\ge n^{3/4}$. Using $\sigma^2=t_2^{-1}+2\delta_2$, $t_2^{-1}<1$, $\sqrt{n/(n-1)}\le\sqrt2$ and the choice of $\eps$, the right side of \eqref{eq:interlace-crit} minus this bound is at least
\[
2\delta_2-\eps\bigl(1+\sqrt2+3S\bigr)-\Bigl(\sqrt{\tfrac n{n-1}}-1\Bigr)-\tau_n\ \ge\ \delta_2-\Bigl(\sqrt{\tfrac n{n-1}}-1\Bigr)-\tau_n ,
\]
which is positive by the choice of $N_2$. Hence \eqref{eq:interlace-crit} holds, so by Lemma~\ref{lem:interlace} the polynomial $Q_n[a^{(n,t_2)}]$ has only real zeros; by \eqref{eq:std-hermite} so has $G_{n,t_2}$, and by \eqref{eq:std-interval} $\Lambda_n^{\rm std}\le t_2$.

Apart from $a_0$, the only coefficient that is not small is $a_2$, and $\alpha_2=-\sqrt{n/(n-1)}$ is close to $-1$; so the test amounts to $t_2^{-1}<\sigma^2=(t_1t_2)^{-1/2}$, that is, to $t_1<t_2$. The coefficient $a_1$ does not enter, because $\alpha_1=0$ at the zeros of $h_{n-1}$.

\emph{Step 5: $\Lambda_n^{\rm std}>t_1$.} The circle $|u-i\sigma|=\eps_0$ lies in $\{|u|<\rho\}$, since $\sigma+\eps_0<\rho$, and on it
\[
|p(u)|=|u-i\sigma|\,|u+i\sigma|\ge\eps_0(2\sigma-\eps_0)\ge\eps_0\sigma=2\delta_1 .
\]
Let $n\ge N_1$. On $E_n\cap F_n$, Step~3(b) gives $|b_n-p|\le\delta_1<|p|$ on this circle, so $|b_n|\ge\delta_1$ there, and by Rouch\'e's theorem $b_n$ has exactly one zero in $\{|u-i\sigma|<\eps_0\}$, as $p$ has (note $|{-i\sigma}-i\sigma|=2\sigma>\eps_0$). Let $B=\{|w-i\sqrt{t_1}\,\sigma|\le\sqrt{t_1}\,\eps_0\}$. Then $a^{(n,t_1)}(w)=b_n(w/\sqrt{t_1})$ has exactly one zero in the interior of $B$, and $|a^{(n,t_1)}|\ge\delta_1$ on $\partial B$. Moreover $B\subset\{|w|\le\sqrt{t_1}(\sigma+\eps_0)\}=\{|w|\le1-d\}$, and $B\subset\{\im w\ge\sqrt{t_1}\,\sigma/2\}$, so $B\cap\R=\emptyset$. By Step~3(a), $|a^{(n,t_1)}_j|\le2\,t_1^{-j/4}$ for all $j$. Lemma~\ref{lem:ext-quant}, with $C_0=2$, $R_0=t_1^{1/4}$, $\delta=\delta_1$ and $d$ as in Step~1, shows that for $n>n_{\rm ext}$ the polynomial $Q_n[a^{(n,t_1)}]$ has a nonreal zero. By \eqref{eq:std-hermite} so has $G_{n,t_1}$, and by \eqref{eq:std-interval} $\Lambda_n^{\rm std}>t_1$.

\emph{Step 6: conclusion.} Let $N_0=\max\{N_1,N_2,\lfloor n_{\rm ext}\rfloor+1\}$. By Steps 4 and 5, $E_n\cap F_n\subset\{t_1<\Lambda_n^{\rm std}\le t_2\}$ for every $n\ge N_0$, up to the null event that $\xi_k=0$ for some $k$. By Step~2, almost surely $E_{n_m}$ occurs for infinitely many $m$ and $F_n$ holds for all large $n$. Hence almost surely $t_1<\Lambda_{n_m}^{\rm std}\le t_2$ for infinitely many $m$.
\end{proof}

Planting a prescribed block of top coefficients is a simple device: for the reversed coupling it also shows directly that $\PP(T_*>t)>0$ for every $t$, since if $\eta_0$ is close to $\delta^2$, $\eta_1$ to $0$, $\eta_2$ to $1$, and the later $\eta_j$ are small enough, then $q$ has a pair of nonreal zeros near $\pm i\delta$, so that $T_*>t$ when $\delta<t^{-1/2}$. What is new here is the almost sure version on disjoint blocks of coordinates, which rules out convergence in probability to any random limit. The events $E_n$ are designed so that the symbols they produce satisfy the hypotheses of Lemma~\ref{lem:interlace} at time $t_2$ and of Lemma~\ref{lem:ext-quant} at time $t_1$ with constants that do not depend on $n$; this gives a single onset $N_0$. Theorem~\ref{thm:main} would not do here, since its onset depends on the symbol, which changes with $n$.

\subsection{Mehler's formula and real-rootedness before time 1}

\begin{lemma}[Mehler's formula]\label{lem:mehler}
For $0<r<1$ and $x,y\in\C$,
\begin{equation}\label{eq:mehler}
\sum_{k\ge0}h_k(x)h_k(y)\,r^k=(1-r^2)^{-1/2}\exp\Bigl(\frac{2rxy-r^2(x^2+y^2)}{2(1-r^2)}\Bigr),
\end{equation}
the series converging absolutely. In particular $\sum_{k\ge0}r^k|h_k(x)|^2=(1-r^2)^{-1/2}\exp\bigl((r|x|^2-r^2\re x^2)/(1-r^2)\bigr)$.
\end{lemma}

\begin{proof}
Let $N,N'$ be independent standard real Gaussian variables. Since $\E e^{\lambda(x+iN)}=e^{\lambda x-\lambda^2/2}=\sum_k\He_k(x)\lambda^k/k!$ for complex $\lambda$, comparing coefficients of $\lambda^k$ gives $\He_k(x)=\E(x+iN)^k$; the exchange of sum and expectation is justified by $\E e^{|\lambda|(|x|+|N|)}<\infty$. Hence
\[
\sum_k\frac{\He_k(x)\He_k(y)}{k!}\,r^k=\E\exp\bigl(r(x+iN)(y+iN')\bigr),
\]
where the series converges absolutely and the exchange is justified because $\E\exp\bigl(r(|x|+|N|)(|y|+|N'|)\bigr)<\infty$ for $r<1$. Integrating first over $N'$, with $\E e^{cN'}=e^{c^2/2}$ for complex $c$, and then over $N$, with $\E e^{bN+r^2N^2/2}=(1-r^2)^{-1/2}e^{b^2/(2(1-r^2))}$ for complex $b$, gives
\[
(1-r^2)^{-1/2}\exp\Bigl(rxy-\frac{r^2x^2}2-\frac{r^2(y-rx)^2}{2(1-r^2)}\Bigr),
\]
which is the right side of \eqref{eq:mehler}. The second statement is the case $y=\bar x$, since the $h_k$ have real coefficients.
\end{proof}

This is Mehler's classical formula \cite{Me}, often stated for real $x,y$. We include the short proof because we use it at complex arguments, with $y=\bar x$ (Steps~1 and~2 of the proof of Proposition~\ref{prop:early}). To extend the real formula by analytic continuation one needs a bound for $|h_k|$ off the real line that makes the series converge locally uniformly on $\C^2$; the representation $\He_k(x)=\E(x+iN)^k$ gives this bound and the formula at once.

\begin{proof}[Proof of Proposition~\ref{prop:early}]
Fix $s\in(0,1)$ and $n\ge1$. For $z=X+iY$ put
\[
g_k(z)=s^{k/2}h_k\bigl(\sqrt{n/s}\,z\bigr),\qquad V_n(z)=\sum_{k=0}^{n}|g_k(z)|^2,\qquad V_\infty(z)=\sum_{k\ge0}|g_k(z)|^2,
\]
and
\[
\varphi(z)=\frac12\Bigl(\frac{X^2}{1+s}+\frac{Y^2}{1-s}\Bigr).
\]
By Lemma~\ref{lem:heat}(a) with $c_k=\xi_k$ and $t=s$, $G_{n,s}=\sum_{k\le n}\xi_kg_k$, so $V_n(z)=\E|G_{n,s}(z)|^2$. The idea is that near the origin $\log|G_{n,s}|$ is close to $n\varphi$ with high probability, that $\varphi$ is strictly subharmonic, and that $\log|G_{n,s}|$ would be harmonic in the upper half-plane if all zeros of $G_{n,s}$ were real.

\emph{Step 1: the variance.} By Lemma~\ref{lem:mehler} with $r=s$ and $x=\sqrt{n/s}\,z$, for which $r|x|^2-r^2\re x^2=n(X^2+Y^2)-ns(X^2-Y^2)$,
\begin{equation}\label{eq:std-Vinf}
V_\infty(z)=(1-s^2)^{-1/2}e^{2n\varphi(z)}\ge e^{2n\varphi(z)}\ge1 .
\end{equation}
(Inside the ellipse of \cite{HHJK}, $(2\pi)^{-1}\Delta(n\varphi)=n/(\pi(1-s^2))$ is the density of $n$ times its uniform distribution; we do not use this.)

\emph{Step 2: the truncation is harmless near $0$.} For $k>n$ we have $s^k\le s^{(n+1)/2}(\sqrt s)^k$. So Lemma~\ref{lem:mehler} with $r=\sqrt s$, together with $-r^2\re x^2\le r^2|x|^2$ and $(\sqrt s+s)/(s(1-s))=1/(\sqrt s(1-\sqrt s))$, gives
\[
V_\infty(z)-V_n(z)\le s^{(n+1)/2}(1-s)^{-1/2}\exp\Bigl(\frac{n|z|^2}{\sqrt s\,(1-\sqrt s)}\Bigr).
\]
For $|z|\le\beta_s$ the exponent is at most $\frac n4\log(1/s)$, so $V_\infty(z)-V_n(z)\le(1-s)^{-1/2}s^{1/2}s^{n/4}$. Let $n_s'$ be the least $n$ with $(1-s)^{-1/2}s^{1/2}s^{n/4}\le\frac12$. For $n\ge n_s'$ and $|z|\le\beta_s$, by \eqref{eq:std-Vinf},
\begin{equation}\label{eq:std-Vn}
V_n(z)\ge V_\infty(z)-\tfrac12\ge\tfrac12V_\infty(z)\ge\tfrac12e^{2n\varphi(z)} .
\end{equation}

\emph{Step 3: geometry.} Let $y_0=\beta_s/2$, $z_0=iy_0$, $r_1=y_0/2$, $r_2=r_1/2$, and $z_l'=z_0+r_2e^{2\pi il/12}$ for $0\le l<12$. The closed disk $\overline D=\overline D(z_0,r_1)$ lies in $\{\im z\ge y_0/2\}\cap\{|z|\le\frac34\beta_s\}$. Since $\varphi$ is a quadratic polynomial in $(X,Y)$, and the averages of $\cos\vartheta$, $\sin\vartheta$, $\cos2\vartheta$ and $\sin2\vartheta$ over $\vartheta=2\pi l/12$ vanish,
\begin{equation}\label{eq:std-gain}
\frac1{12}\sum_{l<12}\varphi(z_l')-\varphi(z_0)=\frac{r_2^2}4\Bigl(\frac1{1+s}+\frac1{1-s}\Bigr)=\frac{r_2^2}{2(1-s^2)}=\frac{y_0^2}{32(1-s^2)}=\gamma_s .
\end{equation}
Moreover $\varphi(z)\le|z|^2/(2(1-s))$, so $\Phi_{\max}:=\max_{\overline D}\varphi\le(3y_0/2)^2/(2(1-s))=36(1+s)\gamma_s$.

\emph{Step 4: two Gaussian estimates.} Let $\eps=\gamma_s/4$.

(U) Let $\Omega_n^1=\{\sum_{k\le n}\xi_k^2\le4(n+1)\}$. If $\chi^2_m$ has the chi-squared distribution with $m$ degrees of freedom and $x>1$, Markov's inequality for $e^{\lambda\chi^2_m}$, with $\E e^{\lambda\chi^2_m}=(1-2\lambda)^{-m/2}$, gives
\[
\PP(\chi^2_m\ge mx)\le\inf_{0<\lambda<1/2}e^{-\lambda mx}(1-2\lambda)^{-m/2}=(xe^{1-x})^{m/2},
\]
the infimum being attained at $\lambda=(1-1/x)/2$. With $m=n+1$ and $x=4$, $\PP((\Omega_n^1)^c)\le(4e^{-3})^{(n+1)/2}\le e^{-0.8(n+1)}$. On $\Omega_n^1$, by the Cauchy--Schwarz inequality and \eqref{eq:std-Vinf}, for all $z\in\C$,
\begin{equation}\label{eq:std-upper}
\log|G_{n,s}(z)|\le\log\bigl(2\sqrt{n+1}\bigr)+\tfrac12\log V_\infty(z)=n\varphi(z)+\log\bigl(2\sqrt{n+1}\bigr)+c_s',
\end{equation}
where $c_s'=\frac14\log\bigl(1/(1-s^2)\bigr)$.

(L) Fix $z\in\C$. For real $\alpha$, $\re(e^{-i\alpha}G_{n,s}(z))=\sum_k\xi_k\re(e^{-i\alpha}g_k(z))$ is a centred Gaussian variable with variance $\frac12V_n(z)+\frac12\re\bigl(e^{-2i\alpha}\sum_kg_k(z)^2\bigr)$, which is at least $\frac12V_n(z)>0$ for a suitable $\alpha$ (note $V_n\ge|g_0|^2=1$). Since $|G_{n,s}(z)|\ge|\re(e^{-i\alpha}G_{n,s}(z))|$, and a centred Gaussian variable with variance $v$ has density at most $(2\pi v)^{-1/2}$,
\[
\PP\Bigl(|G_{n,s}(z)|<e^{-\eps n}\sqrt{V_n(z)/2}\Bigr)\le\sqrt{2/\pi}\,e^{-\eps n}\le e^{-\eps n}.
\]
Let $\Omega_n^2$ be the event that $|G_{n,s}(z)|\ge e^{-\eps n}\sqrt{V_n(z)/2}$ for each of the $13$ points $z\in\{z_0,z_0',\dots,z_{11}'\}$. Then $\PP((\Omega_n^2)^c)\le13e^{-\gamma_sn/4}$. These points lie in $\{|z|\le\beta_s\}$, so on $\Omega_n^2$, for $n\ge n_s'$ and each of them, \eqref{eq:std-Vn} gives
\begin{equation}\label{eq:std-lower}
\log|G_{n,s}(z)|\ge n\varphi(z)-\eps n-\log2 .
\end{equation}

\emph{Step 5: harmonic averaging.} Let $H$ be harmonic on a neighbourhood of $\overline D(z_0,r_1)$, and recall that $r_2/r_1=\frac12$. Then
\begin{equation}\label{eq:std-average}
\Bigl|\frac1{12}\sum_{l<12}H(z_l')-H(z_0)\Bigr|\le c_{12}\cdot\frac1{2\pi}\int_0^{2\pi}\bigl|H(z_0+r_1e^{i\vartheta})\bigr|\,d\vartheta,\qquad c_{12}=\frac{2\cdot2^{-12}}{1-2^{-12}}=\frac2{4095}.
\end{equation}
Indeed, by the Poisson integral, $H(z_0+r_2e^{i\vartheta})=\sum_{k\in\mathbb Z}\widehat H(k)2^{-|k|}e^{ik\vartheta}$ with $\widehat H(k)=\frac1{2\pi}\int_0^{2\pi}H(z_0+r_1e^{i\psi})e^{-ik\psi}d\psi$. Thus $\widehat H(0)=H(z_0)$ and $|\widehat H(k)|\le\frac1{2\pi}\int_0^{2\pi}|H(z_0+r_1e^{i\psi})|\,d\psi$. The average of $e^{ik\vartheta}$ over $\vartheta=2\pi l/12$ is $1$ if $12\mid k$ and $0$ otherwise, so the left side of \eqref{eq:std-average} is at most $2\sum_{m\ge1}2^{-12m}$ times the mean of $|H|$ on the circle.

\emph{Step 6: conclusion.} With $\eps=\gamma_s/4$, $\Phi_s=36(1+s)\gamma_s$ and $c_{12}=2/4095$, let
\[
f_s(n)=n(\gamma_s-\eps)-\log2-\log\bigl(2\sqrt{n+1}\bigr)-c_s'-c_{12}\Bigl(2\bigl(n\Phi_s+\log(2\sqrt{n+1})+c_s'\bigr)+\eps n+\log2\Bigr),
\]
and
\begin{equation}\label{eq:std-onset}
n_s=\min\{n\ge n_s':\ f_s(n)>0\}.
\end{equation}
Here $f_s$ is convex in $n$, $f_s(1)<0$ (as $\gamma_s<1$), and $f_s(n)/n\to\gamma_s\bigl(\frac34-c_{12}(72(1+s)+\frac14)\bigr)>\frac{\gamma_s}2$, so $n_s$ is finite, and by convexity $f_s(n)>0$ for every $n\ge n_s$. For $s=0.3,\ 0.5,\ 0.8,\ 0.95$ one finds $n_s=13\,707,\ 25\,051,\ 93\,949,\ 493\,795$ (and $n_s'=1,4,25,169$).

Let $n\ge n_s$, and suppose that $\Omega_n^1\cap\Omega_n^2$ occurs and that $G_{n,s}$ has only real zeros. On $\Omega_n^2$, $G_{n,s}(z_0)\ne0$, so $H=\log|G_{n,s}|$ is harmonic on the upper half-plane, which contains $\overline D$. Let $I=\frac1{2\pi}\int_0^{2\pi}|H(z_0+r_1e^{i\vartheta})|\,d\vartheta$. By the mean value property, $I=\frac1\pi\int_0^{2\pi}H^+(z_0+r_1e^{i\vartheta})\,d\vartheta-H(z_0)$. By \eqref{eq:std-upper} on the circle, where its right side is positive, and by \eqref{eq:std-lower} at $z_0$ together with $\varphi\ge0$,
\[
I\le2\bigl(n\Phi_{\max}+\log(2\sqrt{n+1})+c_s'\bigr)+\eps n+\log2 .
\]
By \eqref{eq:std-lower} at the points $z_l'$, \eqref{eq:std-upper} at $z_0$, and \eqref{eq:std-gain},
\[
\frac1{12}\sum_{l<12}H(z_l')-H(z_0)\ge n\gamma_s-\eps n-\log2-\log\bigl(2\sqrt{n+1}\bigr)-c_s' .
\]
By \eqref{eq:std-average} the left side is at most $c_{12}I$. Since $\Phi_{\max}\le\Phi_s$, the last three inequalities give $f_s(n)\le0$, a contradiction. Hence, for $n\ge n_s$,
\[
\{G_{n,s}\text{ has only real zeros}\}\subset(\Omega_n^1)^c\cup(\Omega_n^2)^c ,
\]
and \eqref{eq:std-early} follows from Step 4.

For the consequences, note that $\beta_s^2/(1-s^2)=\sqrt s\log(1/\sqrt s)/(2(1+\sqrt s)(1+s))\le1/(2e)$, so $\gamma_s\le1/(256e)$ and $e^{-0.8(n+1)}\le e^{-\gamma_sn/4}$. So the right side of \eqref{eq:std-early} is at most $14e^{-\gamma_sn/4}\le C_se^{-c_sn}$ for $n\ge n_s$, while for $n<n_s$ we have $C_se^{-c_sn}>14$. Finally, the event $\{G_{n,s}$ has only real zeros$\}$ is determined by $(\xi_0,\dots,\xi_n)$, and by \eqref{eq:std-interval} it is the event $\{\Lambda_n^{\rm std}\le s\}$ up to a null set; the same holds for $P_{n,s}$ and $\Lambda_n$, whose coefficient vector has the same law.
\end{proof}

The proof uses neither the convergence of empirical measures of \cite{HHJK} nor Plancherel--Rotach asymptotics: the points are taken near the origin, where the truncation $k\le n$ is exponentially negligible by Mehler's formula alone. Steps 5 and 6 are the standard hole-probability argument for Gaussian analytic functions: a zero-free disk makes $\log|f|$ harmonic, and this is played against lower bounds at finitely many points \cite{ST}. The onset $n_s$ is large, and the right side of \eqref{eq:std-early} is below $1$ only for $n>4\gamma_s^{-1}\log13$, about $2.7\cdot10^4$ at $s=\frac12$. We have not tried to find the true rate of decay.

\subsection{Proof of the limit-point theorem}

\begin{proof}[Proof of Theorem~\ref{thm:std}]
\emph{No limit point below $1$.} By \eqref{eq:std-interval} and Proposition~\ref{prop:early}, $\sum_n\PP(\Lambda_n^{\rm std}\le s)<\infty$ for every $s\in(0,1)$. By the first Borel--Cantelli lemma, applied to each rational $s\in(0,1)$, almost surely $\Lambda_n^{\rm std}>s$ for all large $n$ and every rational $s<1$. Hence almost surely $\liminf_n\Lambda_n^{\rm std}\ge1$, and no point of $[0,1)$ is a limit point of $(\Lambda_n^{\rm std})_{n\in\mathcal N}$, for any infinite $\mathcal N\subset\N$.

\emph{Every point of $[1,\infty]$ is a limit point.} Fix an infinite set $\mathcal N\subset\N$ that does not depend on the $\xi_k$. Applying Lemma~\ref{lem:std-planted} to the countably many pairs of rationals $1<t_1<t_2$, we obtain an almost sure event on which, for every such pair, $t_1<\Lambda_n^{\rm std}\le t_2$ for infinitely many $n\in\mathcal N$. For real $1<t_1<t_2$ choose rationals $t_1<t_1'<t_2'<t_2$; then $(t_1',t_2']\subset(t_1,t_2]$, so the same holds for all real pairs. On this event every $t\in(1,\infty)$ is a limit point of $(\Lambda_n^{\rm std})_{n\in\mathcal N}$, since for every $k$ infinitely many of these values lie in $(t-1/k,t+1/k)$. The set of limit points is closed in $[0,\infty]$, so it also contains $1$ and $\infty$.

Together, the two parts prove (b), and (a) is the case $\mathcal N=\N$; the smallest and largest limit points are $1$ and $\infty$.

(c) Give $[0,\infty]$ the metric $|\arctan x-\arctan y|$. If $\Lambda_n^{\rm std}\to Y$ in probability, then $\Lambda_n^{\rm std}\to Y$ almost surely along some infinite deterministic set $\mathcal N\subset\N$, so that almost surely $(\Lambda_n^{\rm std})_{n\in\mathcal N}$ has exactly one limit point. This contradicts (b).
\end{proof}


\clearpage
\part{The final collision}

\section{The final collision}\label{sec:collision}

In this section $q(w)=\sum_{j\ge0}\eta_jw^j$ is a real power series with radius of convergence $\rho_q\in(0,\infty]$; it is deterministic until Section~\ref{subsec:collision-generic}. The polynomials $P_n$, $P_{n,t}$, $Q_n^{(t)}$ and the real-rooting time $\Lambda_n$ are defined from its coefficients by \eqref{eq:reversed} and \eqref{eq:PQ}, as in Section~\ref{subsec:heatflow}, and we put $P_{n,t}=\exp(-\frac t{2n}\partial_z^2)P_n$ for every real $t$ (a finite sum). Sections~\ref{sec:transfer}--\ref{sec:exterior} call the Hermite variable $z$. Here we write $y$ for it and keep $z$ for the heat variable, so that by \eqref{eq:PQ}
\begin{equation}\label{eq:collision-S1}
P_{n,t}(z)=t^{n/2}\,Q_n^{(t)}(\sqrt n\,y),\qquad z=\sqrt t\,y\qquad(t>0),
\end{equation}
and the zeros of $P_{n,t}$ and of $y\mapsto Q_n^{(t)}(\sqrt n\,y)$ correspond under $z=\sqrt t\,y$, with multiplicities. The argument of the $h_k$ is $x=\sqrt n\,y$, and $v_n(x)=(h_n(x),h_{n-1}(x))^T$ as in Section~\ref{sec:bulk}. Let $N_n(t)$ be the number of nonreal zeros of $P_{n,t}$, counted with multiplicity. We use the following conditions on $q$.
\begin{enumerate}
\item[(G1)] $q(0)\ne0$, and $q$ has a nonreal zero in $\{|w|<\rho_q\}$. Put
\[
R_*=\min\{|u|:\ |u|<\rho_q,\ q(u)=0,\ u\notin\R\},\qquad T_*=R_*^{-2};
\]
the minimum is attained because the zeros of $q$ in $\{|w|<\rho_q\}$ are isolated.
\item[(G2)] The nonreal zeros of $q$ of modulus $R_*$ are exactly one conjugate pair $u_*=R_*e^{i\theta_*}$, $\bar u_*$, with $0<\theta_*<\pi$, and $u_*$ is a simple zero.
\item[(G3)] $q$ has no real zero of modulus $R_*$, and every real zero of $q$ of modulus less than $R_*$ is simple.
\end{enumerate}
Radius $\rho_q\ge1$ is not needed: since $R_*<\rho_q$, the series $a^{(t)}(w)=q(w/\sqrt t)$ has radius $\sqrt t\,\rho_q>1$ for $t$ near $T_*$, and only such $t$ occur. For the Gaussian power series \eqref{eq:kac}, (G1)--(G3) hold almost surely with $\rho_q=1$ (Lemma~\ref{lem:moduli}). Put
\[
y_*=2\cos\theta_*\in(-2,2),\qquad z_*=\sqrt{T_*}\,y_*=2R_*^{-1}\cos\theta_* .
\]

\begin{theorem}[The final collision]\label{thm:collision}
Let $q$ satisfy (G1)--(G3). There are $\delta\in(0,T_*/4)$, $K<\infty$ and $n_0$ such that for every $n\ge n_0$:
\begin{enumerate}
\item[(a)] \emph{(time)} $T_*-\delta<\Lambda_n\le T_*+\delta$ and $|\Lambda_n-T_*|\le K/n$;
\item[(b)] \emph{(the collision)} all zeros of $P_{n,\Lambda_n}$ are real, and exactly one of them, $z_n$, is multiple. Its multiplicity is $2$ or $3$, and, with $y_n=z_n/\sqrt{\Lambda_n}$,
\[
|z_n-z_*|\le\frac Kn,\qquad|y_n-y_*|\le\frac Kn ;
\]
\item[(c)] \emph{(before and after)} for $T_*-\delta\le t<\Lambda_n$ the polynomial $P_{n,t}$ has exactly two nonreal zeros, a conjugate pair, and all its zeros are simple; for $t>\Lambda_n$ all zeros of $P_{n,t}$ are real and simple;
\item[(d)] \emph{(uniqueness and localisation)} $\Lambda_n$ is the only $t\ge T_*-\delta$ at which $P_{n,t}$ has a multiple zero. If $|t-T_*|\le\delta$ and $y$ is a real multiple zero of $y\mapsto Q_n^{(t)}(\sqrt n\,y)$, equivalently $z=\sqrt t\,y$ is a real multiple zero of $P_{n,t}$, then $|t-T_*|+|y-y_*|\le K/n$.
\end{enumerate}
\end{theorem}

For fixed $t<T_*$ close to $T_*$, Theorem~\ref{thm:fixedq}(a) places a nonreal zero of $P_{n,t}$ near $tu_*+u_*^{-1}$, and $T_*u_*+u_*^{-1}=z_*$. So the last pair of zeros becomes real inside the bulk, at the relative position $\cos\theta_*$ of $[-2\sqrt{T_*},2\sqrt{T_*}]$, and within $O(1/n)$ of the time $T_*$; in particular $n(\Lambda_n-T_*)$ is bounded. We do not claim here that $n(\Lambda_n-T_*)$ converges; its first-order behaviour is the subject of Section~\ref{sec:firstorder}. The constants $\delta$, $K$, $n_0$ depend on $q$. The proof below obtains them by compactness; in the companion paper \cite{CompA}, which replaces Theorem~\ref{thm:main} at the two times $T_*\pm\delta$ by a version with an explicit onset, they are bounded by polynomials in condition parameters of $q$, with absolute degrees and, when $T_*\ge1+\eps$, coefficients polynomial in $1/\eps$. % MERGE: was [I, thm:quant] and Theorem thm:collision-explicit of Paper III Part (c) counts the nonreal pair on $[T_*-\delta,\Lambda_n)$ but does not locate it uniformly in $t$.

The proof is in Section~\ref{subsec:collision-proof}. The number of nonreal zeros is carried from $T_*-\delta$ to $T_*+\delta$ by the heat flow (Section~\ref{subsec:collision-heat}): it never increases, and every real multiple zero forces two extra nonreal zeros just before it. This leaves room for exactly one collision. A quantitative form of the bulk estimate of Section~\ref{sec:bulk}, whose lower constant degenerates linearly at $(T_*,y_*)$, then locates it (Sections~\ref{subsec:collision-rate}--\ref{subsec:collision-localise}).

\subsection{The heat flow: the nonreal count and local collisions}\label{subsec:collision-heat}
In this subsection $P$ is a real polynomial of degree $n\ge1$ and $P_t=\exp(-\frac t{2n}\partial_z^2)P$ for $t\in\R$. Then $P_t$ has the leading coefficient of $P$, its coefficients are polynomials in $t$, $P_{t+s}=\exp(-\frac s{2n}\partial_z^2)P_t$, and $\partial_tP_t=-\frac1{2n}\partial_z^2P_t$. Let $Z(p)$ be the number of nonreal zeros of a real polynomial $p$, counted with multiplicity, as in Lemma~\ref{lem:heat}(b); then $N_n(t)=Z(P_{n,t})$. Nothing in this subsection uses the Hermite structure.

\begin{lemma}[The nonreal count does not increase]\label{lem:heatcount}
For $t\in\R$ and $s\ge0$, $Z(P_{t+s})\le Z(P_t)$. Moreover, if real polynomials $p_m$ of degree $n$ converge coefficientwise to a polynomial $p$ of degree $n$, then $Z(p)\le\liminf_mZ(p_m)$. In particular $\{t\ge0:\ Z(P_t)=0\}$ is closed.
\end{lemma}

This is the fact of Lemma~\ref{lem:heat}(b), in the notation of this section.

\begin{proof}
Since $P_{t+s}=\exp(-\frac s{2n}\partial_z^2)P_t$, the first statement is Lemma~\ref{lem:heat}(b) applied to the real polynomial $p=P_t$, for any real $t$, with $\alpha=\frac s{2n}$. For the second, choose around each nonreal zero of $p$, of multiplicity $\mu$, a closed disk that meets neither $\R$ nor the other zeros of $p$; by Hurwitz's theorem it contains exactly $\mu$ zeros of $p_m$ for large $m$, all nonreal. Applied to $p_m=P_{t_m}$ with $t_m\to t$, which have the leading coefficient of $P$, the second statement shows that $\{t\ge0:Z(P_t)=0\}$ is closed.
\end{proof}

\begin{lemma}[Local form of a collision]\label{lem:localform}
Let $t_0\in\R$, and let $x_0\in\R$ be a zero of $P_{t_0}$ of multiplicity $m\ge2$. Let $\rho>0$ be such that $x_0$ is the only zero of $P_{t_0}$ in $\{|z-x_0|\le\rho\}$. Then for all $t\ne t_0$ close enough to $t_0$, $P_t$ has exactly $m$ zeros in $\{|z-x_0|<\rho\}$, and
\begin{enumerate}
\item[(a)] if $t<t_0$, exactly $2\lfloor m/2\rfloor\ (\ge2)$ of them are nonreal;
\item[(b)] if $t>t_0$, all $m$ of them are real and simple.
\end{enumerate}
\end{lemma}

\begin{proof}
Let $b_k=P_{t_0}^{(k)}(x_0)/k!$, so that $b_k=0$ for $k<m$, $b_m\ne0$, and $P_{t_0}(x_0+w)=\sum_kb_kw^k$. Write $t=t_0+\epsilon n\sigma^2$ with $\epsilon\in\{1,-1\}$ and $\sigma>0$, so that $P_t=\exp(-\frac{\epsilon\sigma^2}2\partial^2)P_{t_0}$. Since $\exp(-a\partial^2)w^k=\sum_{i\le k/2}\frac{(-a)^i}{i!}\frac{k!}{(k-2i)!}w^{k-2i}$, we get, exactly,
\[
P_t(x_0+\sigma v)=\sum_{k\ge m}b_k\sigma^k\,\mathcal H_k^\epsilon(v),\qquad\mathcal H_k^{\epsilon}(v)=\sum_{i\le k/2}\Bigl(-\frac\epsilon2\Bigr)^i\frac{k!}{i!\,(k-2i)!}v^{k-2i},
\]
so that $\mathcal H_k^{+}=\He_k$ and $\mathcal H_k^{-}(v)=i^{-k}\He_k(iv)$. Hence $\sigma^{-m}P_t(x_0+\sigma v)\to b_m\mathcal H_m^\epsilon(v)$ uniformly on compact sets as $\sigma\downarrow0$ (the sum is finite). The zeros of $\He_m$ are real and simple, $\lambda_1<\dots<\lambda_m$; they are symmetric about $0$, and $0$ is one of them if and only if $m$ is odd. The zeros of $\mathcal H_m^-$ are $-i\lambda_1,\dots,-i\lambda_m$.

By Hurwitz's theorem, for $t$ close to $t_0$ the disk $\{|z-x_0|<\rho\}$ contains exactly $m$ zeros of $P_t$. Choose pairwise disjoint small disks around the $m$ zeros of $\mathcal H_m^\epsilon$. For small $\sigma$, again by Hurwitz's theorem, each of them contains exactly one zero of the real polynomial $v\mapsto P_t(x_0+\sigma v)$, and these $m$ zeros of $P_t$ lie in $\{|z-x_0|<\rho\}$; so they are all of its zeros there. If $\epsilon=-1$, the disks around $-i\lambda_k$ with $\lambda_k\ne0$ can be taken disjoint from $\R$; this gives $2\lfloor m/2\rfloor$ nonreal zeros. If moreover $m$ is odd, the disk around $0$, taken symmetric about $\R$, contains one zero, which is real because nonreal zeros of a real polynomial come in conjugate pairs. If $\epsilon=+1$, each disk is centred at a real point and taken symmetric about $\R$, so its unique zero is real, and the $m$ zeros are distinct, hence simple.
\end{proof}

\begin{remark}[On Lemma~\ref{lem:localform}]\label{rem:localform}
Lemma~\ref{lem:localform} is classical in substance: $\mathcal H_k^\pm$ are the heat polynomials of \cite{RW}, and the same local form at a zero of order $m$ of the de Bruijn--Newman family is \cite[Proposition~3.1(ii)]{PM15}, by the same computation. The lemma describes the zeros near a multiple zero, but it does not say that multiple zeros are exceptional. In the one-parameter family $t\mapsto P_t$, a real multiple zero at a positive time is forced whenever $P_0$ has nonreal zeros: $P_t$ is real-rooted for large $t$ (Lemma~\ref{lem:heat}(c)), and at its real-rooting time it has a real multiple zero, by the first step of the proof of Lemma~\ref{lem:onedrop}(b) below. A real double zero ($m=2$) persists under small perturbations of the coefficients of $P$: since $\partial_tP_t=-\frac1{2n}\partial_z^2P_t$, the Jacobian of $(t,z)\mapsto(P_t(z),\partial_zP_t(z))$ at $(t_0,x_0)$ has determinant $-P_{t_0}''(x_0)^2/(2n)\ne0$, and the implicit function theorem applies. So multiple zeros cannot be excluded by a general-position argument. Such an argument does predict their type: a real triple zero imposes three conditions on the two parameters $(t,z)$, and Lemma~\ref{lem:notriple} makes this rigorous for random coefficients through Lemma~\ref{lem:conditioning}. That there is only one collision is not a general-position statement; it comes from the monotone count of Lemma~\ref{lem:heatcount}, through Lemma~\ref{lem:onedrop}.
\end{remark}

\begin{lemma}[One drop of two means one collision]\label{lem:onedrop}
Let $0\le T_1<T_2$ with $Z(P_{T_1})=2$ and $Z(P_{T_2})=0$, and let $\Lambda=\inf\{t\ge0:\ Z(P_t)=0\}$. Then:
\begin{enumerate}
\item[(a)] $T_1<\Lambda\le T_2$; $Z(P_t)=2$ for $T_1\le t<\Lambda$, and $Z(P_t)=0$ for $t\ge\Lambda$;
\item[(b)] $P_\Lambda$ has exactly one multiple zero; it is real, of multiplicity $2$ or $3$;
\item[(c)] for every $t\in(T_1,\infty)\setminus\{\Lambda\}$ all zeros of $P_t$ are simple.
\end{enumerate}
In particular, $\Lambda$ is the only time in $(T_1,\infty)$ at which $P_t$ has a multiple zero.
\end{lemma}

\begin{proof}
(a) By Lemma~\ref{lem:heatcount}, $Z(P_t)\le2$ for $t\ge T_1$ and $Z(P_t)=0$ for $t\ge T_2$, and the set $\{t\ge0:Z(P_t)=0\}$ is closed and upward closed. It contains $T_2$, so it is $[\Lambda,\infty)$ with $\Lambda\le T_2$. It does not contain $T_1\ge0$, so $\Lambda>T_1$. For $T_1\le t<\Lambda$ we have $Z(P_t)\ge1$, hence $Z(P_t)\ge2$, because nonreal zeros come in conjugate pairs, hence $Z(P_t)=2$.

(c) Let $t'>T_1$, $t'\ne\Lambda$, and suppose that $P_{t'}$ has a multiple zero. It is not nonreal: by (a), $Z(P_{t'})\le2$, while a nonreal zero of multiplicity at least $2$ and its conjugate would give $Z(P_{t'})\ge4$. So it is a real zero $x_0$ of multiplicity $m\ge2$. Choose $\rho$ as in Lemma~\ref{lem:localform}, so small that the disk about $x_0$ is disjoint from small disks about the nonreal zeros of $P_{t'}$, which keep their nonreal zeros, with multiplicity, for $t$ near $t'$ (Hurwitz's theorem). Lemma~\ref{lem:localform}(a) then gives $Z(P_t)\ge Z(P_{t'})+2$ for $t<t'$ close to $t'$. If $t'<\Lambda$, then $Z(P_{t'})=2$, so $Z(P_t)\ge4$ for some $t\in(T_1,t')$, contradicting (a). If $t'>\Lambda$, then $Z(P_t)\ge2$ for some $t\in(\Lambda,t')$, again contradicting (a).

(b) By (a), $Z(P_\Lambda)=0$. If all zeros of $P_\Lambda$ were simple, hence real, then, as in the proof of Lemma~\ref{lem:constancy}, disjoint disks symmetric about $\R$ around them would each contain exactly one zero of $P_t$, necessarily real, for $t$ close to $\Lambda$; so $Z(P_t)=0$ for $t<\Lambda$ close to $\Lambda$, contradicting (a). So $P_\Lambda$ has a multiple zero, and it is real. If $x_1,\dots,x_l$ are its real multiple zeros, of multiplicities $m_1,\dots,m_l$, Lemma~\ref{lem:localform}(a), applied in disjoint disks, gives $2=Z(P_t)\ge\sum_{i=1}^l2\lfloor m_i/2\rfloor$ for $t<\Lambda$ close to $\Lambda$. Hence $l=1$ and $m_1\in\{2,3\}$.
\end{proof}

\begin{remark}[Simple zeros after real-rootedness]\label{rem:simple-after}
For every real polynomial $P$ of degree $n$, with $\Lambda=\inf\{t\ge0:Z(P_t)=0\}$, all zeros of $P_t$ are simple for every $t>\Lambda$. Indeed, a multiple zero of $P_{t'}$ with $t'>\Lambda$ is real, since $Z(P_{t'})=0$, and by Lemma~\ref{lem:localform}(a) it would give nonreal zeros at times in $(\Lambda,t')$, contrary to Lemma~\ref{lem:heatcount}. This fact is well known.
\end{remark}

\subsection{The bulk inequality at rate \texorpdfstring{$1/n$}{1/n}}\label{subsec:collision-rate}
For a real $2\times2$ matrix $X$ let $s_{\min}(X)=\min_{\|u\|=1}\|Xu\|$, its smallest singular value. Recall that $s_{\min}(X)=|\det X|/\|X\|$ for $X\ne0$, and $s_{\min}(X)\ge s_{\min}(Y)-\|X-Y\|$. The matrices $G_{n,j}$, $G_j$, $H_{n,M}[a]$ and $H_a$ are those of Section~\ref{sec:bulk}, with the variable $y$.

\begin{proposition}[The bulk inequality at rate $1/n$]\label{prop:bulkrate}
Let $(a_s)_{s\in S}$ be a family of real sequences, indexed by an arbitrary set $S$, with $|a_j(s)|\le C_0R_0^{-j}$ for all $j\ge0$ and $s\in S$, where $R_0>1$. Let $\kappa>0$ with $e^\kappa<R_0$, put $L=2\cosh\kappa$ and $\vartheta=e^\kappa/R_0$, and let $M_n=\lceil\ell\log n\rceil$ with $\ell\log(1/\vartheta)\ge2$. There are $C_H,C_W<\infty$ and $n_3$, depending only on $C_0,R_0,\kappa,\ell$, such that for all $n\ge n_3$, $s\in S$ and real $y$ with $|y|\le L$:
\begin{enumerate}
\item[(a)] $\|H_{n,M_n}[a_s](y)-H_{a_s}(y)\|\le C_H/n$;
\item[(b)] with $Q_n^s=\sum_{j=0}^na_j(s)h_{n-j}$ and $x=\sqrt n\,y$,
\[
\Bigl\|\Bigl(Q_n^s(x),\frac{Q_n^s{}'(x)}{\sqrt n}\Bigr)\Bigr\|\ \ge\ \Bigl(s_{\min}\bigl(H_{a_s}(y)\bigr)-\frac{C_W}n\Bigr)\,\|v_n(x)\| .
\]
\end{enumerate}
Hypotheses (H2)--(H4) are not needed.
\end{proposition}

\begin{proof}
Let $\kappa_1=\min(\kappa,1)$, and let $n_3\ge2$ be such that $(M_n+2)^2\le\kappa_1n$ for $n\ge n_3$; then $1\le M_n\le n-1$, and $H_{a_s}$ is defined on $[-L,L]$ by Lemma~\ref{lem:limit}, with $\kappa_0=\kappa$.

(a) We have $H_{n,M_n}[a_s]-H_{a_s}=\sum_{j\le M_n}a_j(s)(G_{n,j}-G_j)-\sum_{j>M_n}a_j(s)G_j$. For $j\le M_n$ we have $(j+2)^2\le\kappa_1n$, and Lemma~\ref{lem:head-rate}, applied to the two rows of $G_{n,j}-G_j$, gives $\|G_{n,j}(y)-G_j(y)\|\le c_2(j+1)^3e^{j\kappa}/n$ with $c_2=2c_1+3\le33/\kappa_1$; here and below the operator norm is bounded by the sum of the absolute values of the entries. So the first sum has norm at most $\frac{C_0c_2}n\sum_{j\ge0}(j+1)^3\vartheta^j$. For $j\ge2$, which covers every $j>M_n$, Lemma~\ref{lem:chebyshev} gives $|A_j|\le(j-1)e^{(j-2)\kappa}$, $|B_j|,|A_{j+1}|\le je^{(j-1)\kappa}$ and $|B_{j+1}|\le(j+1)e^{j\kappa}$, so $\|G_j(y)\|\le4(j+1)e^{j\kappa}$. Hence the second sum has norm at most
\[
4C_0\sum_{j>M_n}(j+1)\vartheta^j\le\frac{4C_0}{(1-\vartheta)^2}(M_n+2)\vartheta^{M_n}\le\frac{4C_0}{(1-\vartheta)^2}\,\frac{M_n+2}{n^2},
\]
since $\vartheta^{M_n}\le n^{-\ell\log(1/\vartheta)}\le n^{-2}$. This proves (a).

(b) Let $E=Q_n^s-\sum_{j\le M_n}a_j(s)h_{n-j}$. By \eqref{eq:head},
\[
\Bigl(Q_n^s(x),\frac{Q_n^s{}'(x)}{\sqrt n}\Bigr)^T=H_{n,M_n}[a_s](y)\,v_n(x)+\Bigl(E(x),\frac{E'(x)}{\sqrt n}\Bigr)^T .
\]
The tail estimate in the proof of Proposition~\ref{prop:bulk} uses only (H1), Lemma~\ref{lem:cd} and the Cauchy--Schwarz inequality, and it gives
\[
\Bigl\|\Bigl(E(x),\frac{E'(x)}{\sqrt n}\Bigr)\Bigr\|\le\eps_n\|v_n(x)\|,\qquad\eps_n=\sqrt2\,C_0(1-R_0^{-2})^{-1/2}R_0^{-M_n}\sqrt{n(1+L/2)} .
\]
Since $\ell\log R_0>\ell\log(1/\vartheta)\ge2$, we have $\eps_n\le C'n^{1/2-2}\le C'/n$ with $C'$ depending only on $C_0,R_0,\kappa$. By (a), $s_{\min}(H_{n,M_n}[a_s](y))\ge s_{\min}(H_{a_s}(y))-C_H/n$. Take $C_W=C_H+C'$.
\end{proof}

The power $(j+1)^3$ in Lemma~\ref{lem:head-rate} is not optimal and does not matter, since any polynomial weight is summable against $\vartheta^j$. Since $\sum_j(j+1)^3\vartheta^j\le6(1-\vartheta)^{-4}$, the proof gives $C_H=C_0(1-\vartheta)^{-4}\bigl(198/\kappa_1+4\sup_{n\ge n_3}(M_n+2)/n\bigr)$, which is polynomial in $1/\kappa_1$ and $(1-\vartheta)^{-1}$.

\subsection{The limit matrix near the collision}\label{subsec:collision-limit}
Assume (G1)--(G3). The constants chosen here are kept until the end of the proof of Theorem~\ref{thm:collision}. The zeros of $q$ in $\{|u|<\rho_q\}$ are isolated and $R_*<\rho_q$; on the circle $\{|u|=R_*\}$ the only zeros are $u_*$ and $\bar u_*$, by (G2) and (G3). Hence there is $\delta_1\in(0,T_*/4)$ such that
\begin{enumerate}
\item[(C1)] $(T_*-2\delta_1)^{-1/2}<\rho_q$, and
\item[(C2)] the only zeros of $q$ in the closed annulus $\mathfrak A=\{(T_*+2\delta_1)^{-1/2}\le|u|\le(T_*-2\delta_1)^{-1/2}\}$ are $u_*$ and $\bar u_*$.
\end{enumerate}
Put $I_1=[T_*-\delta_1,T_*+\delta_1]$ and $a^{(t)}(w)=q(w/\sqrt t)$. Fix $\varrho_1\in((T_*-\delta_1)^{-1/2},\rho_q)$ and put $R_0=\varrho_1(T_*-\delta_1)^{1/2}>1$. By Cauchy's estimate there is $C_0$ with $|\eta_j|\le C_0\varrho_1^{-j}$, hence
\begin{equation}\label{eq:collision-H1}
|a_j^{(t)}|=|\eta_j|\,t^{-j/2}\le C_0R_0^{-j}\qquad(t\in I_1,\ j\ge0).
\end{equation}
Fix $\kappa>0$ with $e^\kappa<R_0$ and
\begin{equation}\label{eq:collision-kappa}
e^{2\kappa}\le\min\Bigl(\frac{T_*-\delta_1}{T_*-2\delta_1},\ \frac{T_*+2\delta_1}{T_*+\delta_1}\Bigr),
\end{equation}
and put $L=2\cosh\kappa>2$. By \eqref{eq:collision-kappa}, if $t\in I_1$ and $e^{-\kappa}\le|w|\le e^\kappa$, then $w/\sqrt t\in\mathfrak A$. Hence, by (C2), for $t\in I_1$ the only zeros of $a^{(t)}$ in $\{e^{-\kappa}\le|w|\le e^\kappa\}$ are $\sqrt t\,u_*$ and $\sqrt t\,\bar u_*$; they have modulus $\sqrt{t/T_*}$, and they lie on the unit circle if and only if $t=T_*$. By \eqref{eq:collision-H1}, Lemma~\ref{lem:limit} applies to every $a^{(t)}$, $t\in I_1$, with $\kappa_0=\kappa$.

\begin{lemma}[Degeneration at the collision point]\label{lem:degenerate}
Let $t\in I_1$ and $y\in[-L,L]$.
\begin{enumerate}
\item[(a)] $\det H_{a^{(t)}}(y)\ge0$, with equality if and only if $(t,y)=(T_*,y_*)$.
\item[(b)] There is $c_*>0$, depending only on $q$, such that
\[
s_{\min}\bigl(H_{a^{(t)}}(y)\bigr)\ \ge\ c_*\bigl(|t-T_*|+|y-y_*|\bigr).
\]
\end{enumerate}
\end{lemma}

\begin{proof}
(a) If $|y|\le2$, write $y=2\cos\theta$ with $\theta\in[0,\pi]$. By Lemma~\ref{lem:limit}, $\det H_{a^{(t)}}(y)=|a^{(t)}(e^{i\theta})|^2=|q(t^{-1/2}e^{i\theta})|^2$. Since $t^{-1/2}e^{i\theta}\in\mathfrak A$, by (C2) this vanishes if and only if $t^{-1/2}e^{i\theta}\in\{u_*,\bar u_*\}$, that is, if and only if $t=T_*$ and $\theta=\theta_*$. (The zero $e^{-i\theta_*}$ of $a^{(T_*)}$ on the circle gives the same point $y_*$, since $y=w+w^{-1}$ does not distinguish $w$ from $w^{-1}=\bar w$.) If $2<|y|\le L$, then $y=w+w^{-1}$ with $w=r(y)$ real and $e^{-\kappa}\le|w|<1$, and $\det H_{a^{(t)}}(y)=a^{(t)}(w)\,a^{(t)}(w^{-1})$ by Lemma~\ref{lem:limit}. Both arguments are real with moduli in $[e^{-\kappa},e^\kappa]$, so neither is a zero of $a^{(t)}$. As a function of $w$ the product is continuous and nonvanishing on $[e^{-\kappa},1]$ and on $[-1,-e^{-\kappa}]$, and at $w=\pm1$ it equals $a^{(t)}(\pm1)^2>0$; so it is positive.

(b) \emph{Local bound.} Let $F(t,\theta)=a^{(t)}(e^{i\theta})=q(t^{-1/2}e^{i\theta})$, a real-analytic function near $(T_*,\theta_*)$ with $F(T_*,\theta_*)=q(u_*)=0$. With $u=t^{-1/2}e^{i\theta}$,
\[
\partial_tF=-\frac{u\,q'(u)}{2t},\qquad\partial_\theta F=i\,u\,q'(u),\qquad\text{so that at }(T_*,\theta_*):\quad\partial_tF=-\frac{\gamma_*}{2T_*},\quad\partial_\theta F=i\gamma_*,
\]
where $\gamma_*=u_*q'(u_*)\ne0$, because $u_*\ne0$ is a simple zero by (G2). The two derivatives are $\R$-linearly independent, one being a real and the other an imaginary multiple of $\gamma_*$; precisely, $|\gamma_*(-\alpha/(2T_*)+i\beta)|\ge2c_F(|\alpha|+|\beta|)$ for real $\alpha,\beta$, with $2c_F=\frac12|\gamma_*|\min(1,\frac1{2T_*})$. By differentiability there is $r_1\in(0,\min(\delta_1,\frac12\theta_*,\frac12(\pi-\theta_*))]$ such that
\[
|F(t,\theta)|\ge c_F\bigl(|t-T_*|+|\theta-\theta_*|\bigr)\qquad(|t-T_*|\le r_1,\ |\theta-\theta_*|\le r_1),
\]
and there $\sin\theta\ge s_0:=\sin(\frac12\min(\theta_*,\pi-\theta_*))>0$.

For $y=2\cos\theta$ with $\theta\in(0,\pi)$ write $H_{a^{(t)}}(y)=\begin{pmatrix}\mathcal A&\mathcal B\\-\mathcal B&\mathcal A+y\mathcal B\end{pmatrix}$ as in Lemma~\ref{lem:limit}, with $\mathcal A,\mathcal B$ real and $\mathcal A+e^{i\theta}\mathcal B=F(t,\theta)$. Then $|\mathcal B|=|\im F|/\sin\theta\le|F|/\sin\theta$ and $|\mathcal A|=|\re F-\mathcal B\cos\theta|\le2|F|/\sin\theta$, so $\|H_{a^{(t)}}(y)\|\le2|\mathcal A|+(2+|y|)|\mathcal B|\le8|F|/\sin\theta$. Since $\det H_{a^{(t)}}(y)=|F|^2$ and $s_{\min}=|\det|/\|\cdot\|$,
\[
s_{\min}\bigl(H_{a^{(t)}}(2\cos\theta)\bigr)\ge\tfrac18|F(t,\theta)|\sin\theta\ge\tfrac18s_0c_F\bigl(|t-T_*|+|\theta-\theta_*|\bigr)\ge\tfrac1{16}s_0c_F\bigl(|t-T_*|+|y-y_*|\bigr)
\]
for $|t-T_*|\le r_1$ and $|\theta-\theta_*|\le r_1$, using $|y-y_*|=2|\cos\theta-\cos\theta_*|\le2|\theta-\theta_*|$. Since $y_*\in(-2,2)$ and $\theta\mapsto2\cos\theta$ is a homeomorphism near $\theta_*$, this bound holds on $V=\{(t,y):|t-T_*|<r_1,\ |y-y_*|<r_2\}$ for some $r_2>0$.

\emph{Global bound.} The map $(t,y)\mapsto H_{a^{(t)}}(y)$ is continuous on $I_1\times[-L,L]$, by \eqref{eq:collision-H1}, the bound for $|A_j|+|B_j|$ in Lemma~\ref{lem:transfer} and the Weierstrass test; hence so is $s_{\min}(H_{a^{(t)}}(y))$. On the compact set $(I_1\times[-L,L])\setminus V$ we have $s_{\min}>0$ by (a), so $s_{\min}\ge c_3$ there for some $c_3>0$, while $|t-T_*|+|y-y_*|\le\delta_1+2L$. Take $c_*=\min\bigl(\frac1{16}s_0c_F,\ c_3/(\delta_1+2L)\bigr)$.
\end{proof}

At $(T_*,y_*)$ the whole limit matrix vanishes: $\mathcal A+e^{i\theta_*}\mathcal B=q(u_*)=0$ with $\mathcal A,\mathcal B$ real and $\sin\theta_*\ne0$ forces $\mathcal A=\mathcal B=0$. So near the collision the local picture is decided by the $O(1/n)$ corrections to the head, which is why Section~\ref{sec:firstorder} needs more than Lemma~\ref{lem:degenerate}.

\subsection{Localisation, uniformly in \texorpdfstring{$t$}{t}}\label{subsec:collision-localise}
Keep the constants of Section~\ref{subsec:collision-limit}.

\begin{proposition}[Localisation of real multiple zeros]\label{prop:localise}
There are $K_1<\infty$ and $n_5$ such that for all $n\ge n_5$ and $t\in I_1$:
\begin{enumerate}
\item[(a)] \emph{(bulk)} for real $y$ with $|y|\le L$ and $x=\sqrt n\,y$,
\[
\Bigl\|\Bigl(Q_n^{(t)}(x),\frac{Q_n^{(t)}{}'(x)}{\sqrt n}\Bigr)\Bigr\|\ \ge\ c_*\Bigl(|t-T_*|+|y-y_*|-\frac{K_1}n\Bigr)\|v_n(x)\| ;
\]
\item[(b)] \emph{(real exterior)} every real multiple zero $y$ of $y\mapsto Q_n^{(t)}(\sqrt n\,y)$ satisfies $|y|<L$;
\item[(c)] every real multiple zero $y$ of $y\mapsto Q_n^{(t)}(\sqrt n\,y)$ satisfies $|t-T_*|+|y-y_*|\le K_1/n$.
\end{enumerate}
\end{proposition}

So the bulk estimate of Proposition~\ref{prop:bulk} holds uniformly in $t\in I_1$, with a lower constant that degenerates linearly at the collision point instead of being bounded below; and in the exterior only real zeros need to be controlled.

\begin{proof}
(a) Fix $\ell$ with $\ell\log(R_0e^{-\kappa})\ge2$. By \eqref{eq:collision-H1}, Proposition~\ref{prop:bulkrate} applies to the family $(a^{(t)})_{t\in I_1}$, whose polynomials $Q_n^s$ are the $Q_n^{(t)}$. Combined with Lemma~\ref{lem:degenerate}(b) it gives (a) for $n\ge n_3$, with $K_1=C_W/c_*$.

(b) By \eqref{eq:collision-H1} and Lemma~\ref{lem:colleague}, for every $t\in I_1$ all zeros of $y\mapsto Q_n^{(t)}(\sqrt n\,y)$ lie in $\{|y|\le\varrho_0\}$, where $\varrho_0=2+|\eta_0|^{-1}C_0(1-R_0^{-2})^{-1/2}$. If $\varrho_0<L$ there is nothing to prove. Otherwise let $J=[-\varrho_0,-L]\cup[L,\varrho_0]$, $\eps'=(L-2)/2$ and $K'=\{y\in\C:\dist(y,J)\le\eps'\}$, a compact subset of $\C\setminus[-2,2]$. Put
\[
\Phi(t,y)=a^{(t)}(r(y)),\qquad\Phi_n(t,y)=\frac{Q_n^{(t)}(\sqrt n\,y)}{h_n(\sqrt n\,y)} .
\]
Proposition~\ref{prop:exterior}, whose only hypothesis is (H1), applies to the family $(a^{(t)})_{t\in I_1}$ (after an affine change of parameter from $I_1$ to $[0,1]$) by \eqref{eq:collision-H1}, and gives $\sup_{I_1\times K'}|\Phi_n-\Phi|\to0$. Both functions are analytic in $y$ on a neighbourhood of $K'$, because $h_n(\sqrt n\,\cdot)$ has no zeros off $[-2,2]$ (Lemma~\ref{lem:resolvent}) and $|r|<1<\sqrt t\,\rho_q$. Every point of $J$ is the centre of a disk of radius $\eps'$ contained in $K'$, so by Cauchy's estimate $\sup_{I_1\times J}|\partial_y\Phi_n-\partial_y\Phi|\to0$.

Let $t\in I_1$ and $y\in J$. Then $w=r(y)$ is real with $0<|w|\le r(L)=e^{-\kappa}$. If $\Phi(t,y)=0$, then $u=w/\sqrt t$ is a real zero of $q$ with $|u|<(T_*-\delta_1)^{-1/2}<(T_*-2\delta_1)^{-1/2}$. By (C2), $u\notin\mathfrak A$, so $|u|<(T_*+2\delta_1)^{-1/2}<R_*$, and $u$ is a simple zero of $q$ by (G3). Hence $(a^{(t)})'(w)=t^{-1/2}q'(u)\ne0$ and $\partial_y\Phi(t,y)=(a^{(t)})'(w)\,r'(y)\ne0$, since $r'$ has no zeros. So $|\Phi|+|\partial_y\Phi|$ is continuous and positive on the compact set $I_1\times J$, with a positive minimum $m_1$. For large $n$, $|\Phi_n-\Phi|+|\partial_y\Phi_n-\partial_y\Phi|<m_1$ on $I_1\times J$, so $\Phi_n(t,\cdot)$ and $\partial_y\Phi_n(t,\cdot)$ have no common zero on $J$. A real multiple zero $y_0\in J$ of $Q_n^{(t)}(\sqrt n\,\cdot)$ would be one, since $h_n(\sqrt n\,y_0)\ne0$.

(c) Let $n_5$ be at least $n_3$ and the threshold in (b). By (b), $|y|<L$. At a multiple zero the left side of (a) vanishes, while $v_n(\sqrt n\,y)\ne0$ by Lemma~\ref{lem:zeros}; so $|t-T_*|+|y-y_*|\le K_1/n$ by (a).
\end{proof}

\subsection{Proof of the collision theorem}\label{subsec:collision-proof}
Let $\delta=\delta_1$ and $t_\pm=T_*\pm\delta$.

\emph{The two anchor times.} The series $a^{(t_\pm)}$ is admissible (Definition~\ref{def:admissible}). Its coefficients are real; its radius of convergence is $\sqrt{t_\pm}\,\rho_q>1$ by (C1); and $a^{(t_\pm)}(0)=\eta_0\ne0$. It has a zero on the unit circle if and only if $q$ has a zero of modulus $t_\pm^{-1/2}$; such a zero would lie in $\mathfrak A$ and have modulus different from $R_*$, which (C2) excludes. Its zeros in $\D$ are the points $\sqrt{t_\pm}\,u$ with $q(u)=0$ and $|u|<t_\pm^{-1/2}$, and such a $u$ is either in $\mathfrak A$, hence $u\in\{u_*,\bar u_*\}$, a simple zero by (G2), or of modulus less than $(T_*+2\delta_1)^{-1/2}<R_*$, hence real by the definition of $R_*$ and simple by (G3). The nonreal zeros of $a^{(t_-)}$ in $\D$ are exactly $\sqrt{t_-}\,u_*$ and $\sqrt{t_-}\,\bar u_*$, since $R_*<t_-^{-1/2}$, and $a^{(t_+)}$ has none, since $t_+^{-1/2}<R_*$. By Theorem~\ref{thm:main} and \eqref{eq:collision-S1} there is $n_7$ such that for $n\ge n_7$
\begin{equation}\label{eq:collision-anchor}
N_n(t_-)=2,\qquad N_n(t_+)=0,\qquad\text{and all zeros of }P_{n,t_-}\text{ are simple.}
\end{equation}

\emph{The collision.} Let $n_0=\max(n_5,n_7)$ and $n\ge n_0$. Apply Lemma~\ref{lem:onedrop} to $P=P_n$, which has degree $n$ because $\eta_0\ne0$, with $T_1=t_->0$ and $T_2=t_+$; its $\Lambda$ is $\Lambda_n$. It gives $t_-<\Lambda_n\le t_+$; $N_n(t)=2$ for $t_-\le t<\Lambda_n$ and $N_n(t)=0$ for $t\ge\Lambda_n$; exactly one multiple zero $z_n$ of $P_{n,\Lambda_n}$, real, of multiplicity $2$ or $3$; and simple zeros for all $t\in(t_-,\infty)\setminus\{\Lambda_n\}$. Together with \eqref{eq:collision-anchor} at $t=t_-$ this proves (c) and the first statement of (d); on $[t_-,\Lambda_n)$ the two nonreal zeros are simple, hence a conjugate pair.

\emph{The place and the time.} By \eqref{eq:collision-S1}, $y_n=z_n/\sqrt{\Lambda_n}$ is a real multiple zero of $Q_n^{(\Lambda_n)}(\sqrt n\,\cdot)$, and $\Lambda_n\in I_1$. By Proposition~\ref{prop:localise}(c), $|\Lambda_n-T_*|+|y_n-y_*|\le K_1/n$. Since $|y_*|<2$ and $|\sqrt{\Lambda_n}-\sqrt{T_*}|\le|\Lambda_n-T_*|/(2\sqrt{T_*-\delta})$,
\[
|z_n-z_*|\le\sqrt{\Lambda_n}\,|y_n-y_*|+|y_*|\,\bigl|\sqrt{\Lambda_n}-\sqrt{T_*}\bigr|\le\Bigl(\sqrt{T_*+\delta}+\frac1{\sqrt{T_*-\delta}}\Bigr)\frac{K_1}n .
\]
This proves (a) and (b) with $K=K_1\bigl(1+\sqrt{T_*+\delta}+(T_*-\delta)^{-1/2}\bigr)$. The last statement of (d) is Proposition~\ref{prop:localise}(c). \qed

\begin{remark}[On the proof]\label{rem:collision-proof}
(1) Proposition~\ref{prop:localise}(c) shows that all real zeros of $P_{n,t}$ are simple for $K_1/n<|t-T_*|\le\delta$. So Lemma~\ref{lem:constancy} on $[T_*-\delta,T_*-2K_1/n]$ and on $[T_*+2K_1/n,T_*+\delta]$, with \eqref{eq:collision-anchor}, already gives $|\Lambda_n-T_*|\le2K_1/n$. Lemma~\ref{lem:heatcount} is needed for (b)--(d): without it, the count could a priori go from $2$ to $4$ and back inside the window.

(2) Theorem~\ref{thm:main} is used only at the two fixed times $t_\pm$, and the bulk estimate only to locate the collision; no uniform control of nonreal zeros near the bulk is needed. Nothing about the size or phase of $h_n(\sqrt n\,y)$ in the bulk is used beyond $v_n(x)\ne0$, since Lemma~\ref{lem:cd} controls the tail relative to $\|v_n(x)\|$. In particular no Plancherel--Rotach asymptotics enter; in Section~\ref{sec:firstorder} they are used only to put the phase of the first-order law in closed form (Proposition~\ref{prop:phase}), on which Theorem~\ref{thm:jointlaw} rests.

(3) For a fixed $q$ the theorem allows a triple zero at $\Lambda_n$. Lemma~\ref{lem:notriple} below excludes it almost surely for random coefficients, using the bottom coefficients $\eta_n,\eta_{n-1},\eta_{n-2}$ of $P_n$, which change with $n$; a deterministic argument under (G1)--(G3) cannot use them.
\end{remark}

\subsection{Genericity for random coefficients}\label{subsec:collision-generic}
In this subsection $q$ is the Gaussian power series \eqref{eq:kac}, and all statements refer to the completed probability space. We show that (G1)--(G3) hold almost surely and that triple collisions do not occur. Apart from the existence of nonreal zeros and the simplicity of the zeros (Proposition~\ref{prop:nonreal} and Lemma~\ref{lem:resample}), the lemmas use only that the coefficients are independent and have densities and that $q$ has radius $1$; see Remark~\ref{rem:collision-laws} for other coefficient laws.

\begin{lemma}[Conditioning]\label{lem:conditioning}
Let $(\eta_j)_{j\ge0}$ be independent real random variables. Let $k\ge1$, let $c=(\eta_{i_1},\dots,\eta_{i_k})$ for distinct indices $i_1,\dots,i_k$, each $\eta_{i_l}$ having a density, and let $g$ be the sequence of the remaining coefficients. Let $\mathcal M$ be a second countable smooth manifold of dimension less than $k$, let $\mathcal G$ be a Borel set with $\PP(g\in\mathcal G)=1$, and let $\Psi:\mathcal M\times\mathcal G\to\R^k$ be jointly Borel on $\mathcal M\times\mathcal G$ and $C^1$ in $m\in\mathcal M$ for each fixed $g$. Then $\PP\bigl(\exists m\in\mathcal M:\ c=\Psi(m,g)\bigr)=0$.
\end{lemma}

\begin{proof}
For fixed $g\in\mathcal G$, the set $\Psi(\mathcal M,g)$ is a countable union of images of compact sets under $C^1$, hence Lipschitz, maps from open subsets of $\R^d$, $d<k$; so it has Lebesgue measure zero in $\R^k$. The set $S=\{(c',g'):c'\in\Psi(\mathcal M,g')\}$ is the projection of the Borel set $\{(m,c',g'):c'=\Psi(m,g')\}\subset\mathcal M\times\R^k\times\mathcal G$, hence analytic and so universally measurable. The vector $c$ has a density on $\R^k$ and is independent of $g$, so Fubini's theorem for the completed product measure gives $\PP((c,g)\in S)=\E\bigl[\PP(c\in\Psi(\mathcal M,g)\mid g)\bigr]=0$.
\end{proof}

\begin{lemma}[Distinct moduli]\label{lem:moduli}
Almost surely:
\begin{enumerate}
\item[(a)] no real zero and no nonreal zero of $q$ in $\D$ have the same modulus;
\item[(b)] if $u,u'$ are nonreal zeros of $q$ in $\D$ with $u'\notin\{u,\bar u\}$, then $|u|\ne|u'|$.
\end{enumerate}
Consequently, almost surely $q$ satisfies (G1)--(G3) with $\rho_q=1$.
\end{lemma}

\begin{proof}
Almost surely $q$ has radius of convergence $1$ and $\eta_0\ne0$: for every $\eps>0$, $\sum_j\PP(|\eta_j|>e^{\eps j})<\infty$, and $|\eta_j|\ge1$ for infinitely many $j$, by the two Borel--Cantelli lemmas. In (a) and (b) let $\mathcal G$ be the Borel set of sequences $(g_j)$ with $\limsup_j|g_j|^{1/j}\le1$, which has full measure. The maps $\Psi$ below are pointwise limits of continuous functions of $(m,g)$ (take partial sums of the series), hence jointly Borel, and they are smooth in $m$.

(a) Let $c=(\eta_0,\eta_1,\eta_2)$, $g_3(w)=\sum_{j\ge3}\eta_jw^j$, and let $\mathcal M_1=\{(x,u):x\in(-1,1)\setminus\{0\},\ u\in\D\setminus\R,\ |u|=|x|\}$, a smooth manifold of dimension $2$, parametrised by $(\pm r,re^{i\varphi})$. For $(x,u)\in\mathcal M_1$ the real-linear map $\Theta_{x,u}:\R^3\to\R\times\C$, $(c_0,c_1,c_2)\mapsto(p(x),p(u))$ with $p(w)=c_0+c_1w+c_2w^2$, is injective: a real polynomial of degree at most $2$ that vanishes at $x$ and $u$ vanishes at the three distinct points $x,u,\bar u$. So $\Theta_{x,u}$ is an isomorphism, depending smoothly on $(x,u)$. If $q(x)=q(u)=0$, then $c=\Psi(x,u,g):=-\Theta_{x,u}^{-1}(g_3(x),g_3(u))$. Lemma~\ref{lem:conditioning} with $k=3$ applies. (A zero at $0$ is excluded by $\eta_0\ne0$.)

(b) Replacing $u,u'$ by their conjugates if necessary, we may assume that $\im u>0$, $\im u'>0$ and $u\ne u'$. Let $c=(\eta_0,\dots,\eta_3)$, $g_4(w)=\sum_{j\ge4}\eta_jw^j$, and let $\mathcal M_2=\{(u,u')\in\D^2:\im u>0,\ \im u'>0,\ u\ne u',\ |u|=|u'|\}$, a smooth manifold of dimension $3$. The real-linear map $\Theta_{u,u'}:\R^4\to\C^2$, $(c_0,\dots,c_3)\mapsto(p(u),p(u'))$ with $p(w)=\sum_{i\le3}c_iw^i$, is injective, because a real cubic in its kernel vanishes at the four distinct points $u,u',\bar u,\bar u'$. If $q(u)=q(u')=0$, then $c=-\Theta_{u,u'}^{-1}(g_4(u),g_4(u'))$, and Lemma~\ref{lem:conditioning} with $k=4$ applies.

\emph{(G1)--(G3).} Almost surely, in addition, $q$ has a nonreal zero in $\D$ (Proposition~\ref{prop:nonreal}) and all zeros of $q$ in $\D\setminus\{0\}$ are simple (Lemma~\ref{lem:resample}). This gives (G1) with $\rho_q=1$; $R_*$ is attained, since the nonreal zeros of modulus at most that of a given nonreal zero $u_0\in\D$ are finitely many. By (b), the nonreal zeros of modulus $R_*$ form exactly one conjugate pair, which is (G2), and by (a) no real zero has modulus $R_*$, which completes (G3).
\end{proof}

\begin{lemma}[No real triple zeros]\label{lem:notriple}
Almost surely, for no $n\ge2$ and no $t\in\R$ does $P_{n,t}$ have a real zero of multiplicity at least $3$.
\end{lemma}

\begin{proof}
Fix $n\ge2$; a countable union of null events is null. The terms $k=0,1,2$ of $P_n$ in \eqref{eq:reversed} are $\eta_n+\eta_{n-1}\sqrt n\,z+\eta_{n-2}nz^2/\sqrt2$, and $\exp(-\frac t{2n}\partial_z^2)$ maps $nz^2$ to $nz^2-t$ and fixes $1$ and $z$. Hence
\[
P_{n,t}(z)=\eta_n+\eta_{n-1}\sqrt n\,z+\eta_{n-2}\,\frac{nz^2-t}{\sqrt2}+R_{n,t}(z),
\]
where $R_{n,t}(z)$ is a polynomial in $(t,z)$ whose coefficients depend only on $\eta_0,\dots,\eta_{n-3}$ ($R_{n,t}=0$ if $n=2$). A real zero $z$ of multiplicity at least $3$ at time $t$ means $P_{n,t}=\partial_zP_{n,t}=\partial_z^2P_{n,t}=0$ at $z$, that is,
\[
\Theta_{t,z}\begin{pmatrix}\eta_n\\\eta_{n-1}\\\eta_{n-2}\end{pmatrix}=-\begin{pmatrix}R_{n,t}(z)\\\partial_zR_{n,t}(z)\\\partial_z^2R_{n,t}(z)\end{pmatrix},\qquad\Theta_{t,z}=\begin{pmatrix}1&\sqrt n\,z&(nz^2-t)/\sqrt2\\0&\sqrt n&\sqrt2\,nz\\0&0&\sqrt2\,n\end{pmatrix}.
\]
The matrix $\Theta_{t,z}$ is invertible, so $(\eta_n,\eta_{n-1},\eta_{n-2})=\Psi(t,z)$ for a map $\Psi$ on $\mathcal M=\R^2$ that is jointly continuous in $(t,z)$ and $(\eta_0,\dots,\eta_{n-3})$, and polynomial in $(t,z)$. Lemma~\ref{lem:conditioning} with $k=3$ applies.
\end{proof}

\emph{The collision point as a fixed rule.} For each $n\ge1$ fix a Borel map $F_n:\R^{n+1}\to\R$ such that $F_n(\eta_0,\dots,\eta_n)$ is a real multiple zero of $P_{n,\Lambda_n}$ whenever $\eta_0\ne0$ and such a zero exists, and $F_n=0$ otherwise. Such a map exists; for instance, take the largest real multiple zero. Indeed, $\Lambda_n$ is a Borel function of $(\eta_0,\dots,\eta_n)$ (proof of Theorem~\ref{cor:weyl}), hence so is the coefficient vector of $P_{n,\Lambda_n}$. Among coefficient vectors of real polynomials of degree $n$, those with a real multiple zero form a relatively closed set, and on it the largest real multiple zero is upper semicontinuous: the zeros stay bounded under small perturbations at fixed degree, and a limit of real multiple zeros is a real multiple zero of the limit. For $n\ge n_0$ in Theorem~\ref{thm:collision}, $F_n(\eta_0,\dots,\eta_n)=z_n$.

\begin{corollary}[The final collision for the Gaussian power series]\label{cor:collision-gauss}
Let $q$ be the Gaussian power series \eqref{eq:kac}, and $z_n=F_n(\eta_0,\dots,\eta_n)$. Almost surely $q$ satisfies (G1)--(G3) with $\rho_q=1$, and the conclusions (a)--(d) of Theorem~\ref{thm:collision} hold with random $\delta$, $K$ and $n_0$, and with $z_n$ a double zero of $P_{n,\Lambda_n}$ for $n\ge n_0$. In particular, almost surely
\[
\sup_{n\ge1}n\,|\Lambda_n-T_*|<\infty\qquad\text{and}\qquad\sup_{n\ge1}n\,\bigl|z_n-2R_*^{-1}\cos\theta_*\bigr|<\infty .
\]
\end{corollary}

\begin{proof}
On the intersection of the almost sure events of Lemmas~\ref{lem:moduli} and~\ref{lem:notriple}, Theorem~\ref{thm:collision} applies to the now fixed $q$, and the multiplicity of $z_n$ is $2$ for $n\ge n_0$. For $n<n_0$, $\Lambda_n<\infty$ by Lemma~\ref{lem:heat}(c), and $z_n$ is a finite number.
\end{proof}

\begin{remark}[Other coefficient laws]\label{rem:collision-laws}
Theorem~\ref{thm:collision} is deterministic. In Lemmas~\ref{lem:conditioning}--\ref{lem:notriple} the Gaussian law enters only through independence, the existence of densities, radius $1$, and, for (G1)--(G3), Proposition~\ref{prop:nonreal} and Lemma~\ref{lem:resample}. So let the $\eta_j$ be independent and identically distributed with a law that has a density and finite variance. Then $q$ has radius $1$ almost surely (by the Borel--Cantelli argument above: $\E\log^+|\eta_0|<\infty$, and $|\eta_j|$ exceeds a fixed positive number for infinitely many $j$), $\eta_0\ne0$ almost surely, $q$ has nonreal zeros by Proposition~\ref{prop:nonreal}, and its zeros in $\D\setminus\{0\}$ are simple by Lemma~\ref{lem:resample}. Lemmas~\ref{lem:moduli} and~\ref{lem:notriple} hold with the same proofs, and so do Corollary~\ref{cor:collision-gauss} and Corollary~\ref{cor:collision-weyl} below, the Weyl polynomials having coefficients of the same law. Finite variance is used only for the existence of nonreal zeros; under a density and $\E\log^+|\eta_0|<\infty$ alone, the conclusions of Corollary~\ref{cor:collision-gauss} hold almost surely on the event that $q$ has a nonreal zero in $\D$. For laws with atoms, such as the Rademacher law, Lemma~\ref{lem:conditioning} is not available. Simplicity of the zeros still holds (Lemma~\ref{lem:resample}), but we do not know whether almost surely the nonreal zeros of modulus $R_*$ form one conjugate pair and no real zero has modulus $R_*$, or whether triple collisions are excluded.
\end{remark}

\subsection{The Weyl polynomials}\label{subsec:collision-weyl}
For the Weyl polynomials \eqref{eq:weyl} with standard Gaussian coefficients, built from a single sequence $(\xi_k)$, only statements about laws can be transferred, as in Theorem~\ref{cor:weyl}. Let $\Lambda_n^{\rm std}$ be the real-rooting time of $G_{n,t}$. Since $G_n$ is the polynomial $P_n$ of \eqref{eq:reversed} for the coefficients $(\eta_0,\dots,\eta_n)=(\xi_n,\dots,\xi_0)$, the number $z_n^{\rm std}=F_n(\xi_n,\dots,\xi_0)$ is a real multiple zero of $G_{n,\Lambda_n^{\rm std}}$ whenever $\xi_n\ne0$ and such a zero exists.

\begin{corollary}[The final collision for Weyl polynomials]\label{cor:collision-weyl}
For each $n\ge1$ there is a Gaussian power series $\tilde q_n$, defined on the probability space of $(\xi_k)$ enlarged by independent standard Gaussian variables, with threshold $\tilde T_{*,n}$ and nonreal zero $\tilde R_{*,n}e^{i\tilde\theta_{*,n}}$, $0<\tilde\theta_{*,n}<\pi$, of minimal modulus, such that the families
\[
n\bigl(\Lambda_n^{\rm std}-\tilde T_{*,n}\bigr),\qquad n\bigl(z_n^{\rm std}-2\tilde R_{*,n}^{-1}\cos\tilde\theta_{*,n}\bigr)\qquad(n\ge1)
\]
are tight. For every $n$ the law of $\tilde T_{*,n}$ is that of $T_*$. With probability tending to one as $n\to\infty$, $z_n^{\rm std}$ is the unique multiple zero of $G_{n,\Lambda_n^{\rm std}}$, and it is a double zero.
\end{corollary}

\begin{proof}
Let $(\xi'_j)_{j\ge1}$ be independent standard Gaussian variables, independent of $(\xi_k)$. For fixed $n$ put $\tilde\eta_j=\xi_{n-j}$ for $0\le j\le n$ and $\tilde\eta_j=\xi'_j$ for $j>n$, and $\tilde q_n(w)=\sum_{j\ge0}\tilde\eta_jw^j$. Then $(\tilde\eta_j)_{j\ge0}$ is an independent standard Gaussian sequence, and the polynomial built from it by \eqref{eq:reversed} in degree $n$ is $\sum_k\tilde\eta_{n-k}(\sqrt nz)^k/\sqrt{k!}=G_n(z)$. Hence $G_{n,t}$ is the heat evolution $\tilde P_{n,t}$ built from $\tilde q_n$. The vector $(\Lambda_n^{\rm std},z_n^{\rm std},\tilde T_{*,n},\tilde R_{*,n},\tilde\theta_{*,n})$ is obtained from $(\tilde\eta_j)$ by the same Borel maps as $(\Lambda_n,z_n,T_*,R_*,\theta_*)$ from $(\eta_j)$, the rule $F_n$ being the same on both sides, so the two vectors have the same law. By Corollary~\ref{cor:collision-gauss}, $D=\sup_mm\bigl(|\Lambda_m-T_*|+|z_m-2R_*^{-1}\cos\theta_*|\bigr)$ is almost surely finite, and therefore
\[
\PP\Bigl(n\bigl(|\Lambda_n^{\rm std}-\tilde T_{*,n}|+|z_n^{\rm std}-2\tilde R_{*,n}^{-1}\cos\tilde\theta_{*,n}|\bigr)>C\Bigr)\le\PP(D>C)\xrightarrow[C\to\infty]{}0
\]
uniformly in $n$. Likewise, the probability that $z_n^{\rm std}$ is the unique multiple zero of $G_{n,\Lambda_n^{\rm std}}$ and is double equals the probability of the same event for $P_n$, which is at least $\PP(n_0\le n)\to1$, with $n_0$ from Corollary~\ref{cor:collision-gauss}.
\end{proof}

Corollary~\ref{cor:collision-weyl} gives tightness only; the first-order law of $n(\Lambda_n-T_*)$ is the subject of Section~\ref{sec:firstorder}. A rate for the convergence of the law of $\Lambda_n^{\rm std}$ to that of $T_*$ on $[1+\eps,\infty)$, which does not use $D$, and a bound for the tail of the variable $D$ of the proof are given in the companion paper \cite{CompA}. % MERGE: was [II, thm:rate] and Section subsec:Dtail of Paper III


\section{The first-order law}\label{sec:firstorder}

Throughout this section $q(w)=\sum_{j\ge0}\eta_jw^j$ is a real power series with radius of convergence $\rho_q$ that
satisfies (G1)--(G3) of Section~\ref{sec:collision}, and $u_*=R_*e^{i\theta_*}$ ($0<\theta_*<\pi$) and $T_*=R_*^{-2}$
are as there. The polynomials $P_n$, $P_{n,t}$ ($t\in\R$) and $Q_n^{(t)}=\sum_{j\le n}a^{(t)}_jh_{n-j}$, with
$a^{(t)}_j=\eta_jt^{-j/2}$, and the real-rooting time $\Lambda_n$ are built from $q$ as in \eqref{eq:reversed} and
\eqref{eq:PQ}. By Theorem~\ref{thm:collision}, for $n\ge n_0$ all zeros of $P_{n,\Lambda_n}$ are real and exactly one of
them, $z_n$, is multiple. We determine $n(\Lambda_n-T_*)$ and $n(z_n-2R_*^{-1}\cos\theta_*)$ up to $o(1)$.

\subsection*{Notation}
Put $\omega=\sin\theta_*$, $w_*=e^{i\theta_*}$, $y_*=2\cos\theta_*$, $x_*=\sqrt n\,y_*$, $z_*=2R_*^{-1}\cos\theta_*$,
$\gamma_*=u_*q'(u_*)$ (nonzero by (G2)), $\rho_*=u_*q''(u_*)/q'(u_*)$, and
\[
W=\frac{T_*}{2\omega^2}=\frac1{2(\im u_*)^2},\qquad
\tau_c=\frac{T_*}2\cot\theta_*\,\im\rho_*=\frac{\re(u_*)\im(\rho_*)}{2|u_*|^2\im u_*},
\]
$\xi_r=-\frac12\cos\theta_*\re\rho_*$ and $\beta=-i\gamma_*/(2\omega)$. The \emph{carrier} and the \emph{carrier phase} are
\begin{align*}
\mathfrak c_n(x)&=h_n(x)+i\,\frac{h_n(x)\cos\theta_*-h_{n-1}(x)}{\sin\theta_*},\\
c_n&=\arg\mathfrak c_n(x_*)+\arg\gamma_*-\frac\pi2-\frac14\sin2\theta_*\re\rho_*\pmod{2\pi};
\end{align*}
$\mathfrak c_n(x_*)\ne0$, since otherwise $h_n(x_*)=h_{n-1}(x_*)=0$, contrary to Lemma~\ref{lem:zeros}.

Let $J(s)=s-\sin s\cos s$. Since $J'(s)=2\sin^2s$ vanishes only on $\pi\Z$, $J$ is an increasing real-analytic
bijection of $\R$; it is odd, and $J(s+\pi)=J(s)+\pi$. For $c\in\R$ put
\[
s_c=J^{-1}(c),\qquad \zeta_c=\sin s_c\cos s_c,\qquad \lambda_c=-\cos^2s_c .
\]
As $s_{c+\pi}=s_c+\pi$, $\zeta_c$ and $\lambda_c$ are $\pi$-periodic in $c$. Let $\tilde c\in[-\frac\pi2,\frac\pi2)$
be congruent to $c$ modulo $\pi$ and $\Psi(c)=2|s_{\tilde c}|\in[0,\pi]$. Since $J$ is odd and $J(\frac\pi2)=\frac\pi2$,
$\Psi=\Psi(c)$ is the root in $[0,\pi]$ of Kepler's equation
\begin{equation}\label{eq:kepler}
\Psi-\sin\Psi=2\dist(c,\pi\Z),
\end{equation}
and $\cos2s_c=\cos\Psi(c)$, $\sin2s_c=\operatorname{sgn}(\tilde c)\sin\Psi(c)$. We write $\Psi_n=\Psi(c_n)$ and
\[
d_n=\max\bigl(n^{-1},\dist(c_n,\pi\Z)\bigr).
\]

The \emph{window} is described as follows: for real $\tau,\xi$ put
\begin{equation}\label{eq:window}
t=T_*+\frac\tau n,\quad y=y_*+\frac\xi n,\quad x=\sqrt n\,y=x_*+\frac\xi{\sqrt n},\quad
\zeta=\omega(\xi-\xi_r),\quad \lambda=\frac{\tau-\tau_c-W}{2W},
\end{equation}
so that $y$ is the Hermite variable and $z=\sqrt t\,y$ the heat variable; $(\tau,\xi)\mapsto(\zeta,\lambda)$ is an affine
bijection that does not depend on $n$. The \emph{model} is
\begin{equation}\label{eq:model}
f_{\lambda,c}(\zeta)=\zeta\cos(\zeta+c)+\lambda\sin(\zeta+c)\qquad(\lambda,c,\zeta\in\R).
\end{equation}
Since $f_{\lambda,c+\pi}=-f_{\lambda,c}$, everything below depends on $c_n$ only modulo $\pi$.

\begin{theorem}[The first-order law]\label{thm:firstorder}
Assume (G1)--(G3). There are $C$ and $n_1\ge n_0$ such that for every $n\ge n_1$ the following hold.
\begin{enumerate}
\item[(a)] \emph{(Time.)} $\bigl|n(\Lambda_n-T_*)-\tau_c+W\cos\Psi_n\bigr|\le Cn^{-1}d_n^{-1/3}\le Cn^{-2/3}$.
\item[(b)] \emph{(Place.)} With $y_n=z_n/\sqrt{\Lambda_n}$, the multiple zero in the Hermite variable,
\[
\Bigl|n(y_n-y_*)-\xi_r-\operatorname{sgn}(\tilde c_n)\frac{\sin\Psi_n}{2\omega}\Bigr|\le Cn^{-1}d_n^{-2/3}\le Cn^{-1/3},
\]
and in the heat variable
\[
n(z_n-z_*)=-\frac{\re(u_*)\re(\rho_*)}{2|u_*|^2}+\operatorname{sgn}(\tilde c_n)\frac{\sin\Psi_n}{2\im u_*}
+\re(u_*)\,n(\Lambda_n-T_*)+O\bigl(n^{-1}d_n^{-2/3}\bigr).
\]
\item[(c)] \emph{(Double zero.)} If $\dist(c_n,\pi\Z)\ge C/n$, then $z_n$ is a double zero of $P_{n,\Lambda_n}$.
\end{enumerate}
\end{theorem}

So, up to $o(1)$, $n(\Lambda_n-T_*)$ is the function $\tau_c-W\cos\Psi_n$ of the single phase $c_n$, with values in
$[\tau_c-W,\tau_c+W]$, and $(n(\Lambda_n-T_*),n(y_n-y_*))$ lies near the ellipse
$((\tau-\tau_c)/W)^2+(2\omega(\xi-\xi_r))^2=1$. By Proposition~\ref{prop:phase}, $c_n=c_\infty-n\pi\nu+O(1/n)$ with
$\pi\nu=J(\theta_*)$, so $c_n$ rotates by $-\pi\nu$ per degree, and $n(\Lambda_n-T_*)$ does not converge in general
(Theorem~\ref{thm:jointlaw}). The upper end $\tau_c+W$ is the time at which the zero of the leading complex envelope
reaches the real line ($\lambda=0$). The error is $O(1/n)$ when $\dist(c_n,\pi\Z)$ is bounded below, where the collision
is a fold of the model; the uniform exponents come from the cusps at $c_n\in\pi\Z$ (Remark~\ref{rem:sharp}).

The proof has three steps. In the window, $Q_n^{(t)}(x)$ is a positive factor times $f_{\lambda,c_n}(\zeta)+O(1/n)$, in
$C^2$ in $\zeta$ (Proposition~\ref{prop:window}). By an exact identity, the model has a unique real multiple zero, at
$\lambda=\lambda_{c}$ (Lemma~\ref{lem:model}). A degree argument and Theorem~\ref{thm:collision} identify the multiple zero
of the perturbed model with the last collision (Proposition~\ref{prop:threshold}).

\subsection{The window expansion}\label{sec:window}

\emph{Constants.} Fix $\delta_1\in(0,T_*/4)$ with $(T_*-2\delta_1)^{-1/2}<\rho_q$, put $I_1=[T_*-\delta_1,T_*+\delta_1]$, fix
$\varrho_1\in((T_*-\delta_1)^{-1/2},\rho_q)$ and put $R_0=\varrho_1(T_*-\delta_1)^{1/2}>1$. By Cauchy's estimate there is $C_0$
with
\begin{equation}\label{eq:F-coeff}
|a^{(t)}_j|\le C_0R_0^{-j}\qquad(t\in I_1,\ j\ge0),
\end{equation}
and $a^{(t)}(w)=q(w/\sqrt t)$ converges for $|w|<R_0$. Fix $\kappa>0$ with $e^\kappa<R_0$, put
$\vartheta=e^\kappa/R_0$ and $S_p=\sum_{j\ge0}(j+1)^p\vartheta^j$, let $M_n=\lceil3\log n/\log(1/\vartheta)\rceil$,
so that $\vartheta^{M_n}\le n^{-3}$, and let $n_W\ge4$ be such that $(M_n+2)^4\le n$ for $n\ge n_W$. Let
$\mathcal D_\kappa=\{w+w^{-1}:e^{-\kappa}\le|w|\le e^\kappa\}$, the closed region bounded by the ellipse with foci
$\pm2$ and semi-axes $2\cosh\kappa$, $2\sinh\kappa$; its interior is the image of the open annulus. Fix $r_0>0$ such
that the disc $D_{2r_0}$, where $D_r=\{y\in\C:|y-y_*|\le r\}$, lies in the interior of $\mathcal D_\kappa$ and does not
meet $(-\infty,-2]\cup[2,\infty)$. Let $w(y)=\frac12(y+i(4-y^2)^{1/2})$ (principal root), which is analytic off
$(-\infty,-2]\cup[2,\infty)$, with $w(y)+w(y)^{-1}=y$, $w(y_*)=w_*$ and $w'=w/(w-w^{-1})$, so that
$w'(y_*)=w_*/(2i\omega)$. For $y\in D_{2r_0}$ the two Joukowski preimages $w(y)^{\pm1}$ of $y$ have modulus in
$(e^{-\kappa},e^\kappa)$. Below, $C$ denotes constants that depend only on $q$ and on these choices, and $C_\Xi$ constants
that depend also on a parameter $\Xi$.

\emph{Continuants.} For complex $y$ let $A_{n,j}(y),B_{n,j}(y)$ be the polynomials of Section~\ref{sec:transfer}
(there the variable is called $z$), $A_j(y)=-U_{j-2}(y/2)$, $B_j(y)=U_{j-1}(y/2)$ as in \eqref{eq:fixed-j}, and
\[
X_{n,j}(y;w)=A_{n,j}(y)+wB_{n,j}(y),\qquad X_j(y;w)=A_j(y)+wB_j(y).
\]
By the proof of Lemma~\ref{lem:limit}, $X_j(w+w^{-1};w)=w^j$ for $w\ne0$.

\begin{lemma}[Second-order continuant expansion]\label{lem:secondorder}
Let $\kappa>0$ and $y\in\mathcal D_\kappa$; write $U_l=U_l(y/2)$.
\begin{enumerate}
\item[(a)] $|U_l|\le U_l(\cosh\kappa)\le(l+1)e^{l\kappa}$ for $l\ge0$; hence $|A_j(y)|+|B_j(y)|\le2(j+1)e^{j\kappa}$.
\item[(b)] Put $A^{(1)}_j=B^{(1)}_j=0$ for $j\le1$ and, for $j\ge2$,
\[
B^{(1)}_j=\sum_{\ell=1}^{j-2}\ell\,U_{\ell-1}U_{j-2-\ell}+\frac{j(j-1)}4U_{j-1},\qquad
A^{(1)}_j=-\sum_{\ell=1}^{j-3}(\ell+1)U_{\ell-1}U_{j-3-\ell}-\frac{j(j-1)}4U_{j-2}.
\]
Then $|A^{(1)}_j(y)|+|B^{(1)}_j(y)|\le j^4e^{j\kappa}$ for every $j\ge0$.
\item[(c)] If $n\ge4$ and $j^2\le n$, then
\[
\Bigl|A_{n,j}(y)-A_j(y)-\frac{A^{(1)}_j(y)}n\Bigr|+\Bigl|B_{n,j}(y)-B_j(y)-\frac{B^{(1)}_j(y)}n\Bigr|
\le\frac{4e^{\coth\kappa}j^6e^{j\kappa}}{n^2}.
\]
\item[(d)] Let $w\in\C\setminus\{0\}$ and $y=w+w^{-1}$. Then $Y_j=A^{(1)}_j(y)+wB^{(1)}_j(y)$ satisfies $Y_0=Y_1=0$ and
\[
Y_j=yY_{j-1}-Y_{j-2}+\tfrac12\bigl((j-1)w^j+(j-2)w^{j-2}\bigr)\qquad(j\ge2).
\]
\end{enumerate}
\end{lemma}

\begin{proof}
(a) Write $y=w+w^{-1}$ with $|w|=e^\sigma$, $|\sigma|\le\kappa$. By \eqref{eq:chebyshev-joukowski} (and by continuity
at $w=\pm1$), $U_l=\sum_{i=0}^lw^{l-2i}$. Pairing $i$ with $l-i$,
$|U_l|\le\sum_ie^{\sigma(l-2i)}=\sum_i\cosh(\sigma(l-2i))\le\sum_i\cosh(\kappa(l-2i))=U_l(\cosh\kappa)$, and
Lemma~\ref{lem:chebyshev} gives the second inequality. For $j\ge1$, $|A_j|+|B_j|\le(j-1)e^{(j-2)\kappa}+je^{(j-1)\kappa}$.

(b) By (a), $\sum_{\ell=1}^{N-1}|U_{\ell-1}U_{N-1-\ell}|\,\ell\le e^{(N-2)\kappa}\sum_{\ell<N}\ell^2(N-\ell)$, and
$\sum_{\ell<N}\ell^2(N-\ell)=N^2(N^2-1)/12\le N^4/8$. With $N=j-1$ for $B^{(1)}_j$ and $N=j-2$, $\ell+1\le2\ell$, for
$A^{(1)}_j$, we get $|B^{(1)}_j|\le(\frac18j^4+\frac14j^3)e^{j\kappa}$ and $|A^{(1)}_j|\le(\frac14j^4+\frac14j^3)e^{j\kappa}$,
and $j^3\le j^4/2$.

(c) For $j\le1$ the left side vanishes, so let $j\ge2$; then $j\le\sqrt n\le n/2$. \emph{A second-order defect bound.}
Let $m\ge2$ and $0\le\delta_\ell\le\bar\delta\le1/m$ for $\ell=1,\dots,m-1$, and put
$T^{(1)}_m=\sum_{\ell=1}^{m-1}\delta_\ell U_{\ell-1}U_{m-\ell-1}$. We claim
\begin{equation}\label{eq:defect2}
\bigl|\Pi_m(y;1-\delta_1,\dots,1-\delta_{m-1})-U_m-T^{(1)}_m\bigr|\le e^{\coth\kappa}m^2(m+1)\bar\delta^2e^{m\kappa}.
\end{equation}
In Lemma~\ref{lem:defect}, $S=\emptyset$ gives $U_m$, and $S=\{\ell\}$ leaves components with $\ell-1$ and $m-\ell-1$
vertices, so the matchings with $|S|=1$ give $T^{(1)}_m$. By (a), the left side of \eqref{eq:defect2} is at most
$\sum_{|S|\ge2}\prod_{\ell\in S}\delta_\ell\prod_CU_{|C|}(\cosh\kappa)\le h(\bar\delta)$, where
$g(s)=\Pi_m(2\cosh\kappa;1-s,\dots,1-s)=\sum_kg_ks^k$ has coefficients
$g_k=\sum_{|S|=k}\prod_CU_{|C|}(\cosh\kappa)\ge0$ and $h(s)=\sum_{k\ge2}g_ks^k$. As $h(s)/s^2$ is nondecreasing,
$h(\bar\delta)\le(m\bar\delta)^2h(1/m)\le m^2\bar\delta^2g(1/m)$. The constant-edge path is $\sqrt{1-s}$ times the unit
path, so $g(s)=(1-s)^{m/2}U_m(\cosh\kappa/\sqrt{1-s})$, and as in the proof of Lemma~\ref{lem:transfer},
$g(1/m)\le(m+1)e^{m\kappa'}$ with $\kappa'\le\kappa+\coth\kappa/m$. This proves \eqref{eq:defect2}.

\emph{Denominators.} Let $L=-\log d_{n,j}=-\frac12\sum_{r=1}^{j-1}\log(1-r/n)$. From $u\le-\log(1-u)\le u+u^2$ on
$[0,\frac12]$, $L=\frac{j(j-1)}{4n}+\eps_d$ with $0\le\eps_d\le\frac{j^3}{6n^2}\le\frac{j^2}{12n}$, so
$0\le L\le\frac{j^2}{3n}\le\frac13$. Write $d_{n,j}^{-1}=e^L=1+D_1+D$ with $D_1=\frac{j(j-1)}{4n}$. Then
$0\le D_1+D=e^L-1\le Le^L\le0.47\,j^2/n$, $0\le D=e^L-1-L+\eps_d\le\frac12L^2e^L+\frac{j^3}{6n^2}\le\frac{j^4}{4n^2}$, and
$d_{n,j}^{-1}\le1.4$.

\emph{The polynomial $B_{n,j}$.} By Lemma~\ref{lem:continuant}, $B_{n,j}=\Pi_{j-1}(y;1-\delta_1,\dots,1-\delta_{j-2})/d_{n,j}$
with $\delta_\ell=\ell/n\le\bar\delta:=(j-2)/n$, and $\bar\delta\le1/(j-1)$ because $(j-1)(j-2)\le n$. Write
$\Pi_{j-1}=U_{j-1}+T_1+T_2$ with $T_1=T^{(1)}_{j-1}$ (and $T_1=T_2=0$ if $j=2$). Then $nT_1=\sum_\ell\ell U_{\ell-1}U_{j-2-\ell}$,
so $B_j+B^{(1)}_j/n=U_{j-1}+T_1+D_1U_{j-1}$ and
\[
B_{n,j}-B_j-\frac{B^{(1)}_j}n=\frac{T_2}{d_{n,j}}+T_1(D_1+D)+U_{j-1}D .
\]
By \eqref{eq:defect2}, $|T_2|\le e^{\coth\kappa}j^5e^{(j-1)\kappa}/n^2$; by the proof of (b), $|T_1|\le j^4e^{(j-3)\kappa}/(8n)$;
and $|U_{j-1}|\le je^{(j-1)\kappa}$. Hence the left side is at most
$(1.4e^{\coth\kappa}j^5+0.06j^6+0.25j^5)e^{j\kappa}/n^2\le2e^{\coth\kappa}j^6e^{j\kappa}/n^2$, as $0.31<0.6e\le0.6e^{\coth\kappa}$.

\emph{The polynomial $A_{n,j}$.} Here $A_{n,j}=-\Pi_{j-2}(y;1-\delta'_1,\dots,1-\delta'_{j-3})/d_{n,j}$ with
$\delta'_\ell=(\ell+1)/n\le(j-2)/n\le1/(j-2)$ (no edges if $j\le3$), and the same algebra gives
$-(A_{n,j}-A_j-A^{(1)}_j/n)=T'_2/d_{n,j}+T'_1(D_1+D)+U_{j-2}D$, with $|T'_2|\le e^{\coth\kappa}j^5e^{(j-2)\kappa}/n^2$,
$|T'_1|\le j^4e^{(j-4)\kappa}/(4n)$ and $|U_{j-2}|\le je^{(j-2)\kappa}$. The bound is now
$(1.4e^{\coth\kappa}j^5+0.12j^6+0.25j^5)e^{j\kappa}/n^2\le2e^{\coth\kappa}j^6e^{j\kappa}/n^2$, as $0.37<0.6e$.
Adding the two bounds gives (c).

(d) The polynomials $X_{n,j}=X_{n,j}(y;w)$ satisfy \eqref{eq:backward}:
$(1-\frac{j-1}n)^{1/2}X_{n,j}=yX_{n,j-1}-(1-\frac{j-2}n)^{1/2}X_{n,j-2}$, with $X_{n,0}=1$, $X_{n,1}=w$. Fix $j\ge2$; $y$
lies in some $\mathcal D_\kappa$, so by (c), $X_{n,i}=w^i+Y_i/n+O(n^{-2})$ for $i\le j$ as $n\to\infty$. Inserting
$(1-k/n)^{1/2}=1-k/(2n)+O(n^{-2})$ and comparing the coefficients of $n^{-1}$ gives
$Y_j-\frac{j-1}2w^j=yY_{j-1}-Y_{j-2}+\frac{j-2}2w^{j-2}$.
\end{proof}

\begin{lemma}[The envelope]\label{lem:envelope}
For $t\in I_1$ and $y\in\mathcal D_\kappa$ let
\[
F_0=\sum_{j\ge0}a^{(t)}_jX_j(y;w_*),\qquad F_1=\sum_{j\ge0}a^{(t)}_j\bigl(A^{(1)}_j(y)+w_*B^{(1)}_j(y)\bigr),\qquad
E_{n,t}=\sum_{j=0}^{M_n}a^{(t)}_jX_{n,j}(y;w_*),
\]
functions of $(t,y)$.
\begin{enumerate}
\item[(a)] The series converge absolutely and uniformly on $I_1\times\mathcal D_\kappa$, also after any number of
$t$-derivatives; $F_0,F_1$ are analytic in $y$ in the interior of $\mathcal D_\kappa$ and smooth in $t$; and
$|E_{n,t}(y)-F_0(t,y)-F_1(t,y)/n|\le C/n^2$ on $I_1\times\mathcal D_\kappa$ for $n\ge n_W$.
\item[(b)] $F_0(T_*,y_*)=0$, $\partial_tF_0(T_*,y_*)=-\gamma_*/(2T_*)$, $\partial_yF_0(T_*,y_*)=\beta$, and
\[
F_1(T_*,y_*)=F_1^*:=\gamma_*\Bigl(\frac1{4\omega^2}-\frac i4\cot\theta_*\,\rho_*\Bigr).
\]
\item[(c)] Let $\Xi>0$. There are $C_\Xi$ and $n'_\Xi$ such that, in the notation \eqref{eq:window}, for $n\ge n'_\Xi$ and
$|\tau|,|\xi|\le \Xi$ the function $\Delta_n=nE_{n,t}(y)-\beta(\zeta-i\lambda)/\omega$ satisfies
$|\partial_\xi^k\Delta_n|\le C_\Xi/n$ for $k=0,1,2$.
\end{enumerate}
\end{lemma}

\begin{proof}
(a) By \eqref{eq:F-coeff} and Lemma~\ref{lem:secondorder}(a),(b), the $j$th terms are at most $2C_0(j+1)\vartheta^j$ and
$C_0j^4\vartheta^j$ on $I_1\times\mathcal D_\kappa$. Since $\partial_t^ka^{(t)}_j$ is $a^{(t)}_j$ times $t^{-k}$ times a
polynomial in $j$ of degree $k$, and $t\ge\frac34T_*$, the $t$-derivatives only add polynomial factors in $j$. This gives
uniform convergence, analyticity by Weierstrass's theorem, and termwise differentiation in $t$. Next,
$E_{n,t}-F_0-F_1/n=\sum_{j\le M_n}a^{(t)}_j\varrho_{n,j}-\sum_{j>M_n}a^{(t)}_j\bigl(X_j+(A^{(1)}_j+w_*B^{(1)}_j)/n\bigr)$ with
$\varrho_{n,j}=X_{n,j}-X_j-(A^{(1)}_j+w_*B^{(1)}_j)/n$. For $n\ge n_W$ and $j\le M_n$ we have $j^2\le n$, so by
Lemma~\ref{lem:secondorder}(c) the first sum is at most $4e^{\coth\kappa}C_0S_6/n^2$. The second is at most
$3C_0\sum_{j>M_n}(j+1)^4\vartheta^j\le3C_0(M_n+2)^4\vartheta^{M_n+1}S_4\le3C_0S_4/n^2$, because
$(M+2+i)^4\le(M+2)^4(1+i)^4$, $\vartheta^{M_n}\le n^{-3}$ and $(M_n+2)^4\le n$.

(b) Let $y\in D_{2r_0}$ and $w=w(y)$. The sequences $A_j+w^{\pm1}B_j$ satisfy the free recurrence with initial values
$1,w^{\pm1}$, so they equal $w^{\pm j}$. Summing against $a^{(t)}_j$ (absolutely, since $|w^{\pm1}|<e^\kappa<R_0$) gives
$\mathcal A_t+w^{\pm1}\mathcal B_t=a^{(t)}(w^{\pm1})$, where $\mathcal A_t=\sum_ja^{(t)}_jA_j$ and
$\mathcal B_t=\sum_ja^{(t)}_jB_j$. Hence
\[
F_0(t,y)=a^{(t)}(w(y))+(w_*-w(y))\mathcal B_t(y),\qquad
\mathcal B_t(y)=\frac{a^{(t)}(w(y))-a^{(t)}(w(y)^{-1})}{w(y)-w(y)^{-1}} .
\]
At $(T_*,y_*)$ we have $w=w_*$, $a^{(T_*)}(w_*)=q(u_*)=0$ and $a^{(T_*)}(w_*^{-1})=q(\bar u_*)=0$; so
$\mathcal B_{T_*}(y_*)=0$ and $F_0(T_*,y_*)=0$. Differentiating, with $w_*T_*^{-1/2}=u_*$,
$\partial_tF_0(T_*,y_*)=\partial_t\bigl[q(w_*t^{-1/2})\bigr]_{t=T_*}=-u_*q'(u_*)/(2T_*)$ and
$\partial_yF_0(T_*,y_*)=(a^{(T_*)})'(w_*)\,w'(y_*)=u_*q'(u_*)/(2i\omega)=\beta$.

Finally $F_1(T_*,y_*)=\sum_ja_jY_j$, with $a_j=a^{(T_*)}_j$ and $Y_j$ as in Lemma~\ref{lem:secondorder}(d) for $w=w_*$.
Let $|w|=1$, $w\ne\pm1$, put
\[
\alpha_2=\frac{w^2+1}{2(w^2-1)},\qquad \alpha_1=\frac{-\frac12(w^2+2)-\alpha_2}{w^2-1},\qquad P(j)=\alpha_2\frac{j(j+1)}2+\alpha_1j,
\]
and $\Delta P(j)=P(j)-P(j-1)=\alpha_2j+\alpha_1$. Then
$w^jP(j)-yw^{j-1}P(j-1)+w^{j-2}P(j-2)=w^j\Delta P(j)-w^{j-2}\Delta P(j-1)$, and the choice of $\alpha_2$ and $\alpha_1$
makes this equal to $\frac12((j-1)w^j+(j-2)w^{j-2})$ (compare the coefficients of $j$ and the constant terms). As the
homogeneous recurrence has the solutions $w^{\pm j}$, $Y_j=w^j(P(j)+k_1)+k_2w^{-j}$ for constants $k_1,k_2$. Since $a$
has radius of convergence greater than $1$,
\[
\sum_ja_jY_j=\frac{\alpha_2}2w^2a''(w)+(\alpha_1+\alpha_2)wa'(w)+k_1a(w)+k_2a(w^{-1}).
\]
At $w=w_*$: $a(w_*^{\pm1})=0$, $w_*a'(w_*)=\gamma_*$, $w_*^2a''(w_*)=u_*^2q''(u_*)=\gamma_*\rho_*$,
$\alpha_2=\cos\theta_*/(2i\sin\theta_*)=-\frac i2\cot\theta_*$, and with $s=w_*^2$,
$\alpha_1+\alpha_2=\frac{(s+1)(s-2)-(s+2)(s-1)}{2(s-1)^2}=\frac{-s}{(s-1)^2}=\frac1{4\omega^2}$, since
$(s-1)^2=-4s\sin^2\theta_*$. This gives $F_1^*$.

(c) Let $n'_\Xi=\max(n_W,\Xi/\delta_1,\Xi/r_0)$ and $n\ge n'_\Xi$. Then $(t,y)$ and the segment joining it to $(T_*,y_*)$ lie in
$I_1\times[y_*-r_0,y_*+r_0]$. By (a) and Cauchy's estimates on discs of radius $r_0$ inside $D_{2r_0}$, the derivatives
of $F_0$ and $F_1$ of order at most $3$ are bounded by $C$ there, and $|\partial_y^k(E_{n,t}-F_0-F_1/n)|\le C/n^2$ for
$k\le2$. Write $nE_{n,t}=nF_0+F_1+n(E_{n,t}-F_0-F_1/n)$. Taylor's formula and (b) give
$nF_0(t,y)=\beta\xi-\gamma_*\tau/(2T_*)+O(\Xi^2/n)$ and $F_1(t,y)=F_1^*+O(\Xi/n)$, so
$nE_{n,t}(y)=\beta\xi-\gamma_*\tau/(2T_*)+F_1^*+O_\Xi(1/n)$. Since $\partial_\xi=n^{-1}\partial_y$,
$\partial_\xi(nE_{n,t})=\partial_yF_0(t,y)+O(1/n)=\beta+O_\Xi(1/n)$ and $\partial_\xi^2(nE_{n,t})=n^{-1}\partial_y^2E_{n,t}=O(1/n)$.
It remains to check that $\beta\xi-\gamma_*\tau/(2T_*)+F_1^*=\beta(\zeta-i\lambda)/\omega$. By \eqref{eq:window},
$\beta(\zeta-i\lambda)/\omega=\beta\xi-\beta\xi_r-i\beta\omega(\tau-\tau_c-W)/T_*$, and $-i\beta\omega/T_*=-\gamma_*/(2T_*)$.
Moreover $-\beta\xi_r=-\frac i4\gamma_*\cot\theta_*\re\rho_*$ and
$\gamma_*(\tau_c+W)/(2T_*)=\gamma_*(\frac14\cot\theta_*\im\rho_*+\frac1{4\omega^2})$, whose sum is $F_1^*$.
\end{proof}

\begin{lemma}[The carrier]\label{lem:carrier}
Let $n\ge2$.
\begin{enumerate}
\item[(a)] For real $x$, $h_n(x)=\re\mathfrak c_n(x)$, $h_{n-1}(x)=\re(w_*\mathfrak c_n(x))$,
$\omega^2|\mathfrak c_n(x)|^2=h_n^2-y_*h_nh_{n-1}+h_{n-1}^2$, and $h_n(x)^2+h_{n-1}(x)^2\le2|\mathfrak c_n(x)|^2$. If
$x=\sqrt n\,(y_*+\xi/n)$ with $\xi$ real, then $\sum_{k<n}h_k(x)^2\le(n\omega^2+|\xi|)|\mathfrak c_n(x)|^2$.
\item[(b)] Let $\Xi>0$ and $n\ge5(1+\Xi)^2/\omega^2$. For real $\xi$ with $|\xi|\le \Xi$ and $x=x_*+\xi/\sqrt n$,
\[
\mathfrak c_n(x)=\mathfrak c_n(x_*)\,e^{(x^2-x_*^2)/4}\,e^{i\omega\xi}\bigl(1+\eps_n(\xi)\bigr),
\]
where $\eps_n$ is smooth, $\eps_n(0)=0$, $|\eps_n|+|\eps_n'|\le5(1+\Xi)^2/(n\omega)\le1$ and $|\eps_n''|\le C_\Xi/n$, with
$C_\Xi$ depending only on $\Xi$ and $\omega$.
\end{enumerate}
\end{lemma}

\begin{proof}
(a) $\re\mathfrak c_n=h_n$ is clear, and $\re(w_*\mathfrak c_n)=h_n\cos\theta_*-(h_n\cos\theta_*-h_{n-1})=h_{n-1}$. Next
$\omega^2|\mathfrak c_n|^2=\omega^2h_n^2+(h_n\cos\theta_*-h_{n-1})^2=h_n^2-y_*h_nh_{n-1}+h_{n-1}^2$. The quadratic form
$a^2-2ab\cos\theta_*+b^2$ has eigenvalues $1\pm\cos\theta_*$, and $\omega^2=(1-|\cos\theta_*|)(1+|\cos\theta_*|)$, so
$h_n^2+h_{n-1}^2\le(1+|\cos\theta_*|)|\mathfrak c_n|^2$. Lemma~\ref{lem:cd} with $z=y_*+\xi/n$ gives
$\sum_{k<n}h_k^2=n(h_n^2-y_*h_nh_{n-1}+h_{n-1}^2)-\xi h_nh_{n-1}=n\omega^2|\mathfrak c_n|^2-\xi h_nh_{n-1}$, and
$|h_nh_{n-1}|\le|\mathfrak c_n|^2$.

(b) Let $u_m(x)=e^{-x^2/4}h_m(x)$. From $\He_m''-x\He_m'+m\He_m=0$ we get $u_n''+(n+\frac12-\frac{x^2}4)u_n=0$, and
\eqref{eq:recurrence} gives $u_n'=\sqrt n\,u_{n-1}-\frac x2u_n$. Put $U(\xi)=u_n(x_*+\xi/\sqrt n)$. Since
$x/(2\sqrt n)=\cos\theta_*+\xi/(2n)$,
\begin{equation}\label{eq:carrier-Z}
e^{-x^2/4}\mathfrak c_n(x)=U-\frac{iU'}\omega-\frac{i\xi U}{2n\omega}.
\end{equation}
Let $Z=U-iU'/\omega$; then $Z(0)=e^{-x_*^2/4}\mathfrak c_n(x_*)\ne0$ by (a). As $x^2/4=n\cos^2\theta_*+\xi\cos\theta_*+\xi^2/(4n)$,
\[
U''+(\omega^2+\epsilon)U=0,\qquad \epsilon(\xi)=\frac1n\Bigl(\frac12-\xi\cos\theta_*-\frac{\xi^2}{4n}\Bigr),\qquad
|\epsilon|\le\frac{1+\Xi}n,\quad|\epsilon'|\le\frac2n .
\]
Since $U=\re Z$, this gives $Z'=i\omega Z+(i\epsilon/\omega)\re Z$. Let $\Upsilon(\xi)=e^{-i\omega\xi}Z(\xi)-Z(0)$, so that
$\Upsilon'=(i\epsilon/\omega)e^{-i\omega\xi}\re Z$. With $\ell=(1+\Xi)\Xi/(n\omega)\le\omega/5$, Gronwall's inequality gives
$|Z|\le e^\ell|Z(0)|\le1.23|Z(0)|$ and then $|\Upsilon|\le1.23\,\ell|Z(0)|$. Define
\[
\eps_n(\xi)=\frac{\Upsilon(\xi)}{Z(0)}-\frac{i\xi e^{-i\omega\xi}U(\xi)}{2n\omega Z(0)} .
\]
By \eqref{eq:carrier-Z}, $Z(0)e^{i\omega\xi}(1+\eps_n)=Z-i\xi U/(2n\omega)=e^{-x^2/4}\mathfrak c_n(x)$, which is the claimed
identity, and $\eps_n(0)=0$. Using $|U|\le|Z|$, $|U'|=\omega|\im Z|\le\omega|Z|$ and $|\Upsilon'|\le1.23(1+\Xi)|Z(0)|/(n\omega)$,
\[
|\eps_n|\le\frac{1.23\Xi^2+1.85\Xi}{n\omega},\qquad|\eps_n'|\le\frac{1.23(1+\Xi)+0.62(1+2\Xi)}{n\omega},
\]
and the sum is at most $5(1+\Xi)^2/(n\omega)\le\omega\le1$. Differentiating once more, with $|Z'|\le2|Z|$ and
$U''=-(\omega^2+\epsilon)U$, gives $|\eps_n''|\le C_\Xi/n$.
\end{proof}

\begin{proposition}[The window expansion]\label{prop:window}
Assume (G1)--(G3). For every $\Xi>0$ there are $n_\Xi$ and $C_\Xi$ such that for $n\ge n_\Xi$ and real $\tau,\xi$ with
$|\tau|,|\xi|\le \Xi$, in the notation \eqref{eq:window},
\[
Q_n^{(t)}(x)=\mathcal P_n(\xi)\bigl[f_{\lambda,c_n}(\zeta)+R_n(\zeta,\lambda)\bigr],\qquad
\mathcal P_n(\xi)=\frac{|\mathfrak c_n(x_*)|\,|\gamma_*|}{2n\omega^2}\,e^{y_*\xi/2+\xi^2/(4n)},
\]
where $R_n$ is a smooth function of $(\zeta,\lambda)$ with $|R_n|+|\partial_\zeta R_n|+|\partial_\zeta^2R_n|\le C_\Xi/n$. The
constants do not depend on $c_n$.
\end{proposition}

\begin{proof}
Let $n_\Xi=\max(n'_\Xi,5(1+\Xi)^2/\omega^2)$ and $n\ge n_\Xi$; then also $M_n+1\le n$ and $\Xi\le n\omega^2$.

\emph{Head representation.} By Lemma~\ref{lem:carrier}(a), $Ah_n+Bh_{n-1}=\re[\mathfrak c_n(A+w_*B)]$ for real $A,B$.
As the $a^{(t)}_j$ and the polynomials $A_{n,j},B_{n,j}$ are real, \eqref{eq:AB} gives
\[
Q_n^{(t)}(x)=\re\bigl[\mathfrak c_n(x)E_{n,t}(y)\bigr]+E^{\rm tail}(x),\qquad E^{\rm tail}(x)=\sum_{j=M_n+1}^na^{(t)}_jh_{n-j}(x).
\]

\emph{The identity.} By Lemmas~\ref{lem:carrier}(b) and~\ref{lem:envelope}(c), with
$(x^2-x_*^2)/4=y_*\xi/2+\xi^2/(4n)$,
\[
\mathfrak c_n(x)=|\mathfrak c_n(x_*)|e^{i\arg\mathfrak c_n(x_*)}e^{\frac{x^2-x_*^2}4}e^{i\omega\xi}(1+\eps_n),\qquad
nE_{n,t}(y)=\frac{|\beta|}\omega e^{i\arg\beta}\Bigl[(\zeta-i\lambda)+\frac{\omega\Delta_n}{\beta}\Bigr].
\]
Since $|\beta|=|\gamma_*|/(2\omega)$, $\arg\beta=\arg\gamma_*-\frac\pi2$ and
$\omega\xi=\zeta+\omega\xi_r=\zeta-\frac14\sin2\theta_*\re\rho_*$, the phases add up to $\zeta+c_n$, and
\[
\mathfrak c_n(x)E_{n,t}(y)=\mathcal P_n(\xi)\,e^{i(\zeta+c_n)}\bigl[(\zeta-i\lambda)+G_n\bigr],\qquad
G_n=\eps_n(\zeta-i\lambda)+(1+\eps_n)\,\omega\Delta_n/\beta .
\]
As $\re[e^{i(\zeta+c)}(\zeta-i\lambda)]=f_{\lambda,c}(\zeta)$, this is the claim with
\[
R_n=\re\bigl[e^{i(\zeta+c_n)}G_n\bigr]+E^{\rm tail}(x)/\mathcal P_n(\xi).
\]
$R_n=Q_n^{(t)}(x)/\mathcal P_n(\xi)-f_{\lambda,c_n}(\zeta)$ is smooth in $(\zeta,\lambda)$, as $Q_n^{(t)}(x)$ is a polynomial in
$(x,t^{-1/2})$ and \eqref{eq:window} is affine.

\emph{Bounds.} At fixed $\lambda$, $\partial_\zeta=\omega^{-1}\partial_\xi$. On the square, $|\zeta-i\lambda|\le C_\Xi$, and
$|\partial_\xi^k\eps_n|,|\partial_\xi^k\Delta_n|\le C_\Xi/n$ for $k\le2$ by Lemmas~\ref{lem:carrier}(b)
and~\ref{lem:envelope}(c); hence $|\partial_\xi^kG_n|\le C_\Xi/n$, and $|\partial_\zeta^k\re[e^{i(\zeta+c_n)}G_n]|\le C_\Xi/n$
uniformly in $c_n$. For the tail, $\partial_\xi=n^{-1/2}\partial_x$ and $h_m'=\sqrt m\,h_{m-1}$ show that
$\partial_\xi^kE^{\rm tail}(x)=\sum_{j>M_n}a^{(t)}_jp_{j,k}h_{n-j-k}(x)$ with weights $p_{j,k}\in[0,1]$ and
all indices below $n$ ($h_{-1}=h_{-2}=0$). By the Cauchy--Schwarz inequality, \eqref{eq:F-coeff} and
Lemma~\ref{lem:carrier}(a),
\[
|\partial_\xi^kE^{\rm tail}(x)|\le\Bigl(\sum_{j>M_n}(a^{(t)}_j)^2\sum_{m<n}h_m(x)^2\Bigr)^{1/2}
\le\frac{C_0R_0^{-M_n}(2n)^{1/2}\omega}{(1-R_0^{-2})^{1/2}}|\mathfrak c_n(x)|\le Cn^{-5/2}|\mathfrak c_n(x)|,
\]
since $R_0^{-M_n}\le\vartheta^{M_n}\le n^{-3}$. By Lemma~\ref{lem:carrier}(b),
$|\mathfrak c_n(x)|\le2|\mathfrak c_n(x_*)|e^{(x^2-x_*^2)/4}=2n\omega\mathcal P_n(\xi)/|\beta|$, and the derivatives of
$1/\mathcal P_n$ are $1/\mathcal P_n$ times polynomials in $y_*/2+\xi/(2n)$ and $1/(2n)$. Hence
$|\partial_\xi^k(E^{\rm tail}/\mathcal P_n)|\le C_\Xi n^{-3/2}$ for $k\le2$.
\end{proof}

The proof uses only that $q$ is real and that $u_*$ is a simple nonreal zero with $|u_*|<\rho_q$; the minimality of $|u_*|$
and (G3) enter only through Theorem~\ref{thm:collision}. No asymptotics of $h_n$ are used: the phase
$\arg\mathfrak c_n(x_*)$ is kept exact. Since the leading envelope $F_0$ vanishes at $(T_*,y_*)$, the window sees its
first derivatives and the first-order envelope $F_1$, together with the carrier phase.

\subsection{The model}\label{sec:model}

For $c\in\R$ let $\Phi_c(\zeta,\lambda)=\bigl(f_{\lambda,c}(\zeta),\partial_\zeta f_{\lambda,c}(\zeta)\bigr)\in\R^2$, with the
Euclidean norm. Let $\Gamma=\{(\zeta,\lambda):\zeta^2+\lambda^2+\lambda=0\}$, the circle of radius $\frac12$ about
$(0,-\frac12)$; the curve $\gamma(s)=(\sin s\cos s,-\cos^2s)$ runs around $\Gamma$ at unit speed with period $\pi$, and
$(\zeta_c,\lambda_c)=\gamma(s_c)$.

\begin{lemma}[The model]\label{lem:model}
Let $c,\zeta,\lambda\in\R$ and $s=\zeta+c$.
\begin{enumerate}
\item[(a)] With $R(s)=\bigl(\begin{smallmatrix}\cos s&\sin s\\-\sin s&\cos s\end{smallmatrix}\bigr)$,
\[
\Phi_c(\zeta,\lambda)=R(s)\bigl[(\zeta,\lambda)-\gamma(s)\bigr],\qquad
|\Phi_c(\zeta,\lambda)|^2=\bigl(J(s)-c\bigr)^2+\bigl(\lambda+\cos^2s\bigr)^2 .
\]
In particular $|\Phi_c(\zeta,\lambda)|\ge\dist((\zeta,\lambda),\Gamma)$.
\item[(b)] $(\zeta_c,\lambda_c)$ is the only zero of $\Phi_c$: $f_{\lambda,c}$ has a real multiple zero if and only if
$\lambda=\lambda_c$, and then only at $\zeta_c$. There $\partial_\zeta^2f_{\lambda,c}=-2\sin s_c$,
$\partial_\zeta^3f_{\lambda,c}=-2\cos s_c$ and $\det D\Phi_c=2\sin^2s_c$; so the zero is double, and regular for
$\Phi_c$, if $c\notin\pi\Z$, and triple if $c\in\pi\Z$.
\item[(c)] All zeros of the entire function $f_{\lambda,c}$ are real if and only if $\lambda\ge\lambda_c$. For
$\lambda>\lambda_c$ they are simple; for $\lambda<\lambda_c$ exactly two are nonreal, a simple conjugate pair, and the
real ones are simple.
\item[(d)] Let $\eps=|\Phi_c(\zeta,\lambda)|\le1$, $d=\max(\eps,\dist(c,\pi\Z))$, $\sigma=|\sin s_c|$ and
$C_J=\sqrt2\,(3\pi^2/2)^{1/6}<2.22$. Then $\sigma^3\ge(12/\pi^3)\dist(c,\pi\Z)$,
\[
|\zeta-\zeta_c|\le\min\bigl\{(\tfrac32\pi^2\eps)^{1/3},\,11.7\,\eps d^{-2/3}\bigr\},\qquad
|\lambda-\lambda_c|\le\min\bigl\{\eps+C_J\eps^{2/3},\,7.1\,\eps d^{-1/3}\bigr\},
\]
and if $\eps<\sigma^3/4$, then $|\zeta-\zeta_c|\le2\eps/\sigma^2$ and $|\lambda-\lambda_c|\le\eps(1+2/\sigma)$.
\end{enumerate}
\end{lemma}

For $c\in\pi\Z$ and $c\in\frac\pi2+\pi\Z$, part (c) is classical: up to the factor $(2/\pi)^{1/2}z^{-1/2}$,
$f_{\lambda,0}(z)$ and $f_{\lambda,\pi/2}(z)$ are the Dini functions $\alpha J_\varrho(z)+zJ_\varrho'(z)$ with $\varrho=\frac12$ and
$\varrho=-\frac12$ and $\alpha=\lambda+\frac12$, whose zeros are all real if and only if $\alpha+\varrho\ge0$ \cite[\S15.25, p.~482]{Wat};
this gives $\lambda_0=-1$ and $\lambda_{\pi/2}=0$. As $c$ runs over $\R/\pi\Z$, the threshold points $(\zeta_c,\lambda_c)$
run once around $\Gamma$, and $(2c,2+2\lambda_c)$ runs along one arch of the cycloid $S\mapsto(S-\sin S,1-\cos S)$, with
$S=2s_c$; its cusps, $c\in\pi\Z$, are the phases with a triple zero. The model is an exact trajectory of the backward
heat flow: $e^uf_{\lambda+2u,c}$ solves $\partial_uF=-\partial_\zeta^2F$, so that (c) describes the real-rooting time of
this flow started from $f_{\lambda,c}$.

\begin{proof}
(a) Since $f=\zeta\cos s+\lambda\sin s$ and $\partial_\zeta f=(1+\lambda)\cos s-\zeta\sin s$, we have
$(f,\partial_\zeta f-\cos s)=R(s)(\zeta,\lambda)$, and $R(s)\gamma(s)=(0,-\cos s)$. As $R(s)$ is orthogonal and
$\zeta-\sin s\cos s=J(s)-c$, the identities follow; and $\gamma(s)\in\Gamma$.

(b) By (a), $\Phi_c=0$ if and only if $J(s)=c$ and $\lambda=-\cos^2s$, that is, $s=s_c$, $\zeta=s_c-c=\zeta_c$ and
$\lambda=\lambda_c$. Inserting these values in $\partial_\zeta^2f=-(2+\lambda)\sin s-\zeta\cos s$ and
$\partial_\zeta^3f=-(3+\lambda)\cos s+\zeta\sin s$ gives $-2\sin s_c$ and $-2\cos s_c$, and the rows of $D\Phi_c$ are
$(\partial_\zeta f,\partial_\lambda f)=(0,\sin s_c)$ and $(\partial_\zeta^2f,\partial_\lambda\partial_\zeta f)=(-2\sin s_c,\cos s_c)$.
Finally $s_c\in\pi\Z$ if and only if $c\in\pi\Z$, since $J(k\pi)=k\pi$.

(c) As $f_{\lambda,c+\pi}=-f_{\lambda,c}$ we may take $0\le c<\pi$. On $|w|=N\pi$, $N\ge1$, one has $|\tan w|\le2$:
for $w=u+iv$ with $|v|\ge1$, $|\tan w|^2\le\coth^2v<2$; if $|v|<1$, then $\dist(u,\pi\Z)<1/\pi$ and
$|\tan w|^2\le(\sin^2u+\sinh^2v)/\cos^2u<2$. So for $N\pi>c+2|\lambda|$, on $|\zeta+c|=N\pi$ we have
$|\lambda\sin(\zeta+c)|<|\zeta\cos(\zeta+c)|$, and by Rouch\'e's theorem $f_{\lambda,c}$ has $2N+1$ zeros in
$\{|\zeta+c|<N\pi\}$, as $\zeta\cos(\zeta+c)$ has. Now count real zeros. Let $0<c<\pi$. On
$I_k=(k\pi-c,(k+1)\pi-c)$ write $w=\zeta+c$ and $f=\sin w\,(h+\lambda)$ with $h(\zeta)=\zeta\cot w$, so that
$h'=(c-J(w))/\sin^2w$; at the endpoints $f=\pm\zeta\ne0$. For $k\ge1$, $\zeta>0$ and
$|\sin w\cos w|\le w-k\pi<\zeta$, so $h$ decreases from $+\infty$ to $-\infty$; for $k\le-1$, similarly $h$ increases from
$-\infty$ to $+\infty$; either way $I_k$ carries exactly one zero, which is simple. On $I_0$, $h\to-\infty$ at both ends
and $h'$ has the sign of $c-J(w)$, so $h$ has the single maximum $h=(s_c-c)\cot s_c=\cos^2s_c=-\lambda_c$, at $\zeta_c$.
Hence $I_0$ carries two simple zeros, one double zero, or none, as $\lambda>\lambda_c$, $\lambda=\lambda_c$ or
$\lambda<\lambda_c$. For $c=0$ the intervals $I_k$, $k\ge1$ and $k\le-2$, are as before; on $(0,\pi)$,
$h=\zeta\cot\zeta$ decreases from $1$ to $-\infty$, $h$ is even, and $\zeta=0$ is a zero with
$f=(1+\lambda)\zeta-(\frac12+\frac\lambda6)\zeta^3+O(\zeta^5)$; so $(-\pi,\pi)$ carries three simple zeros, a triple
zero, or one simple zero, as $\lambda>-1$, $\lambda=-1$ or $\lambda<-1$. In all cases the number of real zeros in
$(-N\pi-c,N\pi-c)$, with multiplicity, is $2N+1$ if $\lambda\ge\lambda_c$ and $2N-1$ if $\lambda<\lambda_c$. Letting
$N\to\infty$ gives (c); a nonreal double zero would come with a conjugate one.

(d) Let $L=|s-s_c|=|\zeta-\zeta_c|$. By (a), $|J(s)-J(s_c)|\le\eps$ and $|\lambda+\cos^2s|\le\eps$. For an interval
$I$ of length $L$, $\int_I2\sin^2u\,du\ge L-\sin L\ge2L^3/(3\pi^2)$ if $L\le\pi$, and $\ge\pi$ if $L\ge\pi$; as $\eps<\pi$,
this gives $L\le(\frac32\pi^2\eps)^{1/3}$. For real $a<b$ with $b-a<\pi$, by the Cauchy--Schwarz inequality,
$|\cos^2b-\cos^2a|\le2\int_a^b|\sin u|\,du\le(2(b-a)(J(b)-J(a)))^{1/2}\le C_J(J(b)-J(a))^{2/3}$, and this also holds
if $b-a\ge\pi$, since then $J(b)-J(a)\ge\pi$; so $|\lambda-\lambda_c|\le\eps+C_J\eps^{2/3}$. If $\eps<\sigma^3/4$: on
the interval between $s_c$ and $s$ of length $\min(L,\sigma/2)$ we have $|\sin u|\ge\sigma/2$, which forces
$L<\sigma/2$ and then $\eps\ge\sigma^2L/2$; moreover $\cos^2s-\cos^2s_c=-\int_{s_c}^s\cot u\cdot2\sin^2u\,du$ with
$|\cot u|\le2/\sigma$, so $|\lambda-\lambda_c|\le\eps+2\eps/\sigma$. Next, with $\tilde s=s_{\tilde c}\in[-\frac\pi2,\frac\pi2)$,
$\sigma=|\sin\tilde s|\ge2|\tilde s|/\pi$ and $\dist(c,\pi\Z)=J(|\tilde s|)\le\frac23|\tilde s|^3$, so
$\sigma^3\ge(12/\pi^3)\dist(c,\pi\Z)$. Finally, if $\dist(c,\pi\Z)>(\pi^3/3)\eps$, then $\sigma^3>4\eps$ and $d=\dist(c,\pi\Z)$,
and the fold bounds give $|\zeta-\zeta_c|\le2(\pi^3/12)^{2/3}\eps d^{-2/3}$ and, as $1\le(\pi/2)^{1/3}d^{-1/3}$,
$|\lambda-\lambda_c|\le((\pi/2)^{1/3}+2(\pi^3/12)^{1/3})\eps d^{-1/3}$; otherwise $d\le(\pi^3/3)\eps$, and the first
bounds give $|\zeta-\zeta_c|\le(\frac32\pi^2)^{1/3}(\pi^3/3)^{2/3}\eps d^{-2/3}$ and
$|\lambda-\lambda_c|\le(1+C_J)\eps^{2/3}\le(1+C_J)(\pi^3/3)^{1/3}\eps d^{-1/3}$. The constants are below $11.7$ and $7.1$.
\end{proof}

\subsection{The threshold}\label{sec:threshold}

Let $B_0=[-1,1]\times[-\frac32,\frac12]$. It is convex and contains $\Gamma\subset[-\frac12,\frac12]\times[-1,0]$, and
$\dist(\partial B_0,\Gamma)=\frac12$. Put $\Xi_0=|\xi_r|+\omega^{-1}+|\tau_c+W|+3W$; then the square $|\tau|,|\xi|\le \Xi_0$ is
mapped by \eqref{eq:window} onto a rectangle containing $B_0$. Let $R_n$ be as in Proposition~\ref{prop:window} for
$\Xi=\Xi_0$, and on $B_0$ put $F_n=f_{\lambda,c_n}+R_n$ and $\Phi_n=(F_n,\partial_\zeta F_n)$, so that
$|\Phi_n-\Phi_{c_n}|\le\bar\eps_n:=C_{\Xi_0}/n$. For $n\ge n_0$ let $\tau_n=n(\Lambda_n-T_*)$ and $\xi_n=n(y_n-y_*)$, with
$y_n=z_n/\sqrt{\Lambda_n}$, and let $(\zeta_n,\lambda_n)$ be the image of $(\tau_n,\xi_n)$ under \eqref{eq:window}.

\begin{proposition}[The last collision in the window]\label{prop:threshold}
For all large $n$, $(\zeta_n,\lambda_n)$ lies in the interior of $B_0$ and $\Phi_n(\zeta_n,\lambda_n)=0$; in particular
$|\Phi_{c_n}(\zeta_n,\lambda_n)|\le\bar\eps_n$.
\end{proposition}

\begin{proof}
\emph{Dictionary.} Let $n\ge n_{\Xi_0}$ and let $(\zeta,\lambda)\in B_0$ be the image of $(\tau,\xi)$. At fixed $\tau$,
$\partial_\xi x=n^{-1/2}$ and $\partial_\xi\zeta=\omega$, so Proposition~\ref{prop:window} gives
$n^{-1/2}Q_n^{(t)}{}'(x)=\mathcal P_n'(\xi)F_n(\zeta,\lambda)+\omega\mathcal P_n(\xi)\partial_\zeta F_n(\zeta,\lambda)$. As
$\mathcal P_n>0$, $\Phi_n(\zeta,\lambda)=0$ if and only if $x$ is a multiple zero of $Q_n^{(t)}$, that is, by
\eqref{eq:PQ}, if and only if $\sqrt t\,y$ is a real multiple zero of $P_{n,t}$.

\emph{Degree.} By Lemma~\ref{lem:model}(a), $|\Phi_c|\ge\frac12$ on $\partial B_0$ for every $c$, and by
Lemma~\ref{lem:model}(b), $(\zeta_c,\lambda_c)\in\Gamma$ is the only zero of $\Phi_c$; for $c\notin\pi\Z$ it is regular
with positive Jacobian, so the Brouwer degree $\deg(\Phi_c,B_0,0)$ is $1$. Since $(c,\zeta,\lambda)\mapsto\Phi_c(\zeta,\lambda)$
is continuous and does not vanish on $\R\times\partial B_0$, homotopy invariance gives $\deg(\Phi_c,B_0,0)=1$ for every
$c$, including $c\in\pi\Z$. The maps $\Phi_{c_n}+\mu(\Phi_n-\Phi_{c_n})$, $0\le\mu\le1$, are continuous and
have modulus at least $\frac12-\bar\eps_n>0$ on $\partial B_0$ for large $n$, so $\deg(\Phi_n,B_0,0)=1$, and $\Phi_n$ has a
zero $(\zeta',\lambda')$ in the interior of $B_0$ \cite[Ch.~1]{De}.

\emph{Identification.} By the dictionary, $P_{n,t'}$ has a real multiple zero, where $t'=T_*+\tau'/n$ and
$|\tau'|\le \Xi_0$. Let $\delta$ be as in Theorem~\ref{thm:collision}. For large $n$, $t'\ge T_*-\delta$, so $t'=\Lambda_n$ by
Theorem~\ref{thm:collision}, and the multiple zero is $z_n$, the only multiple zero of $P_{n,\Lambda_n}$. Hence
$(\zeta',\lambda')=(\zeta_n,\lambda_n)$.
\end{proof}

Proposition~\ref{prop:threshold} uses from Theorem~\ref{thm:collision} only that $\Lambda_n$ is the only time
$t\ge T_*-\delta$ at which $P_{n,t}$ has a real multiple zero, and that $P_{n,\Lambda_n}$ has only one. It does not use
the localisation of the collision within $K/n$ of $(T_*,z_*)$, which it reproves with an explicit window:
$|\tau_n-\tau_c|\le2W$ and $|\xi_n-\xi_r|\le1/\omega$ for large $n$.

\begin{proof}[Proof of Theorem~\ref{thm:firstorder}]
Let $n$ be large, so that $\bar\eps_n\le1$ and Proposition~\ref{prop:threshold} applies. The bounds of
Lemma~\ref{lem:model}(d) increase with $\eps$, and $\max(\bar\eps_n,\dist(c_n,\pi\Z))\ge\min(C_{\Xi_0},1)\,d_n$; so that lemma,
with $c=c_n$, gives
\[
|\lambda_n-\lambda_{c_n}|\le Cn^{-1}d_n^{-1/3},\qquad|\zeta_n-\zeta_{c_n}|\le Cn^{-1}d_n^{-2/3}.
\]
By \eqref{eq:window}, $\tau_n=\tau_c+W+2W\lambda_n$ and $\xi_n=\xi_r+\zeta_n/\omega$, while
$\tau_c+W+2W\lambda_{c_n}=\tau_c-W\cos2s_{c_n}=\tau_c-W\cos\Psi_n$ and
$\zeta_{c_n}=\frac12\sin2s_{c_n}=\frac12\operatorname{sgn}(\tilde c_n)\sin\Psi_n$. This proves (a) and the first part of (b).
Since $z_n=\sqrt{\Lambda_n}\,y_n$, exactly $n(z_n-z_*)=\sqrt{\Lambda_n}\,\xi_n+y_*\tau_n/(\sqrt{\Lambda_n}+\sqrt{T_*})$; as
$\tau_n,\xi_n$ are bounded, this is $\sqrt{T_*}\,\xi_n+y_*\tau_n/(2\sqrt{T_*})+O(1/n)$, and
$\sqrt{T_*}\,\xi_r=-\re(u_*)\re(\rho_*)/(2|u_*|^2)$, $\sqrt{T_*}/(2\omega)=1/(2\im u_*)$, $y_*/(2\sqrt{T_*})=\re u_*$.

(c) Suppose that $\dist(c_n,\pi\Z)\ge3\pi^3C_{\Xi_0}/n$, and let $\sigma=|\sin s_{c_n}|$. Then $\sigma^3\ge36\bar\eps_n$ by
Lemma~\ref{lem:model}(d), and the same lemma gives $|\zeta_n-\zeta_{c_n}|\le2\bar\eps_n/\sigma^2$ and $|\lambda_n-\lambda_{c_n}|\le\bar\eps_n(1+2/\sigma)$.
On $B_0$, which contains $(\zeta_n,\lambda_n)$ and $(\zeta_{c_n},\lambda_{c_n})$, the function
\[
\partial_\zeta^2f_{\lambda,c_n}(\zeta)=-(2+\lambda)\sin(\zeta+c_n)-\zeta\cos(\zeta+c_n)
\]
has partial derivatives bounded by
$\frac92$ in $\zeta$ and by $1$ in $\lambda$. By Lemma~\ref{lem:model}(b) and Proposition~\ref{prop:window}, using
$\sigma\le1$,
\[
|\partial_\zeta^2F_n(\zeta_n,\lambda_n)|\ge2\sigma-5\bar\eps_n\Bigl(\frac2{\sigma^2}+1+\frac2\sigma\Bigr)-\bar\eps_n
\ge2\sigma-\frac{26\bar\eps_n}{\sigma^2}\ge\sigma\Bigl(2-\frac{26}{36}\Bigr)>0 .
\]
So $\zeta_n$ is a zero of $F_n(\cdot,\lambda_n)$ of multiplicity exactly two. As
$Q_n^{(\Lambda_n)}(x)=\mathcal P_n(\xi)F_n(\zeta,\lambda_n)$ with $\mathcal P_n$ smooth and positive and $\zeta$ an affine
function of $x$, $\sqrt n\,y_n$ is a double zero of $Q_n^{(\Lambda_n)}$, and $z_n$ one of $P_{n,\Lambda_n}$.
\end{proof}

\begin{remark}[The exponents]\label{rem:sharp}
For $\upsilon\in\R$ the difference $f_{\lambda,c+\upsilon}-f_{\lambda,c}$ is $O(|\upsilon|)$ on bounded sets, with all its
derivatives and uniformly in $c$; for $\upsilon\asymp1/n$ it is a remainder of the size allowed in
Proposition~\ref{prop:window}. At $c=0$ it moves the multiple zero by $\lambda_\upsilon-\lambda_0=\sin^2s_\upsilon$ and
$\zeta_\upsilon-\zeta_0=\frac12\sin2s_\upsilon$, of orders $|\upsilon|^{2/3}$ and $|\upsilon|^{1/3}$, since $J(s)\sim\frac23s^3$. So no
argument that uses only the conclusion of Proposition~\ref{prop:window} gives better exponents than those of
Theorem~\ref{thm:firstorder} at $d_n\asymp1/n$. In between, $d\lambda_c/dc=\cot s_c$ and
$d\zeta_c/dc=\cos2s_c/(2\sin^2s_c)$; as $\tilde c\to0$ with $n|\tilde c|\to\infty$, a phase shift by $1/n$ moves
$\lambda_c$ and $\zeta_c$ by $(1+o(1))(\frac32|\tilde c|)^{-1/3}/n$ and $(1+o(1))(\frac32|\tilde c|)^{-2/3}/(2n)$. Where
these derivatives vanish, the rate $1/n$ is realised by $\upsilon\sin(\zeta+c)$, which turns $f_{\lambda,c}$ into
$f_{\lambda+\upsilon,c}$, and by translations in $\zeta$. Whether $R_n$ itself has such components at order $1/n$ is a
question about the next order of the expansion.

Numerically, the size of the cusp effect depends on $q$. For eight symbols (the benchmark of
Table~\ref{tab:benchmark}, three other fixed symbols and four samples of the Gaussian power series), with $128$
degrees each in $100\le n\le25607$, we compared the scaled error $n\cdot|n(\Lambda_n-T_*)-(\tau_c-W\cos\Psi(c_n))|$ at
degrees with $\dist(c_n,\pi\Z)<0.05$ with its median at degrees with $\dist(c_n,\pi\Z)>0.3$, and fitted it, for
$n\ge1600$, by $A+B\max(\dist(c_n,\pi\Z),n^{-1})^{-1/3}$. For three symbols, among them the benchmark, the errors near
the cusps are at most three times the median at the folds and the fit gives $B\le0$, so no growth towards the cusps is
detectable for $n\ge1600$. (For the benchmark, the values $1.62$ at $n=110$ and $1.45$ at $n=408$ are about twice the
largest value at the folds, $0.76$; at the three degrees near the cusps above $1600$ they are $0.5$ to $1.0$.) For the
other five, three of them Gaussian samples, $B>0$
and the errors near the cusps are $4.5$ to $14$ times the median at the folds, consistent with the growth
$\dist(c_n,\pi\Z)^{-1/3}$ that Theorem~\ref{thm:firstorder} allows. The samples near the cusps are small ($1$ to $6$
degrees per symbol), and $\dist(c_n,\pi\Z)$ does not go much below $10^{-2}$, so the regime $d_n\asymp1/n$ is not
reached. So for particular symbols the next-order remainder appears to have a negligible component in the direction
that moves the cusp, but not in general.
\end{remark}

\subsection{The carrier phase}\label{sec:phase}

\begin{proposition}[The carrier phase]\label{prop:phase}
Fix $\theta_*\in(0,\pi)$. As $n\to\infty$,
\begin{align*}
\arg\mathfrak c_n(x_*)&=\frac\pi4-\Bigl(n+\frac12\Bigr)\theta_*+\frac n2\sin2\theta_*+O\Bigl(\frac1n\Bigr)\pmod{2\pi},\\
|\mathfrak c_n(x_*)|&=e^{x_*^2/4}\Bigl(\frac2{\pi n}\Bigr)^{1/4}(\sin\theta_*)^{-1/2}\bigl(1+O(1/n)\bigr).
\end{align*}
Consequently $c_n=c_n^\natural+O(1/n)$ modulo $2\pi$, where
\[
c_n^\natural=c_\infty-n\pi\nu,\qquad c_\infty=\arg\gamma_*-\frac\pi4-\frac{\theta_*}2-\frac14\sin2\theta_*\re\rho_*,\qquad
\pi\nu=J(\theta_*)=\theta_*-\sin\theta_*\cos\theta_*,
\]
and Theorem~\ref{thm:firstorder}(a),(b) hold with $(c_n,\Psi_n,d_n)$ replaced by
$(c_n^\natural,\Psi(c_n^\natural),\max(n^{-1},\dist(c_n^\natural,\pi\Z)))$.
\end{proposition}

So the rotation number $\nu$ of the carrier phase is given by the same Kepler function $J$ that gives the threshold of the
model. The $O(1/n)$ terms are locally uniform in $\theta_*$; explicit constants can be obtained from \cite[Theorem~3.1]{SNWW}.

\begin{proof}
By the Plancherel--Rotach asymptotics in the oscillatory region \cite[Theorem~8.22.9(a), (8.22.12)]{Sz} (equivalently, by
formulas 12.7.2, 12.10.14--12.10.16, 12.10.18, 12.10.22 and 12.10.23 of \cite{DLMF}), rewritten for
$h_n=\He_n/\sqrt{n!}=2^{-n/2}H_n(\cdot/\sqrt2)/\sqrt{n!}$, for $x=2(n+\frac12)^{1/2}\cos\varphi$
\begin{equation}\label{eq:PR}
e^{-x^2/4}h_n(x)=\Bigl(\frac2{\pi n}\Bigr)^{1/4}(\sin\varphi)^{-1/2}\Bigl[\cos\Bigl(\bigl(n+\tfrac12\bigr)J(\varphi)-\frac\pi4\Bigr)
+O\Bigl(\frac1n\Bigr)\Bigr],
\end{equation}
uniformly for $\varphi\in[\epsilon,\pi-\epsilon]$. (In Szeg\H{o}'s form the phase is
$(\frac n2+\frac14)(\sin2\varphi-2\varphi)+\frac{3\pi}4$ for $\sin$, which is $\frac{3\pi}4-(n+\frac12)J(\varphi)$.) For real
$m$ near $n$ let $\varphi_m\in(0,\pi)$ be given by $\cos\varphi_m=x_*/(2(m+\frac12)^{1/2})$, and $\mathfrak g(m)=(m+\frac12)J(\varphi_m)$.
Differentiating, $\varphi_m'=\cos\varphi_m/(2(m+\frac12)\sin\varphi_m)$, so
$\mathfrak g'(m)=J(\varphi_m)+2(m+\frac12)\sin^2\varphi_m\,\varphi_m'=\varphi_m$. Since $\varphi_{n-1/2}=\theta_*$ and
$\varphi_m-\theta_*=O(1/n)$ for $m\in[n-1,n]$, we get $\mathfrak g(n)=nJ(\theta_*)+\frac12\theta_*+O(1/n)$ and
$\mathfrak g(n-1)=\mathfrak g(n)-\theta_*+O(1/n)$. Apply \eqref{eq:PR} at $x_*$ to $h_n$ and to $h_{n-1}$; with $\Theta=\mathfrak g(n)-\frac\pi4$ and
$A_n=(2/(\pi n))^{1/4}(\sin\varphi_n)^{-1/2}$,
\[
e^{-\frac{x_*^2}4}h_n(x_*)=A_n\bigl[\cos\Theta+O(\tfrac1n)\bigr],\qquad
e^{-\frac{x_*^2}4}h_{n-1}(x_*)=A_n\bigl[\cos(\Theta-\theta_*)+O(\tfrac1n)\bigr].
\]
As $\cos\Theta\cos\theta_*-\cos(\Theta-\theta_*)=-\sin\Theta\sin\theta_*$, this gives
\[
e^{-x_*^2/4}\mathfrak c_n(x_*)=A_n\bigl[e^{-i\Theta}+O(1/n)\bigr],
\]
hence $\arg\mathfrak c_n(x_*)=-\Theta+O(1/n)$, which is the first claim. The second follows, since
$\sin\varphi_n=\sin\theta_*+O(1/n)$, and the definition of $c_n$ gives $c_n=c_n^\natural+O(1/n)$.

For the last claim, on $B_0$ we have
$|\partial_c\Phi_c|\le|\partial_cf|+|\partial_c\partial_\zeta f|\le1+2(|\zeta|+|\lambda|)\le6$.
Choosing representatives with $|c_n-c_n^\natural|\le C'/n$, Proposition~\ref{prop:threshold} gives
$|\Phi_{c_n^\natural}(\zeta_n,\lambda_n)|\le(C_{\Xi_0}+6C')/n$, and the proof of Theorem~\ref{thm:firstorder} applies with
$c_n^\natural$ in place of $c_n$. The two error bounds are comparable, since
$\max(n^{-1},\dist(c_n^\natural,\pi\Z))\le(1+C')d_n$ and $d_n\le(1+C')\max(n^{-1},\dist(c_n^\natural,\pi\Z))$.
\end{proof}

\subsection{Oscillation and the joint law}\label{sec:jointlaw}

\begin{theorem}[Oscillation and the joint law]\label{thm:jointlaw}
Let $p_W(s)=(\pi W)^{-1}((W+s)/(W-s))^{1/2}$ for $|s|<W$ and $p_W(s)=0$ otherwise.
\begin{enumerate}
\item[(a)] Assume (G1)--(G3) and $\nu=J(\theta_*)/\pi\notin\Q$. Then the set of limit points of $n(\Lambda_n-T_*)$ is
$[\tau_c-W,\tau_c+W]$, and $N^{-1}\sum_{n\le N}\delta_{n(\Lambda_n-T_*)}$ converges weakly to the law with density
$p_W(\tau-\tau_c)$, whose mean is $\tau_c+W/2$. If $\nu=p/r$ in lowest terms, then $n(\Lambda_n-T_*)$ is asymptotically
periodic with period $r$, and it converges if and only if $\theta_*=\pi/2$ and $\dist(c_\infty,\pi\Z)=\pi/4$.
\item[(b)] Let $q$ be the Gaussian power series. Almost surely (G1)--(G3) hold and $\nu\notin\Q$, so that (a) applies.
For every random element $Y$ of a Polish space that is a measurable function of $q$,
\[
\bigl(Y,\,n(\Lambda_n-T_*)\bigr)\longrightarrow\bigl(Y,\,\tau_c-W\cos\Psi(U)\bigr)\qquad\text{in distribution,}
\]
where $U$ is uniform on $[0,\frac\pi2]$ and independent of $q$. In particular
$(T_*,n(\Lambda_n-T_*))\to(T_*,\tau_c-W\cos\Psi(U))$, and $\E[Z\varphi(n(\Lambda_n-T_*))]\to\E[Z\int\varphi(\tau_c+s)p_W(s)\,ds]$
for every bounded measurable function $Z$ of $q$ and bounded continuous $\varphi$: the conditional law of
$n(\Lambda_n-T_*)$ given $q$ converges to the law with density $p_W(\tau-\tau_c)$.
\end{enumerate}
\end{theorem}

The limit law piles up at the upper end $\tau_c+W$, where its density has an inverse square-root singularity, and
vanishes like a square root at $\tau_c-W$.

\begin{proof}
\emph{The law.} Let $V$ be uniform on $\R/\pi\Z$ and $F(c)=\tau_c-W\cos2s_c=\tau_c-W\cos\Psi(c)$, a continuous function on
$\R/\pi\Z$; as $c$ runs once around $\R/\pi\Z$, $2s_c$ runs once around $\R/2\pi\Z$, so $F(\R/\pi\Z)=[\tau_c-W,\tau_c+W]$.
$U=\dist(V,\pi\Z)$ is uniform on $[0,\frac\pi2]$ and $F(V)=\tau_c-W\cos\Psi(U)$. By \eqref{eq:kepler},
$dU=\frac12(1-\cos\Psi)\,d\Psi$, so $\Psi(U)$ has density $(1-\cos\Psi)/\pi$ on $[0,\pi]$, and $s=-W\cos\Psi$ has density
$(1-\cos\Psi)/(\pi W\sin\Psi)=p_W(s)$ and mean $-\frac W\pi\int_0^\pi\cos\Psi(1-\cos\Psi)\,d\Psi=\frac W2$.

(a) By Theorem~\ref{thm:firstorder} and Proposition~\ref{prop:phase}, $\tau_n-F(c_n^\natural)\to0$, where
$\tau_n=n(\Lambda_n-T_*)$. If $\nu\notin\Q$, $(n\nu)$ is equidistributed modulo $1$ \cite[Ch.~1, Example~2.1]{KN}, so
$(c_n^\natural)$ is equidistributed in $\R/\pi\Z$. Every point of $\R/\pi\Z$ is then a limit point of $(c_n^\natural)$, so by
compactness and continuity of $F$ the limit points of $(\tau_n)$ form $F(\R/\pi\Z)$. For bounded continuous $\varphi$,
uniform continuity on a compact set containing all $\tau_n$ gives $N^{-1}\sum_{n\le N}(\varphi(\tau_n)-\varphi(F(c_n^\natural)))\to0$,
and $N^{-1}\sum_{n\le N}\varphi(F(c_n^\natural))\to\E\varphi(F(V))$ by \cite[Ch.~1, Corollary~1.2]{KN}. If
$\nu=p/r$, then $c_n^\natural$ runs periodically through the $r$ points $c_\infty-j\pi/r$ of $\R/\pi\Z$, $0\le j<r$, and
$r\ge2$ as $0<\nu<1$. $F$ depends on $c$ only through $\dist(c,\pi\Z)$, injectively, and this function takes each value in
$(0,\frac\pi2)$ at two points of $\R/\pi\Z$ and each of $0,\frac\pi2$ at one. So $(\tau_n)$ converges if and only if $r=2$,
that is $J(\theta_*)=\frac\pi2$, that is $\theta_*=\frac\pi2$, and $\dist(c_\infty,\pi\Z)=\dist(c_\infty-\frac\pi2,\pi\Z)$,
that is $\dist(c_\infty,\pi\Z)=\frac\pi4$.

(b) By Section~\ref{sec:collision} (Lemma~\ref{lem:moduli}), (G1)--(G3) hold almost surely. \emph{The law of $u_*$ is
absolutely continuous.} Write $q=\eta_0+\eta_1w+g(w)$. For $u\in\D\setminus\R$ the map $\Theta_u(c_0,c_1)=c_0+c_1u$ is
an $\R$-linear isomorphism $\R^2\to\C$ (its determinant is $\im u$), depending smoothly on $u$, and $q(u)=0$ if and only if
$(\eta_0,\eta_1)=S_g(u):=-\Theta_u^{-1}g(u)$. Let $A\subset\D\setminus\R$ be a Borel null set. For fixed $g$, $S_g$ is
locally Lipschitz, so $S_g(A)$ is null; the set of $(c,g)$ with $c\in S_g(A)$ is analytic (a projection of a Borel set),
hence universally measurable; and $(\eta_0,\eta_1)$ has a density and is independent of $g$. By Fubini's theorem,
$q$ has a zero in $A$ with probability $0$. Hence $\PP(u_*\in A)=0$ ($u_*$, the zero of least modulus in $\D\cap\{\im w>0\}$, is a measurable function of $q$),
$\theta_*=\arg u_*$ has a density on $(0,\pi)$, and
so has $\nu=J(\theta_*)/\pi$, since $J$ is a diffeomorphism of $(0,\pi)$ onto itself. In particular $\PP(\nu\in\Q)=0$.

\emph{The phase becomes independent.} Let $p_\nu$ be the density of $\nu$ and $Z$ a bounded complex measurable
function of $q$. For $k\ne0$, $\E[Ze^{-2\pi ikn\nu}]=\int_0^1e^{-2\pi iknv}\E[Z\mid\nu=v]\,p_\nu(v)\,dv\to0$ by the
Riemann--Lebesgue lemma. Applying this to $Ze^{2ikc_\infty}$ gives $\E[Z\psi(c_n^\natural)]\to\E[Z]\,\E\psi(V)$ for
trigonometric polynomials $\psi$ on $\R/\pi\Z$, and by uniform approximation for all continuous $\psi$; this is
the arbitrary functions principle in the form of \cite[eq.~(1)]{Bou}. Let $\bar Y=(Y,T_*,\theta_*,\gamma_*,\rho_*)$. Since $\bar Y$ has a fixed law, $\R/\pi\Z$ is compact,
and products $\phi(\bar y)\psi(c)$ determine laws, $(\bar Y,c_n^\natural)\to(\bar Y,V)$ in distribution, with $V$
independent of $q$ \cite{Bi2}. The map $(\bar y,c)\mapsto F_{\bar y}(c)=\tau_c-W\cos2s_c$, with $\tau_c$ and $W$ computed
from $\bar y$, is continuous, so $(\bar Y,F_{\bar Y}(c_n^\natural))\to(\bar Y,F_{\bar Y}(V))$. Almost surely
$\tau_n-F_{\bar Y}(c_n^\natural)\to0$, by (a) of this proof, so $(\bar Y,\tau_n)\to(\bar Y,F_{\bar Y}(V))$, and
$F_{\bar Y}(V)=\tau_c-W\cos\Psi(U)$ with $U=\dist(V,\pi\Z)$ independent of $q$. Projecting to $(Y,\tau_n)$ gives the
claim; with $Y=(Z,\tau_c,W)$ and the law computed above, it gives the statement about $\E[Z\varphi(\tau_n)]$.
\end{proof}

\begin{corollary}[Weyl polynomials]\label{cor:firstorder-weyl}
Let $\Lambda_n^{\rm std}$ be the real-rooting time of $G_{n,t}$, and $\tilde q_n$, $\tilde T_{*,n}$ as in
Corollary~\ref{cor:collision-weyl}. Then $(\tilde T_{*,n},n(\Lambda_n^{\rm std}-\tilde T_{*,n}))\to(T_*,\tau_c-W\cos\Psi(U))$ in
distribution, with $U$ uniform on $[0,\frac\pi2]$ and independent of $q$. In particular $n(\Lambda_n^{\rm std}-\tilde T_{*,n})$
converges in distribution to the law with density $\tau\mapsto\E[p_W(\tau-\tau_c)]$.
\end{corollary}

\begin{proof}
For each $n$, $G_{n,t}$ is the heat evolution of the reversed polynomial built from $\tilde q_n$ (proof of
Corollary~\ref{cor:collision-weyl}), so $(\Lambda_n^{\rm std},\tilde q_n)$ has the law of $(\Lambda_n,q)$, and
$(\tilde T_{*,n},n(\Lambda_n^{\rm std}-\tilde T_{*,n}))$ has the law of $(T_*,n(\Lambda_n-T_*))$. Apply
Theorem~\ref{thm:jointlaw}(b) with $Y=T_*$; the density of the mixture follows by Fubini's theorem.
\end{proof}

The almost sure statements of Theorem~\ref{thm:jointlaw} do not transfer to $G_n$, since $\tilde q_n$ changes with $n$.

\begin{remark}[A benchmark]\label{rem:benchmark}
Let $q(w)=1+bw+aw^2$ with $b^2<4a$. Then (G1)--(G3) hold, $u_*=a^{-1/2}e^{i\theta_*}$ with $b=-2\sqrt a\cos\theta_*$,
$T_*=a$, $q'(u_*)=2i\sqrt a\sin\theta_*$ and $\rho_*=1-i\cot\theta_*$, so that $\tau_c=-\frac a2\cot^2\theta_*$,
$W=a/(2\sin^2\theta_*)$ and $\tau_c+W=\frac a2$. By Theorem~\ref{thm:jointlaw}(a), if $\nu\notin\Q$, then
$\limsup_nn(\Lambda_n-a)=\frac a2$ and $\liminf_nn(\Lambda_n-a)=-\frac a2(1+\cos^2\theta_*)/\sin^2\theta_*$. On the other
hand, $\Lambda_n\le a\sqrt{n/(n-1)}$ for every $n\ge2$: by \eqref{eq:recurrence},
$\sqrt{n!}\,Q_n^{(t)}=(x+b\sqrt{n/t})\He_{n-1}-\chi\He_{n-2}$ with $\chi=n-1-(a/t)\sqrt{n(n-1)}$, and for
$t\ge a\sqrt{n/(n-1)}$ we have $\chi\ge0$, so that, expanding along the last row, this is $\det(x-J')$ for the real
symmetric matrix $J'$ obtained from $J_n$ by replacing its last diagonal entry by $-b\sqrt{n/t}$ and its last two
off-diagonal entries by $\sqrt\chi$ (recall from the proof of Lemma~\ref{lem:colleague} that $\det(x-J_m)=\He_m$). Hence $Q_n^{(t)}$, and with it $P_{n,t}$, has only
real zeros. Since $n(a\sqrt{n/(n-1)}-a)=\frac a2+O(\frac1n)$, this elementary bound is sharp to first order along a
subsequence whenever $c_n^\natural$ accumulates at $\frac\pi2+\pi\Z$, for instance when $\nu\notin\Q$.

Table~\ref{tab:benchmark} compares, for $a=2$ and $\theta_*=1.1$ (so $\nu=0.2215\ldots$, $\tau_c=-0.2590\ldots$,
$W=1.2590\ldots$), the computed $n(\Lambda_n-2)$ with $\tau_c-W\cos\Psi_n$, where $c_n$ is computed from $h_n(x_*)$ and
$h_{n-1}(x_*)$. $\Lambda_n$ was found by Newton's method for $Q_n^{(t)}=\partial_xQ_n^{(t)}=0$ in $(t,x)$, in double
precision, and for several $n$ confirmed in $50$-digit arithmetic and by counting real zeros just before and after it.
Over the $128$ degrees computed, $100\le n\le25607$, the largest $n(\Lambda_n-2)$ is $1.0017$, at $n=207$, below the bound
$1.0036$ for that $n$.
\end{remark}

\begin{table}[ht]
\caption{The benchmark $q=1+bw+2w^2$, $\theta_*=1.1$: the scaled real-rooting time against the first-order law.}
\label{tab:benchmark}
\begin{tabular}{rcccc}
$n$ & $\dist(c_n,\pi\Z)$ & $n(\Lambda_n-T_*)$ & $\tau_c-W\cos\Psi_n$ & $n\times$ difference\\\hline
$100$ & $0.671$ & $\phantom{-}0.45390$ & $\phantom{-}0.45007$ & $\phantom{-}0.38$\\
$101$ & $0.024$ & $-1.24112$ & $-1.24775$ & $\phantom{-}0.67$\\
$207$ & $1.515$ & $\phantom{-}1.00171$ & $\phantom{-}0.99805$ & $\phantom{-}0.76$\\
$408$ & $0.009$ & $-1.37582$ & $-1.37228$ & $-1.45$\\
$1609$ & $0.076$ & $-0.95726$ & $-0.95684$ & $-0.68$\\
$3203$ & $0.029$ & $-1.21202$ & $-1.21171$ & $-1.00$\\
$12805$ & $1.553$ & $\phantom{-}0.99987$ & $\phantom{-}0.99981$ & $\phantom{-}0.75$\\
$25604$ & $0.065$ & $-1.00556$ & $-1.00558$ & $\phantom{-}0.38$\\
\end{tabular}
\end{table}


\section{Remarks and questions}\label{sec:remarks}

\subsection{The limit law}
As $t\downarrow1$ the event $\{T_*\le t\}$ asks that $q$ has no nonreal zero in a disk of radius $t^{-1/2}$ close to $1$, a hole probability for the nonreal zeros near the whole unit circle.

\begin{question}\label{q:law}
By Theorems~\ref{thm:sharp} and~\ref{thm:lower}, $-\log\PP(T_*\le t)=\frac{\pi^2}{12(t-1)}(1+o(1))$ as $t\downarrow1$, half of the exponent for the Gaussian power series with complex coefficients, whose probability of having no zero in $\{|w|<t^{-1/2}\}$ is $\prod_{k\ge1}(1-t^{-k})=\exp\bigl(-(\pi^2/6+o(1))/(t-1)\bigr)$ \cite[Theorem~2(i) and Corollary~3(i)]{PV}; and $-\log\PP(T_*\le t)\le\frac{\pi^2}{12\log t}+\log\frac1{1-t^{-1/2}}$ for every $t>1$. What is the second-order term of $-\log\PP(T_*\le t)$ as $t\downarrow1$? Numerically, for $0.002\le t-1\le0.08$ (Remark~\ref{rem:fredholm-num}), fits of $(t-1)(-\log\PP(T_*\le t))$ with corrections in powers of $\log\frac1{t-1}$ up to the third give $0.82246$--$0.82247$ for its limit, against $\pi^2/12=0.822467$, while other correction models give $0.820$--$0.826$ with larger residuals. At the computed values of $t$, $-\log\PP(T_*\le t)$ lies below $\frac{\pi^2}{12(t-1)}$, by an amount that grows faster than $\log\frac1{t-1}$; its form is not determined.
\end{question}

\begin{question}\label{q:bernoulli}
Are the eigenvalues of the operators $X(r)$ of Proposition~\ref{prop:fredholm} real and in $[0,1]$ for every $r\in(0,1)$? Since $\E s^{k_r}=\prod_k(1-(1-s)\lambda_k)$ over the eigenvalues $\lambda_k$ of $X(r)$, by Proposition~\ref{prop:fredholm} and \cite[Theorem~14.44]{vN}, this holds if and only if the number $k_r$ of zeros of $q$ in the half disk $\D_r^+$ is distributed as a sum of independent Bernoulli variables, as the number of zeros in $\{|w|<r\}$ is for the complex series \cite{PV}. The constant $\pi^2/12$ no longer depends on this question (Theorem~\ref{thm:lower}). The analogous statement for general regions symmetric under $z\mapsto-\bar z$ is false (Remark~\ref{rem:bernoulli-general}), so an answer must use the shape of the half disk. Numerically, for ten values of $r$ between $0.1$ and $0.99$ and for $r=0.995$ (truncations up to $1401\times1401$), all computed eigenvalues of $X(r)$ are real, and those above $10^{-12}$ lie in $(0,1)$; smaller eigenvalues are not resolved in double precision.
\end{question}

\begin{remark}[Other regions]\label{rem:bernoulli-general}
Let
\[
z_1=-\tfrac{19}{25}+\tfrac{i}{10},\quad z_2=-\bar z_1,\quad z_3=\tfrac{27}{100}+\tfrac{17}{20}\,i,\quad z_4=-\bar z_3,\qquad b_1=b_2=\tfrac{42}5,\quad b_3=b_4=\tfrac9{40},
\]
let $a_j=\pi b_j|1-z_j^2|$, and for $\delta>0$ let $\mathcal A_\delta$ be the union of the disks $B(z_j,\delta\sqrt{a_j/\pi})$, $j\le4$, of areas $\delta^2a_j$; it is symmetric under $z\mapsto-\bar z$. For all sufficiently small $\delta>0$, the number $N_\delta$ of zeros of $q$ in $\mathcal A_\delta$ is not distributed as a sum of independent Bernoulli variables. Indeed, Step~1 of the proof of Proposition~\ref{prop:fredholm} applies to $\mathcal A_\delta$ in place of $\D_r^+$, and with $1-s=\tau/\delta^2$ the $n$-th term of the resulting series \eqref{eq:fpf} for $\E s^{N_\delta}$ is bounded by $(C|\tau|)^n(2n)^{n/2}/n!$, uniformly in small $\delta$. As the correlation functions $\rho_n$ of \eqref{eq:ms-kernel} are continuous and vanish when two arguments coincide, it follows that, locally uniformly in $\tau\in\C$,
\[
\lim_{\delta\to0}\E s^{N_\delta}=P(\tau)=\sum_{S\subset\{1,2,3,4\}}(-\tau)^{|S|}\prod_{j\in S}a_j\cdot\rho_{|S|}\bigl((z_j)_{j\in S}\bigr).
\]
By \eqref{eq:ms-kernel} the coefficients of $P$ are rational numbers; numerically $P(\tau)=1-38.031\tau+534.81\tau^2-3288.5\tau^3+7440.7\tau^4$. Exact rational arithmetic shows that the discriminant of $P$ is negative, so $P$ has two simple real zeros and a pair of simple nonreal zeros, $\tau\approx0.08879\pm0.00176\,i$. By Hurwitz's theorem $\E s^{N_\delta}$, as a function of $\tau$, has a nonreal zero for all small $\delta$, whereas $\E s^N=\prod_k(1-p_k\tau/\delta^2)$ has only real zeros when $N$ is a sum of independent Bernoulli variables with parameters $p_k$. The argument gives no explicit range of $\delta$; numerically, for this configuration, the nonreal zeros persist up to $\delta=0.0275$ and are absent from $\delta=0.03$ on. As the coefficients of $P$ depend continuously on the $z_j$ and $a_j$, the conclusion also holds for all nearby configurations of four points and weights, symmetric or not.

A similar computation shows that the zeros of $q$ in $\D^+$ do not form a determinantal point process with a Hermitian kernel $K$ (with respect to any reference measure, whose density can be absorbed into $K$). For such a process, $a_{ij}=\rho_1(z_i)\rho_1(z_j)-\rho_2(z_i,z_j)=|K(z_i,z_j)|^2$ and $T=\rho_3(z_1,z_2,z_3)-\rho_1(z_1)\rho_1(z_2)\rho_1(z_3)+\sum\rho_1(z_i)a_{jk}=2\re\bigl(K(z_1,z_2)K(z_2,z_3)K(z_3,z_1)\bigr)$, the sum over $\{i,j,k\}=\{1,2,3\}$, so $T^2\le4a_{12}a_{13}a_{23}$ almost everywhere. But the $\rho_n$ of \eqref{eq:ms-kernel} are continuous, and at $(z_1,z_2,z_3)=(-\frac3{20}+\frac i{20},\frac9{10}+\frac{7i}{50},-\frac{93}{100}+\frac{13i}{100})$ the ratio $T^2/(4a_{12}a_{13}a_{23})$ is a rational number equal to $3.9154\ldots$, so the inequality fails on a neighbourhood of this point. Both computations are computer-assisted, in exact rational arithmetic, by the Python program \texttt{part4\_exact\_checks.py} of the verification suite in the ancillary files (see the Statement on the use of AI).
\end{remark}


\subsection{The standard coupling}
Theorem~\ref{thm:std} leaves open how often $\Lambda_n^{\rm std}$ lies below a given level. Heuristically, $\Lambda_n^{\rm std}$ is decided by the top block $(\xi_n,\dots,\xi_{n-M_n})$ of length $O(\log n)$, which has the law of $(\eta_0,\dots,\eta_{M_n})$, and blocks that are more than $M_n$ apart are independent.

\begin{question}\label{q:ergodic}
Is it true that, almost surely,
\[
\frac1N\#\{n\le N:\ \Lambda_n^{\rm std}\le t\}\ \to\ \PP(T_*\le t)\qquad\text{for every }t\ge0\,?
\]
\end{question}

For $t<1$ this holds, with limit $0$, by the first part of the proof of Theorem~\ref{thm:std}. For $t>1$ a proof would need two-sided approximations of $\{\Lambda_n^{\rm std}\le t\}$ by events of the top block whose discrepancy has small frequency, not merely small probability; the one-sided inclusions in the proof of Lemma~\ref{lem:std-planted} do not give this. We also do not know how fast $\limsup_n\Lambda_n^{\rm std}=\infty$ is approached, for instance for which deterministic $\Lambda(n)\to\infty$ the inequality $\Lambda_n^{\rm std}>\Lambda(n)$ holds infinitely often.


\subsection{The cusp phases}
Theorem~\ref{thm:firstorder} gives the time to within $O(n^{-1}d_n^{-1/3})$ and the place to within $O(n^{-1}d_n^{-2/3})$, where $d_n=\max(n^{-1},\dist(c_n,\pi\Z))$: the error is $O(1/n)$ away from the cusp phases $c_n\in\pi\Z$, and $O(n^{-2/3})$, respectively $O(n^{-1/3})$, uniformly. By Remark~\ref{rem:sharp}, no argument that uses only the conclusion of Proposition~\ref{prop:window} does better at the cusps.

\begin{question}\label{q:cusp}
What is the optimal error in Theorem~\ref{thm:firstorder} at the cusp phases, that is, when $\dist(c_n,\pi\Z)=O(1/n)$? Identify the coefficient of the cusp shift: the component of the next-order term of the remainder in Proposition~\ref{prop:window}, at the triple point $(\zeta,\lambda,c)=(0,-1,0)$ of the model, along the direction $\upsilon$ of Remark~\ref{rem:sharp}, which moves the cusp. For which $q$ does it vanish? The numerics in Remark~\ref{rem:sharp} suggest that it vanishes for some symbols, the benchmark among them, and not for others. Is the final collision $z_n$ a triple zero for some deterministic $q$ satisfying (G1)--(G3) and some $n\ge n_0$? By Theorem~\ref{thm:firstorder}(c), for $n\ge n_1$ this requires $\dist(c_n,\pi\Z)<C/n$.
\end{question}


\subsection{Beyond this paper}
The companion papers leave their own questions: whether the rate $n^{-\alpha}$ in Theorem~\ref{cor:weyl} can be improved beyond the small explicit exponent obtained there, and whether the constant in Theorem~\ref{thm:collision-intro}(a) has a power tail without the restriction $T_*\ge1+\eps$ \cite{CompA}; whether the description of Theorem~\ref{thm:heat} persists at times $t_n=1+s_nn^{-1/3}$ whenever $s_n\to\infty$, where the two edge regions remain open \cite{CompB}; and what replaces the Joukowski picture of Theorem~\ref{thm:main} for weights whose equilibrium measure has several intervals \cite{CompC}.

\raggedbottom
\section*{Statement on the use of AI}
The proofs, computations and text of this paper were produced with AI systems from Anthropic, OpenAI and Google, directed by the author. The proofs were checked by independent AI referee rounds; no human expert has yet read them.

Independently of these rounds, the algebraic identities, the explicit constants and several of the quantitative predictions were checked mechanically by a suite of Python programs, deposited with the paper as ancillary files (folder \texttt{anc/}; run \texttt{run\_all.py}). The suite checks the identities of Sections~\ref{sec:prelim}--\ref{sec:transfer} and the window model of Section~\ref{sec:firstorder} symbolically. It encloses the onset $n_s$ and the constants of Proposition~\ref{prop:early}, and the values of $\det(I-X(r))$ reported in Remark~\ref{rem:fredholm-num}, in interval arithmetic. It certifies the zero counts of Theorem~\ref{thm:main} for three admissible symbols, and it carries out the exact computations of Remark~\ref{rem:bernoulli-general}. Finally, it computes $\Lambda_n$ by certified bisection in $t$ and compares $n(\Lambda_n-T_*)$ with Theorem~\ref{thm:firstorder-intro}. All checks pass, and every printed number that the suite recomputes is reproduced to the printed digits. The suite checks identities, constants and predictions, not the proofs themselves.

\begin{thebibliography}{99}
% TODO(author): pinpoints carried over from the companion drafts that were not checked against the sources: the
% theorem numbers of [vN] (Section 14.5, Proposition 14.21, Theorem 14.34, Lemmas 14.38--14.40, Theorem 14.44, as in
% arXiv v7), [HKPV, Section 2.4], [Fe, Theorem 1], [HHJK, Theorem 2.1] and the bibliographic data of [HHJK];
% [Bi, Theorems 27.2 and 29.4] and [Sz, Theorem 3.3.2] were checked only from memory; [Sz, (3.2.4)] and
% [Sz, Thm 3.2.2] likewise; the volume and pages of [Pe]; the journal data of [Du2]; the bibliographic data of [DS],
% [GSC], [KVA], [DKMVZ] and [Fuj]. The pinpoints in [Si], [Ob], [AW] and [Ka] have their own comments in the text.
% Added 22 October 2026 after a referee note on missing references: [Sh], [BD], [Tu], [IS], [Du3], [SV], [LMS], [Shi], [NNS],
% [APS]. Their bibliographic data were checked online (AMS, Crossref, arXiv and RePEc listings; [Tu] through the reference
% lists of Bonciocat--Cipu, Math. Inequal. Appl. 13 (2010), and of [Du2]). Statement numbers quoted: [Du1, Thms 1.2-1.3],
% [Du3, Thms 1.3-1.6], [Du2, Lemma 3.1], [SV, Thm 1.2], [NNS, Thm 1.3] were read in the arXiv/JSTOR texts. [IS] is cited
% without a statement number, because secondary sources disagree on it (Proposition 1 in [Du1], Proposition 6 in Romik's
% Acta Arith. paper); check the original before submission.
\bibitem[A]{CompA} B. Baggins, \emph{Explicit onsets, rates and the tail of the collision constant for heat-evolved Weyl polynomials}, preprint (2026).
\bibitem[B]{CompB} B. Baggins, \emph{Nonreal zeros of heat-evolved Weyl polynomials near the collapse time}, preprint (2026).
\bibitem[C]{CompC} B. Baggins, \emph{Exact counts of nonreal zeros beyond the Weyl weight: varying weights, coefficient profiles and the Laguerre flow}, preprint (2026).
\bibitem{APS} G. Akemann, M. J. Phillips, L. Shifrin, \emph{Gap probabilities in non-Hermitian random matrix theory}, J. Math. Phys. \textbf{50} (2009), 063504; doi:10.1063/1.3133108; arXiv:0901.0897.
\bibitem{AW} J.-M. Aza\"is, M. Wschebor, \emph{Level Sets and Extrema of Random Processes and Fields}, Wiley, Hoboken, NJ, 2009.
\bibitem{BD} A. F. Beardon, K. A. Driver, \emph{The zeros of linear combinations of orthogonal polynomials}, J. Approx. Theory \textbf{137} (2005), 179--186; doi:10.1016/j.jat.2005.09.001.
\bibitem{Bi} P. Billingsley, \emph{Probability and Measure}, 3rd ed., Wiley, New York, 1995.
\bibitem{Bi2} P. Billingsley, \emph{Convergence of Probability Measures}, 2nd ed., Wiley, New York, 1999.
\bibitem{Bou} N. Bouleau, \emph{An extension to the Wiener space of the arbitrary functions principle}, C. R. Math. Acad. Sci. Paris \textbf{343} (2006), no.~5, 329--332; doi:10.1016/j.crma.2006.06.028; arXiv:math/0610509.
\bibitem{dB} N. G. de Bruijn, \emph{The roots of trigonometric integrals}, Duke Math. J. \textbf{17} (1950), 197--226.
\bibitem{CSV} G. Csordas, W. Smith, R. S. Varga, \emph{Lehmer pairs of zeros, the de Bruijn--Newman constant $\Lambda$, and the Riemann hypothesis}, Constr. Approx. \textbf{10} (1994), 107--129.
\bibitem{DS} D. Damanik, B. Simon, \emph{Jost functions and Jost solutions for Jacobi matrices, II. Decay and analyticity}, Int. Math. Res. Not. \textbf{2006}, Art. ID 19396, 1--32; doi:10.1155/IMRN/2006/19396; arXiv:math/0502487.
\bibitem{De} K. Deimling, \emph{Nonlinear Functional Analysis}, Springer, Berlin, 1985.
\bibitem{DKMVZ} P. Deift, T. Kriecherbauer, K. T.-R. McLaughlin, S. Venakides, X. Zhou, \emph{Uniform asymptotics for polynomials orthogonal with respect to varying exponential weights and applications to universality questions in random matrix theory}, Comm. Pure Appl. Math. \textbf{52} (1999), 1335--1425.
\bibitem{DLMF} \emph{NIST Digital Library of Mathematical Functions}, \url{https://dlmf.nist.gov/}, Release 1.2.8 of 2026-09-15, F. W. J. Olver, A. B. Olde Daalhuis, D. W. Lozier, B. I. Schneider, R. F. Boisvert, C. W. Clark, B. R. Miller, B. V. Saunders, H. S. Cohl and M. A. McClain, eds.
\bibitem{Du1} A. J. Dur\'an, \emph{Zeros of linear combinations of Hermite polynomials}, preprint (2025), arXiv:2505.15330.
\bibitem{Du2} A. J. Dur\'an, \emph{Zeros of linear combinations of orthogonal polynomials}, Mediterr. J. Math. \textbf{23} (2026), no.~5, Paper No.~148; doi:10.1007/s00009-026-03145-9; arXiv:2505.11956.
\bibitem{Du3} A. J. Dur\'an, \emph{Zeros of linear combinations of Laguerre polynomials}, preprint (2025), arXiv:2507.22425.
\bibitem{Fe} N. Feldheim, \emph{Zeroes of Gaussian analytic functions with translation-invariant distribution}, Israel J. Math. \textbf{195} (2013), 317--345; doi:10.1007/s11856-012-0130-0; arXiv:1105.3929.
\bibitem{Fo} P. J. Forrester, \emph{The limiting Kac random polynomial and truncated random orthogonal matrices}, J. Stat. Mech. Theory Exp. (2010), P12018; doi:10.1088/1742-5468/2010/12/P12018; arXiv:1009.3066.
\bibitem{GSC} L. Golinskii, S. Serra-Capizzano, \emph{The asymptotic properties of the spectrum of nonsymmetrically perturbed Jacobi matrix sequences}, J. Approx. Theory \textbf{144} (2007), 84--102; arXiv:math/0512222.
\bibitem{HHJK} B. C. Hall, C.-W. Ho, J. Jalowy, Z. Kabluchko, \emph{Zeros of random polynomials undergoing the heat flow}, Electron. J. Probab. \textbf{30} (2025), paper no.~159, 1--55; doi:10.1214/25-EJP1421; arXiv:2308.11685.
\bibitem{HHJK2} B. C. Hall, C.-W. Ho, J. Jalowy, Z. Kabluchko, \emph{The heat flow, GAF, and $SL(2;\mathbb R)$}, Indiana Univ. Math. J. \textbf{74} (2025), 1153--1206; doi:10.1512/iumj.2025.74.60451; arXiv:2304.06665.
\bibitem{HKPV} J. B. Hough, M. Krishnapur, Y. Peres, B. Vir\'ag, \emph{Zeros of Gaussian Analytic Functions and Determinantal Point Processes}, Univ. Lecture Ser. \textbf{51}, Amer. Math. Soc., Providence, RI, 2009.
\bibitem{IS} A. Iserles, E. B. Saff, \emph{Zeros of expansions in orthogonal polynomials}, Math. Proc. Cambridge Philos. Soc. \textbf{105} (1989), 559--573; doi:10.1017/S0305004100077938.
\bibitem{Ka} M. Kac, \emph{On the average number of real roots of a random algebraic equation}, Bull. Amer. Math. Soc. \textbf{49} (1943), 314--320.
\bibitem{KN} L. Kuipers, H. Niederreiter, \emph{Uniform Distribution of Sequences}, Wiley, New York, 1974.
\bibitem{KVA} A. B. J. Kuijlaars, W. Van Assche, \emph{The asymptotic zero distribution of orthogonal polynomials with varying recurrence coefficients}, J. Approx. Theory \textbf{99} (1999), 167--197.
\bibitem{LMS} A. Ledoan, M. Merkli, S. Starr, \emph{A universality property of Gaussian analytic functions}, J. Theoret. Probab. \textbf{25} (2012), 496--504; doi:10.1007/s10959-011-0356-5; arXiv:1003.1951.
\bibitem{MS} S. Matsumoto, T. Shirai, \emph{Correlation functions for zeros of a Gaussian power series and Pfaffians}, Electron. J. Probab. \textbf{18} (2013), no.~49, 1--18; doi:10.1214/EJP.v18-2545; arXiv:1212.6108.
\bibitem{Me} F. G. Mehler, \emph{Ueber die Entwicklung einer Function von beliebig vielen Variablen nach Laplaceschen Functionen h\"oherer Ordnung}, J. Reine Angew. Math. \textbf{66} (1866), 161--176.
\bibitem{NNS} F. Nazarov, A. Nishry, M. Sodin, \emph{Log-integrability of Rademacher Fourier series, with applications to random analytic functions}, Algebra i Analiz \textbf{25} (2013), no.~3, 147--184; arXiv:1301.0529.
\bibitem{vN} J. van Neerven, \emph{Functional Analysis}, Cambridge Stud. Adv. Math. \textbf{201}, Cambridge University Press, Cambridge, 2022; arXiv:2112.11166 (theorem numbers as in v7, 2025).
\bibitem{Ne} C. M. Newman, \emph{Fourier transforms with only real zeros}, Proc. Amer. Math. Soc. \textbf{61} (1976), 245--251.
\bibitem{Ob} N. Obreschkoff, \emph{Verteilung und Berechnung der Nullstellen reeller Polynome}, Hochschulb\"ucher f\"ur Mathematik \textbf{55}, VEB Deutscher Verlag der Wissenschaften, Berlin, 1963.
\bibitem{OT} A. Ostrowski, O. Taussky, \emph{On the variation of the determinant of a positive definite matrix}, Nederl. Akad. Wetensch. Proc. Ser.~A \textbf{54} = Indag. Math. \textbf{13} (1951), 383--385.
\bibitem{Pe} F. Peherstorfer, \emph{Zeros of linear combinations of orthogonal polynomials}, Math. Proc. Cambridge Philos. Soc. \textbf{117} (1995), 533--544; doi:10.1017/S0305004100073357.
\bibitem{PV} Y. Peres, B. Vir\'ag, \emph{Zeros of the i.i.d.\ Gaussian power series: a conformally invariant determinantal process}, Acta Math. \textbf{194} (2005), 1--35; doi:10.1007/BF02392515; arXiv:math/0310297.
\bibitem{PM15} D. H. J. Polymath, \emph{Effective approximation of heat flow evolution of the Riemann $\xi$ function, and a new upper bound for the de Bruijn--Newman constant}, Res. Math. Sci. \textbf{6} (2019), no.~3, Paper No.~31; doi:10.1007/s40687-019-0193-1; arXiv:1904.12438.
\bibitem{Ra} E. M. Rains, \emph{Correlation functions for symmetrized increasing subsequences}, arXiv preprint math/0006097 (2000).
\bibitem{RT} B. Rodgers, T. Tao, \emph{The de Bruijn--Newman constant is non-negative}, Forum Math. Pi \textbf{8} (2020), e6; doi:10.1017/fmp.2020.6; arXiv:1801.05914.
\bibitem{RW} P. C. Rosenbloom, D. V. Widder, \emph{Expansions in terms of heat polynomials and associated functions}, Trans. Amer. Math. Soc. \textbf{92} (1959), 220--266.
\bibitem{SV} L. A. Shepp, R. J. Vanderbei, \emph{The complex zeros of random polynomials}, Trans. Amer. Math. Soc. \textbf{347} (1995), 4365--4384; doi:10.1090/S0002-9947-1995-1308023-8.
\bibitem{SNWW} W. Shi, G. Nemes, X.-S. Wang, R. Wong, \emph{Error bounds for the asymptotic expansions of the Hermite polynomials}, Proc. Roy. Soc. Edinburgh Sect. A \textbf{153} (2023), 417--440; doi:10.1017/prm.2021.90; arXiv:2109.01054.
\bibitem{Shi} T. Shirai, \emph{Limit theorems for random analytic functions and their zeros}, RIMS K\^oky\^uroku Bessatsu \textbf{B34} (2012), 335--359.
\bibitem{Sh} J. Shohat, \emph{On mechanical quadratures, in particular, with positive coefficients}, Trans. Amer. Math. Soc. \textbf{42} (1937), 461--496; doi:10.1090/S0002-9947-1937-1501930-6.
\bibitem{Si} B. Simon, \emph{Trace Ideals and Their Applications}, 2nd ed., Math. Surveys Monogr. \textbf{120}, Amer. Math. Soc., Providence, RI, 2005.
\bibitem{ST} M. Sodin, B. Tsirelson, \emph{Random complex zeroes, III. Decay of the hole probability}, Israel J. Math. \textbf{147} (2005), 371--379; arXiv:math/0312258.
\bibitem{Sz} G. Szeg\H{o}, \emph{Orthogonal Polynomials}, 4th ed., Amer. Math. Soc. Colloq. Publ. \textbf{23}, Amer. Math. Soc., Providence, RI, 1975.
\bibitem{Tu} P. Tur\'an, \emph{Hermite-expansion and strips for zeros of polynomials}, Arch. Math. (Basel) \textbf{5} (1954), 148--152.
\bibitem{Wat} G. N. Watson, \emph{A Treatise on the Theory of Bessel Functions}, 2nd ed., Cambridge University Press, Cambridge, 1944.
\bibitem{We} H. Weyl, \emph{Inequalities between the two kinds of eigenvalues of a linear transformation}, Proc. Nat. Acad. Sci. U.S.A. \textbf{35} (1949), 408--411; doi:10.1073/pnas.35.7.408.
\end{thebibliography}
\end{document}
